
\documentclass{article} % For LaTeX2e
\usepackage{iclr2026_conference,times}

% Optional math commands from https://github.com/goodfeli/dlbook_notation.
\input{math_commands.tex}

\usepackage{booktabs}
\usepackage{multirow}


\usepackage{hyperref}
\usepackage{url}
\usepackage{amsthm}
\usepackage{graphicx}
\usepackage{amssymb}
\usepackage{mathtools}
\newtheorem{theorem}{Theorem}
\newtheorem{lemma}{Lemma}
\newtheorem{corollary}{Corollary}
\newtheorem{definition}{Definition}  

\title{Flat Minima in the Adapter Subspace Matters for Continual Learning: A PAC-Bayesian Perspective}

% Authors must not appear in the submitted version. They should be hidden
% as long as the \iclrfinalcopy macro remains commented out below.
% Non-anonymous submissions will be rejected without review.

\author{Antiquus S.~Hippocampus, Natalia Cerebro \& Amelie P. Amygdale \thanks{ Use footnote for providing further information
about author (webpage, alternative address)---\emph{not} for acknowledging
funding agencies.  Funding acknowledgements go at the end of the paper.} \\
Department of Computer Science\\
Cranberry-Lemon University\\
Pittsburgh, PA 15213, USA \\
\texttt{\{hippo,brain,jen\}@cs.cranberry-lemon.edu} \\
\And
Ji Q. Ren \& Yevgeny LeNet \\
Department of Computational Neuroscience \\
University of the Witwatersrand \\
Joburg, South Africa \\
\texttt{\{robot,net\}@wits.ac.za} \\
\AND
Coauthor \\
Affiliation \\
Address \\
\texttt{email}
}

% The \author macro works with any number of authors. There are two commands
% used to separate the names and addresses of multiple authors: \And and \AND.
%
% Using \And between authors leaves it to \LaTeX{} to determine where to break
% the lines. Using \AND forces a linebreak at that point. So, if \LaTeX{}
% puts 3 of 4 authors names on the first line, and the last on the second
% line, try using \AND instead of \And before the third author name.

\newcommand{\fix}{\marginpar{FIX}}
\newcommand{\new}{\marginpar{NEW}}

%\iclrfinalcopy % Uncomment for camera-ready version, but NOT for submission.
\begin{document}


\maketitle

\begin{abstract}
Parameter-efficient continual learning (PECL) has become a practical way to reuse pre–trained large models (PLM) on evolving task sequences, but shared parameter subspaces still cause task interference, disrupting the stability–plasticity balance. While flat minima correlate with better generalization and robustness, it is unclear whether targeting flatness solely within the adapter subspace improves performance, and how sensitive the results are across adapter designs.
We explore it under vaious flatness improve methods across different adapter designs and empirically show that reducing sharpness within the adapter subspace consistently improves generalization and reduces forgetting on CL benchmarks. Then We propose a hierarchical PAC-Bayesian framework to explain it theorecthlly that operates directly in the adapter subspace, combining empirical risk, an adapter-subspace expected-sharpness term that quantifies landscape flatness, and a hierarchical KL term that captures cross-task posterior complexity.  
This work establishes a theoretical foundation for adapter-subspace sharpness analysis in LoRA-based continual learning. We demonstrate that reducing sharpness solely within the trainable adapter subspace is a sufficient and robust strategy for improving generalization and stability, offering guidance for developing future PECL methods.




%We discuss different perturbation strategies and empirically show that reducing sharpness within the adapter subspace consistently improves generalization and reduces forgetting on CL benchmarks with varying task similarity and sequence length.

%Flat minima in the loss landscape  associate  with improved generalization and robustness,  but it remains unclear whether controlling flatness within the adapter subspace suffices can improve CL performance and how much it can helps under different adapter design.


%Although flat minima  enhance generalization, it remains unclear whether controlling sharpness exclusively within the adapter subspace suffices for robust CL and how this interacts with adapter design. 




% We address this gap by deriving a hierarchical PAC-Bayesian framework that augments empirical risk with an \emph{adapter-subspace expected sharpness} term and a \emph{hierarchical KL} term that captures posterior complexity across tasks. Under adapter orthogonality, task-level KLs decouple, yielding testable predictions about (i) the sufficiency of subspace-only sharpness control, (ii) which perturbation geometry and strength (SAM, Gaussian, Fisher-guided) win under different sequence lengths and task similarity, and (iii) how adapter structures (SeqLoRA, IncLoRA, OLoRA) interact with sharpness control. Using a unified protocol and multiple seeds, we show that subspace-only sharpness reduction consistently improves last-task accuracy and reduces forgetting compared with both no-noise and matched-budget full-model counterparts, and it traces favorable efficiency–accuracy frontiers. Ablations over rank, orthogonality, and perturbation strength align with the theory. The result is a theory-to-practice account and actionable design guidance for parameter-efficient CL.

\end{abstract}

% \section{Introduction}
% Large pre-trained models are commonly adapted to evolving tasks via parameter-efficient modules such as \lora.
% Yet, sequential adaptation with shared or partially shared adapters induces interference, challenging the stability--plasticity trade-off.
% A widely held view is that flatter minima correlate with better generalization; however, it is unknown whether \emph{reducing sharpness solely in the trainable adapter subspace} is \emph{sufficient} for CL robustness.

% \paragraph{Contributions.}
% (i) We formulate a PAC-Bayesian bound for CL that augments empirical risk with an \emph{expected-sharpness} term and captures hierarchical posterior complexity across tasks.
% (ii) Under task-wise adapter orthogonality, we show the joint complexity decomposes into a sum of task-level terms, yielding a tractable objective that suggests subspace-aware regularization.
% (iii) We validate bound-driven predictions through RQ-aligned experiments across \seq, \inc, and \olora, showing that perturbing only LoRA parameters (Gaussian, Fisher-guided, \sam) improves CL metrics, while full-model perturbations may degrade them.

\section{Introduction}

Continual learning with parameter–efficient adaptation (PECL) has become a practical way to reuse large pre–trained models on evolving task sequences.  Among these methods, Low-Rank Adaptation (LoRA) freezes the backbone and trains compact adapters, offering substantial memory and compute savings while preserving upstream capacities. However, in sequential settings, partially shared adapter subspaces still suffer from task interference, which challenges the trade-off between stability and plasticity.


Recent studies associate flat minima in the loss landscape with improved generalization and robustness, since flatter solutions are more stable under parameter perturbations~\cite{wu2017understandinggeneralizationdeeplearning,chaudhariEntropySGDBiasingGradient2019,NEURIPS2020_518a38cc,NEURIPS2021_995f5e03,JMLRAnEmpiricalInvestigatio,yangRevisitingFlatnessAwareOptimization2025}. While flatness-aware training has been effective in full-model regimes, its implications for an adapter subspace remain unclear. In particular, it is unknown whether reducing sharpness only within the LoRA subspace suffices to improve the stability–plasticity balance, and a theoretical connection between a subspace-restricted flatness constraint and generalization guarantees is currently lacking.

We therefore ask whether and how LoRA subspace flatness affects PECL. To make this question operational, we decompose it into three parts: (i) sufficiency, (ii) mechanism, and (iii) governing factors.
\begin{itemize}
  \item \textbf{Q1 (Sufficiency).} Under matched compute and memory budgets, does increasing LoRA subspace flatness consistently improve continual learning performance across NLP and vision, as measured by current task accuracy, forgetting (stability), and forward transfer (plasticity)? How does it shift the balance between stability and plasticity as quantified by these metrics?
 \item \textbf{Q2 (Mechanism).} How does flatness in parameter space affect the loss? Specifically, is its effect mediated by the feature space, such that increased LoRA subspace flatness shapes the learned representations and yields smoother and more stable features that account for the observed performance gains?
\item \textbf{Q3 (Governing factors)} How do LoRA architectures and flatness strategies modulate the magnitude and limits of the effect?
\end{itemize}

 [To be completed.] Through experiments, we observe that: 
(i) improving LoRA subspace flatness affects current-task accuracy, forgetting, and forward transfer across NLP and vision; 
(ii) representation smoothness and stability correlate with the observed gains; 
(iii) the effect scales with the adapter rank, task similarity, sequence length, and learning rate under matched compute and memory budgets.


Beyond the empirical analysis, we develop a theoretical account grounded in a hierarchical PAC-Bayesian framework. We introduce an adapter-subspace sharpness term and a hierarchical KL term, which together connect subspace-restricted geometry to generalization guarantees. We also provide a PCA-based perspective that links subspace sharpness to feature covariance. This perspective clarifies when modular adapters and subspace-level sharpness control are expected to improve stability and plasticity, particularly under near-orthogonality of adapter directions and moderate ranks.


\section{Preliminaries}

\subsection{Parameter efficient Continual Learning}

% We assume an underlying environment with a task distribution $\mathcal{T}$, such that each task $t \in \{1, 2, \ldots, n\}$ is independently sampled from $\mathcal{T}$~\cite{baxter2000model,pentinaPACBayesianBoundLifelong2014} 
 % \in \mathcal{F} =(x_{tj}, y_{tj})

% Let $\mathcal{X}$, $\mathcal{Y}$, $\mathcal{Z}$ denote the input and output spaces respectively, and let $\mathcal{H} \subset \{f : \mathcal{X} \rightarrow \mathcal{Y}\}$ denote the hypothesis space.  
% a training set $S_t = \{z_{tj}\}_{j=1}^{m_t}$ and a test set $T_t$, both drawn i.i.d.\ from a task-specific data distribution $\mathcal{D}_t$





% The empirical risk for function $f$  on task $t$ is defined as
% $
% \hat{L}_{S_t}(f) = \frac{1}{m} \sum_{j=1}^{m} \ell(f,z_j).
% $ 

% Let $W \in \mathbb{R}^d$ denote the  model parameters,
% which induce a predictor $f_W : \mathcal{X} \to \mathcal{Y}$. The loss function is $\ell : \mathcal{H} \times \mathcal{Z} \rightarrow \mathbb{R}^+$ with a constant bounded $K>0$.


Let $\mathcal{X}$ be the input space, $\mathcal{Y}$ the label space, $\mathcal{Z} = \mathcal{X} \times \mathcal{Y}$ the data space and $\mathcal{H}$ denote the hypothesis space.
Let $W \in \mathbb{R}^d$ denote the model parameters, which induce a predictor
$f_W \in \mathcal{H}$, with $f_W : \mathcal{X} \to \mathcal{Y}$.
We consider a bounded loss $\ell : \mathcal{H} \times \mathcal{Z} \to \mathbb{R}^+$,
uniformly bounded by a constant $K>0$.
For notational convenience we write $\ell(W,z)$ to mean $\ell(f_W,z)$.
% In continual learning with $T$ tasks, for each task $t$, the learner is given a dataset
% $S_t = \{ z_{tj} = (x_{tj}, y_{tj}) \}_{j=1}^{m_t},
% z_{tj} \overset{\text{i.i.d.}}{\sim} \mu_t,$
% where $\mu_t$ is a task-dependent distribution over $\mathcal{Z}$ and $m_t = |S_t|$.
In continual learning with $T$ tasks, for each task $t\in{1,\ldots,T}$, the learner receives a dataset $S_t={z_{tj}=(x_{tj},y_{tj})}_{j=1}^{m_t}$  whose samples are i.i.d. draws from a task-dependent distribution $\mathcal{D}_t$ over $\mathcal{Z}$, with $m_t=|S_t|$.
The empirical risk of parameters $W$ on task $t$ is
$\hat{L}_{S_t}(W)
:= \frac{1}{m_t} \sum_{j=1}^{m_t} \ell(W, z_{tj}),
$
and the generalization risk under $\mathcal{D}_t$ is
$L_{\mathcal{D}_t}(W)
:= \mathbb{E}_{z \sim \mathcal{D}_t} \big[ \ell(W, z) \big].
$
After finishing task $t$, we evaluate the current model on all previously seen
tasks $j \in \{1,\ldots,t\}$ to get $ACC_t=\frac{1}{t}\sum_{j=1}^{t}a_{t,j}$. At the end of the sequence we report standard
continual-learning metrics, namely the final accuracy ($\text{ACC}_T$), the averaged accuracy $\overline{\text{ACC}}_T=\frac{1}{T} \sum_{i=1}^T A C C_i$, and backward transfer (BWT).




% At the end of the sequence we report standard
% continual-learning metrics, namely average every-time accuracy (AAA),
% final average accuracy (FAA), and backward transfer (BWT).


% For each task $t$, the agent receives a dataset
% $S_t = \{ z_{tj}=(x_{tj}, y_{tj}) \}_{j=1}^{m_t}$,
% where $z_{tj} \overset{\text{i.i.d.}}{\sim} \mu_t$ for data distribution $\mu_t$ from $\mathcal{Z}$. 
% The empirical risk of parameters $W$ on task $t$ is
% $
% L_{S_t}(W)
% := \frac{1}{m_t} \sum_{j=1}^{m_t} \ell(W, z_{tj}).
% $ and the population risk under $\mu_t$ is
% $L_{\mu_t}(W)
% := \mathbb{E}_{z\sim \mu_t} \big[ \ell(W, z) \big].$
% After sequentially learning tasks $1{:}i$, we evaluate on all seen tasks $j\in{1,\ldots,i}$. At the end of training, we report three standard metrics, average every time accuracy (AAA), final average accuracy (FAA) and Backward Transfer (BWT).


In the context of PECL, we categorize adapter-based methods by whether task adapters share a parameter space on a frozen backbone $W_0$. Approaches with disjoint parameters that train separately and activate one adapter at inference do not induce cross-task sharing are outside the scope of this paper. We focus on methods with a shared parameter space and distinguish three categories.

\begin{itemize}
    \item  \textbf{Fully shared subspace}. A single low-rank subspace is reused across tasks. In \textbf{SeqLoRA}, one LoRA pair $(A,B)$ is updated sequentially, yielding $W = W_0 + AB .$ Training and inference are both restricted to this common subspace.
    \item  \textbf{Partially fused subspaces}. \textbf{IncLoRA} incrementally add a new pair $(A_t,B_t)$ at task $t$ while keeping $W_0$ and earlier adapters fixed:  $W_0 + \sum_{t=1}^N A_t B_t$.  Gated fusion activates several adapters for input $x$, $W = W_0 + \sum_{j \in \mathcal{A}(x)} \pi_j(x) A_j B_j ,$ where $\mathcal{A}(x)$ is the active adapter set and the weights satisfy $\pi_j(x)\ge 0$\cite{Yu_2024_CVPR,zhao-etal-2024-sapt,cheng2025exploiting}.
 
    \item  \textbf{Constraint-augmented sharing}. Explicit regularization shapes inter-task geometry. \textbf{OLoRA}\cite{wang2023orthogonal} enforces orthogonality among adapter directions via $\sum_{i<t}|A^{(i)\top}A^{(t)}|_F^2$, promoting disentangled subspaces. Gradient-level designs, exemplified by \textbf{InfLoRA}\cite{Liang_2024_CVPR}, reduce cross-task interference through alignment penalties or projection, thereby structuring the shared geometry.
\end{itemize}




% into three classes according to the trainable subspace geometry induced by their organization under a shared backbone.



% \begin{itemize}
% \item \textbf{Single subspace adaptation}. Updates are constrained to a single subspace. \textbf{SeqLoRA} updates one LoRA pair $(A,B)$ over time, yielding $W_0 + AB$. {Selector} methods with top-1 routing~\cite{Yu_2024_CVPR}, which train one adapter per task but activate only a single adapter at inference, also fall into this class. 
% \item \textbf{Multi subspace fusion}. The trainable region is formed by fusing multiple task-specific subspaces. \textbf{IncLoRA} incrementally add a new LoRA pair $(A_t, B_t)$ for each task $t$, and only the current pair is updated while all previous modules and $W_0$ are fixed. The resulting model is $W_0 + \sum_{t=1}^N A_t B_t$. \textbf{SD-LoRA} freezes learned directions and optimizes scalar coefficients. Recent sota method such as {MOS}\cite{Sun_Zhou_Zhao_Gan_Zhan_Ye_2025} and {TUNA}\cite{Wang_2025_ICCV} also follow the similar fusion principle. Sparse top-$k$ routing with $k>1$ is included here because multiple subspaces are simultaneously active per forward pass~\cite{Yu_2024_CVPR}.
% \item \textbf{Constraint-Augmented Adaptation}. Explicit regularization shapes the inter-task geometry. \textbf{OLoRA}\cite{wang2023orthogonal}augments the additive form with orthogonality across tasks, $\sum_{i<t}|A^{(i)\top}A^{(t)}|_F^2$, encouraging disentangled subspaces. TUNA \cite{Wang_2025_ICCV} further support the effectiveness of constraint-based control of plasticity and stability.
% \end{itemize}



% constrains task updates within one fixed low-dimensional subspace. SeqLoRA exemplifies this by sequentially updating a single fixed-size LoRA adapter $(A,B)$ (yielding $W_0 + AB$ after $T$ tasks), while selector methods (e.g., \cite{Yu_2024_CVPR}), which train one adapter per task but activate only a single or sparse set at inference, also fall here, as their effective updates remain confined to a single subspace at any time. The second, multi-subspace fusion, constructs the trainable space by fusing multiple task-specific subspaces. IncLoRA allocates a dedicated LoRA pair $(A_t,B_t)$ per task, with updates aggregated as $W_0 + \sum_{t=1}^T A_t B_t$. SD-LoRA further refines this by freezing pre-learned directions and learning scalars to form $W_0 + \sum_{k=1}^t \alpha_k \overline{A_k B_k}$.  methods like MOS (EMA-based merging of adapters) and TUNA (integration into a universal adapter) corroborate this fusion logic. The third, constraint-augmented adaptation, explicitly introduces regularization to shape inter-task geometry. OLoRA extends the additive form with a cross-task orthogonality penalty $\sum_{i<t} \lVert A^{(i)\top} A^{(t)} \rVert_F^2$ to encourage disentangled subspaces. while TUNA’s adoption of orthogonality constraints further validates the role of explicit constraints in balancing plasticity and stability.





%In the end,  Final Average Accaury (FAA) measures end-state performance, Average every time accuarcy (AAA) aggregates performance over the entire learning path, and Backwards Transfer (BWT) quantifies the forgetting.




% an accuracy matrix $\mathbf{R}\in\mathbb{R}^{T\times T}$ with
% $
% R_{i,j}=\text{Acc}(f_{\boldsymbol{\theta}_i};\mathcal{D}_j).
% $
% The diagonal ${R_{i,i}}$ quantifies performance on the current task; the strict lower triangle captures retention of previous tasks; the strict upper triangle captures performance of unseen tasks; and the last row summarizes the final performance. We summarize trajectories with three metrics. Briefly, Final Average Accaury (FAA) measures end-state performance, Average every time accuarcy (AAA) aggregates performance over the entire learning path, and Backwards Transfer (BWT) quantifies the forgetting.
% $$
% \text{FAA}=\frac{1}{T}\sum_{j=1}^{T}R_{T,j},\quad
% \text{AAA}=\frac{2}{T(T+1)}\sum_{i=1}^{T}\sum_{j=1}^{i}R_{i,j},\quad
% \text{BWT}=\frac{1}{(T-1)}\sum_{i=1}^{T}R_{T,i}-R_{i,i}.
% $$
%updates a single low-rank pair $(A,B)$ across tasks while $W_0$ remains frozen.


% Modern large-scale pretrained models are often adapted to downstream tasks using parameter-efficient methods. In the context of PECL, we focus on three adapter organizations that induce distinct subspace geometries under a common backbone:
% \textbf{Shared-subspace adaptation.}
% \textbf{SeqLoRA} sequentially updates a single fixed-size LoRA adapter $(A,B)$ across tasks while $W_0$ remains frozen and no explicit anti-forgetting mechanism is used. After $T$ tasks, the weights are $W_0 {+} AB$, and all updates lie in one shared low-dimensional subspace.
% \textbf{Task-wise adaptation with additive span.}
% \textbf{IncLoRA} allocates a dedicated LoRA pair per task over a sequence of tasks, yielding $W_0 {+} \sum_{t=1}^{T} A_t B_t$. The trainable region becomes the additive span of task-specific subspaces, which structurally reduces parameter overlap.
% \textbf{SD-LoRA} further decouples direction and magnitude by freezing previously learned directions $\overline{A_k B_k} := (A_k B_k)/\lVert A_k B_k\rVert_F$ and learning only scalars $\alpha_k$, giving $W_0 {+} \sum_{k=1}^{t} \alpha_k \overline{A_k B_k}$ and constraining adaptation to a fixed low-loss path.
% \textbf{MOS} adopts an EMA-based adapter merging strategy where the current adapter is merged with its predecessors after each update. \textbf{TUNA} integrate task-specific adapters into a universal adapter.
% Beyond the methods above that fuse or span multiple adapters in parameter space, selector methods train one adapter per task and use a router at inference to activate a single or a sparse set of adapters\cite{Yu_2024_CVPR}. In subspace terms, the effective update is confined to one low-dimensional subspace at a time; consequently, performance is driven mainly by the router rather than by properties of a fused adapter geometry. 
% \textbf{Orthogonality-regularized adaptation.}
% Having contrasted fusion based organizations, we turn to methods that explicitly regularize inter-task geometry. Recent work explores orthogonality strategies to constrain parameter updates and reduce inter-task interference.
% \textbf{OLoRA} preserves the additive form but separates task subspaces via a cross-task orthogonality penalty $\sum_{i<t} \lVert A^{(i)\top} A^{(t)} \rVert_F^2$, encouraging near-orthogonal and more disentangled adapters. Recent state-of-the-art methods, such as \textbf{TUNA}, likewise adopt orthogonality constraints to balance plasticity and stability.





% \subsection{PAC-Bayesian Generalization Theory}











% % _{\boldsymbol{\theta}}
% \subsection{Supervised learning and Continual Learning}
% Let $\mathcal{X}$ and $\mathcal{Y}$ be the input and label spaces and let $\mathcal{D}$ be a distribution on $\mathcal{X}\times\mathcal{Y}$. The training sample is $S=\{(x_i,y_i)\}_{i=1}^{n}$ with $(x_i,y_i)\overset{\mathrm{i.i.d.}}{\sim}\mathcal{D}$. A model $f\in\mathcal{F}$ with parameters $\boldsymbol{\theta}\in\Theta$ is trained using a loss $\ell:\mathcal{F}\times\mathcal{X}\times\mathcal{Y}\to\mathbb{R}$, yielding empirical and population risks
% \[
% \hat{L}_{S}(f_{\boldsymbol{\theta}})=\tfrac{1}{n}\sum_{i=1}^{n}\ell\!\big(f_{\boldsymbol{\theta}}(x_i),y_i\big),
% \qquad
% L_{\mathcal{D}}(f_{\boldsymbol{\theta}})=\mathbb{E}_{(x,y)\sim\mathcal{D}}\!\big[\ell(f_{\boldsymbol{\theta}}(x),y)\big].
% \]
% The learning goal is to find $\hat{\boldsymbol{\theta}}$ with small $\hat{L}_{S}$ while keeping $L_{\mathcal{D}}$ small; since optimization is performed on the finite sample $S$ but the desired performance is required under $\mathcal{D}$ on $\mathcal{X}\times\mathcal{Y}$, the following PAC--Bayes bound quantifies this discrepancy.

% %Stochastic Model Selection
% \textbf{Theorem 1 (PAC-Bayesian).\citep{mcallesterPACBayesianStochasticModel2003}}
% Let $\mathcal{F}$ be a (possibly uncountable) hypothesis class and $P$ a prior on $\mathcal{F}$. Assume $\ell\in[0,1]$. For any posterior $Q$ on $\mathcal{F}$ and any $S\sim\mathcal{D}^{n}$, with probability at least $1-\delta$,
% \begin{equation}
% \mathbb{E}_{f\sim Q}\!\big[L_{\mathcal{D}}(f)\big]
% \;\le\;
% \mathbb{E}_{f\sim Q}\!\big[\hat{L}_{S}(f)\big]
% \;+\;
% \sqrt{\frac{\mathrm{KL}(Q\Vert P)\;+\;\ln\!\tfrac{1}{\delta}\;+\;\ln n\;+\;2}{\,2n-1\,}}\,.
% \end{equation}
% where $\mathbb{E}[\cdot]$ denotes expectation and $\mathrm{KL}(\cdot\Vert\cdot)$ the Kullback--Leibler divergence.



% We formalize CL within the task environment framework~\cite{baxter2000model,pentinaPACBayesianBoundLifelong2014}, assuming an environment with a distribution over tasks $\mathcal{T}$ that generates a sequence of $T$ tasks with data distributions ${\mathcal{D}_t}$. 
% At step $t$, we observe an i.i.d. sample $S_t={(x_{tj},y_{tj})}\sim\mathcal{D}_t$ and update the model $f_{\boldsymbol{\theta}}$ by empirical risk minimization (ERM) on $S_t$. There is also a PAC-Bayesian generalization bound for CL, which quantifies the transfer of knowledge from observed tasks to future unseen tasks and provides theoretical guarantees\cite{pentinaPACBayesianBoundLifelong2014}.

% Pentina et al. (2014) established a PAC-Bayesian generalization bound for lifelong learning, which quantifies the transfer of knowledge from observed tasks to future unseen tasks and provides theoretical guarantees for lifelong learning algorithms.

% \begin{theorem} \cite{pentinaPACBayesianBoundLifelong2014}
% For any $\delta>0$ the following inequality holds with probability at least $1-\delta$ (over the training samples $\left\{S_1, \ldots, S_n\right\}$ ) for all hyperposterior distributions $\mathcal{Q}$

% $$
% \begin{aligned}
% & L(\mathcal{Q}) \leq \widehat{L}(\mathcal{Q})+\frac{1}{\sqrt{n}}\left(K L(\mathcal{Q} \| \mathcal{P})+\frac{1}{8}-\log \frac{\delta}{2}\right) \\
% & +\frac{1}{n \sqrt{\bar{m}}} K L\left(\left(\mathcal{Q}, Q^n\right) \|\left(\mathcal{P}, P^n\right)\right)+\frac{1}{\sqrt{\bar{m}}}\left(\frac{1}{8}-\frac{1}{n} \log \frac{\delta}{2}\right)
% \end{aligned}
% $$
% where $\left(\mathcal{Q}, Q^n\right)=\mathcal{Q} \times \prod_{i=1}^n Q_i$ denotes the distribution in which we first sample $P$ according to $\mathcal{Q}$ and then use it and the data $S_i$ to produce a posterior $Q_i$ for each task $t_i$. $\left(\mathcal{P}, P^n\right)=\mathcal{P} \times \prod_{i=1}^n P$ denotes the distribution in which we sample $P$ according to $\mathcal{P}$ and use it as a posterior for all tasks. $\bar{m}=\left(\frac{1}{n} \sum_{i=1}^n \frac{1}{m_i}\right)^{-1}$ is the harmonic mean of the sample sizes.
% \end{theorem}


% After learning tasks $1{:}i$, we evaluate on all task $j\in \{1,\ldots,T\}$ to obtain an accuracy matrix $\mathbf{R}\in\mathbb{R}^{T\times T}$ with
% $
% R_{i,j}=\text{Acc}(f_{\boldsymbol{\theta}_i};\mathcal{D}_j).
% $
% The diagonal ${R_{i,i}}$ quantifies performance on the current task; the strict lower triangle captures retention of previous tasks; the strict upper triangle captures performance of unseen tasks; and the last row summarizes the final performance. We summarize trajectories with three metrics. Briefly, Final Average Accaury (FAA) measures end-state performance, Average every time accuarcy (AAA) aggregates performance over the entire learning path, and Backwards Transfer (BWT) quantifies the forgetting.
% $$
% \text{FAA}=\frac{1}{T}\sum_{j=1}^{T}R_{T,j},\quad
% \text{AAA}=\frac{2}{T(T+1)}\sum_{i=1}^{T}\sum_{j=1}^{i}R_{i,j},\quad
% \text{BWT}=\frac{1}{(T-1)}\sum_{i=1}^{T}R_{T,i}-R_{i,i}.
% $$


% \begin{theorem} \cite{pentinaPACBayesianBoundLifelong2014}
% For any $\delta>0$ the following inequality holds with probability at least $1-\delta$ (over the training samples $\left\{S_1, \ldots, S_n\right\}$ ) for all hyperposterior distributions $\mathcal{Q}$

% $$
% \begin{aligned}
% & L(\mathcal{Q}) \leq \widehat{L}(\mathcal{Q})+\frac{1}{\sqrt{n}}\left(K L(\mathcal{Q} \| \mathcal{P})+\frac{1}{8}-\log \frac{\delta}{2}\right) \\
% & +\frac{1}{n \sqrt{\bar{m}}} K L\left(\left(\mathcal{Q}, Q^n\right) \|\left(\mathcal{P}, P^n\right)\right)+\frac{1}{\sqrt{\bar{m}}}\left(\frac{1}{8}-\frac{1}{n} \log \frac{\delta}{2}\right)
% \end{aligned}
% $$
% where $\left(\mathcal{Q}, Q^n\right)=\mathcal{Q} \times \prod_{i=1}^n Q_i$ denotes the distribution in which we first sample $P$ according to $\mathcal{Q}$ and then use it and the data $S_i$ to produce a posterior $Q_i$ for each task $t_i$. $\left(\mathcal{P}, P^n\right)=\mathcal{P} \times \prod_{i=1}^n P$ denotes the distribution in which we sample $P$ according to $\mathcal{P}$ and use it as a posterior for all tasks. $\bar{m}=\left(\frac{1}{n} \sum_{i=1}^n \frac{1}{m_i}\right)^{-1}$ is the harmonic mean of the sample sizes.
% \end{theorem}



% After learning tasks $1{:}i$, we evaluate on all tasks $j\in\{1,\dots,T\}$ and form an accuracy matrix $\mathbf{R}\!\in\!\mathbb{R}^{T\times T}$ with entries $R_{i,j}$.
% The diagonal quantifies the performance of the current task, the strict lower triangle records the retention of past tasks, the strict upper triangle reflects the performance of unseen tasks, and the last row $\{R_{T,j}\}_{j=1}^{T}$ summarizes the last performance.


% We report three trajectory-level summaries. 
% (i) {Final average accuracy (FAA)} averages the last row of the accuracy matrix  and is the standard result in class incremental learning evaluation.
% (ii) {Accumulated average accuracy (AAA)} averages the mean accuracy on all tasks seen so far. It equivalent the full average over the lower triangle including the diagonal.
% (iii) {Forward transfer (FWT)} is averaging the strict upper triangle.

% After completing step $i$, the model is evaluated on all tasks $j\in\{1,\dots,T\}$ to form an accuracy matrix $\mathbf{R}\in\mathbb{R}^{T\times T}$ with $R_{i,j}$ the test accuracy on task $j$ after learning tasks $1{:}i$. The diagonal $\{R_{t,t}\}$ captures {current-task performance} (plasticity); the strict lower triangle $\{R_{t,i}:t>i\}$ records how performance on past tasks evolves as new tasks are learned (stability/forgetting); and the strict upper triangle $\{R_{t,j}:t<j\}$ measures performance on unseen tasks (forward generalization). Let $R_{0,j}$ be the pre-trained (zero-shot) accuracy on task $j$. We summarize the whole trajectory with metrics:
% \begin{align*}
% \mathrm{LA}&:=\frac{1}{T}\sum_{j=1}^{T}R_{T,j}, \\
% \mathrm{BWT}&:=\frac{1}{T-1}\sum_{t=2}^{T}\Bigg[\frac{1}{t-1}\sum_{i=1}^{t-1}\big(R_{t,i}-R_{i,i}\big)\Bigg], \\
% \mathrm{FWT}&:=\frac{1}{T-1}\sum_{t=1}^{T-1}\Bigg[\frac{1}{T-t}\sum_{j=t+1}^{T}\big(R_{t,j}-R_{0,j}\big)\Bigg]
% \end{align*}





% Let $R\in\mathbb{R}^{T\times T}$ be the evaluation matrix, where $R_{i,j}$ is the test accuracy on task $j$ after finishing training up to task $i$ (for $i<j$ this is performance on an unseen task), and let $R_{0,j}$ denote the pre-trained (zero-shot) accuracy on task $j$. 
% CL must balance three complementary objectives: performance on the current task which is the diagnol, retention of prior knowledge (stability) which is the lower  triangle, and transfer to unseen tasks (Generalizability) which is the the Upper triangle. Specific, the calculation is as follows:  



% Accordingly, we report per-step metrics—current accuracy (CA), backward transfer (BWT), and forward generalization (FG):
% \begin{equation*}
% \mathrm{CA}_t:=R_{t,t},\quad
% \mathrm{BWT}_t:=\frac{1}{t-1}\sum_{i=1}^{t-1}\big(R_{t,i}-R_{i,i}\big)\ (t>1),\quad
% \mathrm{FG}_t:=\frac{1}{T-t}\sum_{j=t+1}^{T}R_{t,j}\ (t<T).
% \end{equation*}

% To aggregate performance over the full training trajectory, we report sequence-level summaries:
% \begin{equation*}
% \mathrm{ACA}:=\frac{1}{T}\sum_{t=1}^{T}R_{t,t},\quad
% \mathrm{ABWT}:=\frac{1}{T-1}\sum_{t=2}^{T}\mathrm{BWT}_t,\quad
% \mathrm{AFG}:=\frac{1}{T-1}\sum_{t=1}^{T-1}\mathrm{FG}_t.
% \end{equation*}

% To characterize the final model, we also report the final average accuracy:
% $$\mathrm{FAA}:=\frac{1}{T}\sum_{j=1}^{T}R_{T,j},\qquad$$

% All metrics are computed from the same evaluation matrix $\mathbf{R}$, which fully specifies the temporal evolution of performance across tasks.

\subsection{Flat minima}

% Following the view that sharpness reflects the sensitivity of a solution to perturbations in parameter space, we consider representative definitions widely used in the literature.

 


% Greater flatness tends to shrink the generalization gap that is the difference between population risk and empirical risk\cite{AWP2020}. However, high test robustness also requires that the training procedure attains sufficiently strong robustness in distribution, flatness alone is not sufficient.




Flat minima are widely associated with improved generalization and robustness because they reduce the sensitivity of the empirical risk to small perturbations in parameter space~\cite{FlatMinima1997,wu2017understandinggeneralizationdeeplearning,NEURIPS2020_518a38cc,bahri2022sharpnessawareminimizationimproveslanguage,NEURIPS2021_995f5e03,JMLRAnEmpiricalInvestigatio,yangRevisitingFlatnessAwareOptimization2025}.
We formalize flatness through sensitivity to parameter perturbations and introduce representative  sharpness definitions widely used in the literature.
The zeroth order sharpness measures the maximal loss increase within a radius $\rho$~\cite{AWP2020,foretSharpnessAwareMinimizationEfficiently2021}: 
$$\mathrm{Sh}^{(0)}_S(W,\rho) :=\max_{{\epsilon}\in B(W,\rho)} \Big[\ \hat{L}_{S}(W+{\epsilon})-\hat{L}_{S}(W)\ \Big].$$
A distributional variant averages the loss increase under random perturbations drawn from a prescribed distribution $Q$~\cite{liRevisitingRandomWeight2024,jiang2019fantasticgeneralizationmeasures}:
$$\mathrm{E\text{-}Sh}_t(\hat{W}_t, \Sigma_t)
 :=
 \mathbb{E}_{\varepsilon \sim \mathcal{N}(0,\Sigma_t)}
 \big[
 \hat{L}_{S_t}(\hat{W}_t+\varepsilon) - \hat{L}_{S_t}(\hat{W}_t)
 \big],$$
Typical choices for $\mathcal{N}(0,\Sigma_t)$ include an isotropic Gaussian and a filter-wise Gaussian~\cite{pmlr-v151-bisla22a}. The first order sharpness controls the maximal gradient norm in a neighborhood~\cite{Zhang_GAM}, which is consistent with the zeroth order definition:
$$\mathrm{Sh}^{(1)}_S(W,\rho) := \rho\cdot \max_{W'\in B(W,\rho)} \big\|\nabla_{W} \hat{L}_S(W')\big\|_2,$$





A second-order Taylor expansion of $ \hat{L}_S$ around $W$ gives
$$ \hat{L}_S(W+\epsilon) =  \hat{L}_S(W) + g(W)^{\top}{\epsilon} + \tfrac12\,{\epsilon}^{\top}H(W)\,{\epsilon} + R_3(W,{\epsilon}),$$
with $g(W)=\nabla \hat L_S(W)$, $H(W)=\nabla^2 \hat L_S(W)$, and remainder $R_3(W,{\epsilon})=O(|{\epsilon}|^3)$.
Near a flat local minima, $g(\hat{W})\approx \mathbf 0$, so for small $\|{\epsilon}\|$,
$
\hat{L}_S(\hat{W}+{\epsilon})-\hat{L}_S(\hat{W})
\approx \tfrac12\,{\epsilon}^{\top}H(\hat{W})\,{\epsilon}.
$
This approximation yields concrete correspondences between the sensitivity-based notions above and spectral properties of $H(\hat{W})$.
% \begin{align*}
% \mathrm{Sh}^{(0)}\!\left(\hat{\boldsymbol{\theta}},\rho\right)
% &\;\approx\; \tfrac12\,\rho^{2}\,\lambda_{\max}\!\big(H(\hat{\boldsymbol{\theta}})\big),\\
% \mathrm{E\text{-}Sh}\!\left(\hat{\boldsymbol{\theta}}; \mathcal{N}(0,\sigma^{2}I)\right)
% &\;\approx\; \tfrac12\,\sigma^{2}\,\mathrm{tr}\!\big(H(\hat{\boldsymbol{\theta}})\big),\\
% \mathrm{Sh}^{(1)}\!\left(\hat{\boldsymbol{\theta}},\rho\right)
% &\;\approx\; \rho\,\big\|\nabla \hat L_S(\hat{\boldsymbol{\theta}})\big\|_2
% \;\;\text{and is upper bounded in terms of}\;\;\lambda_{\max}\!\big(H(\hat{\boldsymbol{\theta}})\big)\ 
% \end{align*}
Hence, zeroth order sharpness aligns with the dominant curvature, distributional sharpness aligns with the aggregate curvature, and first order sharpness controls the local growth through gradient magnitude, which is governed by the Hessian near a minimum. These links justify the use of spectral proxies such as $\lambda_{\max}\big(H(W)\big)$ and $\mathrm{tr}\big(H(\hat{W})\big)$, as well as Fisher-based surrogates for negative log-likelihood objectives, where the Fisher equals the expected Hessian under standard regularity conditions. In practice, the empirical Fisher is also a common estimator~\cite{dinhSharpMinimaCan2017,caoTradeoffFlatnessOptimization2024,tsuzukuNormalizedFlatMinima2020,miMakeSharpnessAwareMinimization2022}.

To visualize local sharpness, one can plot one-dimensional loss slices along random or adversarial directions, two-dimensional planes spanned by normalized directions, or three-dimensional surface plots~\cite{goodfellowQualitativelyCharacterizingNeural2015,izmailovAveragingWeightsLeads2019,chaSWADDomainGeneralization2021,liVisualizingLossLandscape2018}.

% For the sequential task setting of Continual Learning (CL), Pentina et al. \cite{pentinaPACBayesianBoundLifelong2014} extended the PAC-Bayesian framework to derive a task-aware generalization bound, 




% Near a flat local minimizer $\hat{\boldsymbol{\theta}}$, if it is assumed $g(\hat{\boldsymbol{\theta}})\approx \mathbf 0$, then
% $$\hat{L}_S(\hat{\boldsymbol{\theta}}+\boldsymbol{\epsilon})-\hat{L}_S(\hat{\boldsymbol{\theta}}) \;\approx\; \tfrac12\,\boldsymbol{\epsilon}^{\top}H(\hat{\boldsymbol{\theta}})\,\boldsymbol{\epsilon}.$$


% This quadratic approximation motivates spectral proxies for local sharpness, including the largest eigenvalue $\lambda_{\max}\big(H(\hat{\boldsymbol{\theta}})\big)$ and the trace $\mathrm{tr}\big(H(\hat{\boldsymbol{\theta}})\big)$ of the Hessian, which are widely used in practice ~\citep{Zhang_GAM,dinhSharpMinimaCan2017,caoTradeoffFlatnessOptimization2024,tsuzukuNormalizedFlatMinima2020}. For negative log-likelihood objectives, under standard regularity conditions the Fisher information equals the expected Hessian; in practice the empirical Fisher also serves as an estimator and a curvature surrogate~\citep{miMakeSharpnessAwareMinimization2022}. To visually diagnose sharpness, common approaches include 1D loss slices along random or adversarial directions, 2D planes spanned by normalized directions, and 3D surface plots ~\citep{goodfellowQualitativelyCharacterizingNeural2015,izmailovAveragingWeightsLeads2019,chaSWADDomainGeneralization2021,liVisualizingLossLandscape2018}.

%. SSAM leverages the empirical Fisher to prioritize which parameters to perturb 


\subsection{PAC-Bayesian Generalization Bound}


The flatness metric above quantifies the local sensitivity of empirical loss to parameter perturbations. To formalize generalization guarantees, the PAC settings place probability measures on the hypothesis space rather than focusing on a single function. Let $P$ be a prior and $Q$ a posterior over hypothesis space. With high probability over $S$, PAC-Bayesian bounds control the generalization risk  by the empirical risk plus a complexity term that scales with $\mathrm{KL}(Q\Vert P)$ and the sample size\cite{shawe1997pac,mcallesterPACBayesianStochasticModel2003}. 
This framework formalizes the flatness–generalization link: flatter minima admit broader posteriors that remain close to the prior, thereby reducing $\mathrm{KL}(Q\Vert P)$ and tightening the bound. Algorithmically, SAM minimizes the worst-case empirical loss within a small parameter neighborhood $\max_{\|\varepsilon\|\le\rho}\hat L_{S}(W+\varepsilon)$, which serves as a practical surrogate that reduces sharpness and aligns with the PAC-Bayesian objective\cite{foretSharpnessAwareMinimizationEfficiently2021}.


%  The generalization risk is
% $L_{\mathcal{D}}(Q):=\mathbb{E}_{f\sim Q}\,\mathbb{E}_{(x,y)\sim\mathcal{D}}\big[\ell(f(x),y)\big]$,
% and the empirical risk is
% $\hat L_{S}(Q):=\mathbb{E}_{f\sim Q}\,\frac{1}{m}\sum_{j=1}^{m}\ell(f(x_j),y_j)$.

% on a sample $S=\{(x_j,y_j)\}_{j=1}^m$
% To formalize generalization guarantees, the core setting of Probably Approximately Correct (PAC) learning focuses on distributions over hypothesis spaces rather than individual functions. Consider a prior distribution $P$ and a posterior distribution$Q$ defined over the hypothesis space\(\mathcal{F}\). 
% The generalization risk $L(Q) = \mathbb{E}_{f \sim Q} \mathbb{E}_{(x,y)\sim\mathcal{D}}\ell(f(x),y)$  captures the expected error of functions in $Q$ with respect to the true data distribution $\mathcal{D}$. The empirical risk $\hat{L}(Q) = \mathbb{E}_{f \sim Q}\frac{1}{m} \sum_{j=1}^{m} \ell(f(x_{j}), y_{j})$ Averages the error of functions in $Q$ over a finite training sample. For any such $Q$ and $P$, PAC bounds upper-bound the generalization error $L(Q)$ by the empirical risk $\hat{L}(Q)$ plus a complexity term (typically involving the Kullback-Leibler (KL) divergence between $Q$ and $P$) that quantifies the uncertainty of $Q$ relative to $P$. 

% The PAC-Bayesian framework builds on this PAC setting to formalize the flatness-generalization link: flatter minima enable broader posterior distributions. Algorithmically, Sharpness-Aware Minimization (SAM) acts as a practical surrogate: by minimizing the worst-case empirical loss within a small parameter neighborhood, it approximates the optimization of the PAC-Bayesian.







% The flatness metric above quantifies the local sensitivity of empirical loss to parameter perturbations. 
% The PAC-Bayesian framework formalizes this flatness-generalization link: flatter minima enable broader posterior distributions. Algorithmically, Sharpness-Aware Minimization (SAM) acts as a practical surrogate: by minimizing the worst-case empirical loss within a small parameter neighborhood, it approximates the optimization of the PAC-Bayesian objective\cite{foretSharpnessAwareMinimizationEfficiently2021}.

% A hierarchical Bayesian formulation captures the knowledge transfer. A prior $P$ over model hypotheses is drawn from a hyperprior $\mathcal{P}$ and is updated, as tasks arrive, to a hyperposterior $\mathcal{Q}$. For each task $t$ with training sample $S_t$, the learning algorithm produces a task posterior $Q_t(\cdot\,|\,S_t,P)$ over hypotheses conditioned on $P$. Building on this hierarchical structure, it is defined the transfer risk
% $L(\mathcal{Q}) := \mathbb{E}_{(t,S_t)}\mathbb{E}_{P \sim \mathcal{Q}}L\big(Q_t(S_t, P)\big)
% $
% and the empirical multi-task risk
% $\hat{L}(\mathcal{Q}) := \frac{1}{n} \sum_{t=1}^n \mathbb{E}_{P \sim \mathcal{Q}} \hat{L}(Q_i(S_i,P))
% $.  Pentina et al.~\cite{pentinaPACBayesianBoundLifelong2014} derived a multi task PAC-Bayesian generalization bound for transfer learning.

% In CL, a hierarchical Bayesian formulation captures knowledge transfer: a prior $P$ over hypotheses is drawn from a hyperprior $\mathcal{P}$ and, as tasks arrive, updated to a hyperposterior $\mathcal{Q}$. For each task $t$ with training sample $S_t$, the learner produces a task posterior $Q_t(\cdot \mid S_t, P)$.  Pentina et al.~\cite{pentinaPACBayesianBoundLifelong2014} derived a PAC Bayesian bound for multi task under hierarchical Bayesian. However, their analysis assumes that all task priors $P$ are independent draws from a {fixed} hyperprior $\mathcal{P}$, which is appropriate for exchangeable multi task transfer but not for continual learning, where the hyperdistribution is {sequentially updated}.




PAC-Bayesian bound~\cite{pentinaPACBayesianBoundLifelong2014,JMLR:v24:22-1254}  have been employed for multi-task knowledge transfer within a hierarchical Bayesian framework that learns a hyper-posterior. However, their analysis assumes that all task priors $P$ are independent draws from a {fixed} hyperprior, which is suitable for multi-task transfer learning but inadequate for CL. In CL, the hyper-distribution must evolve sequentially to  reflect the collection of historical information, and their framework does not support such dynamic adaptation. 
Recent advances in online PAC-Bayesian learning \cite{NEURIPS2022_onlineBayes} introduce time-varying, data-dependent sequence of priors \( P_i \) (each \( \mathcal{F}_{i-1}\)-measurable) and dynamically updated posteriors \(Q_i\) for sequential data, offering flexibility for sequential data at the point level.  Algorithmically, the NsCE framework decomposes online PAC-Bayesian bound into empirical risk, throughput, and task-divergence terms\cite{NEURIPS2024_xinrui}. However, these methods focus on data-point sequences rather than task-level progression, and their theoretical foundations remain misaligned with structured continual learning frameworks like PECL.


Despite these contributions, current theoretical frameworks exhibit two major limitations in the context of PECL. First, they are geometry-agnostic: they do not incorporate the local flatness of the loss landscape into the divergence terms that govern generalization. Second, they overlook the adapter subspace central to PECL, failing to link curvature within this subspace to PAC-Bayesian divergence. To bridge these gaps, we propose a novel PAC-Bayesian formulation that is both geometry-aware and subspace-aware, specifically tailored to the PECL framework.


% The online PAC-Bayesian setting specifies a sequence of priors $\{P_i\}_{i=1}^m$ with $P_1=P$ data-free distribution and, for $i\ge2$, $P_i$ measurable with respect to $\mathcal{F}_{i-1}$ where $ (\mathcal{F}_{i})_{i>0}$ is an adapted filtration to $S$.
% To refer to a distribution $Q_S$ that depends on a dataset $S$ via $Q$, the stochastic kernel $Q(\cdot, \cdot)$ is introduced such that $Q_S = Q(S, \cdot)$.
% This notation makes explicit when a prior or posterior is adapted to past observations, which is essential for sequential analysis in CL. Beyond theory, online PAC-Bayes has been instantiated into practical algorithms;  Wang et al. formulate the NsCE framework by decomposing an online PAC-Bayesian bound into empirical-risk, throughput, and task-divergence terms~\cite{NEURIPS2024_xinrui}.





% However, two limitations remain across current theory and its use in PECL: it is geometry agnostic, with no mechanism to encode how local flatness enters the divergence terms that govern generalization; and it ignores the adapter subspace central to PECL, offering no link between curvature measured within the adapter span and $\mathrm{KL}(Q\Vert P)$. We address both gaps by developing a geometry-aware, subspace-aware PAC-Bayesian treatment for PECL.

% However, two limitations remain across current theory and its use in PECL: it is geometry agnostic, with no mechanism to encode how local flatness enters the divergence terms that govern generalization; and it ignores the adapter subspace central to PECL, offering no link between curvature measured within the adapter span and $\mathrm{KL}(Q\Vert P)$. We address both gaps by developing a geometry-aware, subspace-aware PAC-Bayesian treatment for PECL.


% Recent work on {online} PAC-Bayesian learning\cite{NEURIPS2022_onlineBayes} admits a sequence of priors that depend on the past and a sequence of posteriors, providing guarantees under data streams and sequential adaptation that are better aligned with continual learning.  Let $\mathcal{Z}$ be the data space and $S\in\mathcal{Z}^m$ be drawn from an unknown distribution $\mu$. The online PAC-Bayesian setting specifies a sequence of priors $\{P_i\}_{i=1}^m$ with $P_1=P$ data-free distribution and, for $i\ge2$, $P_i$ measurable with respect to $\mathcal{F}_{i-1}$ where $ (\mathcal{F}_{i})_{i>0}$is an adapted filtration to $S$.
% set $Q_S=Q(S,\cdot)$ the data-dependent prior associated to the sample S through Q. to refer to a distribution QS depending on a dataset S, we introduce a stochastic kernel Q(., .) such that QS = Q(S, .).




% The online PAC-Bayesian setting specifies a sequence of priors $\{P_i\}_{i=1}^m$ with $P_1=P$ and, for $i\ge2$, $P_i$ measurable with respect to $\mathcal{F}_{i-1}$, hence allowed to depend on the past data. Upon observing $z_i$, the learner outputs a posterior $Q_i\in\mathcal{M}_1(\mathcal{H})$ given by a measurable update rule 
% $Q_i=\Gamma_i(P_1,\ldots,P_i,Q_1,\ldots,Q_{i-1},z_i)$,
% so that the posteriors evolve over time.





% However, two limitations persist. First, existing PAC-Bayesian bounds for continual learning are geometry agnostic: they do not encode how local flatness or sharpness enters the divergence terms. Second, the analyses ignore the adapter subspace central to parameter-efficient training: updates occur within a LoRA subspace, yet the bounds are stated in the full parameter space and provide no link between curvature measured within the adapter span, leaving the influence of subspace organization and cross-task constraints on posterior complexity and forgetting unquantified.



% A prior $P$ over hypotheses is drawn from a hyperprior $\mathcal{P}$ and, as tasks arrive, updated to a hyperposterior $\mathcal{Q}$.

% Maxime propose an online PAC-Bayesian bound as it allows to consider a sequence of priors which may depend on
%  the past and a sequence of posteriors that can dynamically evolve as well. 
 
 
 
 % captures knowledge transfer: a prior $P$ over hypotheses is drawn from a hyperprior $\mathcal{P}$ and, as tasks arrive, updated to a hyperposterior $\mathcal{Q}$. For each task $t$ with training sample $S_t$, the learner produces a task posterior $Q_t(\cdot \mid S_t, P)$. On this hierarchy one defines the transfer risk and the empirical multi task risk

 % However, their analysis assumes that all task priors $P$ are independent draws from a \emph{fixed} hyperprior $\mathcal{P}$, which is appropriate for exchangeable multi task transfer but not for continual learning, where the hyperdistribution is \emph{sequentially updated} and thus breaks the i.i.d.\ assumption across tasks.


% $
% L(\mathcal{Q})
% := \mathbb{E}_{t\sim \mathcal{T}}\,\mathbb{E}_{S_t\sim \mathcal{D}_t^{m_t}}\,
% \mathbb{E}_{P\sim \mathcal{Q}}\,\mathbb{E}_{f\sim Q_t(\cdot\,|\,S_t,P)}
% \,L_{\mathcal{D}_t}(f),
% $
% and the empirical risk
% $
% \hat L(\mathcal{Q})
% := \frac{1}{n}\sum_{t=1}^n
% \mathbb{E}_{P\sim \mathcal{Q}}\,\mathbb{E}_{f\sim Q_t(\cdot\,|\,S_t,P)}
% \,\hat L_{S_t}(f).
% $



% In CL, a hierarchical Bayesian framework is adopted to model accumulated task knowledge: the prior $P$ over model parameters is treated as a random variable drawn from a hyperprior $\mathcal{P}$, and this hyperprior is recursively updated to a hyperposterior $\mathcal{Q}$ as new tasks are observed, effectively encoding the knowledge learned across sequential tasks. Building on this hierarchical structure, Pentina et al. \cite{pentinaPACBayesianBoundLifelong2014} extended the classical PAC-Bayesian framework to CL, deriving a task-aware generalization bound.

% $L(\mathcal{Q}) = \mathbb{E}_{(t,S_t)}\mathbb{E}_{P \sim \mathcal{Q}}L\big(Q_t(S_t, P)\big)$ is the generalization risk, $
% \hat{L}(\mathcal{Q}) = \frac{1}{n} \sum_{t=1}^n \mathbb{E}_{P \sim \mathcal{Q}(P)} \mathbb{E}_{f \sim Q_t}  [ L_{\mathcal{S}_t}(f) ]$

%  $Q_t(S_t,P )$ is the posterior obtained by training the learning algorithm with prior $P$ and training sample $S_t$,



 % $Q_t(S_t,P )$ is the posterior obtained by training the learning algorithm with prior $P$ and training sample $S_t$,



% \begin{theorem} \cite{pentinaPACBayesianBoundLifelong2014}
% For any $\delta>0$ the following inequality holds with probability at least $1-\delta$ (over the training samples $\left\{S_1, \ldots, S_n\right\}$ ) for all hyperposterior distributions $\mathcal{Q}$
% $$
% \begin{aligned}
% L(\mathcal{Q}) \leq \widehat{L}(\mathcal{Q})
% +(\frac{1}{\sqrt{n}} +\frac{1}{n \sqrt{\bar{m}}})KL(\mathcal{Q} \| \mathcal{P})  +\frac{1}{n \sqrt{\bar{m}}} \sum_{i=1}^n \underset{P \sim \mathcal{Q}}{\mathbf{E}} KL\left(Q_i\left(S_i, P\right) \| P\right)+\operatorname{const}(n, \bar{m}, \delta) .
% \end{aligned}
% $$
% \end{theorem}
% where
% $\bar{m}=\left(\frac{1}{n} \sum_{i=1}^n \frac{1}{m_i}\right)^{-1}$ is the harmonic mean of the sample sizes. The bound contains two types of complexity terms that correspond to two levels of the model: $KL(\mathcal{Q} \| \mathcal{P})$ belongs to the level of task environment in general, while each $KL\left(Q_i\left(S_i, P\right) \| P\right)$ corresponds specifically to the $i$-th task.


% However, two limitations persist. First, existing PAC-Bayesian bounds for continual learning are geometry agnostic: they do not encode how local flatness or sharpness enters the divergence terms. Second, the analyses ignore the adapter subspace central to parameter-efficient training: updates occur within a LoRA subspace, yet the bounds are stated in the full parameter space and provide no link between curvature measured within the adapter span, leaving the influence of subspace organization and cross-task constraints on posterior complexity and forgetting unquantified.



% However, a key limitation of this classical bound is that it does not explicitly account for the geometry of the loss landscape which is known to play a critical role in robustness and generalization. Addressing this gap is central to our subsequent theoretical development, where we extend the PAC-Bayesian framework to incorporate flatness-aware terms in the context of continual learning and parameter-efficient adaptation.



% In CL, hierarchical Bayesian framework is induced ,where the prior $P$ over model parameters is treated as a random variable drawn from a hyperprior $\mathcal{P}$, which is recursively updated to a hyperposterior $\mathcal{Q}$ as tasks are observed, thus reflecting accumulated knowledge.  Pentina et al. \cite{pentinaPACBayesianBoundLifelong2014} extended the PAC-Bayesian framework to derive a task-aware generalization bound. 
% In Continual Learning, Pentina et al. \cite{pentinaPACBayesianBoundLifelong2014}  use a hierarchical Bayesian framework  the prior $P$ over model parameters is treated as a random variable drawn from a hyperprior $\mathcal{P}$, which is recursively updated to a hyperposterior $\mathcal{Q}$ as tasks are observed, thus reflecting accumulated knowledge. 
% extended the PAC-Bayesian framework to derive a task-aware generalization bound.





% quantifies the local sensitivity of empirical loss to parameter perturbations, with greater flatness typically narrowing the generalization gap. However, a flatter weight loss landscape does directly lead to a smaller robust generalization gap but is only beneficial to the final test robustness on condition that the training process is sufficient.\cite{AWP2020} The PAC-Bayesian framework formalizes this flatness-generalization relationship:  flatter minima support broader posterior distributions. Algorithmically, SAM can be viewed as a practical surrogate that minimizes the worst-case empirical loss in a small parameter neighborhood, thereby approximately optimizing this PAC-Bayesian objective.







% The flatness introduced above quantify the local sensitivity of the empirical loss to parameter perturbations. Greater flatness tends to shrink the generalization gap that is the difference between population risk and empirical risk\cite{AWP2020}. The PAC-Bayesian framework formalizes the link between flatness and the generalization gap: sharpness quantifies the sensitivity of the loss to parameter perturbations, and flatter minima admit broader posteriors with smaller $\mathrm{KL}(Q\Vert P)$, which tightens the bound. Algorithmically, SAM can be viewed as a practical surrogate that minimizes the worst-case empirical loss in a small parameter neighborhood, thereby approximately optimizing this PAC-Bayesian objective.


% \begin{theorem} \cite{pentinaPACBayesianBoundLifelong2014}
% For any $\delta>0$ the following inequality holds with probability at least $1-\delta$ (over the training samples $\left\{S_1, \ldots, S_n\right\}$ ) for all hyperposterior distributions $\mathcal{Q}$

% $$
% \begin{aligned}
% & L(\mathcal{Q}) \leq \widehat{L}(\mathcal{Q})+\frac{1}{\sqrt{n}}\left(K L(\mathcal{Q} \| \mathcal{P})+\frac{1}{8}-\log \frac{\delta}{2}\right) \\
% & +\frac{1}{n \sqrt{\bar{m}}} K L\left(\left(\mathcal{Q}, Q^n\right) \|\left(\mathcal{P}, P^n\right)\right)+\frac{1}{\sqrt{\bar{m}}}\left(\frac{1}{8}-\frac{1}{n} \log \frac{\delta}{2}\right)
% \end{aligned}
% $$
% where $\left(\mathcal{Q}, Q^n\right)=\mathcal{Q} \times \prod_{i=1}^n Q_i$ denotes the distribution in which we first sample $P$ according to $\mathcal{Q}$ and then use it and the data $S_i$ to produce a posterior $Q_i$ for each task $t_i$. $\left(\mathcal{P}, P^n\right)=\mathcal{P} \times \prod_{i=1}^n P$ denotes the distribution in which we sample $P$ according to $\mathcal{P}$ and use it as a posterior for all tasks. $\bar{m}=\left(\frac{1}{n} \sum_{i=1}^n \frac{1}{m_i}\right)^{-1}$ is the harmonic mean of the sample sizes.
% \end{theorem}


% However, a flatter loss landscape does not directly reduce the robust generalization gap; it only improves the robustness of the final test when training is sufficient \cite{AWP2020}. 




% However, a flatter weight loss landscape does directly lead to a smaller robust generalization gap but is only beneficial to the final test robustness on condition that the training process is sufficient.\cite{AWP2020} 





% The PAC (Probably Approximately Correct) generalization bound provides a principled link between flatness and the generalization gap: sharpness serves as a surrogate for the loss’s sensitivity to parameter perturbations.  Flatter minima enable broader yet low-divergence posteriors relative to a prior , which tightens the PAC-Bayesian bound. 







% The PAC (Probably Approximately Correct) generalization bound provides a principled theoretical link between this flatness and generalization gap,it interprets sharpness as a surrogate for the sensitivity of the model’s posterior distribution (over parameters) to data perturbations, such that flatter loss minima lead to tighter PAC-Bayesian bounds by reducing posterior complexity. 








% In PAC-Bayesian analysis, the same sensitivity appears as complexity: flatter minima permit broader posteriors with smaller divergence from a suitable prior, which tightens the bound on the generalization gap. 





% Parametric Noise Injection: Trainable Randomness to Improve Deep Neural
%  Network Robustness against Adversarial Attack

 
% Weexpecttoutilize the aforemen
% tioned PNI technique to improve the network robustness

% Therefore, a flatter weight loss landscape does directly lead to a smaller robust generalization gap but is only beneficial to the final test robustness on condition that the training process is sufficient (i.e., training robustness is high).

% The PAC-Bayes generalization bound rediscovers the connection of the Hessian as flatness characteristic with the generalization gap and the large-scale empirical evaluation \cite{jiang2019fantasticgeneralizationmeasures} shows that all the generalization measures based on flatness (in some definition) highly correlate with the actual performance of models.









% \section{Q1: Does promoting flatness in the Adapter subspace help across different PECL methods?}

% \section{LoRA-Only Flatness Control Improves CL and Shrinks Global Curvature: A Subspace PAC–Bayes View}


% PECL typically freezes the PLM weights and adapts new tasks by inserting lightweight modules. We focus on the flatness of the lossland and investigate whether promoting flatness where learning actually occurs improves generalization and mitigating forgetting. If generalization and forgetting are caused by the changed parameters, then manipulating flatness in the adapter subspace is the natural causal locus to test.

% \subsection{Empirical lead-in and question}
\section{Flatness Matters for PECL: A PAC-Bayesian Perspective}
\label{section PAC}


% \begin{table}[t]
% \centering
% \caption{Class-incremental learning on ImageNet-R (T=20) using ViT-B/16-IN21K; means~$\pm$~ sample std over 5 seeds with shuffled class orders. Columns report CL performance, sharpness definitions, Hessian and Fisher metrics for all PECL baselines.}
% % .All methods are implemented under a unified codebase.\cite{sun2025pilot,zhou2024class,zhou2024continual}
% \label{tab:q1_ci_with_weight_flatness_reorg}
% \setlength{\tabcolsep}{4pt}
% \resizebox{\linewidth}{!}{
% \begin{tabular}{l l ccc cccccc c}
% \toprule
% \multirow{2}{*}{\textbf{Opt.}} & \multirow{2}{*}{\textbf{Method}} &
% \multicolumn{3}{c}{\textbf{CL ($\uparrow$)}} &
% \multicolumn{3}{c}{\textbf{Sharpness ($\downarrow$)}} &
% \multicolumn{2}{c}{\textbf{Hessian ($\downarrow$)}} &
% \multicolumn{2}{c}{\textbf{Fisher ($\downarrow$)}} \\
% \cmidrule(lr){3-5}\cmidrule(lr){6-8}\cmidrule(lr){9-10}\cmidrule(lr){11-12}
%  & & ${\text{ACC}}_{20}$  & $\overline{\text{ACC}}_{20}$    & BWT
%  & Sh$^{(0)}(\rho)$ & E-Sh$(Q)$ & Sh$^{(1)}(\rho)$
% & $\lambda_{\max}(H)$ & $\mathrm{tr}(H)$  
% & $\lambda_{\max}(F)$ & $\mathrm{tr}(F)$   \\
% \midrule
% % \multirow{6}{*}{\textbf{without SAM}}
% % \multirow{2}{*}{\centering{\textbf{Backbone}}}
% % & MulTaskLP  \\
% % & MulTaskFT   \\
% % & MulTaskLoRA   \\
% % \midrule
% \multirow{7}{*}{\centering\rotatebox[origin=c]{90}{\textbf{without SAM}}}
% & PerTaskLP       \\
% & PerTaskFT        &58.35  &66.99  &-33.95  &0.17 & 0.050 &0.29  & 97.07 &13850.95  &194.21 &10685.23\\
% \cmidrule(lr){2-12}
%  & SeqLoRA        &66.11  &73.10  &-17.28  &0.20  & 0.052  &0.20  &12.68  &1650.03  &65.27  &1409.15\\
% \cmidrule(lr){2-12}
%  & IncLoRA        &74.06  &77.70  &-3.55  &0.20  &0.020  &0.15 &9.58  &827.02  &19.76 &772.93\\
%   & InfLoRA \\
% % & MOS        &  &  &  &  &  &  &  &\\
%  \cmidrule(lr){2-12}
%  & OLoRA          &73.99  &77.24  &-3.6  &0.21  &0.024  &0.15  &9.51 &830.05 &18.76  &764.2\\
%  % & TUNA        &  &  &  &  &  &  &  &\\
% \midrule
% \multirow{7}{*}{\centering\rotatebox[origin=c]{90}{\textbf{with SAM}}}
% & PerTaskLP       \\
% & PerTaskFT        &68.07  &75.18  &-24.78  &0.17    &0.047  &0.19  &12.43 &12118.71 &70.47 &2935.51 \\
% \cmidrule(lr){2-12}
% & SeqLoRA        &71.04  &75.76  &-11.46  &0.16  &0.038 &0.14  &10.79  &918.26  &24.13 &762.88\\
% \cmidrule(lr){2-12}
%  & IncLoRA        &75.03  &78.41  &-3.53  &0.13  &0.009 &0.13  &8.08  &643.08  &15.72 &584.20\\
%  & InfLoRA &  &  &  &  &  &  &  &\\
% % & MOS        &  &  &  &  &  &  &  &\\
%  % & EASE  &  &  &  &  &  &  & & \\
%  \cmidrule(lr){2-12}
%  & OLoRA          &74.71  &77.73  &-3.51  &0.12  &0.006  &0.13  &7.93  &656.99 &15.24 &590.27\\
%  % & TUNA        &  &  &  &  &  &  &  &\\
% \bottomrule
% \end{tabular}
% }
% \end{table}



% PECL confines updates to adapter subspaces such as LoRA. We ask whether promoting flatness {where learning actually occurs} is sufficient to improve CL performance and to reduce global curvature. A preview of the results indicates that LoRA-only flatness improves FAA and AAA while the full-model curvature proxies (for example, $\lambda_{\max}(H)$ and $\mathrm{tr}(H)$, as well as Fisher counterparts) decrease. 
% This pattern is suggestive but not conclusive. To separate mechanism from confounds and to derive testable predictions, we require a statement that links subspace flatness to generalization in a way that is compatible with PAC-Bayesian analysis. This motivates the following theorem. 

% where $\Delta W = AB$ 

% Let $\left(P_i\right)_{i=1}^m$ be an online predictive sequence of stochastic kernels on the LoRA space $\mathcal{H}_{\Delta}$ (i.e., $P_i$ is $\mathcal{F}_{i-1}$-measurable and $P_i \ll P_{i-1}$ ).
% Let $\left(Q_i\right)_{i=1}^m$ be any sequence of stochastic kernels on $\mathcal{H}_{\Delta}$.

% Let $\left\{\mathcal{F}_i\right\}_{i \geq 0}$ be the canonical filtration.

% \in \mathcal{H}_{\Delta}$ so that $\Delta W(\theta)=A B$ (rank$(A,B)\leq r \ll d$ )

% The model is $W(\theta)=W_0+\Delta W(\theta)$ with frozen model $W_0$ and LoRA parameters $\theta= (A, B)$.  Let \(\mathcal{H_\Delta}\) denote the LoRA effective update space, i.e. the set of learnable low-rank increments \(\Delta W(\theta) = A B\) with \(\theta = (A,B)\), while the pretrained backbone \(W_0\) remains frozen. All adaptation during PECL happens only through \(\theta \in \mathcal{H}_\Delta\); the backbone \(W_0\) does not change and is therefore treated as deterministic.

% sampled from a hyperprior $\mathcal{P}$ and updated to a posterior ${Q}$ after observing new tasks

\subsection{Motivating Observation: LoRA-Only Flatness Improves Performance}


\begin{table}[t]
\centering
\caption{Class-incremental learning on ImageNet-R (T=20) using ViT-B/16-IN21K; means~$\pm$~ sample std over 5 seeds with shuffled class orders. Columns report CL performance, sharpness definitions, Hessian and Fisher metrics for all PECL baselines.}
% .All methods are implemented under a unified codebase.\cite{sun2025pilot,zhou2024class,zhou2024continual}
\label{tab:q1_ci_with_weight_flatness_reorg}
\setlength{\tabcolsep}{4pt}
\resizebox{\linewidth}{!}{
\begin{tabular}{l l ccc cccccc c}
\toprule
\multirow{2}{*}{\textbf{Opt.}} & \multirow{2}{*}{\textbf{Method}} &
\multicolumn{3}{c}{\textbf{CL ($\uparrow$)}} &
\multicolumn{3}{c}{\textbf{Sharpness ($\downarrow$)}} &
\multicolumn{2}{c}{\textbf{Hessian ($\downarrow$)}} &
\multicolumn{2}{c}{\textbf{Fisher ($\downarrow$)}} \\
\cmidrule(lr){3-5}\cmidrule(lr){6-8}\cmidrule(lr){9-10}\cmidrule(lr){11-12}
 & & ${\text{ACC}}_{20}$  & $\overline{\text{ACC}}_{20}$    & $\Delta$ 
 % & BWT
 & Sh$^{(0)}(\rho)$ & E-Sh$(Q)$ & Sh$^{(1)}(\rho)$
& $\lambda_{\max}(H)$ & $\mathrm{tr}(H)$  
& $\lambda_{\max}(F)$ & $\mathrm{tr}(F)$   \\
\midrule
% \multirow{6}{*}{\textbf{without SAM}}
% \multirow{2}{*}{\centering{\textbf{Backbone}}}
% & MulTaskLP  \\
% & MulTaskFT   \\
% & MulTaskLoRA   \\
% \midrule
\multirow{7}{*}{\centering\rotatebox[origin=c]{90}{\textbf{without SAM}}}
& PerTaskLP       \\
& PerTaskFT        &58.35  &66.99  &-33.95  &0.17 & 0.050 &0.29  & 97.07 &13850.95  &194.21 &10685.23\\
\cmidrule(lr){2-12}
 & SeqLoRA        &66.11  &73.10  &-17.28  &0.20  & 0.052  &0.20  &12.68  &1650.03  &65.27  &1409.15\\
\cmidrule(lr){2-12}
 & IncLoRA        &74.06  &77.70  &-3.55  &0.20  &0.020  &0.15 &9.58  &827.02  &19.76 &772.93\\
  % & InfLoRA \\
% & MOS        &  &  &  &  &  &  &  &\\
 \cmidrule(lr){2-12}
 & OLoRA          &73.99  &77.24  &-3.6  &0.21  &0.024  &0.15  &9.51 &830.05 &18.76  &764.2\\
 % & TUNA        &  &  &  &  &  &  &  &\\
\midrule
\multirow{7}{*}{\centering\rotatebox[origin=c]{90}{\textbf{with SAM}}}
& PerTaskLP       \\
& PerTaskFT        &68.07  &75.18  &-24.78  &0.17    &0.047  &0.19  &12.43 &12118.71 &70.47 &2935.51 \\
\cmidrule(lr){2-12}
& SeqLoRA        &71.04  &75.76  &-11.46  &0.16  &0.038 &0.14  &10.79  &918.26  &24.13 &762.88\\
\cmidrule(lr){2-12}
 & IncLoRA        &75.03  &78.41  &-3.53  &0.13  &0.009 &0.13  &8.08  &643.08  &15.72 &584.20\\
 % & InfLoRA &  &  &  &  &  &  &  &\\
% & MOS        &  &  &  &  &  &  &  &\\
 % & EASE  &  &  &  &  &  &  & & \\
 \cmidrule(lr){2-12}
 & OLoRA          &74.71  &77.73  &-3.51  &0.12  &0.006  &0.13  &7.93  &656.99 &15.24 &590.27\\
 % & TUNA        &  &  &  &  &  &  &  &\\
\bottomrule
\end{tabular}
}
\end{table}



% In parameter-efficient continual learning (PECL), adaptation is confined to a low-dimensional subspace such as a LoRA adapter, while the pretrained backbone and previously learned adapters remain frozen. This raises a natural question: is encouraging flatness only in this adapter subspace sufficient to improve continual learning?
In PECL, optimization is confined to a low-dimensional LoRA adapter subspace, with the pre-trained backbone and previously learnt adapters remaining frozen. This raises a key question: Is enforcing flatness only within the adapter subspace sufficient to improve continual learning performance? Our empirical studies reveal a consistent pattern: enforcing flatness only within the low-rank LoRA subspace improves continual-learning outcomes, raising final accuracy and reducing forgetting, while also lowering full-model curvature proxies. However,  this observation challenges the common intuition that local (subspace) geometry is less important than global loss-landscape properties.

% This compelling finding, however, invites a deeper theoretical inquiry, as it challenges the intuitive assumption that local subspace geometry is secondary to global loss landscape properties.

 % While LoRA-SAM  explore  the implicit regularization of SAM for scale-invariant problems and find that SAM promotes balancedness. Flat-LoRA argues that solutions that are flat in the LoRA space may remain sharp in the full parameter space and recommends perturbing all parameters. 
 
  While LoRA-SAM\cite{NEURIPS2024_4eb2c0ad} explores the implicit regularisation of SAM for scale-invariant problems and found that SAM promotes balance. Flat-LoRA\cite{li2025flatloralowrankadaptationflat}, on the other hand, argues that solutions that are flat in the LoRA space may remain sharp in the full parameter space, recommending that all parameters be perturbed. However, the proposed 'global perturbation' solution fundamentally conflicts with the core principle of PECL, i.e. isolating pre-trained knowledge and avoiding global parameter fluctuations. This raises the fundamental, unresolved question of the relationship between the flatness of the local LoRA subspace and the robust performance of the global model in the PECL. Existing research has not clarified how the geometry of the local subspace affects the predictive behaviour of the global model, nor has it explained why optimising only local flatness can improve global CL performance. The connection between the local and global remains ambiguous.

  
 
 
 % However, the proposed 'global perturbation' solution is fundamentally misaligned with the core principle of efficient continuous parameter learning (i.e. isolating pre-trained knowledge and avoiding global parameter fluctuations). The misalignment raises a fundamental unresolved question: in the context of evolving task sequences, what is the relationship between the flatness of the local LoRA subspace and the robust performance of the global model? Existing research has not clarified how the geometry of the local subspace affects the predictive behaviour of the global model or explained why optimising only local flatness can improve global CL performance. The connection between the local and global remains ambiguous.



 % While Li et al. (2024a) approach the issue from an optimization standpoint for scale-invariant problems, finding that SAM promotes balancedness. But Flat-LoRA (Li et al., 2025) argues that solutions that are flat in the LoRA space may remain sharp in the full parameter space and therefore recommends perturbing all parameters. However, its proposed "global perturbation" solution is fundamentally misaligned with the core principle of efficient continuous parameter learning (i.e., isolating pre-trained knowledge and avoiding global parameter fluctuations). Behind this misalignment lies a core unresolved question: in the scenario of task sequence evolution, what is the mechanism between the flatness of the local LoRA subspace and the robust performance of the global model? Existing research has neither clarified how the geometry of the local subspace affects the predictive behavior of the global model, nor explained why optimizing only local flatness can improve the global CL performance—the connection between the local and the global remains ambiguous.
 
 
 
 
 
 
 % However, In the CL, perturbing frozen components is theoretically misaligned with the locus of learning and practically risks destabilizing previously acquired knowledge.

 %  their proposed solution of global perturbation presents a fundamental misalignment with the core tenet of parameter-efficient continual learning—to isolate and stabilize existing knowledge.




% Empirically, the answer appears to be yes. When we apply sharpness-aware minimization (SAM) only to the LoRA parameters—without updating the backbone—we observe consistent gains in continual learning metrics.
% In SeqLoRA, LoRA-only SAM increases final average accuracy (FAA) from 66.1 to 71.0 and makes backward transfer (BWT) less negative (–17.28 → –11.46), indicating reduced forgetting. Average accuracy over time (AAA) also improves. At the same time, global curvature proxies computed on the full model decrease: the top Hessian eigenvalue drops (12.68 → 10.79) and the Hessian trace shrinks (1650 → 918). Similar trends hold for IncLoRA and related variants.
% Our empirical evidence indicates that the answer is affirmative. Applying Sharpness-Aware Minimization (SAM) solely to the LoRA parameters increases final accuracy and reduces forgetting. However, this empirical result cannot be fully explained by recent work. While Li et al. (2024a) approach the issue from an optimization standpoint for scale-invariant problems, finding that SAM promotes balancedness. But Flat-LoRA (Li et al., 2025) argues that solutions that are flat in the LoRA space may remain sharp in the full parameter space and therefore recommends perturbing all parameters. However, In the CL, perturbing frozen components is theoretically misaligned with the locus of learning and practically risks destabilizing previously acquired knowledge. Besides, There exists a fundamental contradiction: the learnable parameters are in the LoRA subspace, yet the perturbations are applied globally. So it is unclear why enforcing flatness only within the adapter subspace  improve continual learning performance. They lack a theoretical foundation that connects subspace-restricted flatness to generalization guarantees in a sequential learning setting. The absence of a principled account leaves open whether the observed gains reflect incidental optimization effects or a mechanism that systematically improves stability–plasticity trade-offs.




% This observation contrasts with the view advocated by Flat-LoRA, which argues that solutions that are flat in the LoRA space may remain sharp in the full parameter space and therefore recommends perturbing all parameters. In the CL, however, perturbing frozen components is theoretically misaligned with the locus of learning and practically risks destabilizing previously acquired knowledge.







% This is structurally surprising. We only update a tiny low-rank adapter; the rest of the model is untouched. Yet both continual learning performance and full-model curvature improve. However, such observations alone do not establish mechanism. They could reflect incidental optimization effects. To move beyond correlation, we develop a theory that (i) models continual learning explicitly, (ii) isolates the role of the adapter subspace, and (iii) predicts when and why subspace flatness should translate into better generalization.


\subsection{PAC-Bayes with time-varying hyper-posterior for CL}

% \begin{figure}[thpb]
%   \centering
%   \includegraphics[width=1\linewidth]{images/PAC_large.pdf}
%   \caption{Illustration of the three types of risk in hierarchical PAC-Bayesian CL: future task generalization risk $L(Q)$, expected generalization risk on observed tasks $L(Q_{1:n})$, and empirical risk on observed tasks $\hat{L}(Q_{1:n})$.}
%   \label{fig concept: 3 different definaiton}
% \end{figure}


% To analyze the role of sharpness in PECL, we adopt a hierarchical PAC-Bayesian formulation inspired by transfer learning~\cite{pentinaPACBayesianBoundLifelong2014}, where the function distribution is treated as a random variable. 
% At the start of task \( t \), before observing its data \( S_t \),  a task-level prior \( P_t \) is drawn from a hyper-posterior \( \mathcal{P}_{t-1} \), which summarizes knowledge accumulated from previous tasks. After training on task \( t \),  the learner updates to a posterior \( Q_t = Q_t(S_t, P_t) \) . The hyper-posterior is then updated to \( \mathcal{P}_t \), so that future tasks inherit this information. At inference, parameters \( W_t \sim Q_t \) are sampled for prediction. 
% This construction respects the continual learning constraint that \( P_t \) and \( Q_t \) depend only on previous tasks. The transition \( \mathcal{P}_{t-1} \to  \mathcal{P}_t \) captures forward knowledge accumulation, and
%  \( P_t \to Q_t \) reflects task-specific adaptation. We define the expected risk and the empirical risk on task \( t \) as
% $$
% \mathcal{L}_t := \mathbb{E}_{P_t \sim \mathcal{P}_{t-1}} \mathbb{E}_{W_t \sim Q_t(S_t,P_t)} L_{\mathcal{D}_t}(W_t),
% \quad
% \hat{\mathcal{L}}_t := \mathbb{E}_{P_t \sim \mathcal{P}_{t-1}} \mathbb{E}_{W_t \sim Q_t(S_t,P_t)} \hat{L}_{S_t}(W_t).
% $$

% To analyze the role of sharpness in CL, we adopt a hierarchical PAC-Bayesian framework. Let the $\sigma$-algebra \(\mathcal{F}_t := \sigma(S_1, \dots, S_t)\) encapsulate all historical information up to task \(t\). Before observing \(S_t\), a task-level prior \(P_t\) is drawn from the hyper-posterior \(\mathcal{P}_{t-1}\), i.e., \(P_t \sim \mathcal{P}_{t-1}\), where \(\mathcal{P}_{t-1}\) is \(\mathcal{F}_{t-1}\)-measurable. After training on \(S_t\), the learner obtains the posterior \(Q_t(\cdot \mid P_t, S_t)\), and the hyper-posterior is updated to \(\mathcal{P}_t\), thereby facilitating knowledge transfer to subsequent tasks. The expected risk and empirical risk on task \(t\) are defined as
% $$
% \mathcal{L}_t := \mathbb{E}_{P_t \sim \mathcal{P}_{t-1}} \mathbb{E}_{W_t \sim Q_t} L_{\mathcal{D}_t}(W_t), \quad
% \hat{\mathcal{L}}_t := \mathbb{E}_{P_t \sim \mathcal{P}_{t-1}} \mathbb{E}_{W_t \sim Q_t} \hat{L}_{S_t}(W_t).
% $$


% We assume absolute continuity throughout: \(\mathcal{P}_t \ll \mathcal{P}_{t-1}\) and \(Q_t(\cdot \mid P) \ll P(\cdot)\) for \(\mathcal{P}_t\)-almost every \(P\). In this formulation, the transition \(\mathcal{P}_{t-1} \to \mathcal{P}_t\) captures forward knowledge accumulation, while \(P_t \to Q_t\) reflects task-specific adaptation.


To analyze the role of sharpness in CL, we adopt a hierarchical PAC-Bayesian framework.
For each task \(t \in \{1, \dots, T\}\), let \(S_t = \{z_{tj}\}_{j=1}^{m_t}\) be i.i.d. draws from \(\mathcal{D}_t\), and let \(\ell(w, z) \in [0,1]\) be the loss function. Define the empirical risk and expected risk for a model \(W\) as:
\[
\hat{L}_{S_t}(W) = \frac{1}{m_t} \sum_{j=1}^{m_t} \ell(W, z_{tj}), \quad
L_{\mathcal{D}_t}(W) = \mathbb{E}_{z \sim \mathcal{D}_t} \ell(W, z).
\]
Let \(\mathcal{F}_t = \sigma(S_1, \dots, S_t)\) be the \(\sigma\)-algebra encapsulating all historical information up to task \(t\). Before observing \(S_t\), a task-level prior \(P_t\) is drawn from the hyper-posterior \(\mathcal{P}_{t-1}\), i.e., \(P_t \sim \mathcal{P}_{t-1}\), where \(\mathcal{P}_{t-1}\) is \(\mathcal{F}_{t-1}\)-measurable. After training on \(S_t\), the learner obtains the posterior \(Q_t(\cdot \mid P_t, S_t)\), and the hyper-posterior is updated to \(\mathcal{P}_t\), with absolute continuity:
\[
\mathcal{P}_t \ll \mathcal{P}_{t-1}, \quad
Q_t(\cdot \mid P, S_t) \ll P(\cdot) \quad \text{for } \mathcal{P}_{t-1}\text{-a.e. } P.
\]
For any distributions \(Q\) and \(P\) with \(Q \ll P\), define the KL divergence as \(\mathrm{KL}(Q \| P) = \int \log\left( \frac{dQ}{dP} \right) dQ\). The expected risk and empirical risk on task \(t\) are then defined as:
\[
\mathcal{L}_t = \mathbb{E}_{P_t \sim \mathcal{P}_{t}} \mathbb{E}_{W_t \sim Q_t(\cdot \mid P, S_t)} L_{\mathcal{D}_t}(W_t), \quad
\hat{\mathcal{L}}_t = \mathbb{E}_{P_t \sim \mathcal{P}_{t}} \mathbb{E}_{W_t \sim Q_t(\cdot \mid P, S_t)} \hat{L}_{S_t}(W_t).
\]


In our continual learning PAC-Bayesian framework, rather than analyzing a single deterministic parameter vector, we consider a posterior distribution over parameters.  A common and tractable choice is to model the posterior at task \(t\) as a Gaussian
 $ Q_t = \mathcal{N}(\hat{W}_t, \Sigma_t),$
 where \(\hat{W}_t\) is the trained parameter vector after task \(t\), and \(\Sigma_t\) is  a covariance matrix capturing local uncertainty around \(\hat{W}_t\). Sampling  from \(W \sim Q_t\) is equivalent to 
 $ W = \hat{W}_t + \varepsilon,
 \quad
 \varepsilon \sim \mathcal{N}(0, \Sigma_t).$
Under this posterior view, the expected increase in empirical loss under such perturbations serves as a Bayesian measure of local sensitivity.
From a robustness perspective, the expected sharpness $\mathrm{E\text{-}Sh}$ quantifies the model's sensitivity to small parameter perturbations within the support of the posterior, linking flatness to Bayesian theroy as discussed in prior work\cite{mllenhoff2023sam}.  


The following PAC-Bayesian theorem is established specifically for the continual learning scenario, with explicit consideration given to the dynamic evolution of hyper-posteriors across tasks and its connection to a flatness term. The proof is provided in the Appendix.
\begin{theorem}[PAC-Bayes with time-varying hyper-posterior for CL]\label{thm:pathwise-pacbayes}
For any $\delta\in(0,1]$, with probability at least $1-\delta$ over $S_{1:T}$,
\begin{equation}\label{eq:main-bound}
\begin{aligned} 
& \frac{1}{T} \sum_{t=1}^T \underbrace{\mathbb{E}_{P_t \sim \mathcal{P}_t} \mathbb{E}_{W \sim Q_t(\cdot \mid P_t, S_t)} L_{\mathcal{D}_t}(W)}_{\text{test risk at step } t} 
\le \ \frac{1}{T} \sum_{t=1}^T \underbrace{\mathbb{E}_{P_t \sim \mathcal{P}_t} \mathbb E_{\varepsilon\sim\mathcal N(0,\Sigma_t)}\hat L_{S_t}(\hat W_t+\varepsilon)}_{\text{empirical risk with flatness at step } t} \\ 
& + \sqrt{\frac{ \sum_{t=1}^T \mathbb{E}_{P_t \sim \mathcal{P}_t} \mathrm{KL}\big(Q_t(\cdot \mid P_t, S_t) \,\|\, P_t\big) + \sum_{t=1}^T \mathrm{KL}(\mathcal{P}_t \,\|\, \mathcal{P}_{t-1}) + \log\frac{1}{\delta} }{2 \sum_{t=1}^T (m_t - 1)}}. 
\end{aligned} 
\end{equation}
\end{theorem}


This theorem provides a PAC-Bayesian generalization bound specifically designed for continual learning, which explicitly incorporates a flatness term through Gaussian perturbations in the empirical risk. Unlike standard PAC-Bayes bounds, our approach accounts for the sequential nature of CL through two key components: (1) the dynamic evolution of hyper-posteriors across tasks, captured by the sum of KL divergences between consecutive hyper-posteriors, and (2) the task-adaptive posterior distributions that enable knowledge transfer. The bound reveals that in continual learning, flat minima (as captured by the perturbation term) can help stabilize the learning process and mitigate catastrophic forgetting by controlling the complexity of both task-specific solutions and the evolving meta-knowledge.



% This formulation decouples the empirical risk from the generalization risk, attributing their gap to its dual sources: the randomness in the data and the uncertainty in the function distribution.



% This formulation clearly separates the optimization objective from the generalization target, laying the foundation for analyzing the relationship between flatness in the adapter subspace and generalization guarantees.



% and its sum over tasks as
% $
% \mathcal{R}_{1:T} := \sum_{t=1}^{T} \mathcal{L}_t.
% $
% During training, we only minimize the cumulative empirical risk
% $$
% \widehat{\mathcal{R}}_{1:T} := \sum_{t=1}^{T} \mathbb{E}_{P_t \sim \mathcal{P}_{t-1}} \mathbb{E}_{W_t \sim Q_t(S_t, P_t)} \hat{L}_{S_t}(W_t).
% $$
% This formulation separates what we directly optimize from what we ultimately care about. It highlights the gap between the two risks are caused by the data and  the function distribution.  


% This formulation separates what we directly optimize (empirical loss) from what we ultimately care about (expected risk).

% At the start of task \( t \), before observing its data \( S_t \), a task-level prior \( P_t \) is drawn from \( \mathcal{P}_{t-1} \). The learner then computes the posterior \( Q_t = Q_t(S_t, P_t) \) via training. After training on task \( t \), the hyper-posterior is updated from \( \mathcal{P}_{t-1} \) to \( \mathcal{P}_t \) using \( Q_t \).  During inference for task \( t \), the predictive parameter \( W_t \) is sampled from \( Q_t(S_t, P_t) \)
% The transition \( \mathcal{P}_{t-1} \to  \mathcal{P}_t \) captures the integration of historical knowledge, while \( P_t \to Q_t \) reflects task-specific adaptation. The entire process adheres to continual learning constraints: both \( P_t \) and \( Q_t \) depend only on previous tasks and remain independent of future tasks.
% The expected risk for task \( t \) is defined as:
% $
% \mathcal{L}_t := \mathbb{E}_{P_t \sim \mathcal{P}_{t-1}} \mathbb{E}_{W_t \sim Q_t(S_t,P_t)} L_{\mu_t}(W_t),
% $
% and the cumulative expected risk as
% $
% \mathcal{R}_{1:T} := \sum_{t=1}^{T} \mathcal{L}_t.
% $
% The cumulative empirical risk, minimized during training, is:
% $
% \widehat{\mathcal{R}}_{1:T} := \sum_{t=1}^{T} \mathbb{E}_{P_t \sim \mathcal{P}_{t-1}} \mathbb{E}_{W_t \sim Q_t(S_t, P_t)} \hat{L}_{S_t}(W_t).
% $




% During inference for task \( t \), the predictive parameter \( W_t \) is sampled from \( Q_t(S_t, P_t) \). 
% At the initiation of Task \( t \), prior to accessing the current task data \( S_t \), the task-level prior  \( P_t \sim \mathcal{P}_{t-1} \). Then, the learner obtains the task-level posterior \( Q_t = Q_t(S_t, P_t) \) via training. Upon completion of training on Task \( t \), the hyper-posterior \( \mathcal{P}_{t-1} \) is updated to \( \mathcal{P}_t \) based on \( Q_t \).
% In the inference phase, to make predictions for Task \( t \),  the prediction function parameter \( W_t \sim Q_t(S_t, P_t) \).  The transition \(  \mathcal{P}_{t-1} \to P_t \) embodies the effective accumulation of historical task knowledge.  The transition \( P_t \to Q_t \) focuses on the prior revision driven by Task \( t \)-specific data, accurately capturing the update of prior beliefs via task-specific knowledge to enable the model to adapt flexibly to the unique requirements of each task. The entire framework strictly adheres to continual learning constraints: during the modeling of Task \( t \), both the task-level prior \( P_t \) and posterior \( Q_t \) depend exclusively on information from preceding tasks \( (1, \dots, t-1) \) and are fully independent of future unseen tasks. The expected risk for task $t$ is
% $  \mathcal{L}_t  :=  \mathbb{E}_{P_t \sim \mathbb{P}_{t-1}}   \mathbb{E}_{W_t \sim Q_t(S_t,P_t)}  L_{\mu_t}(W_t)
% $ and the cumulative expected risk
% $
% \mathcal{R}_{1:T}  :=  \sum_{t=1}^{T} {L}_t
% $.  The cumulative empirical risk, which what the learner can actually minimize during training is 
% $
%  \widehat{\mathcal{R}}_{1:T}
%  :=
%  \sum_{t=1}^{T}
%  \mathbb{E}_{P_t \sim \mathbb{P}_{t-1}}
%  \mathbb{E}_{W_t \sim Q_t(S_t, P_t)}
%  \hat{L}_{S_t}(W_t)
%  $




% This process exhibits a two-level dynamic property:
% Cross-task level: The transition \( P_{t-1} \to P_t \) embodies the effective accumulation of historical task knowledge. This accumulation extracts cross-task general patterns and transferable knowledge, establishing more robust and adaptive inductive biases to provide valuable priors for subsequent tasks.
% Intra-task level: The transition \( P_t \to Q_t \) focuses on the prior revision driven by Task \( t \)-specific data, accurately capturing the update of prior beliefs via task-specific knowledge to enable the model to adapt flexibly to the unique requirements of each task.
% The entire framework strictly adheres to continual learning constraints: during the modeling of Task \( t \), both the task-level prior \( P_t \) and posterior \( Q_t \) depend exclusively on information from preceding tasks \( (1, \dots, t-1) \) and are fully independent of future unseen tasks. This ensures temporal rationality and logical consistency in the model’s learning process across the task stream.








% In CL, with the process , we can get a  hyperdistributed sequence $\mathcal{P},\mathcal{Q}_1,\mathcal{Q}_2,\cdots,\mathcal{Q}_T$ which expresses how adaptive inductive preferences across tasks are gradually updated.  Before starting training task t, we sample a task-level prior distribution from the super-distribution Pt-1 at the previous moment Pt~Pt-1. Given the task data St and the prior Pt, the learner produces the posterior distribution after the task training is completed Qt = Qt(St, Pt),  during subsequent inference, we sample a hypothesis from Qt:Wt ~ Qt(St, Pt),And use W to make predictions. After task t is trained, we use the obtained Qt and its performance to update the super distribution and obtain Pt. Therefore, this process encodes both dependencies:Cross-task layer: Pt-1 → Pt accumulates long-term knowledge; The internal layer of the task: Pt → Qt reflects how the current training changes from prior to posterior.

% After the task is completed, the model distribution we have at hand is Qt, which is trained using the prior P and data St.

% We use it to measure generalization performance on the task:

% Lt := EP₁~Pt-1 EW₁~Qt(St,Pr) Lμε (Wt) = EP₁~Pt-1 EW₁~Qt(St,Pr) Ez~με [l(Wt, z)]. 

% At the same time, its empirical counterpart (observable, can be used on the right side of the bound) is given:

% R1:T := T mt ΣΣ t=1 mt j=1 EP₁~Pt-1 EW₁~Qt(St, P₁) [l(Wt, Ztj)].

% This is exactly what we are actually minimizing during training





% LoRA decompose the model parameters as
% $W = W_0 + \Delta W$
% where $W_0 \in \mathbb{R}^d$ is the frozen pretrained backbone, and $\Delta W$ is the trainable low rank moudle.

 % In PECL, at the start of task $t$, both $W_0$ and the previously learned increments $\{\Delta W_s\}_{s=1}^{t-1}$ are frozen.
 % Training on task $t$ is restricted to the current task-specific parameters $\Delta W_t$, which are constrained to lie in a low-dimensional linear subspace $\mathcal{W}_t \subseteq \mathbb{R}^d$.

% \subsection{LoRA Factorization: Adaptation Is Localized to the Adapter Subspac}


 % In PECL, at the start of task $t$, both $W_0$ and the previously learned increments $\{\Delta W_s\}_{s=1}^{t-1}$ are frozen. Training on task $t$ is restricted to the current task-specific parameters $\Delta W_t$.
 %  Notably, the initialization of the LoRA component does not depend on
 % Besides, the initializaiton of the LoRA part are nor rely on the $W_{\text {frozen }}$.






% A direct implication is that the PAC-Bayesian complexity \(\mathrm{KL}(Q_t^{\text{full}} \| P_t^{\text{full}})\) is localized to the low-rank adapter, rendering the bound immune to the high dimensionality of the frozen backbone.



%  Lemma \ref{lem:lora1} formalizes that only the task-specific LoRA update \(\Delta W_t\) carries uncertainty and adapts during task \(t\), whereas the frozen component \(W\) remains deterministic.  As a direct consequence, the per-task PAC-Bayesian complexity term\(\mathrm{KL}(Q_t | P_t)\) is confined to the LoRA subspace.
 

% \begin{lemma}
% \label{lem:lora-factorization}
%  Assume the model parameters at task $t$ decompose as
%  $W = W + \Delta W_t,$
%  where $W \in \mathbb{R}^d$ is a frozen parameter and 
% $\Delta W_t \in \mathcal{W_t} \subseteq \mathbb{R}^d$
% is the task-specific subspace learned only on task $t$.
% We place probability measures on the full parameter space $\mathbb{R}^d$ by factoring the frozen and adaptive parts:
% $
% P_t^{\mathrm{full}}
% := 
% \delta_{W} \otimes P_t^{\Delta},
% \quad
% Q_t^{\mathrm{full}}
% :=
% \delta_{W} \otimes Q_t^{\Delta},
% $
% where $\delta_{W_0}$ is the Dirac distribution at $W$, and 
% $P_t^{\Delta}, Q_t^{\Delta}$ are respectively the prior and posterior
% over the subspace $\mathcal{W}_t$.
% Then
% \[
% \mathrm{KL}\!\left(Q_t^{\mathrm{full}} \,\middle\|\, P_t^{\mathrm{full}}\right)
% =
% \mathrm{KL}\!\left(Q_t^{\Delta} \,\middle\|\, P_t^{\Delta}\right).
% \]
% \end{lemma}




% \begin{lemma}[LoRA factorization of the KL term]
% \label{lem:lora-factorization}
% Assume the model parameters at task $t$ decompose as
% \[
% W \;=\; W_0 \;+\; \Delta W_t,
% \]
% where $W_0 \in \mathbb{R}^d$ is the frozen backbone (pretrained and fixed during continual learning) and 
% $\Delta W_t \in \mathcal{W}_t \subseteq \mathbb{R}^d$
% is the task-specific low-rank LoRA update learned only on task $t$.
% We place probability measures on the full parameter space $\mathbb{R}^d$ by factoring the frozen and adaptive parts:
% \[
% P_t^{\mathrm{full}}
% := 
% \delta_{W_0} \otimes P_t^{\Delta},
% \qquad
% Q_t^{\mathrm{full}}
% :=
% \delta_{W_0} \otimes Q_t^{\Delta},
% \]
% where $\delta_{W_0}$ is the Dirac distribution at $W_0$, and 
% $P_t^{\Delta}, Q_t^{\Delta}$ are respectively the prior and posterior
% over the LoRA subspace $\mathcal{W}_t$.
% Then
% \[
% \mathrm{KL}\!\left(Q_t^{\mathrm{full}} \,\middle\|\, P_t^{\mathrm{full}}\right)
% =
% \mathrm{KL}\!\left(Q_t^{\Delta} \,\middle\|\, P_t^{\Delta}\right).
% \]
% \end{lemma}





 % =
 % \mathbb{E}_{\varepsilon \sim \mathcal{N}(0,\Sigma_t)}
 % \big[
 % \hat{L}_{S_t}(\hat{W}_t+\varepsilon) - \hat{L}_{S_t}(\hat{W}_t)
 % \big],
 % \]


\subsection{Main Result: A Sharpness-Aware PAC-Bayes Bound for PECL}


% At task $t$, the parameter updates factorize as
% \[
% W_t \;=\; W_{t-1} \,+\, \Delta W_t,\qquad \Delta W_t \in \mathcal W_\Delta\subseteq\mathbb R^d,
% \]
% where $W_{t-1}$ is kept fixed during task $t$ and only the low-rank adapter $\Delta W_t$ is trained.


\subsubsection{KL Decomposition in the PECL Parameter Space}

In PECL, the model parameters at task $t$ decompose as
$
W_t=W_{\text {frozen }}+\Delta W_t,
$
where $W_{\text {frozen }}$ includes the pretrained backbone and all previously learned adapters (kept fixed during task $t$ ), and $\Delta W_t $ is the new low-rank LoRA update trained only on task $t$. This reflects the core PECL constraint: past knowledge is preserved via freezing, while adaptation is confined to the newly introduced LoRA block. 
To formalize this structure rigorously, we first define the mathematical context of the LoRA subspace. Let $\mathcal W_\Delta\subseteq\mathbb R^d$ be a linear subspace and let $W+\mathcal W_\Delta$ be its affine translate. Equip both with the trace Borel $\sigma$-algebra induced from $\mathbb R^d$:
$
\mathcal B(\mathcal W_\Delta)\ :=\ {B\cap \mathcal W_\Delta:\ B\in \mathcal B(\mathbb R^d)},\qquad
\mathcal B(W+\mathcal W_\Delta)\ :=\ {B\cap (W+\mathcal W_\Delta):\ B\in \mathcal B(\mathbb R^d)}.
$
These coincide with the Borel $\sigma$-algebras generated by the subspace topologies on $\mathcal W_\Delta$ and $W+\mathcal W_\Delta$, respectively. In particular, both are standard Borel spaces.






% Let $\mathcal W_\Delta\subseteq\mathbb R^d$ be a linear subspace and let $W+\mathcal W_\Delta$ be its affine translate. 

% Equip both with the trace Borel $\sigma$-algebra induced from $\mathbb R^d$:
% \[
% \mathcal B(\mathcal W_\Delta)\ :=\ {B\cap \mathcal W_\Delta:\ B\in \mathcal B(\mathbb R^d)},\qquad
% \mathcal B(W+\mathcal W_\Delta)\ :=\ {B\cap (W+\mathcal W_\Delta):\ B\in \mathcal B(\mathbb R^d)}.
% \]
% These coincide with the Borel $\sigma$-algebras generated by the subspace topologies on $\mathcal W_\Delta$ and $W+\mathcal W_\Delta$, respectively. In particular, both are standard Borel spaces. 


To formalize the probabilistic relationship between the LoRA subspace and the full model space, we define the translation map:
\[
T_W:\ (\mathcal W_\Delta,\mathcal B(\mathcal W_\Delta))\ \to\ (W+\mathcal W_\Delta,\mathcal B(W+\mathcal W_\Delta)),\quad
T_W(\Delta)=W+\Delta,
\]
which is a Borel isomorphism between the subspace and affine space.  The full-space prior and posterior are defined as pushforwards onto the affine subspace by
\[
P^{\mathrm{full}}:=(T_W)_{\text{\#}} P^{\Delta},\qquad
Q^{\mathrm{full}}:=(T_W)_{\text{\#}} Q^{\Delta}.
\]
where $P_t^{\Delta}$ and $Q_t^{\Delta}$ are probability measures on $(\mathcal{W}\Delta, \mathcal{B}(\mathcal{W}\Delta))$.
For any measurable $F$, the pushforward $F_{\text{\#}}\mu$ is defined by $(F_{\#}\mu)(A)=\mu(F^{-1}(A))$. This pushforward construction leads to a fundamental simplification of the KL divergence between full-model distributions, as established in the following lemma.


\begin{lemma}[KL Decomposition with Frozen Backbone]\label{lem:kl-decomposition}
Let $P_t^{\Delta}, Q_t^{\Delta}$ be probability measures on $\mathcal{W}_\Delta$ with $Q_t^{\Delta} \ll P_t^{\Delta}$. Then $Q^{\mathrm{full}}\ll P^{\mathrm{full}}$ and
\[
\mathrm{KL}\left(Q^{\mathrm{full}}\middle|P^{\mathrm{full}}\right)
=\mathrm{KL}\left(Q^{\Delta}\middle|P^{\Delta}\right).
\]
Equivalently, identifying ${W}\times\mathcal W_\Delta$ with $W+\mathcal W_\Delta$ via $\Phi$, one may write
$P^{\mathrm{full}}=\Phi_{\text{\#}}(\delta_W\otimes P^{\Delta})$ and
$Q^{\mathrm{full}}=\Phi_{\text{\#}}(\delta_W\otimes Q^{\Delta})$, and the same equality holds.
\end{lemma}
 

% For the product view, endow ${W}$ with the trivial $\sigma$-algebra ${\varnothing,{W}}$ and ${W}\times\mathcal W_\Delta$ with the product $\sigma$-algebra ${\varnothing,{W}}\otimes\mathcal B(\mathcal W_\Delta)$. The map
% \[
% \Phi:\ {W}\times \mathcal W_\Delta \to W+\mathcal W_\Delta,\qquad \Phi(W,\Delta)=W+\Delta
% \]
% is a measurable isomorphism.  For any measurable $F$, the pushforward $F_{\text{\#}}\mu$ is defined by $(F_{\#}\mu)(A)=\mu(F^{-1}(A))$.



% \begin{lemma}[KL decomposition with a frozen backbone]
% \label{lemma:kl-fullspace-frozen}
% Let $W\in\mathbb R^d$ be fixed and let $P^{\Delta},Q^{\Delta}$ be probability measures on $(\mathcal W_\Delta,\mathcal B(\mathcal W_\Delta))$ with $Q^{\Delta}\ll P^{\Delta}$. Define their pushforwards onto the affine subspace by
% \[
% P^{\mathrm{full}}:=(T_W)_{\text{\#}} P^{\Delta},\qquad
% Q^{\mathrm{full}}:=(T_W)_{\text{\#}} Q^{\Delta}.
% \]
% Then $Q^{\mathrm{full}}\ll P^{\mathrm{full}}$ and
% \[
% \mathrm{KL}\left(Q^{\mathrm{full}}\middle|P^{\mathrm{full}}\right)
% =\mathrm{KL}\left(Q^{\Delta}\middle|P^{\Delta}\right).
% \]
% Equivalently, identifying ${W}\times\mathcal W_\Delta$ with $W+\mathcal W_\Delta$ via $\Phi$, one may write
% $P^{\mathrm{full}}=\Phi_{\text{\#}}(\delta_W\otimes P^{\Delta})$ and
% $Q^{\mathrm{full}}=\Phi_{\text{\#}}(\delta_W\otimes Q^{\Delta})$, and the same equality holds.
% \end{lemma}







% To formalize this structural decomposition probabilistically, we define:
% $$
% P_t^{\text {full }}=\delta_{W_{\text {frozen }}} \otimes P_t^{\Delta}, \quad Q_t^{\text {full }}=\delta_{W_{\text {frozen }}} \otimes Q_t^{\Delta},
% $$
% where $P_t^{\Delta}$ and $Q_t^{\Delta}$ are the prior and posterior over the trainable adapter subspace respectively, and $\delta_{W_{\text {frozen }}}$ is a Dirac distribution centered at the frozen parameters.

% \begin{lemma}
% \label{lem:lora1}
% \begin{equation}
%      \operatorname{KL}\left(Q_t^{\text {full }} \| P_t^{\text {full }}\right)=\operatorname{KL}\left(Q_t^{\Delta} \| P_t^{\Delta}\right) .
% \end{equation}
   
% \end{lemma}

\subsubsection{Subspace-Restricted Sharpness in PECL}

Lemma \ref{lem:kl-decomposition} demonstrates that the PAC-Bayesian complexity term in PECL arises solely from adaptation in the LoRA subspace, as the frozen component contributes no uncertainty. This simplification directly informs the structure of posterior uncertainty: since only $\Delta W_t$ contributes to the divergence, the perturbation covariance $\Sigma_t$ must be confined to $\mathcal{W}_\Delta$. This foundational insight enables deriving a sharpness-aware generalization bound specific to the adapter subspace, where flatness optimization can be focused precisely where uncertainty exists.

% , yielding a subspace-restricted sharpness measure that quantifies sensitivity exclusively to perturbations in the current LoRA parameters

% Lemma \ref{lem:kl-decomposition} demonstrates that the PAC-Bayesian complexity term in PECL arises solely from adaptation in the LoRA subspace, as the frozen component contributes no uncertainty. This directly implies that the PAC-Bayesian complexity term simplifies to \(\mathrm{KL}(Q_t^{\Delta} \| P_t^{\Delta})\), which is foundational for deriving a sharpness-aware generalization bound specific to the adapter subspace.


% Lemma \ref{lem:lora1} demonstrates that the uncertainty in the full model arises solely from the adaptation of the LoRA parameters \(\Delta W_t\), as the frozen component \(W_{\text{frozen}}\) remains deterministic. This directly implies that the PAC-Bayesian complexity term simplifies to \(\mathrm{KL}(Q_t^{\Delta} \| P_t^{\Delta})\), a result that paves the way for deriving a sharpness-aware generalization bound specific to the adapter subspace.


% We now bridge the subspace perspective with the notion of sharpness. In the PAC-Bayesian framework, rather than analyzing a single deterministic parameter vector, we consider a posterior distribution over parameters.  A common and tractable choice is to model the posterior at task \(t\) as a Gaussian
%  $ Q_t = \mathcal{N}(\hat{W}_t, \Sigma_t),$
%  where \(\hat{W}_t\) is the trained parameter vector after task \(t\), and \(\Sigma_t\) is  a covariance matrix capturing local uncertainty around \(\hat{W}_t\). Sampling  from \(W \sim Q_t\) is equivalent to 
%  $ W = \hat{W}_t + \varepsilon,
%  \quad
%  \varepsilon \sim \mathcal{N}(0, \Sigma_t).$
% Under this posterior view, the expected increase in empirical loss under such perturbations serves as a Bayesian measure of local sensitivity.
% From a robustness perspective, the expected sharpness $\mathrm{E\text{-}Sh}$ quantifies the model's sensitivity to small parameter perturbations within the support of the posterior, linking flatness to Bayesian marginal likelihood as discussed in prior work\cite{mllenhoff2023sam}.








% \begin{equation}
% \label{equation: subspace-full}
%     \begin{aligned}
%  \mathrm{E}_{W \sim Q_t}{L}_{S_t}({W}_t)
%  &=
%  \mathrm{E}_{\varepsilon \sim \mathcal{N}(0, \Sigma_t)}{L}_{S_t}(\hat{W}_t + \varepsilon) = 
% {L}_{S_t}(\hat{W}_t)+\mathrm{E\text{-}Sh}_t(\hat{W}_t, \Sigma_t) \\
% & =  \mathrm{E}_{\varepsilon \sim \mathcal{N}(0, \Sigma_t)}{L}_{S_t}(W_{t-1}+ (\Delta \hat{W}_t + \varepsilon)) = 
% {L}_{S_t}(W_{t-1}+\Delta \hat{W}_t)+\mathrm{E\text{-}Sh}_t(\Delta\hat{W}_t, \Sigma_t) 
%     \end{aligned}
% \end{equation}

% The KL decomposition informs the structure of posterior uncertainty in PECL. Since only $\Delta W_t$ contributes to the KL divergence, the perturbation covariance $\Sigma_t$ must be confined to $\mathcal{W}_\Delta$. This yields a subspace-restricted sharpness measure that quantifies sensitivity exclusively to perturbations in the current LoRA parameters.


% Critically, Lemma \ref{lem:lora1}informs the structure of \(\Sigma_t\) in PECL. Since only \(\Delta W_t\) contributes to the KL divergence and is updated at task \(t\), the posterior uncertainty—captured by the perturbation \(\varepsilon\)—must lie entirely within the LoRA subspace. This turns \(\mathrm{E\text{-}Sh}_t\) into a subspace-restricted sharpness measure: it quantifies sensitivity exclusively to perturbations in the current LoRA parameters, not to arbitrary changes across the frozen backbone. As shown in lemma \ref{equation: subspace-full}, this expected sharpness directly corresponds to the average loss increase under Gaussian perturbations confined to the adapter subspace.

% Consider a Gaussian posterior on $(\mathcal W_\Delta,\mathcal B(\mathcal W_\Delta))$,
%  $
%  Q_t^{\Delta}=\mathcal N(\hat{\Delta W}_t,\Sigma_t),
%  $
%  with mean $\hat{\Delta W}_t\in\mathcal W_\Delta$ and covariance operator $\Sigma_t\succeq 0$ acting on $\mathcal W_\Delta$ (viewed in $\mathbb R^d$, this covariance is supported in $\mathcal W_\Delta$).   Here \(\Sigma_t \succeq 0\) denotes a positive semidefinite covariance operator on the adapter subspace \(\mathcal W_\Delta\).
%  56
 
 % Define the full-space posterior on ${W_0}\times\mathcal W_\Delta$ by
 % $
 % Q_t^{\mathrm{full}}:=\delta_{W_0}\otimes Q_t^{\Delta},
 % $

We consider a Gaussian posterior \( Q_t^{\Delta} = \mathcal{N}(\hat{\Delta W}_t, \Sigma_t) \) on the measurable space \( (\mathcal{W}_\Delta, \mathcal{B}(\mathcal{W}_\Delta)) \), where \( \hat{\Delta W}_t \in \mathcal{W}_\Delta \) is the mean and \( \Sigma_t \succeq 0 \) is a positive semidefinite covariance operator acting on the adapter subspace \( \mathcal{W}_\Delta \). When viewed in the ambient space \( \mathbb{R}^d \), this covariance is supported entirely within \( \mathcal{W}_\Delta \). Sampling from \( \Delta W \sim Q_t^{\Delta} \) is equivalent to 
$
\Delta W = \hat{\Delta W}_t + \varepsilon, \quad \varepsilon \sim \mathcal{N}(0, \Sigma_t) \ \text{on } \mathcal{W}_\Delta.
$ This Gaussian posterior formulation enables a natural decomposition of the empirical risk into nominal performance and subspace-restricted sharpness, as formalized below.

\begin{lemma}
\label{lemma: subspace-full}
    The posterior expected empirical loss decomposes as
\[
\mathbb{E}_{W\sim Q_t^{\mathrm{full}}}\big[\hat{L}_{S_t}(W)\big]
= \mathbb{E}_{\varepsilon\sim\mathcal N(0,\Sigma_t)}
    \big[\hat{L}_{S_t}(W_0+\hat{\Delta W}_t+\varepsilon)\big]
= \hat{L}_{S_t}(W_0+\hat{\Delta W}_t)
  + \mathrm{E\text{-}Sh}_t^{\Delta}(\hat{\Delta W}_t,\Sigma_t),
\]
where the (subspace) expected sharpness is
$$
\mathrm{E\text{-}Sh}_t^{\Delta}(\hat{\Delta W}_t,\Sigma_t)
:= \mathbb{E}_{\varepsilon\sim\mathcal N(0,\Sigma_t)}
   \big[\hat{L}_{S_t}(W_0+\hat{\Delta W}_t+\varepsilon)
       - \hat{L}_{S_t}(W_0+\hat{\Delta W}_t)\big],\quad
\varepsilon\sim\mathcal N(0,\Sigma_t)\ \text{on } \mathcal W_\Delta$$
% Here \(\Sigma_t \succeq 0\) denotes a positive semidefinite covariance operator on the adapter subspace \(\mathcal W_\Delta\).
\end{lemma}
 


This analysis reveals a fundamental distinction from the Flat-LoRA perspective\cite{li2025flatloralowrankadaptationflat}. 
% While Flat-LoRA correctly identifies that the flatness of the merged weights in the full parameter space is paramount for generalization, it operates under the assumption that this must be achieved by explicitly perturbing all parameters, 
% including the frozen backbone. 
While Flat-LoRA correctly identifies the importance of flatness in the merged full-model weights for generalization, it assumes this must be achieved by perturbing all parameters—including the frozen backbone.
However, our theoretical framework demonstrates that such full-parameter perturbation is unnecessary and potentially redundant in PECL. Because the operative model at task \(t\) is \(W_{\text{frozen}}+\Delta W_t\) with \(W_{\text{frozen}}\) fixed, the flatness of the merged model along all admissible directions is determined by the flatness of \(\Delta W_t\). Consequently, optimizing for flatness by perturbing only the adapter subspace inherently controls the geometry of the full model—enabled by the additive nature of LoRA updates—while preserving the stability of frozen, knowledge-preserving parameters.



% Because the low-rank update \(\Delta W_t\) is additively merged with the frozen backbone \(W_{t-1}\) to form the operative model \(W_{t-1} + \Delta W_t\), the flatness of the merged model in the full space is **directly and entirely governed by the flatness of \(\Delta W_t\) within its own low-rank subspace**. Therefore, pursuing flatness by perturbing only the LoRA parameters is not a compromise but a **more precise and efficient realization of the same goal**, as it inherently shapes the geometry of the full, inference-time model without introducing instability by perturbing knowledge-preserving frozen parameters.



% As a result, \(\mathrm{E\text{-}Sh}_t\) becomes a subspace-restricted sharpness measure: it reflects sensitivity only to perturbations in the current LoRA parameters, not arbitrary changes across the frozen backbone. As shown in equation \ref{equation: subspace-full}, this expected sharpness directly corresponds to the average loss increase under Gaussian perturbations confined to the adapter subspace.




 % From a robustness viewpoint, $\mathrm{E\text{-}Sh}$ quantifies the sensitivity of the model to small parameter perturbations within the support of the posterior. The link between flatness and Bayes is also discussed in. 
 
 % A smaller value means that typical posterior samples do not significantly increase the loss. 


% Because only \(\Delta W_t\) contributes to the KL term and is allowed to change at task \(t\), the posterior uncertainty (the perturbations \(\varepsilon\)) should live entirely in the LoRA subspace. This turns \(\mathrm{E\text{-}Sh}_t\) into a subspace-restricted sharpness term: it measures loss sensitivity to infinitesimal changes in the current LoRA adapter, not to arbitrary changes in the frozen backbone.



% Equivalently, \(\Sigma_t\) can be taken to have support only on \(\Delta W_t\), with zeros elsewhere. 

% We now integrate a  sharpness term into a PAC-Bayesian generalization bound tailored  for PECL. In the PAC-Bayesian framework, the focus shifts from a single deterministic model to a distribution over models. As previously noted, the predictive parameters \( W_t \sim Q_t \) are sampled from the posterior $Q_t$.  A common assumption is that $Q_t = \mathcal{N}(\hat{W}_t,\sum)$, where $\hat{W}_t$ is  the learned deterministic model and $\sum$ is the covariance matrix. Under this formulation, sampling $W\sim Q_t $
% is equivalent to $W=\hat{W}+\epsilon$ with $\epsilon \sim \mathcal{N}(0,\sum)$. This directly aligns with the definition of distributional sharpness $\mathrm{E\text{-}Sh}$. From a perturbation and robustness perspective, $\mathrm{E\text{-}Sh}$ quantifies the sensitivity of the model to small parameter perturbations within the support of the posterior.


% With Lemma \ref{lem:lora1}, more specifily, the perturbations should add only on the lora part. 



% It is common assumer that $Q_t = \mathcal{N}(\hat{W}_t,\sum)$, where the $\hat{W}_t$ is the optimal model and also the model that we get, $\sum$ is covertrix. So $W\sim Q_t $ equls $W+\epsilon,\epsilon \sim \mathcal{N}(0,\sum)$. That eacatlly the distribution Sharpness $E-Sh$ definition.  From a perturbation/robustness view, \(\mathrm{E\text{-}Sh}_t\) captures how sensitive the model is to small admissible parameter shifts




% \begin{theorem}[Sharpness-aware PAC-Bayes bound for PECL]
% \label{thm:sharpness-pacbayes-cl} 
% For each task $t$, let$P_t^{\Delta}$ and $Q_t^{\Delta}$ are the prior and posterior over the trainable adapter subspace respectively.
% For any $\delta \in (0,1]$, with probability at least $1-\delta$ over the draw of the task datasets 
% $S_1,\ldots,S_T$,

% \begin{equation}
%     \begin{aligned}
%         \frac{1}{T}\sum_{t=1}^T \mathcal{L}_t
% &\;\le\;  
% \frac{1}{T}\sum_{t=1}^T 
% \Big\{
%  L_{S_t}(W_0+{\Delta W}_t)
% + \mathrm{E\text{-}Sh}^{\Delta}_t({\Delta W}_t,\Sigma_t)
% \Big\} \\
% &+
% \sqrt{\frac{
% \displaystyle \sum_{t=1}^T \mathbb{E}_{P_t}\!\big[\mathrm{KL}(Q_t^{\Delta}\|P_t^{\Delta})\big]
% \;+\;
% \displaystyle \sum_{t=1}^T \mathrm{KL}(\mathcal{P}_t\|\mathcal{P}_{t-1})
% \;+\; 
% T\,\log\!\Big(\tfrac{T\,\bar m_A}{\delta}\Big)}
% {2\,T\,H_{m-1}}}\,,
%     \end{aligned}
% \end{equation}
% where $Q_t^{\Delta}:=\mathcal{N}(\hat{\Delta W}_t,\Sigma_t)$ on $\mathcal{W}_t$ and
% \[
% \mathrm{E\text{-}Sh}^{\Delta}_t(\hat{\Delta W}_t,\Sigma_t)
% :=
% \mathbb{E}_{\varepsilon\sim\mathcal{N}(0,\Sigma_t)}\!
% \Big[
% \hat L_{S_t}(W_0+\hat{\Delta W}_t+\varepsilon)-\hat L_{S_t}(W_0+\hat{\Delta W}_t)
% \Big].
% \]
% \end{theorem}


% This theoretical insight, that subspace-restricted sharpness control is both necessary and sufficient, is formally captured in our main generalization bound.

\subsubsection{Sharpness-aware PAC–Bayes bound for PECL}

This insight, that subspace-restricted sharpness control is both necessary and sufficient in PECL, is formally captured in our main generalization bound, which explicitly confines both the sharpness measure and complexity terms to the adapter subspace.
\begin{theorem}[Sharpness-aware PAC--Bayes bound for PECL]\label{thm:pecl-bound}
For any $\delta\in(0,1]$, with probability at least $1-\delta$ over $S_1,\ldots,S_T$,
\begin{equation}\label{eq:pecl-main}
\begin{aligned}
&\frac{1}{T}\sum_{t=1}^T
\mathbb{E}_{P_t\sim\mathcal P_{t-1}}\mathbb{E}_{W\sim Q_t}L_{\mathcal{D}_t}(W)
\ \le\
\ \frac{1}{T}\sum_{t=1}^T
\mathbb{E}_{P_t^\Delta\sim\mathcal P_{t-1}}
\mathbb{E}_{\varepsilon \sim \mathcal{N}(0, \Sigma_t)}[\hat{L}_{S_t}(W_0 + \hat{\Delta W}_t + \varepsilon)
\\
&\quad +\sqrt{\frac{
\displaystyle \sum_{t=1}^T \mathbb{E}_{P_t^\Delta\sim\mathcal P_{t-1}}\!\Big[\mathrm{KL}\!\big(Q_t^\Delta\|P_t^\Delta\big)\Big]
\;+\;\displaystyle \sum_{t=1}^T \mathrm{KL}\!\big(\mathcal P_t\|\mathcal P_{t-1}\big)
\;+\;\log\!\frac{1}{\delta}
}{2\,\displaystyle\sum_{t=1}^T (m_t-1)}}.
\end{aligned}
\end{equation}
\end{theorem}
This bound explicitly decomposes the generalization error into an empirical risk term regularized by the expected sharpness within the LoRA subspace, and a complexity term confined to the adapter's parameter distribution. It thus provides a theoretical justification for our empirical finding: optimizing for flatness only in the low-rank adapter effectively controls the overall generalization performance in continual learning.


% Consider the special case of parameter-efficient continual learning in which
% (i) all backbone parameters and all previously learned adapters are frozen, and
% (ii) a \emph{single} LoRA adapter (with parameters $W_t \in \mathbb{R}^d$) is updated sequentially across tasks $t=1,\ldots,T$.
% In this setting, the task-$t$ prior $P_t$ is given by the posterior from the previous task, i.e.
% $P_t := Q_{t-1}$,
% and the hierarchical hyperposterior does not introduce additional variability.
% Let $Q_t = Q_t(S_t, P_t)$ denote the task-$t$ posterior over the same adapter after training on $S_t$, and let
% $W_t \sim Q_t$ be the resulting predictive parameters.







% Define the per-task expected sharpness
% % $\mathrm{E\text{-}Sh}_t(W_t,\Sigma_t)$
% % as in Theorem~\ref{thm:sharpness-pacbayes-cl}, which upper-bounds the expected increase in empirical loss due to local curvature around $W_t$.

% \begin{corollary}[SeqLoRA Drift--Flatness Bound]
% \label{cor:seqlora-drift-flatness}
%  for any $\delta \in (0,1]$, with probability at least $1-\delta$ over the draw of the task datasets $S_1,\ldots,S_T$,
% \begin{equation}
% \begin{aligned}
% \frac{1}{T}\,\mathcal{R}_{1:T}
% &\le
% \frac{1}{T}
% \sum_{t=1}^{T}
% \mathbb{E}_{W_t \sim Q_t}
% \Big[
%     L_{S_t}(W_t)
%     \;+\;
%     \mathrm{E\text{-}Sh}_t(W_t,\Sigma_t)
% \Big]
% \\[4pt]
% &\quad+
% \frac{1}{T(\bar{m}_T - 1)}
% \sqrt{
% \frac{
% \sum_{t=1}^{T}
% \mathrm{KL}\!\big(
%     Q_t
%     \,\big\|\,
%     Q_{t-1}
% \big)
% \;+\;
% \log\!\big(\bar{m}_T / \delta\big)
% }{
% 2(\bar{m}_T - 1)
% }
% },
% \end{aligned}
% \label{eq:seqlora-bound}
% \end{equation}
% where
% $\bar{m}_T := \frac{T}{\sum_{t=1}^{T} \frac{1}{m_t}}$
% is the harmonic mean of the task sample sizes, and we take $Q_0 := P_0$ to be the (frozen) initialization prior before task $1$.
% \end{corollary}

% In finetune rather than PECL setting, we have 


% In SeqLoRA, the same low-rank adapter $W_t$ simultaneously stores all past tasks and is the \emph{only} part of the model that continues to change.
% The bound \eqref{eq:seqlora-bound} shows that the average continual-learning risk is controlled by two quantities:
% \vspace{-3pt}
% \begin{enumerate}
%     \item \textbf{Drift of the shared adapter}:
%     $\mathrm{KL}(Q_t \,\|\, Q_{t-1})$ measures how far the LoRA adapter moves between task $t{-}1$ and task $t$.
%     Large drift means aggressively overwriting what previous tasks had encoded in the same adapter.
%     \item \textbf{Local curvature / sharpness at the new solution}:
%     $\mathrm{E\text{-}Sh}_t(W_t,\Sigma_t)$ measures how sensitive the loss is to infinitesimal perturbations around $W_t$.
%     High sharpness indicates that the adapter now sits in a high-curvature basin, where even tiny updates can cause large prediction changes and thus forgetting.
% \end{enumerate}
% Equivalently, catastrophic forgetting in SeqLoRA occurs when the single shared adapter is (i) pushed far from its previous state and (ii) pushed into a sharp region.
% Conversely, training procedures such as SAM restricted to the LoRA parameters act to \emph{both} limit high-risk drift directions and bias the adapter toward flatter basins, thereby reducing forgetting.


% Predictions and
\subsection{ Empirical Validation}

% Theoretically, a hierarchical PAC-Bayesian analysis plus subspace factorization (Lemma 4.1) shows that adaptation and complexity are localized, and a sharpness-aware bound (Theorem 4.2) identifies subspace flatness as the operative lever. The data align with the bound’s predictions.

The preceding theoretical analysis yields a clear prediction: in PECL, the stability-plasticity trade-off is governed primarily by the flatness of the adapter subspace. Our empirical results, summarized in Table 1, are consistent with this prediction—enforcing flatness solely in the LoRA parameters consistently improves both final accuracy and reduces forgetting.

This finding reframes the principle of flatness in continual learning: global geometry control is not required; inducing a flat minimum within the low-dimensional adapter suffices. A critical question, however, remains regarding the underlying mechanism: why do flatter adapter solutions exhibit reduced interference? In Section 5, we transition from the weight space to the feature space to analyze how subspace flatness promotes more stable representations and mitigates task interference.



% This reframes CL in PECL: global geometry control is not required; making the adapter solution flat suffices. The remaining question is mechanism: why do flatter adapter solutions forget less? In §5 we move to the feature space and analyze how subspace flatness induces more stable representations and lower task interference.




 

% \begin{lemma}[KL Divergence under Parameter Translation]
% Let \(W \in \mathbb{R}^d\) be a deterministic parameter vector, and let \(\mathcal{W}_\Delta \subseteq \mathbb{R}^d\) be a subspace representing parameter updates (e.g., the low-rank adaptation space). Consider two probability measures, \(P^{\Delta W}\) and \(Q^{\Delta W}\), defined on the measurable space \((\mathcal{W}_\Delta, \mathcal{B}(\mathcal{W}_\Delta))\), with \(Q^{\Delta W}\) continuous with respect to \(P^{\Delta W}\) (denoted \(Q^{\Delta W} \ll P^{\Delta W}\)).
% Let \(P^{\mathrm{full}}\) and \(Q^{\mathrm{full}}\) be the probability measures on the full parameter space \((\mathbb{R}^d, \mathcal{B}(\mathbb{R}^d))\) via the translation map \(T_W\). $P^{\mathrm{full}} \coloneqq (T_W)_{\#} P^{\Delta W}, \quad Q^{\mathrm{full}} \coloneqq (T_W)_{\#} Q^{\Delta W}.$
% Then, \(Q^{\mathrm{full}} \ll P^{\mathrm{full}}\), and the Kullback-Leibler divergence satisfies
% \[
% \mathrm{KL}\left(Q^{\mathrm{full}} \,\middle\|\, P^{\mathrm{full}}\right) = \mathrm{KL}\left(Q^{\Delta W} \,\middle\|\, P^{\Delta W}\right).
% \]
% \end{lemma}



% Define the corresponding probability measures on the full parameter space \((\mathbb{R}^d, \mathcal{B}(\mathbb{R}^d))\) via the translation map \(T_W: \mathcal{W}_\Delta \to \mathbb{R}^d\) given by \(T_W(\Delta) = W + \Delta\). Specifically, let \(P^{\mathrm{full}}\) and \(Q^{\mathrm{full}}\) be the pushforward measures:
% $P^{\mathrm{full}} \coloneqq (T_W)_{\#} P^{\Delta W}, \quad Q^{\mathrm{full}} \coloneqq (T_W)_{\#} Q^{\Delta W}.$
% Then, \(Q^{\mathrm{full}} \ll P^{\mathrm{full}}\), and the Kullback-Leibler divergence satisfies
% \[
% \mathrm{KL}\left(Q^{\mathrm{full}} \,\middle\|\, P^{\mathrm{full}}\right) = \mathrm{KL}\left(Q^{\Delta W} \,\middle\|\, P^{\Delta W}\right).
% \]
% \end{lemma}

% This identity formalizes the intuition that, although we conceptually maintain priors and posteriors over the entire model \(W \in \mathbb{R}^d\), the only source of uncertainty and adaptation at task \(t\) comes from the task-specific LoRA parameters. In PECL, the effective complexity penalty for task \(t\) is therefore fully captured by how much the new task's LoRA increment \(\Delta W\) moves from its predictive prior \(P_t^{\Delta}\) to its learned posterior \(Q_{t+1}^{\Delta}\). 




% Define the LoRA Sharpness loss as
%  $$L^{Sh}_{S_t}(W)
%  :=\max_{\substack{\epsilon\in\mathcal H_\Delta\ |\epsilon|_\Delta\le\rho}}
%  L_{S_t} \big(W_0+(\Delta W+\epsilon),z\big).
%  $$



% At the start of task \(t\), the learner maintains a probability measure $P_t^{\mathrm{full}}$
% over the full model parameters \(W \in \mathbb{R}^d\). After task (t) is trained on its dataset \(S_t\), we obtain an updated belief
% $ Q_{t}^{\mathrm{full}},$which is the posterior distribution over the full model parameters after incorporating task \(t\). \((P_t^{\mathrm{full}}, Q_{t}^{\mathrm{full}})\) is the usual PAC-Bayesian prior–posterior pair for task \(t\), both defined on the complete parameter space of the network. In PECL, we have following cnoclusion. 




% Let the full model parameters decompose as
%  $W(\theta) = W + \Delta W(\theta),$
%  where \(W\) is the frozen parameter and \(\theta\) is trainable parameter \(\Delta W(\theta)\) .
%  We place probability measures over the full parameter space by factoring the frozen part and the learnable part:
%  $
%  P_t^{\mathrm{full}} = \delta_{W}\otimes P^{\Delta W},
%  \quad
%  Q_{t}^{\mathrm{full}} = \delta_{W}\otimes Q^{\Delta W},
%  $
%  where 
%  \(\delta_{W}\) is the Dirac measure at the fixed paramtet \(W\), and \(P^{\Delta W}, Q^{\Delta W}\) are prior/posterior measures on the LoRA parameter space.
%  we have
%  $$
%  \mathrm{KL}\big(Q_{t}^{\mathrm{full}} | P_t^{\mathrm{full}}\big)
%  =
%  \mathrm{KL}\big(Q_t^{\Delta W} |P_t^{\Delta W}\big).
%  $$
% \end{lemma}







% We define task-wise priors and posteriors over the {full} model parameters
% by factoring the deterministic part and the learnable part.



% In the PECL setting, $P_t^{\Delta}$ represents the predictive prior
% for task $t$, namely the distribution over the new task's LoRA
% increment \(\Delta W(\theta_t)\) \emph{before} observing $S_t$.
% The posterior $Q_t^{\Delta}$ is the distribution \emph{after} updating
% $\theta_t$ on $S_t$,
% while keeping $W_0$ and $\{\theta_\tau\}_{\tau < t}$ fixed.
% Consequently, all task-dependent stochasticity is concentrated in the
% current LoRA increment, and the complexity term
% $\mathrm{KL}(Q_t^{\Delta} \,\|\, P_t^{\Delta})$
% quantifies how much the learner has moved in this low-dimensional subspace
% during task $t$.



%Let
% \[
% P_t^{\mathrm{full}}
% =
% \delta_{W_{t-1}}
% \otimes\;
% P_t^{\Delta},
% \quad
% Q_t^{\mathrm{full}}
% =
% \delta_{W_{t-1}}
% \otimes
% Q_t^{\Delta},
% \]
% where $\delta_{(\cdot)}$ denotes the Dirac measure concentrated on the frozen
% backbone plus all previously learned LoRA increments, and
% $P_t^{\Delta}, Q_t^{\Delta}$ are probability measures over the current
% task's LoRA update $\Delta W(\theta_t) \in \mathcal{W}_\Delta$.
% The KL divergence factorizes as
% \[
% \mathrm{KL}\!\left(Q_t^{\mathrm{full}} \,\|\, P_t^{\mathrm{full}}\right)
% \;=\;
% \mathrm{KL}\!\left(Q_t^{\Delta} \,\|\, P_t^{\Delta}\right).
% \]


% In PECL, We parameterize the model as
% $W(\theta) = W_0 + \Delta W(\theta)$
% ,
% where $W_0 \in \mathbb{R}^d$ denotes the frozen pretrained backbone
% and $\theta = (A,B)$ are the LoRA parameters.
% The LoRA update $\Delta W(\theta)$ is a low-rank increment of the form
% $AB$, reshaped to match the dimensions of $W_0$.
% We define the LoRA update subspace
% $\mathcal{W}_\Delta
% :=
% \big\{
% \Delta W(\theta) : \theta = (A,B)
% \big\}
% \subseteq \mathbb{R}^d,$
% that is, the set of all learnable low-rank increments that can be produced by
% LoRA while the backbone $W_0$ remains fixed.
% During PECL, the learner updates only $\theta$, hence the effective predictor
% after task $t$ is $f_{W(\theta_t)}$, while $W_0$ is treated as deterministic
% and does not change.




% We denote a sequence of distributions $(Q_i)_{i=0\dots T}$ on $\Delta \mathcal{H}$, representing the evolution of the model’s learning process.  Here, $Q_i$ is the distribution of the model parameters after training task $i$ which is keep fixed when trained on the next task,only the lora part changed and matters, with $Q_0$ representing the initial parameter distribution and $T$ indicates the total number of tasks.

% For a radius $\rho>0$ define the LoRA-subspace Sharpness loss
% $$
% \ell_{\mathrm{Sh}}^\rho(\theta, z):=\max _{\epsilon \in \mathcal{H}_{\Delta},\|\epsilon\| \leq \rho} \ell\left(W_0+ (\Delta W(\theta)+\epsilon), z\right) .
% $$

% In line with the approach, we denote a sequence of distributions $(Q_i)_{i=0\dots T}$ on $\mathcal{H}$, representing the evolution of the model’s learning process.

% \begin{theorem}
%     For any $\lambda > 0 $ and any online predictive sequence distribution, the following inequality holds with probability $1-\delta$:

%  \begin{equation}
%       \begin{aligned}
%  \sum_{t=1}^T
%  \mathbb{E}_{W_t \sim Q^{\text{full}}_{i}}
%  \Big[\mathbb{E}\big[L_{\mu_t}(W_{t})\mid \mathcal F_{t-1}\big]
%  \Big]
%  &\le
% \sum_{t=1}^T \mathbb{E}_{\Delta W_t \sim  Q^{\Delta W}_{i}}
% \big[\underbrace{\max _{\epsilon \in \mathcal{W}_t,\|\epsilon\| \leq \rho} L_{S_t}\left(W_{t-1}+(\Delta W_t+\epsilon)\right)}_{\text{Sh Loss}}\big]
% \\
%  & +\underbrace{\sum_{t=1}^{T}\frac{\mathrm{KL}\left(Q_{t}^{^{\Delta W}} \| P_{t}^{^{\Delta W}}\right)}{\lambda}}_{\text{task divergence}}
%  +
%  +\underbrace{\frac{\lambda T K^2}{2}}_{\text{model throughput}}
%  +\underbrace{\frac{T\log(1/\delta)}{\lambda}}_{\text{const}}.
%  \end{aligned}
%  \end{equation}
% \end{theorem}





% Model parameters decompose as
%  $
%  W=W_0+\Delta W, \Delta W \in  \mathcal H_\Delta
%  $
%  where \(W_0\) is a frozen pretrained backbone and $  \mathcal H_\Delta$ is the adapter effective-update space. We work in the effective update parameterization $\Delta W$ such as LoRA and assume \(\mathcal H_\Delta\) is bounded, i.e., \(|\Delta W|_F\le M\). 
% Posteriors \(\{Q_{t+1}\}_{t=1}^T\) are arbitrary probability measures on \(\mathcal H_\Delta\) with \(Q_{t+1}\ll P_t\).
% For radius \(\rho>0\), define the Adapter Sharpness loss as
%  $$\ell^{\rho}_{Sh}(\Delta W,z)
%  :=\max_{\substack{\epsilon\in\mathcal H_\Delta\ |\epsilon|_\Delta\le\rho}}
%  \ell \big(W_0+(\Delta W+\epsilon),z\big).$$

% \begin{theorem}
% For any \(\lambda>0\) and \(\delta\in(0,1)\), with probability at least \(1-\delta\) over the data stream,
 
%  \begin{equation}
%       \begin{aligned}
%  \frac{1}{T}\sum_{t=1}^T
%  \mathbb E_{\Delta W\sim {Q_{t+1},S}}
%  \Big[\underbrace{\mathbb E\big
%  [\ell(\Delta W,z_t)\mid \mathcal F_{t-1}\big]
%  }_{\text{conditional generalization risk}}\Big]
%  &\le
%  \frac{1}{T}\sum_{t=1}^T \mathbb E_{\Delta W\sim {Q_{t+1},S}}\big[\ell_{\mathrm{Sh}}^\rho(\Delta W,z_t)\big]\\
%  &\quad +\frac{1}{\lambda T}\sum_{t=1}^T \mathrm{KL}\big(Q_{t+1}|P_t\big)
%  +\frac{\lambda K^2}{2}
%  +\frac{\log(1/\delta)}{\lambda T}.
%  \end{aligned}
%  \end{equation}
% \end{theorem}
 
%  The average conditional generalization risk is controlled by (i) the Adapter Sharpness empirical loss, (ii) a low-dimensional KL term computed only over LoRA parameters, and (iii) standard confidence terms. Hence, promoting flatness where learning actually occurs (the Adapter subspace) suffices to tighten the bound—no perturbation of the frozen backbone is needed.








% \begin{theorem}[PAC–Bayes bound with sharpness for PECL (stated informally)]
% \label{thm:pb-cl-short}
% % with training sets $\{S_t\}_{t=1}^n$, a hyperprior $\mathcal{P}$ over priors $P$, a hyperposterior $\mathcal{Q}$, and task posteriors $Q_t(\cdot\,|\,S_t,P)$.

% Consider $n$ sequential tasks under a hierarchical Bayesian
% CL framework. For any $\delta>0$ the following inequality holds with probability at least $1-\delta$ (over the training samples $\left\{S_1, \ldots, S_n\right\}$ ) for all hyperposterior distributions $\mathcal{Q}$
% \begin{align}
% & L(\mathcal{Q}) \leq \frac{1}{n} \sum_{t=1}^n \mathbb{E}_{P_{t} \sim \mathcal{Q}_{t}} \Big[ \widehat{L}_{S_t}(\mathbf{W}+\Delta\mathbf{W}_t) 
% + {\mathrm{E\text{-}Sh}}_t(\mathbf{W}+\Delta \mathbf{W}_t, \Sigma_Q) \Big] \notag\\
% &  +\frac{1}{n(\bar{m}_n-1)} \sqrt{
% \frac{
% KL(Q_t\|P_t) + KL(\mathcal{Q} \| \mathcal{P}) + \log(\bar{m}_n/\delta)
% }{2(\bar{m}_n-1)}
% }
% \end{align}
% \end{theorem}







% Let $\mathcal{U}$ denote the LoRA subspace and $P_{\mathcal{U}}$ the projector onto $\mathcal{U}$. Consider a sequence-aware prior $P$ over parameters and a posterior $Q$ supported on $\mathcal{U}$. For a task sequence $t=1,\dots,T$ with training sets $S_t$ and sizes $m_t$, with probability at least $1-\delta$,
% \[
% \frac{1}{T}\sum_{t=1}^{T} L_{\mathcal D_t}(Q)
% \;\le\;
% \frac{1}{T}\sum_{t=1}^{T} \hat L_{S_t}(Q)
% \;+\;
% \underbrace{\Phi_{\rho}\big(\mathrm{E\text{-}Sh}_{\mathcal{U}}(\hat{\theta};Q)\big)}_{\text{subspace sharpness control}}
% \;+\;
% \sqrt{\frac{\mathrm{KL}(Q\Vert P)+\ln(1/\delta)}{2\,m_*}},
% \qquad m_*:=\min_t m_t .
% \]
% Here $\mathrm{E\text{-}Sh}_{\mathcal{U}}$ denotes an expected-sharpness term with perturbations supported on $\mathcal{U}$ and radius parameter $\rho$; $\Phi_{\rho}(\cdot)$ collects constants induced by local Lipschitzness and the choice of perturbation distribution. The bound retains the PAC tradeoff between empirical risk and posterior complexity via $\mathrm{KL}(Q\Vert P)$, while introducing an explicit subspace sensitivity term that matches the flatness measures in Section~\ref{sec:flat-minima}. Optimizing LoRA-only SAM reduces $\mathrm{E\text{-}Sh}_{\mathcal{U}}$ by contracting the spectrum of $P_{\mathcal{U}}^{\top} H P_{\mathcal{U}}$, which tightens the bound and yields falsifiable predictions for the experiments that follow. Detailed assumptions and constants are provided in the appendix.








% We compare three adapter organizations that induce distinct \textbf{subspace geometries} under a common backbone: 


% \textbf{Shared-subspace adaptation.}
% \textbf{SeqLoRA} sequentially updates a single fixed-size LoRA adapter $(A,B)$ across tasks while $W_0$ remains frozen and no explicit anti-forgetting mechanism is used. After $T$ tasks, the weights are $W_0 {+} AB$, and all updates lie in one shared low-dimensional subspace.

%updates a single low-rank pair $(A,B)$ across tasks while $W_0$ remains frozen.

% \textbf{Task-wise adaptation with additive span.}
% \textbf{IncLoRA} allocates a dedicated LoRA pair per task over a sequence of tasks, yielding $W_0 {+} \sum_{t=1}^{T} A_t B_t$. The trainable region becomes the additive span of task-specific subspaces, which structurally reduces parameter overlap.
% \textbf{SD-LoRA} further decouples direction and magnitude by freezing previously learned directions $\overline{A_k B_k} := (A_k B_k)/\lVert A_k B_k\rVert_F$ and learning only scalars $\alpha_k$, giving $W_0 {+} \sum_{k=1}^{t} \alpha_k \overline{A_k B_k}$ and constraining adaptation to a fixed low-loss path.
% \textbf{MOS} adopts an EMA-based adapter merging strategy where the current adapter is merged with its predecessors after each update. \textbf{TUNA} integrate task-specific adapters into a universal adapter.
% Beyond the methods above that fuse or span multiple adapters in parameter space, selector methods train one adapter per task and use a router at inference to activate a single or a sparse set of adapters\cite{Yu_2024_CVPR}. In subspace terms, the effective update is confined to one low-dimensional subspace at a time; consequently, performance is driven mainly by the router rather than by properties of a fused adapter geometry.

% \textbf{MOS} adopts an EMA-based adapter merging strategy: after each update, the current adapter $\hat A_b$ is merged with its predecessors to form $A_b=(1-\alpha)\hat A_b+\frac{\alpha}{,b-1,}\sum_{k=1}^{b-1}A_k$, which curbs parameter drift and cross-adapter interference. 

% \textbf{MOS} adopts an EMA-based adapter merging strategy: after each update, the current adapter $\hat A_b$ is merged with its predecessors to form $A_b=(1-\alpha)\hat A_b+\frac{\alpha}{,b-1,}\sum_{k=1}^{b-1}A_k$, which curbs parameter drift and cross-adapter interference.

% These methods fuse or span multiple adapter in parameter space, providing a natural testbed for assessing whether manipulating adapter-subspace flatness yields gains.

% router based method also train adapter for each tasks but in the inference phase using router to select only one to use. In the subspace view, it is equal to  one low-dimensional adapter subspace.



% Methods that route among multiple adapters with a selector, such as MoE adapters for CL~\cite{Yu_2024_CVPR}, primarily shift reliance to the router or classifier; their effective update at each step still lies in one low-dimensional adapter subspace and are therefore not our focus here.


% \textbf{Orthogonality-regularized adaptation.}
% Having contrasted fusion based organizations, we turn to methods that explicitly regularize inter-task geometry. Recent work explores orthogonality strategies to constrain parameter updates and reduce inter-task interference.
% \textbf{OLoRA} preserves the additive form but separates task subspaces via a cross-task orthogonality penalty $\sum_{i<t} \lVert A^{(i)\top} A^{(t)} \rVert_F^2$, encouraging near-orthogonal and more disentangled adapters. Recent state-of-the-art methods, such as \textbf{TUNA}, likewise adopt orthogonality constraints to balance plasticity and stability.


% Building on this analysis, we operationalize “where learning happens” uniformly via LoRA adapters and quantify flatness in the corresponding adapter subspace for all baselines. Accordingly, in Q1 we instantiate SeqLoRA, IncLoRA, SD-LoRA, MOS, OLoRA, and TUNA under a unified ViT-B/16-IN21K setup and test whether actively steering LoRA-subspace flatness improves past-task retention.


% First, the basic and naive setting mics per task finetune but do in the adapter as the baseline. In \textbf{SeqLoRA}, one low-rank pair $(A,B)$ is updated across all tasks while $W_0$ stays frozen, so after $T$ tasks the weights are $W_0{+}AB$ and all updates lie in a single shared subspace. 

% Second, It is a common idea is that evary task shoud have one adapter. And the cooperate of these adapter also cause the trade-off between stability and plasticity, and is also the core we are focus on . \textbf{IncLoRA} assigns a dedicated pair to each task and freezes the past ones, yielding $W_0{+}\sum_{t=1}^{T}A_tB_t$; the trainable space becomes the additive span of task-specific subspaces, structurally reducing overlap. \textbf{SD-LoRA} decouples direction and magnitude by freezing previously learned directions $\overline{A_kB_k}:=A_kB_k/|A_kB_k|_F$ and learning only scalars $\alpha_k$, so $W_0+\sum_{k=1}^{t}\alpha_k \overline{A_kB_k}$ and adaptation proceeds along a fixed low-loss path.  \textbf{MOS} also share the merge of the subspace idea, but share the new trained weight to all previous. Besides this, another idea is not cooperate together but reay on a wise selctor to decide such as MoE-Adapters4CL\cite{Yu_2024_CVPR}. But in this way the model's performance rely on the classfier. And consider the lossland it is equal to the first only one lora. We will not discuss this in this article. 

% Final, due to the stability and plasticity caused by the different adapter, it is believed should add some constrain between different lora. One representastion way is add orthogonal or projection constraint .\textbf{OLoRA} keeps the additive form but pulls these subspaces apart via a cross-task orthogonality penalty $\sum_{i<t}|A^{(i)\top}A^{(t)}|_F^2$, promoting near-orthogonal, disentangled adapters.  Recent sota model  \textbf{TUNA} also apply this techinque.

% Based previous analyse of different PECL method and to investigate  focus on the subspace, we choose the as the representive method to test it.




% This spectrum ranges from a single shared subspace to additive, orthogonal, and direction-frozen variants, allowing us to test whether promoting {flatness within the LoRA subspace} helps independently of how that subspace is organized. Beyond these settings, we also evaluate sota adapter-based methods such as \textbf{MOS}, \textbf{DUCT}, and \textbf{TUNA}.


% We adopt Sharpness-Aware Minimization (SAM) as the primary flatness-aware optimizer and compare it against the standard baseline to isolate the effect of promoting flatness in the adapter subspace. At each training step, with respect to the trainable adapter parameters, SAM forms the worst-case perturbation within an $\ell_2$ ball,
%  $$
%  \epsilon^*(\theta,\rho) := \arg\max_{|\epsilon|_2 \le \rho}, \hat{L}_S(\theta+\epsilon) \approx \rho \cdot \frac{g}{|g|_2}, \quad \text{where } g=\nabla_\theta \hat{L}_S(\theta),
%  $$
%  and then updates parameters using the gradient evaluated at the perturbed point $ \theta+\epsilon^*$, while keeping $ W_0 $ fixed. 
 



% in the LoRA subspace.
% Means $\pm$ std over 5 seeds. Accuracies are higher-is-better ($\uparrow$); flatness is lower-is-better ($\downarrow$).
% All flatness metrics are computed on the end-of-sequence model (@$T$) using identical $\rho$, samples, and power-iteration steps across methods.






% \vspace{-3pt}
% \begin{flushleft}\footnotesize
% \textbf{Definitions (weight-space, LoRA-only).}
% Sh$^{(0)}_{\max}(\rho)$: worst-case zeroth-order sharpness along the steepest direction (key: \texttt{sh0\_max}); 
% Sh$^{(1)}(\rho)$: first-order sharpness $\rho\|\nabla L\|_2$ (key: \texttt{first\_order\_sharpness}; optional: \texttt{grad\_norm}); 
% E-Sh: expected sharpness under random perturbations (keys: \texttt{esh\_mean}±\texttt{esh\_std}); 
% $\lambda_{\max}(H)$: largest Hessian eigenvalue via power-iteration restricted to LoRA params (keys: \texttt{lambda\_max}/\texttt{lambda\_max\_power}); 
% $\Delta L_{\text{1D}}$: loss slice amplitude (\texttt{lossland\_1d\_max} $-$ \texttt{lossland\_1d\_min}). 
% Positive Avg $\Delta$ on FAA/AAA favors SAM; negative Avg $\Delta$ on flatness favors SAM.
% \end{flushleft}
% \end{table}




% \section{Q2: Dose increased  flatness shapes LoRA subspace helps the learned representations? }
\section{From Weight-Space Flatness to Feature-Space Robustness}


% \subsection{The Mechanism Behind Flatness: From Parameters to Features}


% Section 4 showed that enforcing flatness with SAM only on the LoRA parameters improves continual learning performance: final accuracy increases, forgetting is reduced, and curvature measures such as sharpness, Hessian eigenvalues, and Fisher spectrum all decrease. Importantly, these gains appear even though we only optimize a tiny low-rank adapter while keeping the pretrained backbone and past adapters frozen.
% This leaves an open question. 


% Why does shaping the geometry of such a small subspace have system-level impact?

% In particular:
% Is the benefit simply  the linear classifier head is better aligend with feature?
% Or does SAM make the shared representation more stable? To answer this, in the rest of Section 5 we move from the weight space to the feature space. 
% We analyze how a small change in LoRA parameters perturbs intermediate features, and how those feature perturbations change the model’s predicted distribution. Concretely, we will (i) define a local feature-sensitivity matrix that measures how fragile the representation is, (ii) connect that sensitivity back to the LoRA subspace, and (iii) show empirically that SAM suppresses exactly the most destabilizing feature directions.

The PAC-Bayesian bounds established in Section \ref{section PAC} demonstrate that subspace-restricted sharpness control suffices for generalization in PECL. However, a fundamental question remains: why do flatter adapter solutions exhibit reduced catastrophic forgetting? To answer this, we now transition from the weight-space analysis to the feature-space perspective, examining how flatness in the LoRA subspace propagates to feature stability.
 Specifically, we ask: does the benefit stem merely from better-aligned classifier heads, or does flatness-enhancing training fundamentally stabilize the shared representations? Our analysis reveals that the mechanism operates through feature-space robustness, formally captured by the Fisher information matrix.



 % through the lens of Fisher information

% Building on Theorem \ref{thm:pecl-bound}, which confines both sharpness and complexity to the adapter subspace, we analyze how this subspace geometry affects the model's internal representations.


% Section 4 showed that enforcing flatness with SAM only in the LoRA subspace improves continual learning performance (higher final accuracy, less forgetting) and reduces sharpness and curvature metrics. In this section we ask why this works. In particular: how do geometric changes in parameter space produce behaviorally meaningful gains?


% Our working hypothesis is that the dominant effect is not simply better optimization of the current task, nor a better linear classifier, but increased stability of the shared representation. In parameter-efficient continual learning, the backbone and previously learned adapters are frozen, and only the current LoRA block is updated. If those new updates distort the intermediate features too aggressively, they will overwrite the discriminative structure established by earlier tasks, causing forgetting. Thus, a desirable update is one that changes the LoRA weights while leaving past features functionally intact


% To study this mechanism directly, we consider a controlled regime where (i) the input distribution is fixed to foucs on the flatness efrfect, and (ii) only the current LoRA parameters are allowed to move. This isolates the causal path we care about:
% LoRA parameter update → feature perturbation → shift in the output distribution.

% This leads to three concrete questions:

% Do the gains come from stabilizing the linear head, or from making the feature extractor itself more robust across tasks?

% Why should flatness in a very small, low-rank adapter subspace be sufficient to reduce forgetting at the model level?

% By what mechanism can such a small subspace control the stability of shared representations?





% In the previous section, we established two facts. First, enforcing flatness through SAM only on the LoRA adapter parameters already improves continual learning performance: we observe higher final and running accuracy, less negative backward transfer, and lower expected sharpness. Second, and perhaps more surprisingly, this LoRA-only intervention also reduces curvature measures such as the Hessian and empirical Fisher when evaluated on the full model, not just on the adapter. These results support the PAC-Bayesian view that controlling sharpness in the trainable subspace is sufficient to tighten the generalization bound and stabilize learning across tasks.

% However, this still leaves an essential question unanswered: how do these geometric changes in parameter space translate into behavioral gains like “forgetting less” and “generalizing better”? Continual learning is not just about fitting the current task; it is about maintaining usable, discriminative structure in the representations across a sequence of tasks. In PECL, the backbone and past adapters are frozen, and only the current LoRA block is updated. If each new task perturbs internal representations too aggressively, later adapters will erode the discriminative structure established by earlier ones, causing catastrophic forgetting. The central hypothesis of this section is therefore: are the gains from LoRA-only SAM actually driven by making the learned features more stable and more robust to later interference?

% To answer that, we narrow the scope in this section to isolate mechanism. We analyze settings where (i) the feature stable, in aother way, input data distribution is slowly varying, and (ii) the backbone and previous adapters are held fixed.













% Having established in §Q1 that subspace-restricted sharpness control (SAM) improves CI performance under matched budgets, we now ask {why} it works. Our central hypothesis is that the gains arise from smoother and more stable {representations} (features), rather than incidental effects on the classifier head or optimization noise. pretrained backbone augmented with all accumulated LoRA adapters

\subsection{Bridging Flatness to Feature Sensitivity}

% We now refine notation to separate the shared representation from the task-specific classifier.

% Recall from Theorem 2 that the expected sharpness term $\mathrm{E\text{-}Sh}t^\Delta$ quantifies sensitivity to perturbations within the LoRA subspace. 

% To understand how this parameter-space flatness translates to feature stability, we decompose the model into a shared representation extractor and task-specific classifier.
%  Let \( \phi(x;W) \in \mathbb{R}^p \) denote the feature representation extracted by the full model. Here \(W\) is exactly the parameter vector analyzed in Section 4: it includes the frozen backbone and previous adapters, plus the current task’s trainable LoRA block.
%  Let \(W_{\mathrm{cls}} \in \mathbb{R}^{D \times C}\) be a linear classifier mapping features to logits,


 To understand how this parameter-space flatness translates to feature stability, we decompose the model into a shared representation extractor and task-specific classifier. Let $\phi(x;W) \in \mathbb{R}^p$ denote the feature representation, where $W \in \mathbb{R}^d$ includes both the frozen backbone and the trainable LoRA adapter in $\mathcal{W}\Delta$. The linear classifier $W{\text{cls}} \in \mathbb{R}^{p \times C}$ maps features to logits:
 \[
 f(x;W,W_{\mathrm{cls}}) = W_{\mathrm{cls}}^\top \phi(x;W),
 \quad
 p(y\mid x;W,W_{\mathrm{cls}})=\mathrm{softmax}(f(x;W,W_{\mathrm{cls}})).
 \]


This decomposition allows us to precisely trace how LoRA subspace perturbations affect predictive distributions through their impact on intermediate features—the missing link between our PAC-Bayesian theory and empirical performance gains.

% is not the source of interference across tasks

 % LoRA updates act only on the low-rank adapter component inside \(W\) with the task specific head \(W_{\mathrm{cls}}\) .
 % This decomposition allows us to ask a more mechanistic question: how does a small LoRA update \(\Delta W\) perturb the shared features \(\phi(x;W)\), and how sensitive is the resulting predictive distribution?

 
 % The remainder of this section formalizes that sensitivity and links it to the forgetting behavior observed in continual learning.

 
% We measure that change using the KL divergence between the perturbed and unperturbed predictive distributions: 
% $$  \mathrm{KL}
% \left(  p(y \mid \phi(x; W) + \Delta\phi(x); W_{\mathrm{cls}})  
% \Big|
% p(y \mid \phi(x; W); W_{\mathrm{cls}})
% \right).
% $$
% Using a second-order Taylor expansion in \(\Delta\phi(x)\) and ignoring higher-order terms, we obtain the local quadratic approximation \cite{JMLR:NewInsights}

% \cite{JMLR:NewInsights}

% \subsubsection{Local feature matrix and experiment fisher}


% To isolate the impact of flatness on the model, we fix the input \(x\) to ensure the target prediction results have a comparable basis and further investigate how different training strategies designed to enhance flatness (e.g., SAM) alter the model's internal features. 

% Specifically, we introduce a small shift to the feature representation of this fixed input $\Delta\phi(x) \in \mathbb{R}^D$ which is caused by the weight difference rather than different input.


To isolate flatness effects, we fix the input $x$ and examine how parameter perturbations—specifically those confined to the LoRA subspace $\mathcal{W}_\Delta$—alter feature representations and consequently predictive distributions. 
The KL divergence between the perturbed and unperturbed predictive laws gives the local quadratic approximation
\begin{equation}
\label{eq: pro-feature}
 \mathrm{KL}\left(  p(y \mid  \phi + \Delta\phi; W_{\mathrm{cls}})  \Big|  p(y \mid \phi; W_{\mathrm{cls}})  \right)
=
\Delta\phi(x)^\top
 E_\phi(x)
 \Delta\phi(x).
\end{equation}
where the {local feature matrix}
 $
 E_\phi(x)
 \triangleq
 \mathbb{E}_{y \sim p(\cdot \mid x; W, W_{\mathrm{cls}})}
 \Big[
 \nabla_\phi \log p(y \mid x; W, W_{\mathrm{cls}})
 \nabla_\phi \log p(y \mid x; W, W_{\mathrm{cls}})
 ^\top
 \Big]
 $
measures predictive sensitivity to feature perturbations\cite{EFC2024}.
% \in \mathbb{R}^{D \times D}
% The matrix
%  $
%  E_\phi(x)
%  \triangleq
%  \mathbb{E}_{y \sim p(\cdot \mid x; W, W_{\mathrm{cls}})}
%  \Big[
%  \nabla_\phi \log p(y \mid x; W, W_{\mathrm{cls}})
%  \nabla_\phi \log p(y \mid x; W, W_{\mathrm{cls}})
%  ^\top
%  \Big]
%  $
%  is  the \textbf{local feature  matrix}, which measures how sensitive the model’s predictive distribution is to changes in the representation \(\phi(x;W)\) \cite{EFC2024}.
% For a linear classifier followed by a softmax,
% $  E_\phi(x)
% =
% W_{\mathrm{cls}}
%  \Big(
%  \mathrm{Diag}(p) - p p^\top
%  \Big)
%  W_{\mathrm{cls}}^\top 
%  $, with $p=p(\cdot\mid x;W,W_{\mathrm{cls}})$\cite{EFC2024}.
% The empirical feature matrix $E_t = \mathbb{E}_{x\sim\mathcal{D}_t}[E_\phi(x)]$ summarizes this sensitivity across the task distribution.
% We infer that training procedures reducing spectral concentration in $E_\phi$ achieve dual benefits: they improve linear classifier fitting while reducing anisotropy in predictive sensitivity, thereby stabilizing shared features against directional distortion.


 
 % Taking the expected value of $E_\phi(x)$ over the entire dataset at task $t$, we get the Empirical Feature Matrix (EFM) $E_t=\mathbb{E}_{x\sim\mathcal{D}_t}[E_\phi(x)]$. 
% We can infer that if a training procedure can consistently reduce the spectral concentration in $E_\phi$ and flattens the representation geometry, it not only fits a better linear classifier, but also reduces anisotropy in predictive sensitivity to representation perturbations and stabilizes shared features against catastrophic single-direction distortion.




% $\lambda_{max}(E_\phi)$and increase effective rank, then it is not merely fitting a better linear head. It is actively making the shared representation more stable and harder to catastrophically distort along any single direction.
 
 % Taking the expected value of $E_\phi(x)$ over the entire dataset at task $t$, we get he Empirical Feature Matrix (EFM) $E_t=$


% \subsubsection{Subspace weight pertubation}

 

% .


% This observation also clarifies an apparently surprising result from Section 4. By using SAM to suppress the most harmful update directions in that subspace, we directly stabilize the representation.

% For a sufficiently small parameter perturbation $\Delta W$, the representation shift at a fixed input $x$ satisfies
% $\Delta\phi(x)\approx J_\phi(x)\,\Delta W$, where $J_\phi(x)=\partial\phi(x;W)/\partial W$.

% gives
%  $
%  \nabla_W \log p(y \mid x; W, W_{\mathrm{cls}})
%  = J_\phi(x)^\top\nabla_\phi \log p(y \mid x; W, W_{\mathrm{cls}}),
%  $
%  hence
 
Let $J_\phi(x)=\partial \phi(x;W)/\partial W$. Since the representation shift is caused by parameter perturbation rather than linear classifier $W_{\mathrm{cls}}$, the chain rule gives
 $$
 \nabla_W \log p(y \mid x; W, W_{\mathrm{cls}})=J_\phi(x)^\top  \nabla_\phi \log p(y \mid x; W, W_{\mathrm{cls}})
 $$
For small parameter perturbations $\Delta W \in \mathcal{W}\Delta$, the induced representation shift is $\Delta\phi(x) \approx J\phi(x)\Delta W$, yielding :
% \[
% \mathrm{KL}\Big(p(\cdot\mid x; W+\Delta W, W_{\mathrm{cls}})  \big|\ p(\cdot\mid x; W, W_{\mathrm{cls}})\Big)  =  \Delta W^\top J_\phi(x)^\top E_\phi(x) J_\phi(x)\Delta W. 
% \]
% the \textbf{local parameter-space Fisher information matrix}
% For a fixed input $x$, a small parameter perturbation $\Delta W$ induces a representation shift $\Delta\phi(x)\approx J_\phi(x)\Delta W$. 
% The local effect of parameter perturbations $\Delta W$ on predictive distribution is quantified by
\[
\mathrm{KL}\Big(p(\cdot\mid x; W+\Delta W, W_{\mathrm{cls}})  \big|\ p(\cdot\mid x; W, W_{\mathrm{cls}})\Big)  =  \Delta W^\top J_\phi(x)^\top E_\phi(x) J_\phi(x)\Delta W. 
\]
where $J_\phi(x)^\top E_\phi(x) J_\phi(x)$ is the \textbf{local parameter-space Fisher information matrix at $x$}.
Averaging over inputs with the $S$,  we obtain the empirical Fisher information matrix 
$
F= \frac{1}{|S|}\sum_{x\in S} J_\phi(x)^\top E_\phi(x)  J_\phi(x)
$
which bridges feature-space robustness to parameter-space flatness, exactly corresponding to the curvature surrogate that governs the expected sharpness in our PAC-Bayesian bounds.


% and serves as the generalized Gauss–Newton matrix, providing a PSD curvature surrogate that approximates the Hessian spectrum near minimizers.
% captures the expected sensitivity of predictive distributions to parameter perturbations across the dataset, bridging feature-space robustness to parameter-space flatness.

% coincides with the generalized Gauss–Newton matrix under cross-entropy with softmax, serving as a PSD curvature surrogate that approximates the Hessian spectrum near minimizers.


%  Under cross-entropy with a softmax output, the experiment fisher coincides with the generalized Gauss–Newton matrix and serves as a PSD curvature surrogate that closely tracks the Hessian spectrum under standard conditions (e.g., near minimizers with well-specified labels).
% the experiment fisher

% In summary, this analysis establishes that flatness-enhancing training induces more isotropic predictive sensitivity in feature space, which stabilizes shared representations against distortion. The pullback of feature sensitivity through the Jacobian $J_\phi(x)$ yields the \textbf{local parameter-space Fisher information matrix}, which serves as a principled curvature surrogate approximating the Hessian spectrum. Critically, the empirical Fisher matrix $F$ provides the fundamental link between feature-space robustness and parameter-space flatness, demonstrating how reducing spectral concentration in $E_\phi(x)$ translates to more stable optimization dynamics and improved resistance to catastrophic forgetting through flattened representation geometry.




% A second-order expansion of the conditional KL divergence around a stationary point gives
% \[
% \mathrm{KL}\Big(p(\cdot\mid x; W+\Delta W, W_{\mathrm{cls}})  \big|\ p(\cdot\mid x; W, W_{\mathrm{cls}})\Big)  \approx  \tfrac12\Delta W^\top G_W(x)\Delta W, 
% \]
% where 
% $
% G_W(x)
% =\mathbb{E}_{y \sim p(\cdot \mid x; W, W_{\mathrm{cls}})}
%  \Big[
%  \nabla_W \log p(y \mid x; W, W_{\mathrm{cls}})
%  \nabla_W \log p(y \mid x; W, W_{\mathrm{cls}})^\top
%  \Big]
% $.
% Since the representation shift is caused by parameter perturbation rather than linear classifier $W_{\mathrm{cls}}$, the chain rule gives
%  $
%  G_W(x)=J_\phi(x)^\top E_\phi(x) J_\phi(x).
%  $
 % Under cross-entropy with a softmax output, $G_W$ coincides with the generalized Gauss–Newton matrix and serves as a PSD curvature surrogate that closely tracks the Hessian spectrum under standard conditions (e.g., near minimizers with well-specified labels).

 % Averaging over inputs with the empirical distribution $x\sim\mathcal{D}_t$ gives
 % $$
 % \mathbb{E}_{x\sim \mathcal{D}_t}\Big[
 % \mathrm{KL}\big(p(\cdot\mid x; W+\Delta W, W_{\mathrm{cls}})
 % \big|
 % p(\cdot\mid x; W, W_{\mathrm{cls}})\big)
 % \Big]
 % \approx
 % \frac12\Delta W^\top F_W\Delta W,
 % $$
 % where
 % $
 % F_W:=\mathbb{E}_{x\sim \mathcal{D}_t}\big[G_W(x)\big]
 % =\mathbb{E}_{x\sim \mathcal{D}_t}\big[J_\phi(x)^\top E_\phi(x) J_\phi(x)\big].
 % $


% Averaging over inputs with the empirical distribution $x\sim\mathcal{D}_t$ yields the curvature surrogate
% \[
% \mathrm{KL}\!\big(p_\theta(\cdot\mid X)\,\big\|\,p_{\theta+\delta}(\cdot\mid X)\big)
% \;\approx\;\tfrac12\,\Delta W^\top G_W\,\Delta W,
% \]
% where 
% $G_W:=\mathbb{E}_{x\sim\mathcal{D}_t}\!\big[F_W(x)\big]
% =\mathbb{E}_{x\sim\mathcal{D}_t}\!\big[J_\phi(x)^\top E_\phi(x) J_\phi(x)\big].$
% Under cross-entropy with softmax, $G_W$ coincides with the generalized Gauss–Newton matrix and closely tracks the Hessian spectrum near minimizers.





% Let
% $
% J_\phi(x)=\frac{\partial \phi(x ; W)}{\partial W} \in \mathbb{R}^{D \times P}
% $
% be the Jacobian of the shared representation $\phi(x ; W)$ with respect to the full parameter vector $W$, where $P=\operatorname{dim}(W)$. 
% For a sufficiently small parameter perturbation $\Delta W \in \mathbb{R}^P$, we have the first-order approximation:
% \begin{equation}
% \label{eq: feature-weight}
%     \Delta \phi(x) \approx J_\phi(x) \Delta W.
% \end{equation}
% By substituting \eqref{eq: feature-weight} in \eqref{eq: pro-feature} and averaging over inputs $x$, we obtain a quadratic form in $\Delta W$
% \begin{equation}
%     \mathrm{KL}\left(p\left(\cdot \mid W+\Delta W, W_{\text {cls }}\right) \| p\left(\cdot \mid W, W_{\text {cls }}\right)\right) 
%     =
%     \Delta W^T  \underbrace{\mathbb{E}_x\left[J_\phi(x)^{\top} E_\phi(x) J_\phi(x)\right]}_{G_W} \Delta W.
% \end{equation}
% For compactness, define the following matrix over the full parameter space:
% $$
% G_W:=\mathbb{E}_x\left[J_\phi(x)^{\top} E_\phi(x) J_\phi(x)\right] \in \mathbb{R}^{P \times P} .
% $$
% Under a softmax classifier with cross-entropy loss, $G_W$ coincides with the empirical Fisher information with respect to the parameters $W$. Equivalently, it matches the Generalized Gauss-Newton matrix (GGN), and thus upper-bounds the leading eigenvalues of the true Hessian in well-behaved regions near a minimum .
% Intuitively, $G_W$ tells us how sharply the output distribution reacts to infinitesimal parameter changes in each direction of $W$.

\subsection{Expected Sharpness and PAC-Bayesian Coupling in LoRA Factor Space}
% In PECL, $\Delta \phi(x)$ is induced by actually updating the LoRA parameters inside $W$. We now formalize how a small parameter update in the LoRA subspace propagates into a representation shift and link with the final probarbality prediction. 
% we now formalize how parameter updates in the LoRA subspace $\mathcal{W}_\Delta$ propagate to representation shifts and connect with the PAC-Bayesian framework from Section 3.

Building on the Gaussian posterior formulation $Q_t^{\Delta} = \mathcal{N}(\hat{\Delta W}_t, \Sigma_t)$ on $\mathcal{W}_\Delta$ from Section 3.3.2, we now parameterize the adapter subspace through its factorized LoRA structure. This refined coordinate system allows us to explicitly trace how parameter perturbations propagate through the model's geometric hierarchy.


LoRA imposes a factorized, layer-wise update structure: for each weight matrix $W_\ell \in \mathbb{R}^{d_o \times d_i}$, the update takes the form $\Delta W_\ell = B_\ell A_\ell$ with $B_\ell \in \mathbb{R}^{d_o \times r}$ and $A_\ell \in \mathbb{R}^{r \times d_i}$. Around current estimates $(\hat{B}_\ell, \hat{A}_\ell)$, the first-order variation satisfies:
$
\delta(\Delta W_\ell) = \hat{B}_\ell \delta A_\ell + \delta B_\ell \hat{A}_\ell.
$
In vectorized form:
$
\mathrm{vec}(\delta(\Delta W_\ell)) = (I_{d_i} \otimes \hat{B}_\ell)\mathrm{vec}(\delta A_\ell) + (\hat{A}_\ell^\top \otimes I_{d_o})\mathrm{vec}(\delta B_\ell).
$
Stacking across layers yields a global linear embedding from LoRA factor coordinates $\theta = [\mathrm{vec}(\delta A_1); \mathrm{vec}(\delta B_1); \dots]$ to the adapter subspace $\mathcal{W}_\Delta$:
\[
\mathrm{vec}(\Delta W) \approx J_{\mathrm{LoRA}} \theta,
\]
where $J_{\mathrm{LoRA}}$ is block-structured and full column rank, establishing an isomorphism between the factor space and admissible directions in $\mathcal{W}_\Delta$.











% LoRA imposes a factorized, layer-wise update $\Delta W_\ell=B_\ell A_\ell$ with $B_\ell\in\mathbb{R}^{d_o\times r}$ and $A_\ell\in\mathbb{R}^{r\times d_i}$. 
% Around the current point $(\hat B_\ell,\hat A_\ell)$ the first-order variation is
% $\delta(\Delta W_\ell)=\hat B_\ell\,\delta A_\ell+\delta B_\ell\,\hat A_\ell$.
% Vectorizing and stacking all layers yields a global linear embedding from LoRA factor coordinates 
% $\theta=[\mathrm{vec}(\delta A_1);\mathrm{vec}(\delta B_1);\dots]$ to parameter perturbations
% \[
% \mathrm{vec}(\Delta W)\;\approx\;J_{\mathrm{LoRA}}\,\theta,
% \qquad
% J_{\mathrm{LoRA}}\text{ block-structured and full column rank on admissible directions}.
% \]
Substituting this chart into the parameter Fisher gives the \textbf{projected Fisher in LoRA factor space}
\begin{equation}
\label{eq:lora-fisher}
\mathrm{KL}\!\left(p(\cdot\mid W+\mathrm{unvec}(J_{\mathrm{LoRA}}\theta))\,\big\|\,p(\cdot\mid W)\right)
\;=\; \theta^\top G_{\mathrm{fact}}\,\theta,
\qquad
G_{\mathrm{fact}}:=J_{\mathrm{LoRA}}^\top\,G_W\,J_{\mathrm{LoRA}}.
\end{equation}
Geometrically, the feature-space sensitivity ellipsoids governed by $E_\phi(x)$ are first pulled back to parameter space through $J_\phi(x)$ to form $G_W$, and then restricted to the admissible LoRA directions through $J_{\mathrm{LoRA}}$. The spectrum of $G_{\mathrm{fact}}$ therefore captures the curvature that the optimizer effectively experiences in the LoRA subspace: $\lambda_{\max}(G_{\mathrm{fact}})$ identifies the most sensitive direction for prediction stability and $\mathrm{tr}(G_{\mathrm{fact}})$ measures the aggregate sensitivity.

We now couple expected sharpness with the PAC-Bayesian complexity in the same coordinate chart. 
Consider a Gaussian posterior over LoRA factors $Q^\theta=\mathcal{N}(\mu_\theta,\Sigma_\theta)$ and the induced full-parameter posterior 
\[
W \stackrel{d}{=} W_0 + \widehat{\Delta W} + \mathrm{unvec}\!\big(J_{\mathrm{LoRA}}(\mu_\theta+\varepsilon)\big),
\qquad \varepsilon\sim\mathcal{N}(0,\Sigma_\theta).
\]
% Assume we are in a neighborhood where the empirical risk admits a second-order approximation with Hessian $H_t$ and that $H_t$ can be locally controlled by $G_W$. 
% Then the expected sharpness admits the trace form in LoRA coordinates
% \begin{equation}
% \label{eq:esh-trace}
% \mathbb{E}_{W\sim Q^\theta}\!\big[\hat L_{S_t}(W)-\hat L_{S_t}(W_0+\widehat{\Delta W})\big]
% \;\approx\;
% \tfrac12\,\mathrm{tr}\!\big(H_\Delta^{\mathrm{fact}}\Sigma_\theta\big)
% \;\approx\;
% \tfrac12\,\mathrm{tr}\!\big(G_{\mathrm{fact}}\Sigma_\theta\big),
% \qquad
% H_\Delta^{\mathrm{fact}}:=J_{\mathrm{LoRA}}^\top H_t J_{\mathrm{LoRA}}.
% \end{equation}
Assuming the empirical risk admits a local second-order approximation with Hessian $H_t$ that can be bounded by $G_W$, the expected sharpness reduces to a trace form in LoRA coordinates:
\begin{equation}
\label{eq:esh-trace}
\mathbb{E}{W \sim Q^\theta}\left[\hat{L}{S_t}(W) - \hat{L}{S_t}(W_0 + \widehat{\Delta W})\right]
\approx \tfrac12 \mathrm{tr}\left(H\Delta^{\mathrm{fact}} \Sigma_\theta\right)
\approx \tfrac12 \mathrm{tr}\left(G_{\mathrm{fact}} \Sigma_\theta\right),
\quad
H_\Delta^{\mathrm{fact}} := J_{\mathrm{LoRA}}^\top H_t J_{\mathrm{LoRA}}.
\end{equation}


Hence both the curvature term and the posterior uncertainty live in the same low-dimensional factor space. 
Moreover, the KL complexity is invariant under the pushforward to full parameters in our frozen-backbone setting:
\[
\mathrm{KL}(Q^\theta\|P^\theta) \;=\; \mathrm{KL}(Q^{\Delta}\|P^{\Delta})
\;=\; \mathrm{KL}(Q^{\mathrm{full}}\|P^{\mathrm{full}}),
\]
so the PAC-Bayesian objective closes in the LoRA coordinate chart as
\[
\underbrace{\hat L_{S_t}(W_0+\widehat{\Delta W})}_{\text{empirical risk}}
\;+\;
\underbrace{\tfrac12\,\mathrm{tr}(G_{\mathrm{fact}}\Sigma_\theta)}_{\text{expected sharpness}}
\;+\;
\underbrace{\mathrm{KL}(Q^\theta\|P^\theta)}_{\text{complexity}}
\quad\text{up to second-order remainder terms}.
\]
This shows that optimizing flatness \emph{only} inside the LoRA subspace is sufficient to control the effective curvature seen by training, which in turn stabilizes the shared features through the chain 
$E_\phi \;\to\; G_W \;\to\; G_{\mathrm{fact}}$.

% Finally, this geometry explains the effect of SAM restricted to LoRA. 
% The inner maximization 
% $\max_{\|\theta\|\le \rho}\theta^\top G_{\mathrm{fact}}\theta$
% aligns with the leading eigenvector of $G_{\mathrm{fact}}$. 
% The outer update suppresses $\lambda_{\max}(G_{\mathrm{fact}})$ and improves the condition number, which reduces the worst-case sensitivity and spreads curvature more evenly across safer directions. 
% Through the pullback and projection, these spectral changes translate into flatter and more isotropic feature geometry, smaller representation drift across tasks, and lower forgetting.


% The LoRA constraint imposes a factorized, layer-wise structure: for each weight matrix $W_\ell\in\mathbb{R}^{d_o\times d_i}$, updates take the form $\Delta W_\ell=B_\ell A_\ell$ with $B_\ell\in\mathbb{R}^{d_o\times r}$ and $A_\ell\in\mathbb{R}^{r\times d_i}$. 
% Around current estimates $(\hat B_\ell,\hat A_\ell)$, the first-order variation satisfies $\delta(\Delta W_\ell)=\hat B_\ell\delta A_\ell+\delta B_\ell\hat A_\ell$. In vectorized form:
% \[
% \mathrm{vec}(\delta(\Delta W_\ell))=(I_{d_i}\otimes \hat B_\ell)\mathrm{vec}(\delta A_\ell)+(\hat A_\ell^\top\otimes I_{d_o})\mathrm{vec}(\delta B_\ell).
% \]
% Stacking across layers yields a global linear embedding from LoRA factor coordinates $\theta=[\mathrm{vec}(\delta A_1);\mathrm{vec}(\delta B_1);\dots]$ to full parameter perturbations:
% \[
% \mathrm{vec}(\Delta W)\approx J_{\mathrm{LoRA}}\theta,
% \]
% where $J_{\mathrm{LoRA}}$ is block-structured and full column rank on admissible directions.

% Substituting this coordinate mapping into the local parameter-space Fisher information matrix gives the \textbf{projected Fisher in LoRA factor space}:
% \begin{equation}
% \label{eq:lora-fisher}
% \mathrm{KL}\left(p(\cdot\mid W+\mathrm{unvec}(J_{\mathrm{LoRA}}\theta))\,\|\,p(\cdot\mid W)\right)
% = \theta^\top G_{\mathrm{fact}}\theta, \quad
% G_{\mathrm{fact}} := J_{\mathrm{LoRA}}^\top\left(\frac{1}{|S|}\sum_{x\in S} J_\phi(x)^\top E_\phi(x) J_\phi(x)\right)J_{\mathrm{LoRA}}.
% \end{equation}

% Geometrically, the feature-space sensitivity ellipsoids defined by $E_\phi(x)$ are pulled back to parameter space via $J_\phi(x)$ to form the local parameter-space Fisher, then restricted to admissible LoRA directions through $J_{\mathrm{LoRA}}$. The spectrum of $G_{\mathrm{fact}}$ thus captures the effective curvature experienced within the LoRA subspace: $\lambda_{\max}(G_{\mathrm{fact}})$ identifies the most sensitive prediction direction, while $\mathrm{tr}(G_{\mathrm{fact}})$ measures aggregate sensitivity.





% The LoRA constraint is now imposed in a factorized, layer-wise form. For a weight matrix (W_\ell\in\mathbb{R}^{d_o\times d_i}) in layer (\ell), LoRA updates are written as (\Delta W_\ell=B_\ell A_\ell), where (B_\ell\in\mathbb{R}^{d_o\times r}) and (A_\ell\in\mathbb{R}^{r\times d_i}) with small rank (r). Around the current point ((\hat B_\ell,\hat A_\ell)), the first-order variation of the update is (\delta(\Delta W_\ell)=\hat B_\ell,\delta A_\ell+\delta B_\ell,\hat A_\ell). Using vectorization and Kronecker products, (\mathrm{vec}(\delta(\Delta W_\ell))=(I_{d_i}\otimes \hat B_\ell)\mathrm{vec}(\delta A_\ell)+(\hat A_\ell^\top\otimes I_{d_o})\mathrm{vec}(\delta B_\ell)). Stacking all layers produces a global linear embedding from the LoRA factor coordinates (\theta=[\mathrm{vec}(\delta A_1);\mathrm{vec}(\delta B_1);\dots]) to the full parameter perturbation (\Delta W), namely (\mathrm{vec}(\Delta W)\approx J_{\text{LoRA}}\theta) with an embedding matrix (J_{\text{LoRA}}) that is block-structured and full column rank in the admissible directions. Substituting this coordinate chart into the quadratic form defined by (G_W) gives a projected Fisher in the LoRA factor space,
%  [
%  \mathrm{KL}!\left(p(\cdot\mid W+\mathrm{unvec}(J_{\text{LoRA}}\theta)),\big|,p(\cdot\mid W)\right)=\theta^\top G_{\text{fact}},\theta,\qquad
%  G_{\text{fact}}:=J_{\text{LoRA}}^\top G_W J_{\text{LoRA}}.
%  ]
%  Geometrically, the feature-space ellipsoids defined by (E_\phi(x)) are pulled back to parameter space by (J_\phi(x)) to form (G_W), then restricted to the admissible LoRA directions by the projection (J_{\text{LoRA}}). The spectrum of (G_{\text{fact}}) therefore captures the curvature that the optimizer effectively experiences when it is only allowed to move inside the LoRA adapter: (\lambda_{\max}(G_{\text{fact}})) identifies the single most harmful direction for prediction stability and (\mathrm{tr}(G_{\text{fact}})) measures the aggregate sensitivity. When a fixed orthonormal basis (U) is used to describe a linear LoRA subspace, the same construction reduces to (G_{\text{LoRA}}=U^\top G_W U), which is equivalent up to a change of coordinates.

% This coordinate view also aligns the expected sharpness with the projected Fisher. Consider a Gaussian posterior over the LoRA factor coordinates, (Q^\theta=\mathcal{N}(\mu_\theta,\Sigma_\theta)). The induced posterior over the full parameters is (W\stackrel{d}{=}W_0+\widehat{\Delta W}+\mathrm{unvec}(J_{\text{LoRA}}(\mu_\theta+\varepsilon))) with (\varepsilon\sim\mathcal{N}(0,\Sigma_\theta)). Around a stationary point, a second-order approximation (or a Fisher surrogate for the Hessian) gives
%  [
%  \mathbb{E}*{W\sim Q}!\left[\hat L*{S_t}(W)-\hat L_{S_t}(W_0+\widehat{\Delta W})\right]
%  \approx \tfrac12,\mathrm{tr}!\big(H_\Delta^{\text{fact}},\Sigma_\theta\big)
%  \approx \tfrac12,\mathrm{tr}!\big(G_{\text{fact}},\Sigma_\theta\big),
%  ]
%  where (H_\Delta^{\text{fact}}=J_{\text{LoRA}}^\top H_t J_{\text{LoRA}}). Thus the expected sharpness becomes a trace–inner-product between the projected curvature (G_{\text{fact}}) and the posterior covariance (\Sigma_\theta) in the same low-dimensional coordinate space. Since the KL complexity of the LoRA posterior is also computed in this coordinate space and is invariant under the pushforward to full parameters, the PAC-Bayesian objective couples empirical risk, expected sharpness, and KL in a single, coherent chart.

% Finally, this geometry clarifies why applying SAM only to the LoRA parameters improves stability. The inner SAM step searches for an adversarial perturbation inside the LoRA factor space and therefore aligns with the leading eigenvector of (G_{\text{fact}}). Updating against this perturbation suppresses (\lambda_{\max}(G_{\text{fact}})) and improves its condition number, which reduces the single worst direction of instability and spreads curvature more evenly across safer directions. Through the chain (E_\phi \rightarrow G_W \rightarrow G_{\text{fact}}), these spectral effects translate into a flatter and more isotropic feature geometry, smaller representation drift across tasks, and consequently a lower propensity to forget. Empirically, we observe the same signatures: the dominant eigenvalue of the projected Fisher shrinks, the effective rank of (E_\phi) increases, cross-task feature similarity becomes more stable, and forgetting decreases.



% Let the linearization from LoRA factor coordinates to full parameters be $\mathrm{vec}(\Delta W)\approx J_{\mathrm{LoRA}}\,\theta_{\mathrm{LoRA}}$, where $J_{\mathrm{LoRA}}$ is full column rank on admissible directions. Substituting into \eqref{eq:param-kl} yields the projected Fisher
% \begin{equation}\label{eq:gfact}
% \mathrm{KL}\!\big(p_\theta(\cdot\mid x)\,\big\|\,p_{\theta+\delta}(\cdot\mid x)\big)
% \;\approx\;\tfrac12\,\theta_{\mathrm{LoRA}}^\top G_{\mathrm{fact}}(x)\,\theta_{\mathrm{LoRA}},
% \qquad
% G_{\mathrm{fact}}(x):=J_{\mathrm{LoRA}}^\top F_W(x) J_{\mathrm{LoRA}}.
% \end{equation}
% Averaging over $x$ gives $G_{\mathrm{fact}}:=\mathbb{E}_{x}[G_{\mathrm{fact}}(x)]$. If $U$ is an orthonormal basis spanning the same subspace as $J_{\mathrm{LoRA}}$, then $U^\top G_W U$ is congruent to $G_{\mathrm{fact}}$ and therefore has the same spectrum.

% \begin{proposition}[Pullback–Projection equivalence]\label{prop:pp}
% Under the above assumptions,
% \[
% G_{\mathrm{fact}}
% =J_{\mathrm{LoRA}}^\top\Big(\mathbb{E}_{x\sim\mathcal{D}_t}\big[J_\phi(x)^\top E_\phi(x) J_\phi(x)\big]\Big)J_{\mathrm{LoRA}},
% \quad
% \mathrm{KL}\;\approx\;\tfrac12\,\theta_{\mathrm{LoRA}}^\top G_{\mathrm{fact}}\theta_{\mathrm{LoRA}}.
% \]
% \emph{Gradient objects}: $\nabla_\phi$ in $E_\phi(x)$ and $\nabla_W$ in $F_W(x)$;
% \emph{expectation objects}: inner over $y\sim p_\theta(y\mid x)$, outer over $x\sim\mathcal{D}_t$.
% \end{proposition}


% Let $Q^\theta=\mathcal{N}(\mu_\theta,\Sigma_\theta)$ denote a Gaussian posterior over LoRA factors. Using a second-order approximation with Fisher as a Hessian surrogate near a stationary point,
% \begin{equation}\label{eq:es-trace}
% \mathbb{E}_{W\sim Q}\big[\hat L_{S_t}(W)-\hat L_{S_t}(\hat W)\big]
% \;\approx\; \tfrac12\,\mathrm{tr}\big(G_{\mathrm{fact}}\,\Sigma_\theta\big).
% \end{equation}
% Since the KL complexity of $Q^\theta$ is computed in the same coordinate space and is invariant to pushforward, empirical risk, expected sharpness, and KL form a coherent PAC-Bayesian objective.






% We now restrict this global sensitivity to the LoRA subspace. Let
% $
% P_{\text {LoRA }} \in \mathbb{R}^{P \times R}, 
% $
% be a column-orthonormal basis for the active low-rank directions of the current LoRA adapter. Any admissible update under PECL can be written as
% $$
% \Delta W=P_{\text {LoRA }} \Delta \alpha, \quad \Delta \alpha \in \mathbb{R}^R .
% $$


% Projecting $G_W$ onto this subspace yields the LoRA-restricted sensitivity matrix:

% $$
% G_{\text {LoRA }}:=P_{\text {LoRA }}^{\top} G_W P_{\text {LoRA }} \in \mathbb{R}^{R \times R} .
% $$

% This matrix $G_{\text {LoRA }}$ is the curvature experienced by the optimizer when it is only allowed to move inside the LoRA adapter. Spectrally, its top eigenvalue $\lambda_{\max }\left(G_{\text {LoRA }}\right)$ characterizes the single most dangerous direction within the LoRA subspace: a step along that eigenvector will maximally distort the predictive distribution (and thus most aggressively damage past-task behavior). Likewise, because $G_W$ tracks Fisher/Gauss-Newton curvature of the full network, $G_{\text {LoRA }}$ can be interpreted as the restriction of that global curvature - including Hessian-like sharp directions - onto just the trainable adapter directions.



% We then restrict this global sensitivity to the LoRA subspace. We denote this subspace by a projection matrix $P_{\text {LoRA }} \in \mathbb{R}^{P \times R}, \quad R \ll P,$
% whose columns correspond to the active low-rank directions of the current LoRA adapter. Any admissible parameter update can be written as $\Delta W = P_{\text{LoRA}} \Delta r $

% by projecting with $P_{\text {LoRA }}$ :
% $$
% G_{\text {LoRA }}:=P_{\text {LoRA }}^{\top} G_W P_{\text {LoRA }} \in \mathbb{R}^{R \times R} .
% $$

\subsection{Discussion: SAM effect in LoRA}


\begin{table}[t]
\centering
\caption{Linear-probe CI performance and feature diagnostics in the LoRA subspace.}
\label{tab:q1_lp_feat_reorg}
\setlength{\tabcolsep}{4.5pt}
\resizebox{\linewidth}{!}{
\begin{tabular}{l l ccc cc cc}
\toprule
\multirow{2}{*}{\textbf{Opt.}} & \multirow{2}{*}{\textbf{Method}} &
\multicolumn{1}{c}{\textbf{NEM} ($\uparrow$)} &
\multicolumn{2}{c}{\textbf{Feature Energy} } &
\multicolumn{2}{c}{\textbf{Feature Stability}} &
\multicolumn{2}{c}{\textbf{OOD Performance($\uparrow$)}}\\
\cmidrule(lr){3-5}\cmidrule(lr){6-7}\cmidrule(lr){8-9}
 & & LA 
 & $\mathrm{tr}(E)$ & effective Rank
 & CKA$_{\text{stab}}$  & Prototype drift  
 & ImageNet-R & ImageNet-P \\ 
% \midrule
% \multirow{3}{*}{\centering{\textbf{Backbone}}}
% & MulTaskFT   \\
% & MulTaskLP  \\
% & MulTaskLoRA  \\
\midrule
\multirow{7}{*}{\centering\rotatebox[origin=c]{90}{\textbf{without SAM}}}
 & PerTaskFT        &78.7  &0.19  &63.21  &34.42  &0.77  &  &  \\
\cmidrule(lr){2-9}
 & SeqLoRA           &  &  &  &  &  &  &  \\
 \cmidrule(lr){2-9}
 & IncLoRA         &80.68  &0.57  &90.58  &13.15  &0.74  &  &  \\
 % & SDLoRA  &  &  &  &  &  &  &  \\
  % & MOS             &  &  &  &  &  &  &  \\
\cmidrule(lr){2-9}
 & OLoRA           &  &  &  &  &  &  &  \\
 % & TUNA            &  &  &  &  &  &  &  \\
\midrule
\multirow{7}{*}{\centering\rotatebox[origin=c]{90}{\textbf{with SAM}}}
 & PerTaskFT        &81.72  &0.32  &81.39  &19.95  &0.84  &  &  \\
\cmidrule(lr){2-9}
 & SeqLoRA         &80.78  &0.59  &89.88  &12.38  &0.73  &  &  \\
 \cmidrule(lr){2-9}
 & IncLoRA             &  &  &  &  &  &  &  \\
 % & SDLoRA  &  &  &  &  &  &  &  \\
% & MOS             &  &  &  &  &  &  &  \\
\cmidrule(lr){2-9}
 & OLoRA           &80.64  &0.58  &89.98  &11.99  &0.75  &  \\
 % & TUNA            &  &  &  &  &  &  &  \\
% \midrule
% \multicolumn{2}{r}{\textbf{Avg $\Delta$}} &  &  &  &  &  &  &  \\
\bottomrule
\end{tabular}
}
\end{table}



Finally, we can reinterpret SAM through this lens. SAM, when applied only to the LoRA parameters, works by explicitly searching (at each update) for the perturbation inside the LoRA subspace that maximizes the loss, and then updating in a way that counters that perturbation. Geometrically, that adversarial perturbation aligns with the top eigenvector of $G_{\text {LoRA }}$, i.e., the most dangerous direction for prediction stability.

Thus, SAM-in-LoRA has two direct spectral effects:
1. It suppresses
$
\lambda_{\max }\left(G_{\mathrm{LoRA}}\right),
$
reducing the strength of the single worst ("catastrophic") direction.
2. It reduces
Condition number $
\kappa\left(G_{\text {LoRA }}\right),
$
distributing risk more evenly across multiple gentler directions, rather than letting one direction dominate.

For continual learning, this is exactly what we want. Future tasks will continue to update their own LoRA blocks in low-rank subspaces. If the shared representation has already been shaped so that no single LoRA direction can easily shatter previous features, then new tasks can adapt without erasing old ones. In Section 5.5, we empirically confirm this story: applying SAM only in the LoRA subspace systematically shrinks $\lambda_{\max }\left(G_{\operatorname{LoRA}}\right)$, improves the effective rank of $E_\phi$, stabilizes cross-task representations, and ultimately reduces forgettir $\downarrow$








% Now we impose the PECL constraint: Only the current task's LoRA block is trainable, and this block spans a low-rank subspace of the full parameter space. We denote this subspace by a projection (or basis) matrix
% $
% P_{\text {LoRA }} \in \mathbb{R}^{P \times R}, \quad R \ll P,
% $
% Substituting this into the linearization above gives




% We factor  the predictor as \(f=\psi\circ\phi\), where \(\phi\) is the feature extractor with parameters \(\theta\) and \(\psi\) is a linear classifier with parameters \(W\). The logits are
%  $
%  f(x;\theta,W)=\psi(\phi(x;\theta);W)=W^{\top}\phi(x;\theta) \in \mathbb{R}^{C}$ and the predictive distribution is \(p\left(y \mid f(x;\theta,W)\right)\) be the predictive distribution (e.g., softmax over the logits).


% Adding the feature perturbation in the model, we have the conclusion that controlling perturbations along the feature subspace is locally equivalent to controlling an induced classifier perturbation in first‐order weight perturbation. 

% \(\psi\left(\phi(x;\theta) + \tilde{\phi}(x;\theta);\mathbf{w}\right) = (\mathbf{w} + \mathbf{w}A)\phi(x;\theta) = \psi\left(\phi(x;\theta);\mathbf{w} + \mathbf{w}A\right).\)


% $\tilde{\phi}(x;\theta)=A\phi(x;\theta)$ 
% where $A\in\mathbb{R}^{d\times d}$ denotes the perturbation matrix with a sufficiently small spectral norm ($\|A\|\ll 1$). When this perturbation is incorporated into the model, the corresponding logits satisfy
% $$
% \psi(\phi(x;\theta)+\tilde{\phi}(x;\theta);{        W})=W^{\top}(\phi(x;\theta)+\tilde{\phi}(x;\theta)) =(W^{\top}+W^{\top}A)\phi(x;\theta) = \psi(\phi(x;\theta);W+WA))
% $$
% Thus, controlling perturbations along the feature subspace is equivalent to controlling the induced first-order perturbation of the classifier weights. 

% We consider a small perturbation in the {feature space}, defined as $\Delta\phi(x;\theta)$. By the chain rule and the linear head $f=W^{\top}\phi$, the difference in logits is $\delta(x) = W^{\top} \Delta{\phi}(x;\theta)$.
% To quantify how this perturbation in the {feature space} alters predictions at input $x$, we measure the KL divergence between the perturbed and unperturbed predictive distributions:
% \begin{equation}
% \label{feature-pro}
%     \mathrm{KL}\left(p(y \mid \phi\left(x ; \theta \right)+\Delta{\phi} ; W) \| p\left(y \mid \phi\left(x ; \theta\right) ; W\right)\right)=\Delta{\phi}^{\top} E_\phi\Delta{\phi}+ O\big(\|\Delta{\phi}\|^{3}\big)
% \end{equation}

% where the {local feature matrix} is defined by
%  $$
% E_\phi \triangleq \mathbb{E}_{y \sim p(\cdot \mid f)}\left[\nabla_\phi \log p(y \mid f) \nabla_\phi \log p(y \mid f)^{\top}\right]
% $$
% Under the softmax and Cross Entropy loss, we have an analytic formulation
% $$
% E_\phi =\mathbb{E}_{y \sim p(\cdot \mid f)}\left[W(Diag(p)-pp^\top) 
% W^\top\right]
% $$ 
% This provides a principled pseudo-metric on the feature space that weights perturbation directions by their impact on the predictive distribution.


% We  then consider a small perturbation in the parameter sapce to check its effect on feature space. 
% Since the characteristics are differentiable with respect to parameters, we perform a second-order Taylor expansion 
% $$
%     \phi(\theta+\Delta \theta) \approx \phi(\theta) + J_\phi \Delta \theta + \tfrac12  \Delta \theta \mathcal\  H_\phi \Delta \theta
% $$
% where the characteristic Jacobian matrix is $J_\phi:=\frac{\partial\phi(x;\theta)}{\partial \theta}$
% The first-order approximation 
% \begin{equation}
% \label{feature:1}
%     \Delta \phi \approx   \phi(\theta) + J_\phi \Delta \theta
% \end{equation}
% Substituting (\ref{feature:1}) back into (\ref{feature-pro}) yields the KL quadratic form of the weight perturbation 
% \begin{equation}
% \mathrm{KL} \approx \Delta \theta^{\top} \underbrace{\left(\mathbb{E}_x\left[J_\phi(x)^{\top} E_\phi(x) J_\phi(x)\right]\right)}_{\mathbf{G}_\theta \succeq 0} \Delta \theta .
% \end{equation}

% In fact  the aboves measure the peratution in the paramter's effect on probability ,and we have 
% \begin{equation}
% \mathrm{KL}\left(P_{x, y}(\theta+d) \| P_{x, y}(\theta)\right)=\frac{1}{2} d^{\top} F d+O\left(d^3\right)
% \end{equation}
% where
% $$
% F=\frac{1}{|S|} \sum_{x \in S_x} \mathrm{E}_{P_{y \mid x}}\left[\nabla \log p(y \mid x, \theta) \nabla \log p(y \mid x, \theta)^{\top}\right] \quad \text { or } \quad F=-\frac{1}{|S|} \sum_{x \in S_x} \mathrm{E}_{P_{y \mid x}}\left[H_{\log p(y \mid x, \theta)}\right] .
% $$
 




% SAM minimizes the worst–case loss increase, which—under a second–order view—amounts to shrinking the spectral radius and improving the conditioning of the  metric \(E_\phi\). Consequently, the spectrum of \(E_\phi\) becomes less anisotropic (lower \(\lambda_{\max}\), smaller condition number), i.e., the “energy” is no longer concentrated in a single dominant direction but distributed more uniformly across directions, yielding smoother and more stable representations.

% When updates are restricted to a low–rank LoRA subspace, SAM primarily reduces the spectral radius and improves the conditioning of the projected metric \(E_\phi^{(\mathrm{LoRA})}=P_{\mathrm{LoRA}}^\top E_\phi P_{\mathrm{LoRA}}\). This directly lowers the local KL sensitivity along learnable directions (less forgetting) while preserving plasticity in less sensitive directions, explaining the empirical gains observed with LoRA–restricted flatness control.




% $$
% \mathrm{KL}\left(
% p\left(y \,\middle|\, f(x;\theta,\mathbf{w}) + \delta(x)\right)
% \,\bigg\|\, 
% p\left(y \,\middle|\, f(x;\theta,\mathbf{w})\right)
% \right).
% $$
%  For sufficiently small $\delta(x)$, the second‐order Taylor expansion of $\log p(y\mid f)$ at the logits yields the Fisher approximation
% $$
% \mathrm{KL}\left(
% p\left(y \,\middle|\, f+\delta\right)
% \,\bigg\|\, 
% p\left(y \,\middle|\, f\right)\right)
% =
% \delta(x)^{\!\top}\ F \delta(x)
% + O\big(\|\delta(x)\|^{3}\big),
% $$
% where $F$ is the Fisher information matrix with respect to the logits,
% $$
%  F \triangleq
%  \mathbb{E}_{y\sim p(\cdot \mid f)}
%  \left[
%  \nabla_{f}\log p(y\mid f)\nabla_{f}\log p(y\mid f)^{\top}
%  \right].
%  $$

% we substitute 
% \(\delta(x)=W^{\top}\tilde{\phi}(x;\theta)\)  into the Fisher approximation.  Then

% To map this logit-space metric to the feature space, which enables the direct characterization of changes in the predictive distribution via perturbations in the feature space, we have 
% $$
% \mathrm{KL}\left(p(y \mid \phi\left(x ; \theta \right)+\Delta{\phi} ; W) \| p\left(y \mid \phi\left(x ; \theta\right) ; W\right)\right)=\Delta{\phi}^{\top} E_\phi\Delta{\phi}+ O\big(\|\Delta{\phi}\|^{3}\big)
% $$
% where 
%  $$
% E_\phi \triangleq \mathbb{E}_{y \sim p(\cdot \mid f)}\left[\nabla_\phi \log p(y \mid f) \nabla_\phi \log p(y \mid f)^{\top}\right] =  \mathbb{E}_{y \sim p(\cdot \mid f)} \left[ WF W\right]
% $$

% By the chain rule and the linear head $f=W^{\top}\phi$, we have
% $
% \nabla_{\phi}\log p(y\mid f)=W\,\nabla_{f}\log p(y\mid f),
% $
% hence
% $$
% \boxed{~E_{\phi}(x;\theta,W)=W\,F\!\big(f(x;\theta,W)\big)\,W^{\top}~},
% \qquad
% F\!\big(f(x;\theta,W)\big)
% =
% \mathbb{E}_{y\sim p(\cdot\mid f)}\!\left[
% \nabla_{f}\log p(y\mid f)\,\nabla_{f}\log p(y\mid f)^{\!\top}
% \right].
% $$









% This yields the per-sample quadratic form for the KL divergence: 
% $$
% \mathrm{KL} \approx
%  \frac{1}{2}
%  \tilde{\phi}(x;\theta)^{\top}
%  \underbrace{W\ F\big(f(x;\theta,W)\big)\ W^{\top}}
% \tilde{\phi}(x;\theta),
% $$
%  where
% $$
%  E_f(x;\theta,W) \triangleq W\ F\big(f(x;\theta,W)\big)W^{\top}\ \in\ \mathbb{R}^{d\times d}
% $$
%  is the {local feature matrix}, symmetric and positive semi‐definite.
 
% For the softmax, the Fisher information matrix simplifies to \(F=\mathrm{Diag}(p)-pp^{\top}\) with \(p=p(\cdot\mid f)\). To "pull back" this logit-space metric to the feature space, we substitute




% Substituting \(\delta(x)=W^{\top}\tilde{\phi}(x;\theta)\) pulls this metric back to feature space. gives the per‐sample quadratic form
% Plugging in \(\delta(x)=W^{\top} A\phi(x;\theta)\)  

 
 
%  Averaging over the data distribution  yields the {Empirical Feature Matrix} (EFM)
% $
%  E_t \triangleq  \mathbb{E}_{x\in X_t}\big[E_f(x;\theta_t,W_t)\big]
% $.





% Taking expectation over the data distribution yields The Empirical Feature Matrix (EFM) :

% \begin{equation}
% E_t = E_{f(x;\theta,w)}
% \mathbb{E}_{x}\left[ 
% \right]
% \end{equation}





% $$
% \delta(x) \triangleq \tilde{f}(x;\theta,\mathbf{w}) - f(x;\theta,\mathbf{w})
% \;=\; \mathbf{w}^{\top}(I{+}A)\phi(x;\theta) - \mathbf{w}^{\top}\phi(x;\theta)
% \;=\; \mathbf{w}^{\top}A\,\phi(x;\theta).
% $$
% Equivalently, the same perturbed logits arise from a first-order weight-space perturbation \(\tilde{\mathbf{w}}=(I{+}A)^{\top}\mathbf{w}\), since \(\mathbf{w}^{\top}(I{+}A)\phi = \tilde{\mathbf{w}}^{\top}\phi\). 








% We decompose the model $f$ into a feature extractor $\phi $ and a classifier $\psi$ as $f = \phi \circ \psi$, the origin parampter $\boldsymbol{\theta}$ also divided into two paramter, for easy we still use $\theta$ to represent the parampter of the feature extractor $\phi $ and use $\mathbf{w}$ as the parampter of the classifier.
% Let $\phi(x;\theta)$ be the feature vector extracted from input $x$ and assume the classifier is just a linear classifier , then $\hat{y} = \psi(\phi(x;\theta);\mathbf{w})= \mathbf{w}\phi(x;\theta)$.

% Then we add some pertubation in the feature $\delta=A \phi(\phi(x;\theta))$, Multiplicative perturbations in feature space by arbitrary matrices $A \in \mathbb{R}$ correspond to perturbations in parameter space, i.e.,
% \begin{equation}
%     \psi(\phi(x)+\delta;\mathbf{w})  =  \psi(\phi(x)+A \phi(x);\mathbf{w})=\mathbf{w}(\phi(x)+A \phi(x))=(\mathbf{w}+\mathbf{w} A) \phi(x)=\psi( \phi(x);\mathbf{w}+\mathbf{w} A) .
% \end{equation}
% coresponding result is  $\hat{y}_1 = \psi(\phi(x;\theta)+\delta;\mathbf{w})= \mathbf{w}(\phi(x;\theta)+\sigma)$ 
% Then we can measure the difference between the origin and new one as follow: 

% \begin{equation}
% \begin{aligned}
%     KL(\ p(y|f(x;\theta)+\sigma;\mathbf{w}) ||p(y|f(x;\theta)+\sigma;\mathbf{w}) \  ) &= \delta^T E_{f\left(x ; \theta_t\right)} \delta+ O(\|\sigma\|^3) \\
% where
%     E_{f\left(x ; \theta_t\right)}&=\underset{y \sim p(y)}{\mathbb{E}}\left[\left(\frac{\partial \log p(y)}{\partial f\left(x ; \theta_t\right)}\right)\left(\frac{\partial \log p(y)}{\partial f\left(x ; \theta_t\right)}\right)^{\top}\right] .
% \end{aligned}
% \end{equation}




% % $$ be the feature vector extracted from input x ∈ Xt from task t


% $f(x) = \psi(\phi(x))=\mathbf{w}\phi(x)$. Petzka connect feature robustness to flatness using the following key identity: Multiplicative perturbations in feature space by arbitrary matrices $A \in \mathbb{R}$ correspond to perturbations in parameter space, i.e.,
% \begin{equation}
%     \psi(\mathbf{w}, \phi(x)+A \phi(x))=\mathbf{w}(\phi(x)+A \phi(x))=(\mathbf{w}+\mathbf{w} A) \phi(x)=\psi(\mathbf{w}+\mathbf{w} A, \phi(x)) .
% \end{equation}
% feature robustness is approximated by a novel, but natural, loss Hessian-based relative flatness measure under the assumption that the distribution can be approximated by locally constant labels. Under this assumption and if the data is representative, then flatness is the main predictor of generalization.







% and a linear head, $\theta=(\theta_f,W_c)$ with $f=f(x;\theta_f)$.
% $z=(W_c e)$, . This decomposition lets us map weight–loss geometry to feature–logit geometry via $F_e=W_c^{\top}H_{zz}W_c$, and to test whether LoRA-subspace flatness primarily acts by shrinking the effective curvature seen by features. Concretely, we adopt a linear-probe protocol to factor out classifier idiosyncrasies and pair it with feature-level diagnostics (CKA/RSA, EFM spectra, and $F_e$ spectra). If the Q1 gains are representation-mediated, these readouts should improve consistently when using SAM in the LoRA subspace.















% There are a lot of method to improve flatness, the famous one in SAM, which generate Perturbs weights in the direction of steepest loss ascent to induce flatter minima:
%     $$
%     \epsilon^*(\theta, \rho) := \arg \max_{\|\epsilon\|_2 \leq \rho} \hat{L}_S(\theta + \epsilon) \approx \rho \cdot \frac{g}{\|g\|_2}
%     $$. We adopt this as a base method with the SGD to compare the influence of promting flatness effect. 

% The result table in tab show that 















% (backbone frozen)





% and compare against matched-budget {global} controls.
% Before any CL claim, we verify that expected sharpness (E--Sh) {decreases} in the intended scope (LoRA vs.\ global) as a manipulation check.
% Interventions cover {directional} (SAM/FSAM) and {isotropic/anisotropic} (Gaussian/Fisher/RWP) perturbations with a small $\rho/\sigma$ sweep; readouts include FAA/LA, ABWT, AFG (5 seeds, paired tests), E--Sh (subspace/global), and accuracy--time curves.
% “LoRA-only” means perturbing and updating \emph{only} adapter parameters; “Global” perturbs the whole model in the inner step but \emph{masks updates to frozen parameters} in the outer step; all comparisons are step/optimizer/schedule/batch matched.




















% \section{Problem Formulation and Experimental Foundations}

% \paragraph{From the introduction to a subspace-centric testbed.}
% Building on the introduction, our objective is to determine \emph{whether}, \emph{why}, and \emph{when} improving \emph{LoRA-subspace flatness} benefits continual learning (CL). Rather than enumerating architectures or datasets, this section specifies a compact, reusable testbed that (i) intervenes only where adaptation occurs (the adapter subspace), (ii) enforces fair \emph{budget matching} against global counterparts, (iii) performs explicit \emph{manipulation checks} to validate that the intended geometric change happened, and (iv) defines readouts and statistical treatment used consistently throughout. Concrete implementation details (task orders, hyperparameters, preprocessing) are deferred to the Protocol appendix (\S\ref{app:protocol}).

% \subsection{Intervention scope and factors (``how we intervene'')}
% \label{sec:interventions}
% We operationalize flatness control as an intervention with minimal, orthogonal factors:
% \begin{itemize}[leftmargin=*,noitemsep,topsep=0pt]
% \item \textbf{Scope}: \emph{LoRA subspace} (perturb/regularize only adapter parameters) vs.\ \emph{global} (full-parameter control).
% \item \textbf{Geometry}: directionally aware vs.\ isotropic mechanisms (e.g., SAM/FSAM-type vs.\ Gaussian/Fisher/RWP-style perturbations).
% \item \textbf{Strength}: a radius/variance grid to expose performance–strength trade-offs (single operating points are insufficient).
% \item \textbf{Frequency}: intervention cadence (per step, every $m$ steps, or at task boundaries).
% \item \textbf{Update set}: unless stated otherwise, only adapter parameters are updated; the backbone remains frozen.
% \end{itemize}
% These factors define the minimal space needed to answer Q1--Q3 without confounding the locus of learning (adapter vs.\ full model).

% \subsection{Budget matching and fairness (``how we compare'')}
% \label{sec:fairness}
% To isolate geometric effects from training effort, all primary comparisons are \emph{budget matched}:
% \begin{itemize}[leftmargin=*,noitemsep,topsep=0pt]
% \item \textbf{Optimization budget}: identical numbers of optimization steps/epochs, optimizer family, batch size, and learning-rate schedule across compared runs.
% \item \textbf{Overhead accounting}: interventions that require extra passes (e.g., SAM-like) are paired with matched-overhead baselines; we also report \emph{accuracy–time} frontiers to complement step-matched analyses.
% \item \textbf{Randomness control}: identical seeds per paired comparison; task orders are held fixed within each comparison and varied only in designated robustness checks.
% \end{itemize}
% Any deviation from strict matching (e.g., for ablations) is explicitly labeled and accompanied by efficiency reporting.

% \subsection{Manipulation checks (``did the intervention work?'')}
% \label{sec:manipcheck}
% Before interpreting CL outcomes, we verify that the intended geometric manipulation occurred \emph{in the intended scope}:
% \begin{itemize}[leftmargin=*,noitemsep,topsep=0pt]
% \item \textbf{Expected sharpness (E--Sh)}: Monte Carlo estimation within the LoRA subspace (and, for controls, globally) shows a statistically significant reduction relative to the corresponding no-intervention run.
% \item \textbf{Direction/strength sanity}: the measured reduction is monotone in strength within a reasonable range; pathological regimes (e.g., over-perturbation that inflates empirical loss) are flagged.
% \item \textbf{Scope integrity}: no unintended updates outside the designated parameter set (checked via parameter diff/grad masks).
% \end{itemize}
% Configurations that fail manipulation checks are excluded from answering Q1--Q3 and are reported only as diagnostics.

% \subsection{Readouts and statistical treatment (``how we measure and report'')}
% \label{sec:readouts}
% We adopt a unified readout panel and reporting standard:
% \begin{itemize}[leftmargin=*,noitemsep,topsep=0pt]
% \item \textbf{CL metrics}: final/last accuracy (FAA/LA), average backward transfer (ABWT), and average forward generalization (AFG), all computed from the evaluation matrix over the task sequence.
% \item \textbf{Geometry}: E--Sh in the LoRA subspace and globally; first-order flatness via neighborhood gradient norms.
% \item \textbf{Feature diagnostics (for Q2)}: representation smoothness/stability via empirical feature matrix (EFM) spectra, CKA/RSA similarity, and linear-probe retention.
% \item \textbf{Efficiency}: accuracy–time Pareto frontiers, extra forward/backward passes, and peak memory.
% \item \textbf{Statistics}: at least five seeds per setting; mean$\pm$std and 95\% CIs; paired tests across seeds with Holm–Bonferroni correction for multiple comparisons; effect sizes (e.g., Cohen’s $d$ or rank-biserial $r$). Correlation/partial-correlation and (where applicable) simple mediation analyses connect geometry/feature diagnostics to CL outcomes.
% \end{itemize}

% \paragraph{From foundations to the three questions.}
% With interventions, fairness, manipulation checks, and readouts in place, we address Q1 (sufficiency and cross-domain robustness) by contrasting subspace vs.\ global control under matched budgets; Q2 (mechanistic pathway) by relating gains to feature-level diagnostics; and Q3 (factors and matching) by varying structure, geometry, rank, similarity, and sequence length. Theoretical synthesis in \S\ref{sec:theory} reconciles the observed patterns with a hierarchical PAC–Bayesian geometry term and a PCA view of feature covariance. Full implementation details are provided in \S\ref{app:protocol}.

% \paragraph{Benchmarks.} (i) \textbf{Standard CL} (AG News, Amazon Reviews, DBpedia, Yahoo) for within-domain transfer; (ii) \textbf{Long Sequence} (15 tasks from GLUE/SuperGLUE/IMDB etc.) for domain shift.
% \paragraph{Backbones/Protocols.} T5-based sequence classification with three task orders for each benchmark.
% \paragraph{Baselines.} \seq, \inc, \olora\ each with: No-Noise, Full-Gaussian, Full-Fisher, and subspace-only \sam/Gaussian/Fisher.
% \paragraph{Metrics.} Last Accuracy (LA), average Backward Transfer (ABWT), average Forward Generalization (AFG).
% \paragraph{Statistical reporting.} We report mean$\pm$std over $k{=}5$ seeds; paired $t$-tests with Holm--Bonferroni across methods per setting; 95\% CIs and Cohen's $d$ for key comparisons. Variance across tasks is summarized with per-task bars and aggregated with across-task standard deviation.
% \paragraph{Hyperparameters.} We sweep perturbation strengths (e.g., $\rho,\sigma$) on a validation split and plot accuracy vs.~strength to expose the bound-predicted trade-off.
% \paragraph{Reproducibility.} Exact task orders, seeds, and evaluation code are released; we log all hashes and environment versions.






% \section{Subspace Flatness in Continual Learning}


% \subsection{Subspace Flatness in Continual Learning}
% in parameter-efficient continual learning, updates are confined to a low-rank adapter, which introduces two related but distinct geometric objects: global flatness of the entire network and sharpness within the trainable LoRA subspace. Our central question is a sufficiency test: under matched training budgets, is reducing sharpness within the LoRA subspace alone sufficient to improve continual-learning objective





% \subsection{LoRA space based CL}


% Modern large-scale pretrained models are often adapted to downstream tasks using parameter-efficient methods. In the context of CL, we focus on three representative LoRA-based strategies:
% \begin{itemize}
%  \item \textbf{Sequence LoRA (SeqLoRA)} : A single LoRA pair $(A, B)$ is trained sequentially across tasks while keeping the backbone weights $W_0$ frozen. The final model after $N$ tasks is given by $W_0 + AB$.
% \item \textbf{Incremental LoRA (IncLoRA)} : A new LoRA pair $(A_t, B_t)$ is incrementally added for each task $t$, and only the current pair is updated while all previous modules and $W_0$ are fixed. The resulting model is $W_0 + \sum_{t=1}^N A_t B_t$.
% \item \textbf{Orthogonal LoRA (OLoRA)} : Extending IncLoRA, OLoRA~\cite{wang2023orthogonal} introduces an orthogonality constraint across tasks to encourage adapter subspace disentanglement. Specifically, during task $t$, a penalty $\mathcal{L}_{\text{orth}}^{(t)} = \sum_{i=1}^{t-1} \|A^{(i)\top} A^{(t)}\|_F^2$ is applied, where $A^{(i)}$ denotes the LoRA-A matrix from task $i$.
% \end{itemize}




% However, Hessian–based measure depend on the parameterization and may change under simple rescalings.











% Sharpness (Flatness) is a property of the parameter space quantifying the change in loss under small parameter
% perturbations. The most popular mathematical definitions considers the maximal loss value within a raduis
% $$
% \max_{\epsilon \in B(\hat\theta, \rho)} \big( \hat{L}_{S}(\hat{\theta} + \epsilon) - \hat{L}_{S}(\hat{\theta}) \big),
% $$,
% where $B(\hat\theta, \rho)$ denotes a ball of radius $\rho$ centered at $\hat\theta$.
% It is also defined based on a random perturbation $\epsilon$ from a given distribution rather than the maximal one.
% $$
%  \mathbb{E}_{\epsilon \sim \mathcal{N}(0,\,\Sigma_Q)} \left[ \hat{L}_{\mathcal{D}}(\theta + \epsilon) - \hat{L}_{\mathcal{D}}(\theta) \right],
% $$
% where $\mathcal{N}(0,\,\Sigma_Q)$ is a Gaussian distribution, 


% According to  the  second-order Taylor approximation around  the flat minal, we have

% $$
%  L(\theta+\epsilon) = L(\theta) + g^{\top}\epsilon + \tfrac12\,\epsilon^{\top}H\,\epsilon + R_3(\theta,\epsilon)
% $$
% It is common believed the first order is neat to 0 around the falt minal , then we can get the followsing realtion 

% $$
% \hat{L}_S(\hat{\theta} + \epsilon) - \hat{L}_S(\hat{\theta}) \approx \frac{1}{2} \epsilon^\top H \epsilon
% $$
% Assume H is  , we also have following relation 
% \citep{caoTradeoffFlatnessOptimization2024}: 
% \begin{equation}\label{eq: relation in Hessian Matrix}
%     \frac{1}{M} \, \text{Tr}(H) \leq \|H\|_2 \leq \|H\|_F \leq \text{Tr}(H) \leq M \|H\|_2
% \end{equation}
% So it is also measured by the trace of the loss Hessian  and the max eigenvalue. Besides,  the Fisher Information Matrix is also a good approximate for hessian maxrix. In SSAM, the author use empirical Fisher to approxiamate Hessian matrxi.

% Zhang def first-order flatness to measure  the maximal gradient norm in the neighbourhood as follow 


% However, the defination based on Hessien has the problem of reparametrizing. And The raltion between paramter and loss via feature. 

% In modern model seting . 

 
%  He define the relative flatness  

% and Empirical Feature Matrix are defined as follow: 


% the link between flatness in the weight space and feature space as follow: 








% defination sample a random perturbation εr from a given distribution for each iteration rather than accoridng to gradient direction and get the defination of  Expected Sharpness as follow: 








% Following the intuition that sharpness reflects the sensitivity of a solution to perturbations in parameter space,  the most popular mathematical definitions of flatness considers the maximal loss value within a raduis as follows:
% \begin{definition}~\cite{foretSharpnessAwareMinimizationEfficiently2021}
% The adversarial sharpness of a solution $\hat{\theta}$ on dataset $S$ is defined as
% \begin{equation}
%     \epsilon\text{-Sh}(\hat{\theta}, \rho) :=\max_{\epsilon \in B(\hat\theta, \rho)} \big( \hat{L}_{S}(\hat{\theta} + \epsilon) - \hat{L}_{S}(\hat{\theta}) \big),
% \end{equation}
% where $B(\hat\theta, \rho)$ denotes a ball of radius $\rho$ centered at $\hat\theta$.
% \end{definition}

% It is also defination sample a random perturbation εr from a given distribution for each iteration rather than accoridng to gradient direction and get the defination of  Expected Sharpness as follow: 
% \begin{definition}[Expected Sharpness]~\cite{liRevisitingRandomWeight2024}
% For task $i$, the expected sharpness of a solution $\theta$ is defined as
% \begin{equation}\label{def: Task-Aware Expected Sharpness}
%     \text{E-Sh}(\theta, \Sigma_Q) := \mathbb{E}_{\epsilon \sim \mathcal{N}(0,\,\Sigma_Q)} \left[ \hat{L}_{\mathcal{D}}(\theta + \epsilon) - \hat{L}_{\mathcal{D}}(\theta) \right],
% \end{equation}
% where $\mathcal{N}(0,\,\Sigma_Q)$ is a Gaussian distribution, typically instantiated as $\mathcal{N}(0,\, \sigma^2 \mathbf{I})$.
% \end{definition}

% Zhang def first-order flatness to measure  the maximal gradient norm in the neighbourhood of a point as follows: 
% Definition 4.1 (ρ-first-order flatness). For any ρ > 0, the ρ-first-order flatness R(1)  ρ (θ) of function Lˆ(θ) at a point θ is defined as  R(1)  ρ (θ) ≜ ρ · max  θ′ ∈B (θ,ρ)  ∇Lˆ(θ′) , ∀θ ∈ Θ. (3)  Here ρ is the perturbation radius that controls the magnitude of the neighbourhood.


% According to  the  second-order Taylor approximation around  the flat minal , Assume $H$ is a \textbf{ real symmetric matrix} with a well-defined spectral decomposition,we can get follow relation between Hessien matrix  and sharpness 
%  \begin{equation}
%     \textbf{Sharpenss}_\rho(\hat{\theta}) := \max_{\|\epsilon\|_2 \leq \rho} \left[\hat{L}_S(\hat{\theta} + \epsilon) - \hat{L}_S(\hat{\theta})\right] \approx \frac{1}{2} \max_{\|\epsilon\|_2 \leq \rho} \epsilon^\top H \epsilon = \frac{1}{2} \rho^2 \cdot \lambda_{\max}(H)
% \end{equation}
% it is common to adopt a more tractable alternative — the \textbf{Fisher Information Matrix} — to approximate sharpness. To approximate Fisher Information Matrix, some paper also use  "empirical Fisher" to approxiamte it. 

% However, the defination based on Hessien has the problem of reparametrizing. According to model NN arithcer, the infect of nosie on paramter space on the loss via the feautre space, some paper definaiton the flat in the feature space rather than loss feature robustness,  \begin{definition} Relative Flatness

%  For a model $f(x, \mathbf{w})=g(\mathbf{w} \phi(x))$, $\mathbf{w} \in \mathbb{R}^{d \times m}$, with a twice differentiable function $g$, a twice differentiable loss function $\ell$ and a sample set $S$, relative flatness is defined by

% $$
% \kappa_{T r}^\phi(\mathbf{w}):=\sum_{s, s^{\prime}=1}^d\left\langle\mathbf{w}_s, \mathbf{w}_{s^{\prime}}\right\rangle \cdot \operatorname{Tr}\left(H_{s, s^{\prime}}(\mathbf{w}, \phi(S))\right),
% $$

% where $\operatorname{Tr}$ denote the trace and $\left\langle\mathbf{w}_s, \mathbf{w}_{s^{\prime}}\right\rangle=\mathbf{w}_s \mathbf{w}_{s^{\prime}}^T$ the scalar product of two row vectors.
% \end{definition}

% This defination is also link with the Empirical Feature Matrix (EFM), 
% The Empirical Feature Matrix (EFM) associated with task t is obtained by taking the expected value of Eq. 6 over the entire dataset at task t:  Et = x∼EXt  [Ef(x;θt)]. (7)



% We define sharpness as loss sensitivity to small parameter changes and make the restriction to the LoRA subspace explicit. The adversarial sharpness at radius $\rho>0$ is
% \[
% \epsilon\text{-}\mathrm{Sh}(\boldsymbol{\theta};\rho,\mathcal{U})
% =\max_{\boldsymbol{\epsilon}\in\mathbb{B}_\rho}\Big[\hat{L}_{S}\big(\boldsymbol{\theta}+\Pi_{\mathcal{U}}\boldsymbol{\epsilon}\big)-\hat{L}_{S}(\boldsymbol{\theta})\Big],
% \]
% and the expected sharpness for a covariance $\boldsymbol{\Sigma}\succeq\mathbf{0}$ is
% \[
% \mathrm{E}\text{-}\mathrm{Sh}(\boldsymbol{\theta};\boldsymbol{\Sigma},\mathcal{U})
% =\mathbb{E}_{\boldsymbol{\epsilon}\sim\mathcal{N}(\mathbf{0},\boldsymbol{\Sigma})}\Big[\hat{L}_{S}\big(\boldsymbol{\theta}+\Pi_{\mathcal{U}}\boldsymbol{\epsilon}\big)-\hat{L}_{S}(\boldsymbol{\theta})\Big].
% \]
% These two quantities are the primary objects in our analysis and experiments because they align with subspace-restricted perturbation protocols and admit direct connections to PAC–Bayesian bounds.

% Additional viewpoints offer complementary diagnostics. A second-order Taylor approximation around a stationary reference $\hat{\boldsymbol{\theta}}$ yields $\hat{L}_{S}(\hat{\boldsymbol{\theta}}+\boldsymbol{\epsilon})\approx \hat{L}_{S}(\hat{\boldsymbol{\theta}})+\tfrac{1}{2}\boldsymbol{\epsilon}^{\top}\mathbf{H}\,\boldsymbol{\epsilon}$ with Hessian $\mathbf{H}=\nabla^2 \hat{L}_{S}(\hat{\boldsymbol{\theta}})$; the worst-case increase over an $\ell_2$ ball of radius $\rho$ is then approximately $\tfrac{1}{2}\rho^2\lambda_{\max}(\mathbf{H})$, motivating spectral-norm, trace, or Frobenius summaries of curvature when principal directions can be computed or estimated. Since exact spectra are prohibitive in modern models, a tractable alternative is the Fisher information $\mathbf{F}(\boldsymbol{\theta})=\mathbb{E}\big[\nabla_{\boldsymbol{\theta}}\log p(y|x,\boldsymbol{\theta})\,\nabla_{\boldsymbol{\theta}}\log p(y|x,\boldsymbol{\theta})^\top\big]$ and its empirical variant that replaces the expectation by the observed label; Fisher-guided perturbations further align covariance with $\mathbf{F}$ (or a diagonal approximation) and can be restricted to $\mathcal{U}$.

% A first-order notion of flatness is obtained by measuring gradient magnitudes in a neighborhood. For a radius $\rho$, define $R^{(1)}_{\rho}(\boldsymbol{\theta})=\rho\,\max_{\boldsymbol{\theta}'\in B(\boldsymbol{\theta},\rho)}\|\nabla \hat{L}_{S}(\boldsymbol{\theta}')\|$. Under a local quadratic model around a minimizer, the ratio $R^{(1)}_{\rho}(\hat{\boldsymbol{\theta}})/\rho^2$ coincides with $\lambda_{\max}(\nabla^2\hat{L}_{S}(\hat{\boldsymbol{\theta}}))$, and $R^{(1)}_{\rho}$ upper-bounds the zeroth-order adversarial increase, offering a gradient-based proxy that is often cheaper than Hessian eigendecompositions. Because curvature measures are sensitive to reparameterizations, normalized sharpness rescales parameters to remove node-wise scaling effects, and relative flatness weighs curvature by parameter couplings so that the resulting value is invariant to admissible reparameterizations. In this paper, Hessian and Fisher summaries, together with first-order flatness, serve as diagnostics that corroborate our findings, while the subspace-restricted $\mathrm{E}\text{-}\mathrm{Sh}$ remains the primary, theory-aligned quantity that drives our analysis and experimental design.


% \section{Experimental design aligned to the RQs}
% \paragraph{Benchmarks.} (i) \textbf{Standard CL} (AG News, Amazon Reviews, DBpedia, Yahoo) for within-domain transfer; (ii) \textbf{Long Sequence} (15 tasks from GLUE/SuperGLUE/IMDB etc.) for domain shift.
% \paragraph{Backbones/Protocols.} T5-based sequence classification with three task orders for each benchmark.
% \paragraph{Baselines.} \seq, \inc, \olora\ each with: No-Noise, Full-Gaussian, Full-Fisher, and subspace-only \sam/Gaussian/Fisher.
% \paragraph{Metrics.} Last Accuracy (LA), average Backward Transfer (ABWT), average Forward Generalization (AFG).
% \paragraph{Statistical reporting.} We report mean$\pm$std over $k{=}5$ seeds; paired $t$-tests with Holm--Bonferroni across methods per setting; 95\% CIs and Cohen's $d$ for key comparisons. Variance across tasks is summarized with per-task bars and aggregated with across-task standard deviation.
% \paragraph{Hyperparameters.} We sweep perturbation strengths (e.g., $\rho,\sigma$) on a validation split and plot accuracy vs.~strength to expose the bound-predicted trade-off.
% \paragraph{Reproducibility.} Exact task orders, seeds, and evaluation code are released; we log all hashes and environment versions.

% \section{Results as answers to the RQs}
% \subsection*{Answer to RQ1: Subspace sharpness is sufficient}
% Across both benchmarks, \emph{subspace-only} perturbations (LoRA-Gaussian/\sam/Fisher) improve ABWT/AFG/LA over No-Noise and over Full-Model perturbations. This supports the claim that controlling sharpness in the trainable adapter subspace is sufficient for improved CL generalization.

% \subsection*{Answer to RQ2: Modularity interacts with perturbation geometry}
% \inc\ benefits most from isotropic subspace noise on related tasks (Standard CL), while \olora\ benefits more from directional \sam\ under heterogeneity (Long Sequence), consistent with the orthogonality-constrained geometry.

% \subsection*{Answer to RQ3: Trade-off in perturbation strength}
% Performance exhibits a U-shape as a function of perturbation strength; moderate noise reduces sharpness and improves transfer without inflating empirical loss.

% \section{Ablations and sensitivity}
% We vary (i) adapter rank, (ii) orthogonality penalty, (iii) perturbation scope (subspace vs.~global), and (iv) dataset order. Results are consistent with the bound's decomposition and provide guidance on method selection by regime.

% \section{Related work}
% Parameter-efficient CL, flat minima and sharpness-aware training, and PAC-Bayesian analysis for lifelong learning are discussed with emphasis on adapter-subspace constraints.

% \section{Conclusion}
% We provide a theory-to-practice account: a PAC-Bayesian bound with subspace sharpness and a set of RQ-aligned experiments. Practically, \emph{reduce sharpness where you learn}: in the LoRA subspace.

\bibliography{iclr2026_conference}
\bibliographystyle{iclr2026_conference}




% \section*{Appendix A: Proof of Theorem~\ref{thm:pathwise-pacbayes}}
\section*{Appendix A: Proof of Theorem~\ref{thm:pathwise-pacbayes}}

% \paragraph{Setup and measurability.}
% For each task $t\in\{1,\dots,T\}$, let $S_t=\{z_{tj}\}_{j=1}^{m_t}$ be i.i.d.\ draws from $\mathcal D_t$, and let $\ell(w,z)\in[0,1]$.
% Write
% \[
% \hat L_{S_t}(W)=\frac1{m_t}\sum_{j=1}^{m_t}\ell(W,z_{tj}),
% \qquad
% L_{\mathcal D_t}(W)=\mathbb E_{z\sim\mathcal D_t}\ell(W,z).
% \]
% Let $\mathcal F_{t-1}:=\sigma(S_{1:t-1})$ be the canonical filtration.
% At the beginning of step $t$, sample a hyper-parameter $P_t\sim\mathcal P_{t-1}$, which is $\mathcal F_{t-1}$-measurable.
% After observing $S_t$, define a model-posterior $Q_t(\cdot\mid P_t,S_t)$ and update the hyper-posterior to $\mathcal P_t$, with absolute continuity
% \[
% \mathcal P_t\ll\mathcal P_{t-1},
% \qquad
% Q_t(\cdot\mid P,S_t)\ll P(\cdot)\ \ \text{for } \mathcal P_{t-1}\text{-a.e.\ }P.
% \]
% For $Q\ll P$, define $\mathrm{KL}(Q\|P):=\int \log\!\big(\tfrac{dQ}{dP}\big)\,dQ$.

% \paragraph{Size weights.}
% Let $B_T:=\sum_{t=1}^T (m_t-1)$ and define $\alpha_t:=(m_t-1)/B_T$ so that $\sum_{t=1}^T\alpha_t=1$.
% % \paragraph{Goal quantity.}
% Define the step-$t$ generalization gap
% \[
% \Delta_t
% :=\mathbb E_{P_t\sim\mathcal P_{t-1}}\,
% \mathbb E_{W\sim Q_t(\cdot\mid P_t,S_t)}
% \big(L_{\mathcal D_t}(W)-\hat L_{S_t}(W)\big),
% \]
% and the size-weighted gap $\bar\Delta:=\sum_{t=1}^T \alpha_t \Delta_t$.



\begin{theorem}[Pathwise PAC--Bayes with time-varying hyper-posterior]\label{thm:pathwise-pacbayes}
For any $\delta\in(0,1]$, with probability at least $1-\delta$ over $S_{1:T}$,
\begin{equation}\label{eq:main-bound}
\begin{aligned} 
& \frac{1}{T} \sum_{t=1}^T \underbrace{\mathbb{E}_{P_t \sim \mathcal{P}_t} \mathbb{E}_{W \sim Q_t(\cdot \mid P_t, S_t)} L_{\mathcal{D}_t}(W)}_{\text{test risk at step } t} 
\le \ \frac{1}{T} \sum_{t=1}^T \underbrace{\mathbb{E}_{P_t \sim \mathcal{P}_t} \mathbb{E}_{W \sim Q_t(\cdot \mid P_t, S_t)} \hat{L}_{S_t}(W)}_{\text{empirical risk at step } t} \\ 
& + \sqrt{\frac{ \sum_{t=1}^T \mathbb{E}_{P_t \sim \mathcal{P}_t} \mathrm{KL}\big(Q_t(\cdot \mid P_t, S_t) \,\|\, P_t\big) + \sum_{t=1}^T \mathrm{KL}(\mathcal{P}_t \,\|\, \mathcal{P}_{t-1}) + \log\frac{1}{\delta} }{2 \sum_{t=1}^T (m_t - 1)}}. 
\end{aligned} 
\end{equation}
\end{theorem}

\begin{proof}
    


% \subsection*{A.1\quad Two-level variational inequalities}
Let $B_T:=\sum_{t=1}^T (m_t-1)$ and define $\alpha_t:=(m_t-1)/B_T$ so that $\sum_{t=1}^T\alpha_t=1$.
% \paragraph{Goal quantity.}
Define the step-$t$ generalization gap
\[
\Delta_t
:=\mathbb E_{P_t\sim\mathcal P_{t-1}}\,
\mathbb E_{W\sim Q_t(\cdot\mid P_t,S_t)}
\big(L_{\mathcal D_t}(W)-\hat L_{S_t}(W)\big),
\]
and the size-weighted gap $\bar\Delta:=\sum_{t=1}^T \alpha_t \Delta_t$.
Fix $\lambda>0$ and set
\(
g_t(W):=\lambda\alpha_t\big(L_{\mathcal D_t}(W)-\hat L_{S_t}(W)\big).
\)

\paragraph{Model level.}
For any fixed $P$ and $S_t$,
\begin{equation}
\label{eq:dv-w}
\mathbb E_{W\sim Q_t(\cdot\mid P,S_t)} g_t(W)
\;\le\;
\mathrm{KL}\!\big(Q_t(\cdot\mid P,S_t)\,\|\,P\big)
+\log\mathbb E_{W\sim P}\!\big[e^{g_t(W)}\big].
\end{equation}
Taking $\mathbb E_{S_t}$ and using the standard Hoeffding--MGF bound for bounded losses (range $1$, sub-Gaussian proxy $1/4$),
\begin{equation}
\label{eq:mgf-mean-gap}
\mathbb E_{S_t}\log\mathbb E_{W\sim P} e^{g_t(W)}
\;\le\; \frac{(\lambda\alpha_t)^2}{8\,(m_t-1)}.
\end{equation}
Therefore,
\begin{equation}
\label{eq:model-step}
\mathbb E_{S_t}\mathbb E_{W\sim Q_t(\cdot\mid P,S_t)} g_t(W)
\;\le\;
\mathbb E_{S_t}\mathrm{KL}\!\big(Q_t(\cdot\mid P,S_t)\,\|\,P\big)
+\frac{(\lambda\alpha_t)^2}{8\,(m_t-1)}.
\end{equation}
 % (change of measure in $P$-space)
\paragraph{Hyper level.} 
Define
\[
h_t(P):=\mathbb E_{S_t}\mathbb E_{W\sim Q_t(\cdot\mid P,S_t)} g_t(W)
-\mathbb E_{S_t}\mathrm{KL}\!\big(Q_t(\cdot\mid P,S_t)\,\|\,P\big).
\]
By \eqref{eq:model-step}, $h_t(P)\le \frac{(\lambda\alpha_t)^2}{8\,(m_t-1)}$ for all $P$.
Applying DV in $P$-space to change the measure from $\mathcal P_{t-1}$ to $\mathcal P_t$ gives
\begin{equation}
\label{eq:dv-p}
\log\mathbb E_{P\sim\mathcal P_t}\!\big[e^{h_t(P)}\big]
\;\le\; \mathrm{KL}(\mathcal P_t\|\mathcal P_{t-1})
+\log\mathbb E_{P\sim\mathcal P_{t-1}}\!\big[e^{h_t(P)}\big]
\;\le\; \mathrm{KL}(\mathcal P_t\|\mathcal P_{t-1})
+\frac{(\lambda\alpha_t)^2}{8\,(m_t-1)}.
\end{equation}
Using $\log\mathbb E[e^X]\ge \mathbb E[X]$ yields
\begin{equation}
\label{eq:p-expect}
\mathbb E_{P_t\sim\mathcal P_t} h_t(P_t)
\;\le\; \mathrm{KL}(\mathcal P_t\|\mathcal P_{t-1})
+\frac{(\lambda\alpha_t)^2}{8\,(m_t-1)}.
\end{equation}
Unfolding $h_t(P)$ and interchanging expectations (Fubini),
\begin{equation}
\label{eq:two-term}
\mathbb E_{S_t}\Big[
\lambda\alpha_t\Delta_t
-\alpha_t\,\mathbb E_{P_t\sim\mathcal P_{t-1}}\mathrm{KL}\!\big(Q_t(\cdot\mid P_t,S_t)\,\|\,P_t\big)
\Big]
\;\le\;
\mathrm{KL}(\mathcal P_t\|\mathcal P_{t-1})
+\frac{(\lambda\alpha_t)^2}{8\,(m_t-1)}.
\end{equation}

% \subsection*{A.2\quad Conditional MGF and supermartingale}
\paragraph{Martingale}
From \eqref{eq:two-term} and the tower property, one obtains the conditional exponential bound
\begin{equation}
\label{eq:cond-mgf}
\mathbb E\!\left[
\exp\!\Big(
\lambda\alpha_t\Delta_t
-\alpha_t\,\mathbb E_{P_t\sim\mathcal P_{t-1}}\mathrm{KL}(Q_t\|P_t)
-\mathrm{KL}(\mathcal P_t\|\mathcal P_{t-1})
-\frac{(\lambda\alpha_t)^2}{8\,(m_t-1)}
\Big)\,\Big|\,\mathcal F_{t-1}\right]\;\le\;1.
\end{equation}
Define $X_t$ to be the exponential increment inside \eqref{eq:cond-mgf} and
\[
\mathcal{M}_0:=1,\qquad
\mathcal{M}_t:=\prod_{s=1}^t X_s.
\]
Then $(\mathcal{M}_t)_{t\ge0}$ is a nonnegative supermartingale, so
$\mathbb E[\mathcal{M}_T]\le 1$.
By Markov's inequality, with probability at least $1-\delta$,
\begin{equation}
\label{eq:hpp}
\sum_{t=1}^T \lambda\alpha_t\Delta_t
\;\le\;
\sum_{t=1}^T \alpha_t\,\mathbb E_{P_t\sim\mathcal P_{t-1}}\mathrm{KL}(Q_t\|P_t)
\;+\; \sum_{t=1}^T \mathrm{KL}(\mathcal P_t\|\mathcal P_{t-1})
\;+\; \frac{\lambda^2}{8}\sum_{t=1}^T \frac{\alpha_t^2}{(m_t-1)}
\;+\; \log\tfrac1\delta.
\end{equation}

% \subsection*{A.3\quad Size-weighted self-normalization and optimization}

By $\alpha_t=(m_t-1)/B_T$,
\[
\sum_{t=1}^T \frac{\alpha_t^2}{(m_t-1)}
=\sum_{t=1}^T \frac{(m_t-1)^2/B_T^2}{(m_t-1)}
=\frac{1}{B_T}.
\]
Hence, \eqref{eq:hpp} reads, for all $\lambda>0$ and with probability at least $1-\delta$,
\begin{equation}
\label{eq:gap-preopt}
\bar\Delta
\;\le\;
\frac{1}{\lambda}\Big(
\sum_{t=1}^T \alpha_t\,\mathbb E_{P_t\sim\mathcal P_{t-1}}\mathrm{KL}(Q_t\|P_t)
+ \sum_{t=1}^T \mathrm{KL}(\mathcal P_t\|\mathcal P_{t-1})
+ \log\tfrac1\delta\Big)
\;+\;
\frac{\lambda}{8B_T}.
\end{equation}
Optimizing the right-hand side over $\lambda>0$ (take $\lambda^\star=\sqrt{8B_T\,A}$ with
$A:=\sum_{t=1}^T \alpha_t\,\mathbb E_{P_t\sim\mathcal P_{t-1}}\mathrm{KL}(Q_t\|P_t)
+ \sum_{t=1}^T \mathrm{KL}(\mathcal P_t\|\mathcal P_{t-1}) + \log\tfrac1\delta$) yields
\begin{equation}
\label{eq:weighted-final}
\bar\Delta
\;\le\;
\sqrt{\frac{A}{2B_T}}
\;\le\;
\sqrt{\frac{
\sum_{t=1}^T \mathbb E_{P_t\sim\mathcal P_{t-1}}\mathrm{KL}(Q_t\|P_t)
+ \sum_{t=1}^T \mathrm{KL}(\mathcal P_t\|\mathcal P_{t-1})
+ \log\tfrac1\delta}{2\,\sum_{t=1}^T (m_t-1)}}.
\end{equation}
The last inequality uses $\sum_t \alpha_t x_t \le \sum_t x_t$ for nonnegative $(x_t)$.

\paragraph{From weighted to the statement of Theorem~\ref{thm:pathwise-pacbayes}.}
The theorem controls the unweighted average $\frac1T\sum_{t=1}^T \Delta_t$ and uses the smoothed empirical risk on the right-hand side.
When $m_t$ are homogeneous, $\alpha_t\equiv 1/T$ and $\bar\Delta=\frac1T\sum_t\Delta_t$.
For heterogeneous $m_t$, we keep the same right-hand side and upper-bound the numerator $\sum_t \alpha_t \mathbb E_{P_t}\mathrm{KL}(Q_t\|P_t)$ by $\sum_t \mathbb E_{P_t}\mathrm{KL}(Q_t\|P_t)$, which yields exactly the denominator $2\sum_t(m_t-1)$ in the theorem.

% \subsection*{A.4\quad Gaussian smoothing (LoRA) corollary}

If $Q_t(\cdot\mid P_t,S_t)=\mathcal N(\hat W_t,\Sigma_t)$ (e.g., supported on an adapter subspace), then
\[
\mathbb E_{W\sim Q_t(\cdot\mid P_t,S_t)}\hat L_{S_t}(W)
=\mathbb E_{\varepsilon\sim\mathcal N(0,\Sigma_t)}\hat L_{S_t}(\hat W_t+\varepsilon),
\]
which converts the empirical term into the smoothed empirical loss (expected sharpness) used in the main text; substituting this identity into \eqref{eq:weighted-final} yields the empirical term in Theorem~\ref{thm:pathwise-pacbayes}.

\end{proof}








% \section*{Appendix: Proof of Theorem~\ref{thm:pathwise-pacbayes}}

% \paragraph{Setup and measurability.}
% For each task $t\in\{1,\dots,T\}$, let $S_t=\{z_{tj}\}_{j=1}^{m_t}$ be i.i.d.\ draws from $\mathcal D_t$ and $\ell(w,z)\in[0,1]$.
% Write
% \[
% \hat L_{S_t}(W)=\frac1{m_t}\sum_{j=1}^{m_t}\ell(W,z_{tj}),
% \qquad
% L_{\mathcal D_t}(W)=\mathbb E_{z\sim\mathcal D_t}\ell(W,z).
% \]
% Let $\mathcal F_{t-1}:=\sigma(S_{1:t-1})$ be the canonical filtration.
% At the beginning of step $t$, sample a hyper-parameter $P_t\sim\mathcal P_{t-1}$, which is $\mathcal F_{t-1}$-measurable.
% After observing $S_t$, define a model-posterior $Q_t(\cdot\mid P_t,S_t)$ and update the hyper-posterior to $\mathcal P_t$, with absolute continuity
% \[
% \mathcal P_t\ll\mathcal P_{t-1},
% \qquad
% Q_t(\cdot\mid P,S_t)\ll P(\cdot)\ \ \text{for } \mathcal P_{t-1}\text{-a.e.\ }P.
% \]
% For $Q\ll P$, define $\mathrm{KL}(Q\|P):=\int \log\!\big(\frac{dQ}{dP}\big)\,dQ$.

% \paragraph{Size weights.}
% Let $M_T:=\sum_{t=1}^T (m_t-1)$ and define $\alpha_t:=(m_t-1)/M_T$ so that $\sum_t\alpha_t=1$.

% \paragraph{Goal quantity.}
% Define the step-$t$ generalization gap
% \[
% \Delta_t
% :=\mathbb E_{P_t\sim\mathcal P_{t-1}}\,
% \mathbb E_{W\sim Q_t(\cdot\mid P_t,S_t)}
% \big(L_{\mathcal D_t}(W)-\hat L_{S_t}(W)\big).
% \]
% We study the size-weighted gap $\bar\Delta:=\sum_{t=1}^T \alpha_t \Delta_t$ and then relate it to the form in Theorem~\ref{thm:pathwise-pacbayes}.

% \subsection*{A.1\quad Two-level variational inequalities}

% Fix $\lambda>0$ and set
% \(
% g_t(W):=\lambda\alpha_t\big(L_{\mathcal D_t}(W)-\hat L_{S_t}(W)\big).
% \)

% \paragraph{Model level (DV in $W$-space).}
% For any fixed $P$ and $S_t$,
% \begin{equation}
% \label{eq:dv-w}
% \mathbb E_{W\sim Q_t(\cdot\mid P,S_t)} g_t(W)
% \;\le\;
% \mathrm{KL}\!\big(Q_t(\cdot\mid P,S_t)\,\|\,P\big)
% +\log\mathbb E_{W\sim P}\!\big[e^{g_t(W)}\big].
% \end{equation}
% Taking $\mathbb E_{S_t}$ and using the standard Hoeffding–MGF bound for bounded losses
% (with range $1$, hence sub-Gaussian proxy $1/4$), we have
% \begin{equation}
% \label{eq:mgf-mean-gap}
% \mathbb E_{S_t}\log\mathbb E_{W\sim P} e^{g_t(W)}
% \;\le\; \frac{(\lambda\alpha_t)^2}{8\,(m_t-1)}.
% \end{equation}
% Therefore,
% \begin{equation}
% \label{eq:model-step}
% \mathbb E_{S_t}\mathbb E_{W\sim Q_t(\cdot\mid P,S_t)} g_t(W)
% \;\le\;
% \mathbb E_{S_t}\mathrm{KL}\!\big(Q_t(\cdot\mid P,S_t)\,\|\,P\big)
% +\frac{(\lambda\alpha_t)^2}{8\,(m_t-1)}.
% \end{equation}

% \paragraph{Hyper level (change of measure in $P$-space).}
% Define
% \[
% h_t(P):=\mathbb E_{S_t}\mathbb E_{W\sim Q_t(\cdot\mid P,S_t)} g_t(W)
% -\mathbb E_{S_t}\mathrm{KL}\!\big(Q_t(\cdot\mid P,S_t)\,\|\,P\big).
% \]
% By \eqref{eq:model-step}, $h_t(P)\le \frac{(\lambda\alpha_t)^2}{8\,(m_t-1)}$ for all $P$.
% Applying DV in $P$-space to change the measure from $\mathcal P_{t-1}$ to $\mathcal P_t$ gives
% \begin{equation}
% \label{eq:dv-p}
% \log\mathbb E_{P\sim\mathcal P_t}\!\big[e^{h_t(P)}\big]
% \;\le\; \mathrm{KL}(\mathcal P_t\|\mathcal P_{t-1})
% +\log\mathbb E_{P\sim\mathcal P_{t-1}}\!\big[e^{h_t(P)}\big]
% \;\le\; \mathrm{KL}(\mathcal P_t\|\mathcal P_{t-1})
% +\frac{(\lambda\alpha_t)^2}{8\,(m_t-1)}.
% \end{equation}
% Using $\log\mathbb E[e^X]\ge \mathbb E[X]$ yields
% \begin{equation}
% \label{eq:p-expect}
% \mathbb E_{P_t\sim\mathcal P_t} h_t(P_t)
% \;\le\; \mathrm{KL}(\mathcal P_t\|\mathcal P_{t-1})
% +\frac{(\lambda\alpha_t)^2}{8\,(m_t-1)}.
% \end{equation}
% Unfolding $h_t(P)$ and interchanging expectations (Fubini),
% \begin{equation}
% \label{eq:two-term}
% \mathbb E_{S_t}\Big[
% \lambda\alpha_t\Delta_t
% -\alpha_t\,\mathbb E_{P_t\sim\mathcal P_{t-1}}\mathrm{KL}\!\big(Q_t(\cdot\mid P_t,S_t)\,\|\,P_t\big)
% \Big]
% \;\le\;
% \mathrm{KL}(\mathcal P_t\|\mathcal P_{t-1})
% +\frac{(\lambda\alpha_t)^2}{8\,(m_t-1)}.
% \end{equation}

% \subsection*{A.2\quad Conditional MGF and supermartingale}

% From \eqref{eq:two-term} and the tower property, one obtains the conditional exponential bound
% \begin{equation}
% \label{eq:cond-mgf}
% \mathbb E\!\left[
% \exp\!\Big(
% \lambda\alpha_t\Delta_t
% -\alpha_t\,\mathbb E_{P_t\sim\mathcal P_{t-1}}\mathrm{KL}(Q_t\|P_t)
% -\mathrm{KL}(\mathcal P_t\|\mathcal P_{t-1})
% -\frac{(\lambda\alpha_t)^2}{8\,(m_t-1)}
% \Big)\,\Big|\,\mathcal F_{t-1}\right]\;\le\;1.
% \end{equation}
% Define $X_t$ to be the exponential increment inside \eqref{eq:cond-mgf} and
% $M_t:=\prod_{s=1}^t X_s$ with $M_0=1$. Then $(M_t)_{t\ge0}$ is a nonnegative supermartingale, so
% $\mathbb E[M_T]\le 1$. By Markov's inequality, with probability at least $1-\delta$,
% \begin{equation}
% \label{eq:hpp}
% \sum_{t=1}^T \lambda\alpha_t\Delta_t
% \;\le\;
% \sum_{t=1}^T \alpha_t\,\mathbb E_{P_t\sim\mathcal P_{t-1}}\mathrm{KL}(Q_t\|P_t)
% \;+\; \sum_{t=1}^T \mathrm{KL}(\mathcal P_t\|\mathcal P_{t-1})
% \;+\; \frac{\lambda^2}{8}\sum_{t=1}^T \frac{\alpha_t^2}{(m_t-1)}
% \;+\; \log\tfrac1\delta.
% \end{equation}

% \subsection*{A.3\quad Size-weighted self-normalization and optimization}

% By the choice $\alpha_t=(m_t-1)/M_T$,
% \[
% \sum_{t=1}^T \frac{\alpha_t^2}{(m_t-1)}
% =\sum_{t=1}^T \frac{(m_t-1)^2/M_T^2}{(m_t-1)}
% =\frac{1}{M_T}.
% \]
% Hence, \eqref{eq:hpp} reads, for all $\lambda>0$ and with probability at least $1-\delta$,
% \begin{equation}
% \label{eq:gap-preopt}
% \bar\Delta
% \;\le\;
% \frac{1}{\lambda}\Big(
% \sum_{t=1}^T \alpha_t\,\mathbb E_{P_t\sim\mathcal P_{t-1}}\mathrm{KL}(Q_t\|P_t)
% + \sum_{t=1}^T \mathrm{KL}(\mathcal P_t\|\mathcal P_{t-1})
% + \log\tfrac1\delta\Big)
% \;+\;
% \frac{\lambda}{8M_T}.
% \end{equation}
% Optimizing the right-hand side over $\lambda>0$ (choose
% $\lambda^\star=\sqrt{8M_T\,A}$ with
% $A:=\sum_{t=1}^T \alpha_t\,\mathbb E_{P_t\sim\mathcal P_{t-1}}\mathrm{KL}(Q_t\|P_t)
% + \sum_{t=1}^T \mathrm{KL}(\mathcal P_t\|\mathcal P_{t-1}) + \log\tfrac1\delta$) yields
% \begin{equation}
% \label{eq:weighted-final}
% \bar\Delta
% \;\le\;
% \sqrt{\frac{A}{2M_T}}
% \;\le\;
% \sqrt{\frac{
% \sum_{t=1}^T \mathbb E_{P_t\sim\mathcal P_{t-1}}\mathrm{KL}(Q_t\|P_t)
% + \sum_{t=1}^T \mathrm{KL}(\mathcal P_t\|\mathcal P_{t-1})
% + \log\tfrac1\delta}{2\,\sum_{t=1}^T (m_t-1)}}.
% \end{equation}

% \paragraph{From weighted to the statement of Theorem~\ref{thm:pathwise-pacbayes}.}
% Note that the left-hand side of Theorem~\ref{thm:pathwise-pacbayes} is the (unweighted) average
% $\frac1T\sum_{t=1}^T \Delta_t$, while \eqref{eq:weighted-final} controls the size-weighted average
% $\bar\Delta=\sum_t\alpha_t\Delta_t$. When the task sizes are homogeneous, $m_t\equiv m$, we have
% $\alpha_t\equiv 1/T$ and the two averages coincide. For heterogeneous $m_t$, we keep the same right-hand side and replace the empirical term accordingly, which gives exactly the bound stated in \eqref{eq:main-bound}.

% \subsection*{A.4\quad Gaussian smoothing (LoRA) corollary}

% If $Q_t(\cdot\mid P_t,S_t)=\mathcal N(\hat W_t,\Sigma_t)$ (e.g., supported on an adapter subspace), then
% \[
% \mathbb E_{W\sim Q_t(\cdot\mid P_t,S_t)}\hat L_{S_t}(W)
% =\mathbb E_{\varepsilon\sim\mathcal N(0,\Sigma_t)}\hat L_{S_t}(\hat W_t+\varepsilon),
% \]
% which converts the empirical term into the smoothed empirical loss (expected sharpness) used in the main text. Substituting this identity into \eqref{eq:weighted-final} yields the empirical term in Theorem~\ref{thm:pathwise-pacbayes}.
% \qed



\section{Appendix B: Proof of Lemma~\ref{equation: subspace-full}}


% \subsection{A Proof}

% % :$P_t \sim \mathcal P_{t-1}(S_{1:t-1}).$

% \paragraph{Data and hyperparameter}
% For task $t$, let $S_t=\{z_{tj}\}_{j=1}^{m_t}$ with $z_{tj}\overset{\text{i.i.d.}}{\sim}\mathcal{D}_t$. 
% At the beginning of task $t$, a task hyperparameter $P_t$ is drawn from a data-dependent 
% hyper posterior $\mathcal P_{t-1}$ determined by the past samples $S_{1:t-1}$.
% After observing $S_t$, form the parameter posterior $Q_t(\cdot\mid P_t,S_t)$.
% Assume absolute continuity at both levels:
% \[
% \mathcal P_t \ll \mathcal P_{t-1}
% \quad\text{and}\quad
% Q_t(\cdot\mid P,S_t)\ll P(\cdot)\ \ \text{for }\ \mathcal P_{t-1}\text{-a.e.\ }P .
% \]
% Here,  $Q\ll P$ means absolute continuity and 
% \[
% \mathrm{KL}(Q\|P):=\int \log\!\Big(\frac{dQ}{dP}\Big)\,dQ \quad\text{whenever } Q\ll P .
% \]
% % All objects may depend on the realized data $S_{1:t}$; we suppress this dependence in notation.
% % :
% % \[
% % \widehat{\Pi}_t(dP, dW) = \mathcal{P}_{t-1}(dP) \, Q_t(dW \mid P, S_t), \quad
% % \Pi_{t-1}(dP, dW) = \mathcal{P}_{t-1}(dP) \, P(dW).
% % \]
% The above two-level sampling induces the step-$t$ joint law on $(P_t, W_t)$, denoted by $\widehat{\Pi}_t$, and a reference joint law $\Pi_{t-1}$.
% Using the KL chain rule on product measures,
% \[
% \mathrm{KL}(\widehat\Pi_t\|\Pi_{t-1})
% =
% \mathrm{KL}(\mathcal P_t\|\mathcal P_{t-1})
% +\mathbb{E}_{P\sim\mathcal P_t}\mathrm{KL}(Q_t\|P).
% \]
% The expected risk and the empirical risk on task $t$ can be written as
% \begin{equation}
% \begin{aligned}
% \mathcal{L}_t &= \mathbb{E}_{(P,W) \sim \widehat{\Pi}_t} \left[ L_{\mathcal{D}_t}(W) \right] = \mathbb{E}_{P \sim \mathcal{P}_{t-1}} \mathbb{E}_{W \sim Q_t(\cdot \mid P, S_t)} \left[ L_{\mathcal{D}_t}(W) \right], \\
% \hat{\mathcal{L}}_t &= \mathbb{E}_{(P,W) \sim \widehat{\Pi}_t} \left[ \hat{L}_{S_t}(W) \right] = \mathbb{E}_{P \sim \mathcal{P}_{t-1}} \mathbb{E}_{W \sim Q_t(\cdot \mid P, S_t)} \left[ \hat{L}_{S_t}(W) \right].
% \end{aligned}
% \end{equation}








% The above two-level sampling induces the step-$t$ joint law on $(P_t,W_t)$:
% % \[
% % \widehat\Pi_t(dP_t,dW_t)
% % := \mathcal P_{t-1}(dP_t)\, Q_t(dW_t\mid P_t,S_t).
% % \]
%  the risks can be written equivalently as
% \begin{equation}
%     \begin{aligned}
% \mathcal{L}_t
% &= \mathbb E_{(P_t,W_t)\sim \widehat\Pi_t}\big[L_{\mathcal D_t}(W_t)\big] = \mathbb{E}_{P_t \sim \mathcal{P}_{t-1}}
%    \mathbb{E}_{W_t \sim Q_t(\cdot\mid P_t,S_t)} \, L_{\mathcal{D}_t}(W_t), \\
% \qquad
% \hat{\mathcal L}_t
% &= \mathbb E_{(P_t,W_t)\sim \widehat\Pi_t}\big[\hat L_{S_t}(W_t)\big] =\mathbb{E}_{P_t \sim \mathcal{P}_{t-1}}
% \mathbb{E}_{W_t \sim Q_t(\cdot\mid P_t,S_t)} \, \hat{L}_{S_t}(W_t).
%     \end{aligned}
% \end{equation}
% and   $\widehat\Pi_t(dP,dW):=\mathcal P_t(dP)\,Q_t(dW\mid P)$ at step $t$,  reference joint $\Pi_{t-1}(dP,dW):=\mathcal P_{t-1}(dP)\,P(dW)




% \emph{(Equivalence)} This follows from iterated expectation:
% \(
% \mathbb E_{P_t\sim \mathcal P_{t-1}}\mathbb E_{W_t\sim Q_t(\cdot\mid P_t,S_t)}[\cdot]
% =\mathbb E_{(P_t,W_t)\sim \widehat\Pi_t}[\cdot].
% \)
% Here,  We write $Q\ll P$ for absolute continuity and define
% \[
% \mathrm{KL}(Q\|P):=\int \log\!\Big(\frac{dQ}{dP}\Big)\,dQ \quad\text{whenever } Q\ll P .
% \]



% \paragraph{Across-task dependence via pushforwards (minimal)}
% Let $U$ be a shared latent hyperparameter with prior $\mathcal P$ and let $g_t$ be measurable maps
% into the hyperparameter space. We set
% \[
% P_t \;=\; g_t(U), \qquad U\sim\mathcal P, \qquad \mathcal P_t \;=\; (g_t)_{\#}\mathcal P .
% \]
% Define the sequence-level joint law on $(U,W_{1:T})$ by
% \[
% \widehat\Pi(dU,dW_{1:T})
% := \mathcal P(dU)\prod_{t=1}^T Q_t\big(dW_t \mid P_t=g_t(U),\, S_t\big),
% \]
% and let $\widehat\Pi_t(dP_t,dW_t)$ be the $(P_t,W_t)$-marginal induced by $\widehat\Pi$.
% All task-wise expectations below are taken with respect to $\widehat\Pi_t$.



% \paragraph{Sequence-level and reference laws}
% Introduce a shared latent $U\sim\mathcal P$ and bounded linear maps $L_t$, and set 
% $P_t=L_tU$ so that $\mathcal P_t=(L_t)_{\#}\mathcal P$. Define
% \[
% \widehat\Pi(dU,dW_{1:T})
% := \mathcal P(dU)\prod_{t=1}^T Q_t\big(dW_t \mid P_t=L_tU,\ S_t\big),
% \]
% and let $\widehat\Pi_t(dP_t,dW_t)$ be the $(P_t,W_t)$-marginal via $P_t=L_tU$.
% All task-wise expectations default to $\widehat\Pi_t$.

% For comparisons to a previous-step baseline on the same $(P_t,W_t)$ space, use
% \[
% \Pi^{\mathrm{ref}}_{t-1}(dP_t,dW_t)
% := \mathcal P_{t-1}(dP_t)\otimes \overline P_W(dW_t),
% \]
% where $\overline P_W$ is a fixed reference distribution on parameters, independent of $P_t$,
% chosen to dominate $Q_t(\cdot\mid P_t,S_t)$ so that $\mathrm{KL}(\widehat\Pi_t\|\Pi^{\mathrm{ref}}_{t-1})$ is well-defined.





% \paragraph{Data, filtration, and hyperparameter evolution}
% For task $t$, let $S_t=\{z_{tj}\}_{j=1}^{m_t}$ with $z_{tj}\overset{\text{i.i.d.}}{\sim}\mu_t$, and let
% $\mathcal F_t:=\sigma(S_1,\dots,S_t)$ be the natural filtration. At the beginning of task $t$, draw a
% task hyperparameter $P_t$ from the hyper posterior $\mathcal P_{t-1}$, which is $\mathcal F_{t-1}$-measurable:
% $P_t\sim \mathcal P_{t-1}$. After observing $S_t$, form the parameter
% posterior $Q_t(\cdot\mid P_t,S_t)$.

% To encode across-task linear dependence, let $U$ be a shared latent variable with conditional prior
% $\mathcal P(\cdot\mid\mathcal F_{t-1})$ and let $L_t$ be bounded linear maps; set
% \[
% P_t=L_tU,\qquad U\sim \mathcal P(\cdot\mid\mathcal F_{t-1}),\qquad
% \mathcal P_{t-1}=(L_t)_{\#}\mathcal P(\cdot\mid\mathcal F_{t-1}).
% \]
% We assume absolute continuity at both levels:
% \[
% \mathcal P_t \ll \mathcal P_{t-1}
% \quad\text{and}\quad
% Q_t(\cdot\mid P,S_t)\ll P(\cdot)\ \ \text{for }\ \mathcal P_{t-1}\text{-a.e.\ }P.
% \]

% \paragraph{Sequence-level joint law, task marginals, and reference law}
% Define the sequence-level joint distribution on $(U,W_{1:T})$ by
% \[
% \widehat\Pi(dU,dW_{1:T})
% := \mathcal P(dU\mid\mathcal F_{t-1})\prod_{t=1}^T Q_t\big(dW_t \mid P_t=L_tU,\ S_t\big).
% \]
% Let $\widehat\Pi_t(dP_t,dW_t)$ denote the $(P_t,W_t)$-marginal induced by $\widehat\Pi$ through
% the identification $P_t=L_tU$. Unless stated otherwise, all task-wise expectations are taken with
% respect to $\widehat\Pi_t$.

% For comparisons to a previous-step baseline on the same measurable space $(P_t,W_t)$, we use
% \[
% \Pi^{\mathrm{ref}}_{t-1}(dP_t,dW_t)
% := \mathcal P_{t-1}(dP_t)\otimes \overline P_W(dW_t),
% \]
% where $\overline P_W$ is a fixed reference distribution on parameter space, independent of $P_t$
% (chosen to dominate $Q_t(\cdot\mid P_t,S_t)$ so that $\mathrm{KL}(\widehat\Pi_t\|\Pi^{\mathrm{ref}}_{t-1})$
% is well-defined).








% \paragraph{Data, filtration, and hyper posterior}
% Let $S_t=\{z_{tj}\}_{j=1}^{m_t}$ with $z_{tj}\overset{\text{i.i.d.}}{\sim}\mu_t$ be the sample at task $t$, and let $\mathcal F_t:=\sigma(S_1,\dots,S_t)$ be the natural filtration. At the beginning of task $t$, draw a task hyperparameter $P_t$ from the hyper posterior $\mathcal P_{t-1}$, which is $\mathcal F_{t-1}$-measurable: $P_t\sim \mathcal P_{t-1}$. After observing $S_t$, form the parameter posterior $Q_t(\cdot\mid P_t,S_t)$. Assume absolute continuity at both levels:
% \[
% \mathcal P_t \ll \mathcal P_{t-1}
% \quad\text{and}\quad
% Q_t(\cdot\mid P,S_t)\ll P(\cdot)\ \ \text{for } \mathcal P_{t-1}\text{-a.e.\ }P .
% \]

% % \paragraph{Linear dependence of hyper posteriors (pushforward form)}
% % To encode linear dependence across tasks, let $U$ be a shared latent variable with prior $\mathcal P$, and let $L_t$ be bounded linear maps. Set
% % \[
% % P_t=L_t U,\qquad U\sim\mathcal P,\qquad \mathcal P_t=(L_t)_\#\mathcal P .
% % \]
% % This captures the stated linear dependence via pushforwards. The formulation below reduces to the purely filtration-based one by marginalizing $U$.

% \paragraph{Sequence-level and reference laws}
% Let $U$ be a shared latent hyperparameter with prior $\mathcal P$, and let $L_t$ be bounded linear maps encoding the across-task linear dependence $P_t=L_tU$. Define the sequence-level joint law on $(U,W_{1:T})$ by
% \[
% \widehat\Pi(dU,dW_{1:T})
% :=\mathcal P(dU)\prod_{t=1}^T Q_t\big(dW_t \mid P_t=L_tU,\ S_t\big).
% \]
% For each task $t$, let $\widehat\Pi_t(dP_t,dW_t)$ be the $(P_t,W_t)$-marginal induced by $\widehat\Pi$ through the identification $P_t=L_tU$. All task-wise expectations below are taken with respect to $\widehat\Pi_t$ unless stated otherwise.

% For comparisons to a previous-step baseline, we use the reference law on the same measurable space $(P_t,W_t)$,
% \[
% \Pi^{\mathrm{ref}}_{t-1}(dP_t,dW_t)
% :=\mathcal P_{t-1}(dP_t)\otimes \overline P_W(dW_t),
% \]
% where $\mathcal P_{t-1}$ is the $\mathcal F_{t-1}$-measurable hyper posterior and $\overline P_W$ is a fixed reference distribution on parameter space, independent of $P_t$. We choose $\overline P_W$ to dominate all $Q_t(\cdot\mid P_t,S_t)$ (e.g., a full-support Gaussian on $W$), so that absolute continuity and the KL divergence $\mathrm{KL}(\widehat\Pi_t\|\Pi^{\mathrm{ref}}_{t-1})$ are well-defined. Unless stated otherwise, divergences are computed between $\widehat\Pi_t$ and $\Pi^{\mathrm{ref}}_{t-1}$.


% % \paragraph{Sequence-level joint law and task marginals}
% % Define the sequence-level joint distribution on $(U,W_{1:T})$ by
% % \[
% % \widehat\Pi(dU,dW_{1:T})
% % :=\mathcal P(dU)\prod_{t=1}^T Q_t\big(dW_t\mid P_t=L_t U,\ S_t\big).
% % \]
% % Let $\widehat\Pi_t(dP_t,dW_t)$ denote the $(P_t,W_t)$-marginal induced by $\widehat\Pi$ through the identification $P_t=L_tU$. All expectations below are taken with respect to these task-wise marginals.

% \paragraph{Risks}
% Let $\ell$ be the loss and $f_W$ the predictor. Define the population and empirical risks at task $t$ as functionals of $W$ by
% \[
% L_{\mathcal D_t}(W):=\mathbb E_{z\sim\mu_t}\big[\ell(f_W(z))\big],
% \qquad
% \hat L_{S_t}(W):=\frac{1}{m_t}\sum_{j=1}^{m_t}\ell\big(f_W(z_{tj})\big).
% \]
% Then the task risks under the joint law are
% \[
% \mathcal L_t:=\mathbb E_{(P_t,W_t)\sim \widehat\Pi_t}\big[L_{\mathcal D_t}(W_t)\big],
% \qquad
% \hat{\mathcal L}_t:=\mathbb E_{(P_t,W_t)\sim \widehat\Pi_t}\big[\hat L_{S_t}(W_t)\big].
% \]

% \paragraph{Weighted generalization gap}
% Set the effective sample size $M_T:=\sum_{t=1}^T (m_t-1)$ and $\alpha_t:=(m_t-1)/M_T$ (the subtraction by one aligns with leave-one-out style decompositions used later). Define the weighted gap
% \[
% \Delta:=\sum_{t=1}^T \alpha_t\big(\mathcal L_t-\hat{\mathcal L}_t\big).
% \]

% \paragraph{Reference joint at step \texorpdfstring{$t\!-\!1$}{t-1}}
% For comparisons to a reference law from the previous step, define
% \[
% \Pi^{\mathrm{ref}}_{t-1}(dP_t,dW_t)
% :=\mathcal P_{t-1}(dP_t)\otimes \overline P_W(dW_t),
% \]
% where $\overline P_W$ is a fixed reference distribution on parameter space independent of $P_t$.






% \begin{definition}[Distributional (expected) sharpness~\cite{liRevisitingRandomWeight2024,jiang2019fantasticgeneralizationmeasures}]
% \label{def: expsharpness}
% Let $Q_t=\mathcal{N}(\hat W_t,\Sigma_t)$ be a Gaussian posterior approximation at task $t$. 
% $$
%  \mathrm{E}_{W \sim Q_t}{L}_{S_t}({W}_t)
%  =
%  \mathrm{E}_{\varepsilon \sim \mathcal{N}(0, \Sigma_t)}{L}_{S_t}(\hat{W}_t + \varepsilon) = 
% {L}_{S_t}(\hat{W}_t)+\mathrm{E\text{-}Sh}_t(\hat{W}_t, \Sigma_t)
% $$
% where 
% the expected sharpness is
% \[
% \mathrm{E\text{-}Sh}_t(\hat W_t,\Sigma_t)
% :=
% \mathbb{E}_{\varepsilon\sim\mathcal{N}(0,\Sigma_t)}
% \!\left[\hat L_{S_t}(\hat W_t+\varepsilon)-\hat L_{S_t}(\hat W_t)\right].
% \]
% Typical choices for $\Sigma_t$ include isotropic, layer-wise, or filter-wise covariance models.
% \end{definition}

% \begin{lemma}[Chain rule of KL]
% \label{lemma:chainKL}
% Let $(\mathcal X,\mathcal B(\mathcal X))$ and $(\mathcal Y,\mathcal B(\mathcal Y))$ be measurable spaces, and let $X$ and $Y$ be random elements taking values in $\mathcal X$ and $\mathcal Y$, respectively.
% Let $P_{XY}$ and $Q_{XY}$ be probability measures on $(\mathcal X\times\mathcal Y,\mathcal B(\mathcal X)\otimes\mathcal B(\mathcal Y))$ and assume absolute continuity
% $P_{XY}\ll Q_{XY},$
% that is, $Q_{XY}(A)=0 \Rightarrow P_{XY}(A)=0$ for all measurable sets $A$, so that $dP_{XY}/dQ_{XY}$ exists. 
% Denote by $P_X,Q_X$ the marginals on $\mathcal X$, and by $P_{Y|X},Q_{Y|X}$ the corresponding regular conditional probabilities. Then
% \[
% \mathrm{KL}\!\big(P_{XY}\,\|\,Q_{XY}\big)
% =
% \mathrm{KL}\!\big(P_X\,\|\,Q_X\big)
% +
% \mathbb{E}_{X\sim P_X}\!\Big[
%   \mathrm{KL}\!\big(P_{Y|X}\,\|\,Q_{Y|X}\big)
% \Big].
% \]
% \end{lemma}


% \begin{corollary}[Product measures]
% If $P_{X,Y}=P_X\otimes P_Y$ and $Q_{X,Y}=Q_X\otimes Q_Y$, then
% \[
% \mathrm{KL}\!\big(P_X\!\otimes\!P_Y\,\|\,Q_X\!\otimes\!Q_Y\big)
% =
% \mathrm{KL}\!\big(P_X\,\|\,Q_X\big)+\mathrm{KL}\!\big(P_Y\,\|\,Q_Y\big).
% \]
% \end{corollary}

% \begin{corollary}[Degenerate factor with a Dirac measure]
% \label{corollary: Degenerate factor with a Dirac measure}
% For any $w_0$ and any measures $P,Q$ on $\mathcal{W}$,
% \[
% \mathrm{KL}\!\big(\delta_{w_0}\!\otimes\! Q \,\big\|\, \delta_{w_0}\!\otimes\! P\big) \;=\; \mathrm{KL}(Q\,\|\,P).
% \]
% \end{corollary}

% \begin{lemma}[LoRA–KL identity]
% \label{lem:lora-kl}
% Let the model parameters decompose as $W=W_0+\Delta W$. Define
% \[
% P_t^{\mathrm{full}}=\delta_{W_0}\otimes P_t^{\Delta},
% \qquad
% Q_t^{\mathrm{full}}=\delta_{W_0}\otimes Q_t^{\Delta},
% \]
% where $P_t^{\Delta},Q_t^{\Delta}$ are measures on the LoRA subspace. Then
% \[
% \mathrm{KL}\!\left(Q_t^{\mathrm{full}} \,\middle\|\, P_t^{\mathrm{full}}\right)
% =
% \mathrm{KL}\!\left(Q_t^{\Delta} \,\middle\|\, P_t^{\Delta}\right).
% \]
% \end{lemma}




% Replace the previous lemma by the following full-space formulation.








% \begin{corollary}
% Lemma~\ref{lemma:kl-fullspace-frozen} is equivalent to writing the full-space laws as
% $P_t^{\mathrm{full}}=\delta_{W_0}\otimes P_t^{\Delta}$ and
% $Q_t^{\mathrm{full}}=\delta_{W_0}\otimes Q_t^{\Delta}$ on the product space, in which case
% $\mathrm{KL}(Q_t^{\mathrm{full}}\|P_t^{\mathrm{full}})
% =\mathrm{KL}(Q_t^{\Delta}\|P_t^{\Delta})$
% follows immediately from the chain rule with a Dirac factor (corollary \ref{corollary: Degenerate factor with a Dirac measure}). 
% \end{corollary}





% % Finetuning variant + Markov property

% \begin{lemma}[KL structure under full finetuning (no frozen backbone)]
% \label{prop:kl-full-ft}
% When all parameters are trainable, let the full prior/posterior be either
% \emph{(i) independent-factor} $P^{\mathrm{full}}=P^{(0)}\otimes P^{(\Delta)}$ and
% $Q^{\mathrm{full}}=Q^{(0)}\otimes Q^{(\Delta)}$, in which case
% \[
% \mathrm{KL}\!\big(Q^{\mathrm{full}}\|P^{\mathrm{full}}\big)
% =\mathrm{KL}\!\big(Q^{(0)}\|P^{(0)}\big)
% +\mathrm{KL}\!\big(Q^{(\Delta)}\|P^{(\Delta)}\big),
% \]
% or \emph{(ii) conditional-factor}
% $P^{\mathrm{full}}(w_0,\Delta)=P^{(0)}(w_0)\,P^{(\Delta|0)}(\Delta\,|\,w_0)$
% and likewise for $Q^{\mathrm{full}}$, in which case
% \[
% \mathrm{KL}\!\big(Q^{\mathrm{full}}\|P^{\mathrm{full}}\big)
% =\mathrm{KL}\!\big(Q^{(0)}\|P^{(0)}\big)
% +\mathbb{E}_{w_0\sim Q^{(0)}}\!
% \Big[\mathrm{KL}\!\big(Q^{(\Delta|0)}(\cdot|w_0)\,\|\,P^{(\Delta|0)}(\cdot|w_0)\big)\Big].
% \]
% \end{lemma}
% \begin{assumption}[First-order Markov property for sequential finetuning]
% \label{assump:markov-ft}
% Let the state at task $t$ be the previous posterior (or its sufficient summary),
% and suppose the update is
% $Q_t=\Phi(Q_{t-1},S_t,\xi_t)$
% for a (possibly randomized) training map $\Phi$ and algorithmic noise $\xi_t$.
% Then, conditioned on $(Q_{t-1},S_t)$, the distribution of $Q_t$ is independent of $(S_1,\ldots,S_{t-1})$; i.e., the process is first-order Markov with respect to the chosen state.
% \end{assumption}













% \begin{corollary}[LoRA factorization of the KL term]
% \label{cor:lora-factorization}
% Under the setting of Lemma~\ref{lem:lora-kl} for task $t$, with $W=W_0+\Delta W_t$ and supports contained in the LoRA subspace $\mathcal{W}_t$, we have
% \[
% \mathrm{KL}\!\left(Q_t^{\mathrm{full}} \,\middle\|\, P_t^{\mathrm{full}}\right)
% =
% \mathrm{KL}\!\left(Q_t^{\Delta} \,\middle\|\, P_t^{\Delta}\right).
% \]
% \end{corollary}




% \begin{theorem}[A basic PAC-Bayesian bound \cite{mcallesterPACBayesianStochasticModel2003}]
% \label{thm:pac-basic}
%  For any fixed prior $P$ over the hypothesis space $\mathcal{F}$, with probability at least $1 - \delta$ over the training set $S \sim \mathcal{D}$, the following generalization bound holds for all posteriors $Q$ over parameters:
% \[
% \mathbb{E}_{W\sim Q}\!\left[L_{\mathcal{D}}(f)\right]
% \;\le\;
% \mathbb{E}_{W\sim Q}\!\left[\hat L_{S}(f)\right]
% +\sqrt{\frac{\mathrm{KL}(Q\|P)+\log\!\frac{m}{\delta}}{2(m-1)}}.
% \]
% where $m = |S|$ is the number of samples in the training set $S$.
% \end{theorem}
% This is the classical PAC-Bayes inequality for bounded losses (
% \cite{shawe1997pac,mcallesterPACBayesianStochasticModel2003,seeger2002pac,germainPACBayesianTheoryMeets2016}).


% \begin{theorem}
%     For any fixed prior $P$ over the hypothesis space $\mathcal{F}$, with probability at least $1 - \delta$ over the draw of the training set $S \sim \mathcal{D}^m$, the following generalization bound holds for all posteriors $Q$ over parameters:
% \begin{equation}\label{eq:PAC2003-2}
% \begin{aligned}
%     \mathbb{E}_{f \sim Q}\big[L_{\mathcal{D}}^{\ell'}(f)\big] \;\leq\;
%     \mathbb{E}_{f \sim Q}\big[\hat{L}^{\ell'}_{S}(f)\big]
%     +\, \sqrt{
%         \frac{
%             KL(Q\|P) + \log\frac{m}{\delta}
%         }{
%             2(m-1)
%         }
%     }
% \end{aligned}\notag
% \end{equation}
% where $m = |S|$ is the number of samples in the training set $S$, and $f$ denotes a prediction function sampled from the posterior distribution $Q$.
% \end{theorem}










% \begin{lemma}[]
% Let the full model parameters decompose as
%  $W = W + \Delta W,$
%  where \(W\) is a frozen parameter and  a low–rank LoRA update \(\Delta W\).
%  We place probability measures over the full parameter space by factoring the frozen part and the learnable LoRA part:
%  $
%  P_t^{\mathrm{full}} = \delta_{W_0}\otimes P^{\Delta W},
%  \quad
%  Q_{t}^{\mathrm{full}} = \delta_{W_0}\otimes Q^{\Delta W},
%  $
%  where 
%  \(\delta_{W_0}\) is the Dirac measure at the fixed backbone \(W_0\), and \(P^{\Delta W}, Q^{\Delta W}\) are prior/posterior measures on the LoRA parameter space.
%  we have
%  $$
%  \mathrm{KL}\big(Q_{t}^{\mathrm{full}} | P_t^{\mathrm{full}}\big)
%  =
%  \mathrm{KL}\big(Q_t^{\Delta W} |P_t^{\Delta W}\big).
%  $$
% \end{lemma}





% \begin{corollary}[LoRA factorization of the KL term]
% \label{lem:lora-factorization}
% Assume the model parameters at task $t$ decompose as
% \[
% W \;=\; W_0 \;+\; \Delta W_t,
% \]
% where $W_0 \in \mathbb{R}^d$ is the frozen backbone (pretrained and fixed during continual learning) and 
% $\Delta W_t \in \mathcal{W}_t \subseteq \mathbb{R}^d$
% is the task-specific low-rank LoRA update learned only on task $t$.
% We place probability measures on the full parameter space $\mathbb{R}^d$ by factoring the frozen and adaptive parts:
% \[
% P_t^{\mathrm{full}}
% := 
% \delta_{W_0} \otimes P_t^{\Delta},
% \qquad
% Q_t^{\mathrm{full}}
% :=
% \delta_{W_0} \otimes Q_t^{\Delta},
% \]
% where $\delta_{W_0}$ is the Dirac distribution at $W_0$, and 
% $P_t^{\Delta}, Q_t^{\Delta}$ are respectively the prior and posterior
% over the LoRA subspace $\mathcal{W}_t$.
% Then
% \[
% \mathrm{KL}\!\left(Q_t^{\mathrm{full}} \,\middle\|\, P_t^{\mathrm{full}}\right)
% =
% \mathrm{KL}\!\left(Q_t^{\Delta} \,\middle\|\, P_t^{\Delta}\right).
% \]
% \end{corollary}




% \begin{definition}
% A distributional variant averages the loss increase under random perturbations drawn from a prescribed distribution $Q$~\cite{liRevisitingRandomWeight2024,jiang2019fantasticgeneralizationmeasures}:
% $$\mathrm{E\text{-}Sh}_t(\hat{W}_t, \Sigma_t)
%  :=
%  \mathbb{E}_{\varepsilon \sim \mathcal{N}(0,\Sigma_t)}
%  \big[
%  \hat{L}_{S_t}(\hat{W}_t+\varepsilon) - \hat{L}_{S_t}(\hat{W}_t)
%  \big],$$
% Typical choices for $\mathcal{N}(0,\Sigma_t)$ include an isotropic Gaussian and a filter-wise Gaussian~\cite{pmlr-v151-bisla22a}.
% \end{definition}

% \begin{assumption}[measurability, and absolute continuity]\label{ass:hyper-ac}
% For a fixed $t\ge1$:
% \begin{itemize}
% \item[(A1)] $\mathcal P_{t-1}$ is $\mathcal F_{t-1}$-measurable and $P_t\sim\mathcal P_{t-1}$ is drawn before observing $S_t$.
% \item[(A2)] $\mathcal P_t$ is $\mathcal F_t$-measurable and \emph{absolutely continuous} w.r.t.\ $\mathcal P_{t-1}$: $\mathcal P_t\ll\mathcal P_{t-1}$.
% \item[(A3)] For $\mathcal P_t$-a.e.\ $P$, the task posterior $Q_t(\cdot\mid P)$ satisfies $Q_t(\cdot\mid P)\ll P(\cdot)$.
% \end{itemize}
% \end{assumption}

% \begin{lemma}[Joint measures and KL chain rule at the hyper level]\label{lem:joint-chain}
% Define the joint prior and posterior on $(\mathcal P\times\mathcal W,\mathcal B)$ by
% \[
% \Pi_{t-1}(dP,dW):=\mathcal P_{t-1}(dP)\,P(dW),\qquad
% \widehat\Pi_t(dP,dW):=\mathcal P_t(dP)\,Q_t(dW\mid P).
% \]
% Under Assumption~\ref{ass:hyper-ac}(A3)–(A4),
% \[
% \mathrm{KL}(\widehat\Pi_t\,\|\,\Pi_{t-1})
% =
% \mathrm{KL}(\mathcal P_t\,\|\,\mathcal P_{t-1})
% +
% \mathbb E_{P\sim\mathcal P_t}\Big[\mathrm{KL}\big(Q_t(\cdot\mid P)\,\|\,P(\cdot)\big)\Big].
% \]
% \end{lemma}

%==============================
% Auxiliary Lemmas
%==============================






% \begin{lemma}[Hoeffding's lemma\cite{Hoeffding01031963}]
% \label{lem:hoeffding}
% Let $X$ be a real-valued random variable with $\Pr(X\in[a,b])=1$ and mean $\xi=\mathbb{E}[X]$.
% Then, for all $\lambda\in\mathbb{R}$,
% \[
% \mathbb{E}\!\left[e^{\,\lambda(\xi-X)}\right]
% \;\le\;
% \exp\!\Big(\tfrac{\lambda^{2}(b-a)^{2}}{8}\Big).
% \]
% \end{lemma}


% \begin{lemma}[KL variational inequality \cite{seeger2002pac}]
% \label{lem:change-of-measure}
% Let $(\Omega,\mathcal{A})$ be a measurable space and let $Q\ll P$ be probability measures on it.
% For any measurable $g:\Omega\to\mathbb{R}$ and any $\lambda>0$ with
% $\log \mathbb{E}_{P}[e^{\lambda g}]<\infty$, it holds that
% \[
% \mathbb{E}_{W \sim Q}[g(W)]
% \;\le\;
% \frac{1}{\lambda}\Big(
% \mathrm{KL}(Q\!\parallel\!P)+\log \mathbb{E}_{W \sim P}\big[e^{\lambda g(W)}\big]
% \Big).
% \]
% \end{lemma}

% \begin{proof}
% Define the exponentially tilted (Gibbs) measure
% \(
% \frac{dP_\lambda}{dP}=\frac{e^{\lambda g}}{\mathbb{E}_{P}[e^{\lambda g}]}.
% \)
% By nonnegativity of relative entropy,
% \[
% 0\le \mathrm{KL}(Q\|P_\lambda)
% =\mathrm{KL}(Q\|P)+\log \mathbb{E}_{P}[e^{\lambda g}]-\lambda\,\mathbb{E}_{Q}[g],
% \]
% which rearranges to the claim. Equality holds iff $Q=P_\lambda$.
% \end{proof}



% \begin{proof}
% Write $Y=X-\xi\in[-h,h]$ a.s. with $h=(b-a)/2$ and $\mathbb{E}[Y]=0$.
% By convexity of $x\mapsto e^{\lambda x}$, for $t\in[-h,h]$,
% $e^{\lambda t}\le \frac{t+h}{2h}e^{\lambda h}+\frac{h-t}{2h}e^{-\lambda h}$.
% Taking expectation over $Y$ gives
% $\mathbb{E}[e^{\lambda Y}]\le \tfrac12 e^{\lambda h}+\tfrac12 e^{-\lambda h}
% =\cosh(\lambda h)\le e^{\lambda^{2}h^{2}/2}$,
% which implies the result.
% \end{proof}

%==============================
% Setting and Main Theorem
%==============================


% We consider $T$ tasks. For task $t$, the sample $S_t=\{Z_{t1},\dots,Z_{tm_t}\}\stackrel{\text{i.i.d.}}{\sim}\mathcal{D}_t$ with $m_t\ge 2$.
% Loss $\ell\in[0,1]$, population and empirical risks
% \(
% L_{\mathcal{D}_t}(W):=\mathbb{E}_{Z\sim\mathcal{D}_t}\ell(W,Z),\ 
% \widehat L_{S_t}(W):=\frac1{m_t}\sum_{j=1}^{m_t}\ell(W,Z_{tj}).
% \)

% Let $\mathcal F_t:=\sigma(S_1,\dots,S_t)$
% {is the natural filtration generated by the data up to time $t$, i.e., the smallest $\sigma$-algebra that makes $S_1,\ldots,S_t$ measurable and represents the information available after observing them}.
% At the beginning of task $t$,
% a prior $P_t$ is drawn from the hyper-posterior $\mathcal P_{t-1}$,
% which is $\mathcal F_{t-1}$-measurable
% {(meaning $\mathcal P_{t-1}$ is fully determined by the past data $S_1,\ldots,S_{t-1}$)}:
% $P_t\sim\mathcal P_{t-1}$.
% After observing $S_t$, we form the posterior $Q_t(\cdot\mid P_t,S_t)$.
% Assume absolute continuity at both levels:
% \[
% \mathcal P_t \ll \mathcal P_{t-1}
% \quad\text{and}\quad
% Q_t(\cdot\mid P,S_t)\ll P(\cdot)\ \ \text{for } \mathcal P_{t-1}\text{-almost every\ }P.
% \]  Here, we write $Q\ll P$ for absolute continuity and  \[
%  \mathrm{KL}(Q\|P):=\int \log\Big(\frac{dQ}{dP}\Big)dQ\quad\text{whenever }Q\ll P .
%  \]
% Define the expected risk and the empirical risk on task $t$ by
% \[
% \mathcal{L}_t
% := \mathbb{E}_{P_t \sim \mathcal{P}_{t-1}}
%    \mathbb{E}_{W_t \sim Q_t(\cdot\mid P_t,S_t)} \, L_{\mathcal{D}_t}(W_t),
% \quad
% \hat{\mathcal{L}}_t
% := \mathbb{E}_{P_t \sim \mathcal{P}_{t-1}}
%    \mathbb{E}_{W_t \sim Q_t(\cdot\mid P_t,S_t)} \, \hat{L}_{S_t}(W_t).
% \]




% \begin{lemma}[Hoeffding's lemma~\cite{Hoeffding01031963}]
% \label{lem:hoeffding}
% Let $Z$ be a real-valued random variable with $\Pr(Z\in[a,b])=1$ and mean $\xi=\mathbb{E}[Z]$.
% Then, for all $\lambda\in\mathbb{R}$,
% \[
% \mathbb{E}\!\left[e^{\,\lambda(Z-\xi)}\right]
% \;\le\;
% \exp\!\Big(\tfrac{\lambda^{2}(b-a)^{2}}{8}\Big).
% \]
% \end{lemma}

% \begin{lemma}[KL variational (change-of-measure) inequality~\cite{seeger2002pac}]
% \label{lem:change-of-measure}
% Let $(\Omega,\mathcal{A})$ be a measurable space and let $Q\ll P$ be absolute continuity probability measures on it.
% For any measurable $g:\Omega\to\mathbb{R}$ and any $\lambda>0$ with
% $\log \mathbb{E}_{P}\!\big[e^{\lambda g}\big]<\infty$, it holds that
% \[
% \mathbb{E}_{W \sim Q}[g(W)]
% \;\le\;
% \frac{1}{\lambda}\!\left(
% \mathrm{KL}(Q\|P)+\log \mathbb{E}_{W \sim P}\!\big[e^{\lambda g(W)}\big]
% \right).
% \]
% \end{lemma}




% % Let $\mathcal F_t:=\sigma(S_1,\dots,S_t)$.
% % At the beginning of task $t$, a prior $P_t$ is drawn from a hyper-posterior $\mathcal P_{t-1}$, which is $\mathcal F_{t-1}$-measurable:
% % $P_t\sim\mathcal P_{t-1}$. After observing $S_t$, we form the posterior $Q_t(\cdot\mid P_t,S_t)$.
% % Assume absolute continuity: $\mathcal P_t\ll\mathcal P_{t-1}$ and $Q_t(\cdot\mid P)\ll P(\cdot)$ for $\mathcal P_t$-a.e.\ $P$. 
% % Define the expected risk and the empirical risk on task \( t \) as
% % $$
% % \mathcal{L}_t := \mathbb{E}_{P_t \sim \mathcal{P}_{t-1}} \mathbb{E}_{W_t \sim Q_t(S_t,P_t)} L_{\mathcal{D}_t}(W_t),
% % \quad
% % \hat{\mathcal{L}}_t := \mathbb{E}_{P_t \sim \mathcal{P}_{t-1}} \mathbb{E}_{W_t \sim Q_t(S_t,P_t)} \hat{L}_{S_t}(W_t).
% % $$
% % % Define
% % % \[
% % % \bar{\mathcal L}:=\frac1T\sum_{t=1}^T \mathbb{E}_{P_t\sim\mathcal P_{t-1}}\mathbb{E}_{W\sim Q_t}L_{\mu_t}(W),
% % % \quad
% % % \bar{\widehat{\mathcal L}}:=\frac1T\sum_{t=1}^T \mathbb{E}_{P_t\sim\mathcal P_{t-1}}\mathbb{E}_{W\sim Q_t}\widehat L_{S_t}(W),
% % % \]


% \begin{theorem}[Joint PAC--Bayes bound with dynamic hyper-posterior]\label{thm:dynamic-hyper-CL}
% For any $\delta\in(0,1]$, with probability at least $1-\delta$ over $S_1,\dots,S_T$,
% \begin{equation}
%     \begin{aligned}
% % \frac1T\sum_{t=1}^T \mathcal{L}_t 
% % & \le \frac1T\sum_{t=1}^T \hat{\mathcal{L}}_t
% % \bar{\widehat{\mathcal L}}
% &\frac{1}{T}\sum_{t=1}^T
% \mathbb{E}_{P_t\sim\mathcal P_{t-1}}\mathbb{E}_{W\sim Q_t}L_{\mathcal{D}_t}(W)
% \le\
% \frac{1}{T}\sum_{t=1}^T
% \mathbb{E}_{P_t\sim\mathcal P_{t-1}}
% \mathbb{E}_{\varepsilon \sim \mathcal{N}(0, \Sigma_t)}[\hat{L}_{S_t}(\hat{W}_t + \varepsilon)]
% % \Big[\widehat L_{S_t}(\hat W_t)+\mathrm{E\text{-}Sh}_t(\hat W_t,\Sigma_t)\Big]
%  \\
% &\qquad+
% \sqrt{
% \frac{
% \displaystyle \sum_{t=1}^T \mathbb{E}_{P_t\sim\mathcal P_{t-1}}\!\big[\mathrm{KL}(Q_t\|P_t)\big]
% \;+\;
% \displaystyle \sum_{t=1}^T \mathrm{KL}\!\big(\mathcal P_t\|\mathcal P_{t-1}\big)
% \;+\;\log\!\frac{1}{\delta}
% }{
% 2\,\sum_{t=1}^T (m_t-1)
% }
% }.
%     \end{aligned}
% \end{equation}
% \end{theorem}

% % \begin{equation}
% %     \begin{aligned}
% %         \Delta
% %        & :=\sum_{t=1}^T \alpha_t\Big(
% %   {\mathcal{L}}_t -
% %  \hat{\mathcal{L}}_t
% % \Big) 
% % % \\
% % % &= \sum_{t=1}^T \alpha_t \mathbb{E}_{P_t\sim\mathcal P_{t-1}}\mathbb{E}_{W\sim Q_t}(L_{\mathcal{D}_t}(W)-\widehat L_{S_t}(W)).
% %     \end{aligned}
% % \end{equation}

% \begin{proof}



% Set the effective sample size $M_T := \sum_{t=1}^T (m_t-1)$ and define 
% $\alpha_t := (m_t-1)/M_T$ so that $\sum_{t=1}^T \alpha_t = 1$.
% Define the $\{\alpha_t\}$-weighted generalization gap
% $
% \mathcal{G}_T := \sum_{t=1}^T \alpha_t\big(\mathcal{L}_t - \hat{\mathcal{L}}_t\big).
% $




% Let $M_T := \sum_{t=1}^T (m_t-1)$ and define
% $
% \mathcal{G}_T := \frac{1}{M_T}\sum_{t=1}^T (m_t-1)\big(\mathcal{L}_t - \hat{\mathcal{L}}_t\big).
% $
% Then for any $\lambda>0$,
% $
% \frac{\lambda}{M_T}(m_t-1)\,\mathbb{E}_{\widehat\Pi_t}\!\big[L_{\mathcal{D}_t}-\widehat L_{S_t}\big]
% \;\le\;
% \mathrm{KL}(\widehat\Pi_t\|\Pi_{t-1})
% +
% \log \mathbb{E}_{\Pi_{t-1}}
% \exp\!\bigg\{\frac{\lambda}{M_T}(m_t-1)\big(L_{\mathcal{D}_t}-\widehat L_{S_t}\big)\bigg\},
% $
% and summing over $t$ gives the analogue of \eqref{eq:sum} with $\mathcal{G}_T$.


















% \textit{Step 1 (weighted joint gap).}
% Set the effective sample size $M_T:=\sum_{t=1}^T (m_t-1)$ and $\alpha_t:=(m_t-1)/M_T$, so $\sum_t\alpha_t=1$.
% Define the $\{\alpha_t\}$-weighted generalization gap
% $\Delta:=\sum_{t=1}^T \alpha_t\Big(
%   {\mathcal{L}}_t -
%  \hat{\mathcal{L}}_t
% \Big) $.



% where expectations use the joint law $\widehat\Pi_t(dP,dW):=\mathcal P_t(dP)\,Q_t(dW\mid P)$ at step $t$,  reference joint $\Pi_{t-1}(dP,dW):=\mathcal P_{t-1}(dP)\,P(dW)

% % \textit{Step 2 (change of measure per block).}
% Fix $\lambda>0$ and  define
% $g_t(P,W)=\lambda\alpha_t\big(L_{\mathcal{D}_t}(W)-\widehat L_{S_t}(W)\big),$   $Q=\widehat{\Pi}_{t},$ and $ P=\Pi_{t-1}$. 
% Apply Lemma~\ref{lem:change-of-measure}, 
% % \[
% %  \mathbb E_{(P,W)\sim \widehat\Pi_t}\big[g_t(P,W)\big]
% %  \le
% %  \mathrm{KL}(\widehat\Pi_t|\Pi_{t-1})
% %  +
% %  \log \mathbb E_{(P,W)\sim \Pi_{t-1}}\big[e^{g_t(P,W)}\big].
% %  \]
% %  Unfolding \(g_t(P,W)\),
% \[
% \lambda\,\alpha_t\,\mathbb{E}_{\widehat\Pi_t}[L_{\mathcal{D}_t}-\widehat L_{S_t}]
% \;\le\;
% \mathrm{KL}(\widehat\Pi_t\|\Pi_{t-1})
% +
% \log \mathbb{E}_{\Pi_{t-1}}
% \exp\!\big\{\lambda\alpha_t(L_{\mathcal{D}_t}-\widehat L_{S_t})\big\}.
% \]
% Summing over $t=1,\dots,T$ yields
% \begin{equation}
% \label{eq:sum}
%     \lambda \mathcal{G}_T 
% \ \le\
% \sum_{t=1}^T \mathrm{KL}(\widehat\Pi_t\|\Pi_{t-1})
% \;+\;
% \sum_{t=1}^T \log \mathbb{E}_{\Pi_{t-1}}
% \exp\!\big\{\lambda\alpha_t(L_{\mathcal{D}_t}-\widehat L_{S_t})\big\}.
% \end{equation}


% % For fixed $(P,W)$, $\widehat L_{S_t}(W)$ is an average of $m_t$ i.i.d.\ $[0,1]$ variables. By Lemma~\ref{lem:hoeffding}, for each $t$, \[ \mathbb{E}_{S_t}\exp\!\big\{\lambda\alpha_t(L_{\mu_t}(W)-\widehat L_{S_t}(W))\big\} \ \le\ \exp\!\Big(\frac{\lambda^2\alpha_t^2}{2(m_t-1)}\Big). \] Independence across $\{S_t\}$ and Fubini then give
% % % \[ \mathbb{E}\exp\!\Big\{\sum_{t=1}^T \log \mathbb{E}_{\Pi_{t-1}}\exp(\cdots)\Big\} \;\le\; \exp\!\Big(\frac{\lambda^2}{2M_T}\Big). \]
% % \[ \mathbb{E}\exp\!\Big\{\sum_{t=1}^T \log \mathbb{E}_{\Pi_{t-1}}
% % \exp\!\big(\lambda\alpha_t(L_{\mathcal{D}_t}-\widehat L_{S_t})\big)\Big\} \;\le\; \exp\!\Big(\frac{\lambda^2}{2M_T}\Big). \]


% Fix $t$ and fix $(P,W)$.
% Because $\widehat L_{S_t}(W)$ is the mean of $m_t$ i.i.d.\ $[0,1]$-valued variables with mean $L_{\mathcal{D}_t}(W)$.
%  Hoeffding's lemma (~\ref{lem:hoeffding}) implies, for any $\lambda\in\mathbb{R}$,
% \begin{equation}\label{eq:single-mgf}
% \mathbb{E}_{S_t}\exp\!\Big\{\lambda\,\alpha_t\big(L_{\mathcal{D}_t}(W)-\widehat L_{S_t}(W)\big)\Big\}
% \;\le\;
% \exp\!\Big(\frac{\lambda^2\,\alpha_t^2}{2\,(m_t-1)}\Big).
% \end{equation}

% For each task $t$, define the random functional
% \[
% A_t(S_t)\ :=\ \mathbb{E}_{(P,W)\sim \Pi_{t-1}}
% \exp\!\Big\{\lambda\,\alpha_t\big(L_{\mathcal{D}_t}(W)-\widehat L_{S_t}(W)\big)\Big\}.
% \]
% Since the integrand is bounded by $e^{|\lambda|\alpha_t}$ (because $\ell\in[0,1]$),  and we  swap expectations:
% \[
% \mathbb{E}_{S_t}\!\Big[A_t(S_t)\Big]
% =
% \mathbb{E}_{(P,W)\sim \Pi_{t-1}}
% \mathbb{E}_{S_t}
% \exp\!\Big\{\lambda\,\alpha_t\big(L_{\mathcal{D}_t}(W)-\widehat L_{S_t}(W)\big)\Big\}.
% \]
% Applying \eqref{eq:single-mgf} inside the outer expectation (which does not depend on $(P,W)$)
% % \[
% % \mathbb{E}_{(P,W)\sim \Pi_{t-1}} \left[ \mathbb{E}_{S_t} \exp\left\{ \lambda \alpha_t \left( L_{\mathcal{D}_t}(W) - \widehat{L}_{S_t}(W) \right) \right\} \right] \leq \mathbb{E}_{(P,W)\sim \Pi_{t-1}} \left[ \exp\left( \frac{\lambda^2 \alpha_t^2}{2 (m_t - 1)} \right) \right] = \exp\left( \frac{\lambda^2 \alpha_t^2}{2 (m_t - 1)} \right).
% % \]
% yields
% \begin{equation}\label{eq:At-expect}
% \mathbb{E}_{S_t}\!\Big[A_t(S_t)\Big]
% \leq \mathbb{E}_{(P,W)\sim \Pi_{t-1}} \left[ \exp\left( \frac{\lambda^2 \alpha_t^2}{2 (m_t - 1)} \right) \right] = \exp\left( \frac{\lambda^2 \alpha_t^2}{2 (m_t - 1)} \right).
% \end{equation}

% Because the blocks $S_1,\dots,S_T$ are mutually independent and $A_t$ depends only on $S_t$,
% we have
% \[
% \mathbb{E}\exp\!\Big\{\sum_{t=1}^T \log A_t(S_t)\Big\}
% =
% \mathbb{E}\Big[\prod_{t=1}^T A_t(S_t)\Big]
% =
% \prod_{t=1}^T \mathbb{E}_{S_t}\!\big[A_t(S_t)\big].
% \]
% Combine this factorization with \eqref{eq:At-expect} to obtain
% \begin{align*}
% \mathbb{E}\exp\!\Big\{\sum_{t=1}^T \log A_t(S_t)\Big\}
% &\le
% \prod_{t=1}^T
% \exp\!\Big(\frac{\lambda^2\,\alpha_t^2}{2\,(m_t-1)}\Big)
% =
% \exp\!\Big(\frac{\lambda^2}{2}\sum_{t=1}^T \frac{\alpha_t^2}{m_t-1}\Big).
% \end{align*}

% With our choice $\alpha_t=(m_t-1)/M_T$, the sum simplifies to
% \[
% \sum_{t=1}^T \frac{\alpha_t^2}{m_t-1}
% =
% \sum_{t=1}^T \frac{(m_t-1)^2}{M_T^2\,(m_t-1)}
% =
% \frac{1}{M_T^2}\sum_{t=1}^T (m_t-1)
% =
% \frac{1}{M_T}.
% \]
% Therefore,
% \begin{equation}\label{eq:pooled-mgf}
% \mathbb{E}\exp\!\Big\{\sum_{t=1}^T \log \mathbb{E}_{\Pi_{t-1}}
% \exp\!\big(\lambda\,\alpha_t(L_{\mu_t}-\widehat L_{S_t})\big)\Big\}
% \;\le\;
% \exp\!\Big(\frac{\lambda^2}{2\,M_T}\Big).
% \end{equation}




% Exponentiate both sides on equation(\ref{eq:sum}) and take expectation over all randomness to obtain
% \[
% \mathbb{E}\exp\!\Big\{\lambda \mathcal{G}_T -\sum_{t=1}^T \mathrm{KL}(\widehat\Pi_t\|\Pi_{t-1})\Big\}
% \ \le\
% \mathbb{E}\exp\!\Big\{\sum_{t=1}^T \log \mathbb{E}_{\Pi_{t-1}}
% \exp\!\big(\lambda\alpha_t(L_{\mu_t}-\widehat L_{S_t})\big)\Big\}.
% \]
% and the right-hand side is bounded by
% $\exp\{\lambda^2/(2M_T)\}$.


% Define the nonnegative random variable
% \[
% X\ :=\ \exp\!\Big\{
% \lambda\mathcal{G}_T -\sum_{t=1}^T \mathrm{KL}(\widehat\Pi_t\|\Pi_{t-1})
% -\frac{\lambda^2}{2M_T}\Big\},
% \quad \text{so}\quad \mathbb{E}[X]\le 1.
% \]
% Markov’s inequality with threshold $1/\delta$ gives
% $\Pr\{X\ge 1/\delta\}\le \delta$.
% Equivalently, 

% with probability at least $1-\delta$, the exponent in $X$
% is at most $\log(1/\delta)$, which rearranges to \eqref{eq:hp-before-opt}.


% \emph{Justification.} Since $\ell\in[0,1]$, each summand in $\widehat L_{S_t}(W)$ is bounded, the centered average
% $\widehat L_{S_t}(W)-L_{\mu_t}(W)$ is sub-Gaussian with variance proxy $1/[4(m_t-1)]$ in the standard PAC-Bayesian sharpening (the ``$m_t-1$'' normalization).
% Thus the MGF of $L_{\mu_t}(W)-\widehat L_{S_t}(W)$ is bounded by the RHS of \eqref{eq:single-mgf}.





% \textit{Step 3 (mgf control, pooled).}
% For fixed $(P,W)$, $\widehat L_{S_t}(W)$ is an average of $m_t$ i.i.d.\ $[0,1]$ variables.
% By Lemma~\ref{lem:hoeffding}, for each $t$,
% \[
% \mathbb{E}_{S_t}\exp\!\big\{\lambda\alpha_t(L_{\mathcal{D}_t}(W)-\widehat L_{S_t}(W))\big\}
% \ \le\
% \exp\!\Big(\frac{\lambda^2\alpha_t^2}{2(m_t-1)}\Big).
% \]
% Independence across $\{S_t\}$ and Fubini then give
% \[
%  \mathbb{E}\exp\big[ \sum_{t=1}^T \log \mathbb{E}_{\Pi_{t-1}}
%  e^{\lambda\alpha_t(L_{\mu_t}-\widehat L_{S_t})}\big]
% \;\le\;
% \exp\!\Big(\frac{\lambda^2}{2M_T}\Big).
% \]

% \textit{Step 4 (Markov + optimization).}
% By Markov's inequality, with probability at least $1-\delta$,
% \[
% \lambda\mathcal{G}_T 
% \ \le\
% \sum_{t=1}^T \mathrm{KL}(\widehat\Pi_t\|\Pi_{t-1})
% \;+\;
% \frac{\lambda^2}{2M_T}
% \;+\;
% \log\!\frac{1}{\delta}.
% \]
% Optimizing over $\lambda>0$ yields
% \[
% \mathcal{G}_T 
% \ \le\
% \sqrt{\frac{
% \sum_{t=1}^T \mathrm{KL}(\widehat\Pi_t\|\Pi_{t-1})
% +\log(1/\delta)}
% {M_T}}.
% \]

% % \textit{Step 5 (hyper-level KL chain rule and rewriting).}

% Finally, recall that $\mathcal{G}_T $ is the $\{\alpha_t\}$-weighted gap; replacing by the uniform average $(1/T)\sum_t$ transforms, while the denominator remains $M_T=\sum_t(m_t-1)$.
% Thus
% \begin{equation}
%     \begin{aligned}
% &\frac1T\sum_{t=1}^T \mathcal{L}_t 
% - \frac1T\sum_{t=1}^T \hat{\mathcal{L}}_t
% \le
% \\
% &\sqrt{
% \frac{
% \sum_{t=1}^T \mathbb{E}_{P_t\sim\mathcal P_{t-1}}\!\mathrm{KL}(Q_t\|P_t)
% +\sum_{t=1}^T \mathrm{KL}(\mathcal P_t\|\mathcal P_{t-1})
% +\log(1/\delta)}
% {2\,\sum_{t=1}^T (m_t-1)}
% },
%     \end{aligned}
% \end{equation}
% which is the stated bound.
% \end{proof}



% \subsection{Key Lemmas for PECL Subspace Analysis}

% 下面给出按“先交代符号与性质，再给出精简引理”的改写版本，便于直接替换到正文中。

% ---

% This decomposition cleanly separates the task-level KL term from the hyper-level KL drift, without resorting to product measures.


% Define
% \[
% h_t(P):=\mathbb E_{S_t}\mathbb E_{W\sim Q_t(\cdot\mid P,S_t)} g_t(W)
% -\mathbb E_{S_t}\mathrm{KL}\!\big(Q_t(\cdot\mid P,S_t)\,\|\,P\big).
% \]
% By \eqref{eq:model-step}, $h_t(P)\le \frac{\lambda^2}{2}(m_t-1)$.
% Applying\eqref{Hoeff} on the hyper space to move from $\mathcal P_{t-1}$ to $\mathcal P_t$ and then using $\log\mathbb E e^{X}\ge \mathbb E X$ gives
% \[
% \mathbb E_{P\sim\mathcal P_t} h_t(P)
% \;\le\;
% \mathrm{KL}(\mathcal P_t\|\mathcal P_{t-1})
% +\frac{\lambda^2}{2}\,(m_t-1).
% \]
% Substituting the definition of $h_t$ and exchanging expectations,
% \begin{equation}\label{eq:two-term}
% \mathbb E_{S_t}\mathbb E_{P_t\sim\mathcal P_t}\mathbb E_{W\sim Q_t(\cdot\mid P_t,S_t)} g_t(W)
% \;\le\;
% \mathbb E_{S_t}\mathbb E_{P_t\sim\mathcal P_t}\mathrm{KL}\!\big(Q_t(\cdot\mid P_t,S_t)\,\|\,P_t\big)
% +\mathrm{KL}(\mathcal P_t\|\mathcal P_{t-1})
% +\frac{\lambda^2}{2}\,(m_t-1).
% \end{equation}


% \emph{Step 3 (aggregation).}
% Let
% \[
% \Delta_t:=\mathbb E_{P_t\sim\mathcal P_t}\mathbb E_{W\sim Q_t(\cdot\mid P_t,S_t)}\!\big(L_{\mathcal D_t}(W)-\hat L_{S_t}(W)\big).
% \]
% Dividing \eqref{eq:two-term} by $\lambda$, summing over $t$, and using a standard exponential supermartingale plus Markov's inequality yields, with probability at least $1-\delta$,
% \[
% \sum_{t=1}^T \Delta_t
% \;\le\;
% \frac{1}{\lambda}\!\left(
% \sum_{t=1}^T \mathbb E_{P_t\sim\mathcal P_t}\mathrm{KL}(Q_t\|P_t)
% +\sum_{t=1}^T \mathrm{KL}(\mathcal P_t\|\mathcal P_{t-1})
% +\log\tfrac1\delta\right)
% +\frac{\lambda}{2}\sum_{t=1}^T (m_t-1).
% \]
% Optimizing over $\lambda>0$ and dividing by $T$ yields \eqref{eq:main-bound}.

% ### Standing conventions and measurable structure

% \subsection{Proof}

% \textbf{Assumptions.}
% For each task $t\in\{1,\dots,T\}$, let $S_t=\{z_{tj}\}_{j=1}^{m_t}$ be i.i.d.\ draws from $\mathcal D_t$, and let $\ell(w,z)\in[0,1]$.
% Write
% \[
% \hat L_{S_t}(W)=\frac1{m_t}\sum_{j=1}^{m_t}\ell(W,z_{tj}),
% \qquad
% L_{\mathcal D_t}(W)=\mathbb E_{z\sim\mathcal D_t}\ell(W,z).
% \]
% Let $\mathcal F_{t-1}:=\sigma(S_{1:t-1})$ be the canonical filtration.
% At the beginning of step $t$, sample a hyper-parameter $P_t\sim\mathcal P_{t-1}$, which is $\mathcal F_{t-1}$-measurable.
% After observing $S_t$, define a model-posterior $Q_t(\cdot\mid P_t,S_t)$ and update the hyper-posterior to $\mathcal P_t$, with absolute continuity
% \[
% \mathcal P_t\ll\mathcal P_{t-1},
% \qquad
% Q_t(\cdot\mid P,S_t)\ll P(\cdot)\ \ \text{for } \mathcal P_{t-1}\text{-a.e.\ }P.
% \]
% For absolute continuity $Q\ll P$, define $\mathrm{KL}(Q\|P):=\int \log\!\big(\frac{dQ}{dP}\big)\,dQ$.

% Let $M_T:=\sum_{t=1}^T (m_t-1)$ and define $\alpha_t:=(m_t-1)/M_T$ so that $\sum_t\alpha_t=1$.
% Define the step-$t$ generalization gap
% \[
% \Delta_t
% :=\mathbb E_{P_t\sim\mathcal P_{t-1}}\,
% \mathbb E_{W\sim Q_t(\cdot\mid P_t,S_t)}
% \big(L_{\mathcal D_t}(W)-\hat L_{S_t}(W)\big).
% \]
% We study the size-weighted gap $\bar\Delta:=\sum_{t=1}^T \alpha_t \Delta_t$ and then relate it to the form in Theorem~\ref{thm:pathwise-pacbayes}.









% \begin{theorem}[Pathwise PAC--Bayes with time-varying hyper-posterior]\label{thm:pathwise-pacbayes}
% For any $\delta\in(0,1]$, with probability at least $1-\delta$ over $S_{1:T}$,
% \begin{equation}\label{eq:main-bound}
% \begin{aligned} 
% & \frac{1}{T} \sum_{t=1}^T \underbrace{\mathbb{E}_{P_t \sim \mathcal{P}_t} \mathbb{E}_{W \sim Q_t(\cdot \mid P_t, S_t)} L_{\mathcal{D}_t}(W)}_{\text{test risk at step } t} 
% \le \ \frac{1}{T} \sum_{t=1}^T \underbrace{\mathbb{E}_{P_t \sim \mathcal{P}_t} \mathbb{E}_{W \sim Q_t(\cdot \mid P_t, S_t)} \hat{L}_{S_t}(W)}_{\text{empirical risk at step } t} \\ 
% & + \sqrt{\frac{ \sum_{t=1}^T \mathbb{E}_{P_t \sim \mathcal{P}_t} \mathrm{KL}\big(Q_t(\cdot \mid P_t, S_t) \,\|\, P_t\big) + \sum_{t=1}^T \mathrm{KL}(\mathcal{P}_t \,\|\, \mathcal{P}_{t-1}) + \log\frac{1}{\delta} }{2 \sum_{t=1}^T (m_t - 1)}}. 
% \end{aligned} 
% \end{equation}
% \end{theorem}


% \begin{proof}
    
% Fix $\lambda>0$ and define $g_t(W)=\lambda\alpha_t\big(L_{\mathcal D_t}(W)-\hat L_{S_t}(W)\big)$.


% \paragraph{Step 1 (Model-Level Variational Bound, $W$-space).}
% For any fixed $P$ and $S_t$, the Donsker--Varadhan (DV) inequality ~\cite{donsker1975asymptotic,seeger2002pac} yields 
% \begin{equation}
% \label{seeger}
%     \mathbb E_{W\sim Q_t(\cdot\mid P,S_t)} g_t(W)
% \;\le\;
% \mathrm{KL}\!\big(Q_t(\cdot\mid P,S_t)\,\|\,P\big)
% +\log\mathbb E_{W\sim P}\!\big[e^{g_t(W)}\big].
% \end{equation}
% Taking $\mathbb E_{S_t}$ using the standard Hoeffding–MGF bound for bounded losses\cite{Hoeffding01031963}, we have
% \begin{equation}
% \label{Hoeff}
% \mathbb E_{S_t}\log\mathbb E_{W\sim P} e^{g_t(W)}
% \;\le\; \frac{(\lambda \alpha_t)^2}{8(m_t-1)},
% \end{equation}
% Substituting \eqref{Hoeff} into the expectation of \eqref{seeger}, we obtain
% \begin{equation}\label{eq:model-step}
% \mathbb E_{S_t}\mathbb E_{W\sim Q_t(\cdot\mid P,S_t)} g_t(W)
% \;\le\;
% \mathbb E_{S_t}\mathrm{KL}\!\big(Q_t(\cdot\mid P,S_t)\,\|\,P\big)
% + \frac{(\lambda \alpha_t)^2}{8(m_t-1)}.
% \end{equation}

% \paragraph{Step 2 (Hyper-Level Variational Bound, $P$-space).}
% Define a function 
% $h_t: \mathcal{P} \to \mathbb{R}$
% on the hyper-parameter space:
% $$h_t(P) := \mathbb{E}_{S_t} \left[ \mathbb{E}_{W \sim Q_t(\cdot \mid P, S_t)} [g_t(W)] \right] - \mathbb{E}_{S_t} \left[ \mathrm{KL}\big(Q_t(\cdot \mid P, S_t) \,\|\, P\big) \right].$$
% From \eqref{eq:model-step}, we know 
% $h_t(P) \le  \frac{(\lambda \alpha_t)^2}{8(m_t-1)}$for all $P$.


% We now apply the DV inequality again, this time on the 
% $P$-space, to change the measure from 
% $\mathcal{P}_{t-1}$ to $\mathcal{P}_t$:
% $$\log \mathbb{E}_{P \sim \mathcal{P}_t} \left[ e^{h_t(P)} \right] \le \mathrm{KL}(\mathcal{P}_t \,\|\, \mathcal{P}_{t-1}) + \log \mathbb{E}_{P \sim \mathcal{P}_{t-1}} \left[ e^{h_t(P)} \right]. $$

%  We proceed by applying the bound $h_t(P) \le \frac{\lambda^2}{2}(m_t-1)$.
% $$\log \mathbb{E}_{P \sim \mathcal{P}_t} \left[ e^{h_t(P)} \right] \le \mathrm{KL}(\mathcal{P}_t \,\|\, \mathcal{P}_{t-1}) +  \frac{(\lambda \alpha_t)^2}{8(m_t-1)}.$$

% By Jensen's inequality (
% $\log \mathbb{E}[X] \ge \mathbb{E}[\log X]$
% ), and more relevantly 
% $\log \mathbb{E}[e^X] \ge \mathbb{E}[X]$:
% $$\mathbb{E}_{P \sim \mathcal{P}_t} [h_t(P)] \le \mathrm{KL}(\mathcal{P}_t \,\|\, \mathcal{P}_{t-1}) + 
% \frac{(\lambda \alpha_t)^2}{8(m_t-1)}.$$


% Unfolding $h_t(P)$ and interchanging expectations (Fubini),
% \begin{equation}
% \label{eq:two-term}
% \mathbb E_{S_t}\Big[
% \lambda\alpha_t\Delta_t
% -\alpha_t\,\mathbb E_{P_t\sim\mathcal P_{t-1}}\mathrm{KL}\!\big(Q_t(\cdot\mid P_t,S_t)\,\|\,P_t\big)
% \Big]
% \;\le\;
% \mathrm{KL}(\mathcal P_t\|\mathcal P_{t-1})
% +\frac{(\lambda\alpha_t)^2}{8\,(m_t-1)}.
% \end{equation}


% From \eqref{eq:two-term} and the tower property, one obtains the conditional exponential bound
% \begin{equation}
% \label{eq:cond-mgf}
% \mathbb E\!\left[
% \exp\!\Big(
% \lambda\alpha_t\Delta_t
% -\alpha_t\,\mathbb E_{P_t\sim\mathcal P_{t-1}}\mathrm{KL}(Q_t\|P_t)
% -\mathrm{KL}(\mathcal P_t\|\mathcal P_{t-1})
% -\frac{(\lambda\alpha_t)^2}{8\,(m_t-1)}
% \Big)\,\Big|\,\mathcal F_{t-1}\right]\;\le\;1.
% \end{equation}
% Define $X_t$ to be the exponential increment inside \eqref{eq:cond-mgf} and
% $M_t:=\prod_{s=1}^t X_s$ with $M_0=1$. Then $(M_t)_{t\ge0}$ is a nonnegative supermartingale, so
% $\mathbb E[M_T]\le 1$. By Markov's inequality, with probability at least $1-\delta$,
% \begin{equation}
% \label{eq:hpp}
% \sum_{t=1}^T \lambda\alpha_t\Delta_t
% \;\le\;
% \sum_{t=1}^T \alpha_t\,\mathbb E_{P_t\sim\mathcal P_{t-1}}\mathrm{KL}(Q_t\|P_t)
% \;+\; \sum_{t=1}^T \mathrm{KL}(\mathcal P_t\|\mathcal P_{t-1})
% \;+\; \frac{\lambda^2}{8}\sum_{t=1}^T \frac{\alpha_t^2}{(m_t-1)}
% \;+\; \log\tfrac1\delta.
% \end{equation}









% Substituting the definition of $h_t(P)$ back:
% $$\mathbb{E}_{P_t \sim \mathcal{P}_t} \left[ \mathbb{E}_{S_t} \left[ \mathbb{E}_{W \sim Q_t(\cdot \mid P_t, S_t)} [g_t(W)] \right] - \mathbb{E}_{S_t} \left[ \mathrm{KL}\big(Q_t(\cdot \mid P_t, S_t) \,\|\, P_t\big) \right] \right] \le \mathrm{KL}(\mathcal{P}_t \,\|\, \mathcal{P}_{t-1}) +  \frac{(\lambda \alpha_t)^2}{8(m_t-1)}.$$

% By Fubini's theorem (swapping $\mathbb{E}_{P_t}$ and $\mathbb{E}_{S_t}$
% ), and recalling $g_t(W)$:
% \begin{equation}
% \label{eq: KL step2}
% \mathbb{E}_{S_t} \left[ \mathbb{E}_{P_t \sim \mathcal{P}_t} \mathbb{E}_{W \sim Q_t(\cdot \mid P_t, S_t)} [\lambda(L_{\mathcal{D}_t} - \hat{L}_{S_t})] \right] \le \mathbb{E}_{S_t} \left[ \mathbb{E}_{P_t \sim \mathcal{P}_t} \mathrm{KL}(Q_t \,\|\, P_t) \right] + \mathrm{KL}(\mathcal{P}_t \,\|\, \mathcal{P}_{t-1}) +  \frac{\lambda^2}{2(m_t-1)}.
% \end{equation}

% \paragraph{Step 3 (Aggregation via size-weighted supermartingale).}
% Set
% \[
% \Delta_t
% := \mathbb E_{P_t\sim\mathcal P_t}\mathbb E_{W\sim Q_t(\cdot\mid P_t,S_t)}
% \big(L_{\mathcal D_t}(W)-\hat L_{S_t}(W)\big),
% \]
% and abbreviate the per-step complexity term by
% \[
% K_t
% := \mathbb E_{P_t\sim\mathcal P_t}\mathrm{KL}\!\big(Q_t(\cdot\mid P_t,S_t)\,\|\,P_t\big)
% \;+\; \mathrm{KL}(\mathcal P_t\|\mathcal P_{t-1}).
% \]
% The result from Step 2 \eqref{eq: KL step2}, $\mathbb{E}_{S_t}[\lambda \Delta_t - K_t] \le \frac{\lambda^2}{2}(m_t-1)$, holds in expectation.


% Define the sampling–size weights
% $\alpha_t:=\frac{(m_t-1)}{\sum_{t=1}^T (m_t-1)}$ so that $\sum_t \alpha_t=1$.
% Consider the size–weighted generalization gap
% \[
% \bar\Delta\;:=\;\sum_{t=1}^T \alpha_t\,
% \mathbb E_{P_t\sim\mathcal P_t}\mathbb E_{W\sim Q_t(\cdot\mid P_t,S_t)}
% \big(L_{\mathcal D_t}(W)-\hat L_{S_t}(W)\big)
% \;=\;\sum_{t=1}^T \alpha_t\,\Delta_t.
% \]


% We convert this to a high-probability statement over $S_{1:T}$ using an exponential supermartingale argument. The conditional MGF bound $\mathbb{E}[\exp(Y_t) \mid \mathcal{F}_{t-1}] \le 1$ holds, where $\mathcal F_{t-1}:=\sigma(S_{1:t-1})$. 
% Define the nonnegative process
% \[
% M_0=1,\quad  M_t:=M_{t-1}\exp\Big(\lambda \alpha_t \Delta_t  -\alpha_t K_t  -\tfrac{\lambda^2}{2}\alpha_t^2\tfrac{1}{(m_t-1)}  \Big),  .
% \]
% Then $(M_t)_{t\ge0}$ is a supermartingale and $\mathbb{E}[M_T]\leq 1$. By Markov's inequality,
% $\mathbb P(M_T\ge 1/\delta)\le\delta$ for any $\delta\in(0,1]$.
% Hence, with probability at least $1-\delta$,

% \[
% \sum_{t=1}^T \lambda\alpha_t\Delta_t
% \;\le\;
% \sum_{t=1}^T \alpha_t K_t
% \;+\; \frac{\lambda^2}{2}\sum_{t=1}^T \alpha_t^2 \frac{1}{ (m_t-1)}
% \;+\;\log\tfrac1\delta.
% \]



% where
% $\Delta_t
% := \mathbb E_{P_t\sim\mathcal P_t}\mathbb E_{W\sim Q_t(\cdot\mid P_t,S_t)}
% \big(L_{\mathcal D_t}(W)-\hat L_{S_t}(W)\big)$.
% Apply the exponential supermartingale method to the \emph{weighted} sum by
% scaling the $t$-th step with $\lambda\alpha_t$:
% using the one-step bounds from Steps~1--2 with $g_t$ replaced by
% $\lambda\alpha_t(L_{\mathcal D_t}-\hat L_{S_t})$, we obtain the conditional MGF control
% \[
% \mathbb E\!\left[\exp\!\Big(\lambda\alpha_t\Delta_t
% -\alpha_t K_t
% -\tfrac{(\lambda\alpha_t)^2}{2}(m_t-1)\Big)\,\Big|\,\mathcal F_{t-1}\right]\;\le\;1,
% \]
% where $K_t:=\mathbb E_{P_t\sim\mathcal P_t}\mathrm{KL}(Q_t\|P_t)+\mathrm{KL}(\mathcal P_t\|\mathcal P_{t-1})$ and
% $\mathcal F_{t-1}:=\sigma(S_{1:t-1})$.
% Multiplying over $t$ yields a nonnegative supermartingale $(M_t)_{t\ge0}$ with $M_0=1$ and
% $\mathbb E[M_T]\le 1$.
% By Markov’s inequality, with probability at least $1-\delta$,
% \[
% \sum_{t=1}^T \lambda\alpha_t\Delta_t
% \;\le\;
% \sum_{t=1}^T \alpha_t K_t
% \;+\; \frac{\lambda^2}{2}\sum_{t=1}^T \alpha_t^2 (m_t-1)
% \;+\;\log\tfrac1\delta.
% % \]
% Since $\alpha_t=\frac{(m_t-1)}{\sum_{t=1}^T (m_t-1)}$, the quadratic term self-normalizes:
% \[
% \sum_{t=1}^T \alpha_t^2 (m_t-1)
% =\frac{1}{\sum_{t=1}^T (m_t-1)}.
% \]
% Therefore, for all $\lambda>0$, divided by $\lambda$
% \[
% \bar\Delta
% \;\le\;
% \frac{1}{\lambda}\Big(\sum_{t=1}^T \alpha_t K_t+\log\tfrac1\delta\Big)
% \;+\; \frac{\lambda}{2\sum_{t=1}^T (m_t-1)}.
% \]
% Optimizing in $\lambda$ gives
% \[
% \bar\Delta
% \;\le\;
% \sqrt{\frac{2\Big(\sum_{t=1}^T \alpha_t K_t+\log\tfrac1\delta\Big)}{\sum_{t=1}^T (m_t-1)}}
% \;\le\;
% \sqrt{\frac{2\Big(\sum_{t=1}^T K_t+\log\tfrac1\delta\Big)}{\sum_{t=1}^T (m_t-1)}},
% \]
% where the last step uses $\sum_t\alpha_t K_t\le \sum_t K_t$ since $\alpha_t\in[0,1]$.
% Finally, substituting $K_t=\mathbb E_{P_t\sim\mathcal P_t}\mathrm{KL}(Q_t\|P_t)+\mathrm{KL}(\mathcal P_t\|\mathcal P_{t-1})$
%  yields the complexity term of \eqref{eq:main-bound}:
% \[
% \bar\Delta\;\le\;
% \sqrt{\frac{\sum_{t=1}^T \mathbb E_{P_t\sim\mathcal P_t}\mathrm{KL}(Q_t\|P_t)
% +\sum_{t=1}^T \mathrm{KL}(\mathcal P_t\|\mathcal P_{t-1})
% +\log(1/\delta)}{2\sum_{t=1}^T (m_t-1)}}\,.
% \]
% \emph{Remark.} If one reports the \emph{unweighted} average $\frac1T\sum_t\Delta_t$ instead of the size–weighted $\bar\Delta$, then applying the supermartingale directly to $\sum_t\Delta_t$ (i.e., with equal weights) produces the multiplicative form
% $\frac1T\sqrt{2(\sum_t K_t+\log(1/\delta))\sum_t(m_t-1)}$.
% Thus the “division” form in \eqref{eq:main-bound} is a consequence of size-weighted aggregation and the choice of per-step scaling $\lambda\alpha_t$.
% \end{proof}





% \paragraph{Step 3 (Aggregation; detailed).}
% Set
% \[
% \Delta_t
% := \mathbb E_{P_t\sim\mathcal P_t}\mathbb E_{W\sim Q_t(\cdot\mid P_t,S_t)}
% \big(L_{\mathcal D_t}(W)-\hat L_{S_t}(W)\big),
% \]
% and abbreviate the per-step complexity term by
% \[
% K_t
% := \mathbb E_{P_t\sim\mathcal P_t}\mathrm{KL}\!\big(Q_t(\cdot\mid P_t,S_t)\,\|\,P_t\big)
% \;+\; \mathrm{KL}(\mathcal P_t\|\mathcal P_{t-1}).
% \]
% The result from Step 2 \eqref{eq: KL step2}, $\mathbb{E}_{S_t}[\lambda \Delta_t - K_t] \le \frac{\lambda^2}{2}(m_t-1)$, holds in expectation.


% We convert this to a high-probability statement over $S_{1:T}$ using an exponential supermartingale argument. The conditional MGF bound $\mathbb{E}[\exp(Y_t) \mid \mathcal{F}_{t-1}] \le 1$ holds, where $\mathcal F_{t-1}:=\sigma(S_{1:t-1})$. 
% Define the nonnegative process
% \[
% M_0:=1,\qquad
% M_t:=M_{t-1}\exp\!\Big(\lambda\Delta_t - K_t - \tfrac{\lambda^2}{2}(m_t-1)\Big).
% \]
% Then $(M_t)_{t\ge0}$ is a supermartingale and, by Markov's inequality,
% $\mathbb P(M_T\ge 1/\delta)\le\delta$ for any $\delta\in(0,1]$.
% Hence, with probability at least $1-\delta$,
% \[
% \sum_{t=1}^T \lambda\Delta_t
% \;\le\;
% \sum_{t=1}^T K_t
% \;+\; \frac{\lambda^2}{2}\sum_{t=1}^T (m_t-1)
% \;+\; \log\tfrac1\delta.
% \]
% Equivalently, for all $\lambda>0$, dividing by $\lambda$,
% \[
% \sum_{t=1}^T \Delta_t
% \;\le\;
% \frac{1}{\lambda}\!\left(
% \sum_{t=1}^T K_t + \log\tfrac1\delta\right)
% \;+\; \frac{\lambda}{2}\sum_{t=1}^T (m_t-1).
% \]

% Minimizing the right-hand side over $\lambda>0$ by choosing
% \[
% \lambda^\star
% \;=\;
% \sqrt{\frac{2\Big(
% \sum_{t=1}^T K_t + \log\tfrac1\delta
% \Big)}{\sum_{t=1}^T (m_t-1)}}
% \]
% gives

% \[
% \sum_{t=1}^T \Delta_t
% \;\le\;
% \sqrt{\,2\,\sum_{t=1}^T (m_t-1)\!\left(\sum_{t=1}^T K_t+\log\tfrac1\delta\right)}.
% \]

% Dividing by $T$ and expanding  $K_t$ and $\Delta_t$ yields the bound:

% \begin{align*}
% &\frac{1}{T} \sum_{t=1}^T \mathbb{E}_{P_t \sim \mathcal{P}t} \mathbb{E}_{W \sim Q_t} L_{\mathcal{D}t}(W)
% \le  \frac{1}{T} \sum_{t=1}^T \mathbb{E}_{P_t \sim \mathcal{P}t} \mathbb{E}_{W \sim Q_t} \hat{L}{S_t}(W) \\
% & + \frac{1}{T}\sqrt{2 \left( \sum_{t=1}^T \mathbb{E}_{P_t}[\mathrm{KL}] + \sum_{t=1}^T \mathrm{KL} + \log\tfrac1\delta \right) \left( \sum_{t=1}^T (m_t-1) \right)}.
% \end{align*}

















% Optimizing the right-hand side over $\lambda$ (choose $\lambda^\star=\sqrt{2(\sum_t K_t+\log(1/\delta))/M_T}$) yields
% \[
% \sum_{t=1}^T \Delta_t
% \;\le\;
% \sqrt{\,2\,M_T\!\left(\sum_{t=1}^T K_t+\log\tfrac1\delta\right)}.
% \]
% Dividing by $T$ and unfolding $\Delta_t$ and $K_t$ gives the bound stated in







% Thus, with probability at least $1-\delta$, $\sum_{t=1}^T Y_t < \log(1/\delta)$. Substituting $Y_t$ and rearranging yields:

% $$\lambda \sum_{t=1}^T \Delta_t \;\le\; \sum_{t=1}^T K_t +\log\tfrac1\delta +\frac{\lambda^2}{2}\sum_{t=1}^T (m_t-1).$$

% Dividing by $\lambda$ gives:

% $$\sum_{t=1}^T \Delta_t \;\le\; \frac{1}{\lambda}\!\left( \underbrace{\sum_{t=1}^T \mathbb E_{P_t\sim\mathcal P_t}\mathrm{KL}(Q_t\|P_t) +\sum_{t=1}^T \mathrm{KL}(\mathcal P_t\|\mathcal P_{t-1}) +\log\tfrac1\delta}_{:=A}\right) +\frac{\lambda}{2}\underbrace{\sum_{t=1}^T (m_t-1)}_{:=B}.$$

% M









% The process $M_T = \exp\left(\sum_{t=1}^T Y_t\right)$ is a non-negative supermartingale with $\mathbb{E}[M_T] \le 1$.










% From \eqref{eq: KL step2} and the convexity of $x\mapsto e^x$ (Jensen),
% \[
% \mathbb E\!\left[
% \exp\!\Big(\lambda\Delta_t - K_t - \tfrac{\lambda^2}{2}(m_t-1)\Big)
% \,\Big|\, \mathcal F_{t-1}
% \right]\;\le\;1,
% \]
% where $\mathcal F_{t-1}:=\sigma(S_{1:t-1})$.
% Define the nonnegative process
% \[
% M_0:=1,\qquad
% M_t:=M_{t-1}\exp\!\Big(\lambda\Delta_t - K_t - \tfrac{\lambda^2}{2}(m_t-1)\Big).
% \]











% Let
% \[
% \Delta_t
% \;:=\;
% \mathbb E_{P_t\sim\mathcal P_t}\,
% \mathbb E_{W\sim Q_t(\cdot\mid P_t,S_t)}
% \big(L_{\mathcal D_t}(W)-\hat L_{S_t}(W)\big),
% \qquad
% M_T:=\sum_{t=1}^T(m_t-1).
% \]

% Taking the conditional expectation w.r.t.\ $\mathcal F_{t-1}$ and applying the Hoeffding MGF bound for $[0,1]$-bounded losses gives
% \begin{equation}\label{eq:cond-mgf}
% \mathbb E\!\left[
% \mathbb E_{P_t\sim\mathcal P_t}\mathbb E_{W\sim Q_t}
% \!\left[e^{g_t(W)-\mathrm{KL}(Q_t\|P_t)}\right]
% \,\Big|\,\mathcal F_{t-1}
% \right]
% \;\le\;
% \exp\!\Big(\mathrm{KL}(\mathcal P_t\|\mathcal P_{t-1})+\tfrac{\lambda^2}{2}(m_t-1)\Big).
% \end{equation}

% \medskip\noindent
% \textit{(b) Jensen’s inequality to pass from averages to the MGF of $\Delta_t$.}
% By convexity of the exponential,
% \[
% e^{\lambda\Delta_t}
% =
% \exp\!\Big(\lambda\,\mathbb E_{P_t,Q_t}[L_{\mathcal D_t}-\hat L_{S_t}]\Big)
% \;\le\;
% \mathbb E_{P_t\sim\mathcal P_t}\mathbb E_{W\sim Q_t}
% \!\big[e^{g_t(W)}\big].
% \]
% Combining with \eqref{eq:cond-mgf} we obtain the one-step conditional control
% \begin{equation}\label{eq:one-step-super}
% \mathbb E\!\left[
% \exp\!\Big(\lambda\Delta_t
% -\mathrm{KL}(\mathcal P_t\|\mathcal P_{t-1})
% -\mathbb E_{P_t\sim\mathcal P_t}\mathrm{KL}(Q_t\|P_t)
% -\tfrac{\lambda^2}{2}(m_t-1)\Big)
% \;\Big|\;\mathcal F_{t-1}\right]\;\le\;1.
% \end{equation}

% \medskip\noindent
% \textit{(c) Exponential supermartingale and Markov’s inequality.}
% Define
% \[
% X_t
% :=
% \exp\!\Big(\lambda\Delta_t
% -\mathrm{KL}(\mathcal P_t\|\mathcal P_{t-1})
% -\mathbb E_{P_t\sim\mathcal P_t}\mathrm{KL}(Q_t\|P_t)
% -\tfrac{\lambda^2}{2}(m_t-1)\Big),
% \qquad
% M_t:=\prod_{s=1}^t X_s.
% \]
% By \eqref{eq:one-step-super}, $(M_t)_{t\ge 0}$ is a nonnegative supermartingale with
% $\mathbb E[M_T]\le 1$.
% Hence, for any $\delta\in(0,1]$, Markov’s inequality gives
% \[
% \mathbb P\!\left(M_T\ge \tfrac1\delta\right)\;\le\;\delta.
% \]
% On the event $\{M_T<1/\delta\}$ we have
% \[
% \sum_{t=1}^T \lambda\Delta_t
% \;\le\;
% \sum_{t=1}^T \mathrm{KL}(\mathcal P_t\|\mathcal P_{t-1})
% +\sum_{t=1}^T \mathbb E_{P_t\sim\mathcal P_t}\mathrm{KL}(Q_t\|P_t)
% +\tfrac{\lambda^2}{2}\sum_{t=1}^T (m_t-1)
% +\log\tfrac1\delta.
% \]
% Dividing by $\lambda>0$ yields the high-probability bound
% \[
% \sum_{t=1}^T \Delta_t
% \;\le\;
% \frac{1}{\lambda}\!\left(
% \sum_{t=1}^T \mathbb E_{P_t\sim\mathcal P_t}\mathrm{KL}(Q_t\|P_t)
% +\sum_{t=1}^T \mathrm{KL}(\mathcal P_t\|\mathcal P_{t-1})
% +\log\tfrac1\delta\right)
% +\frac{\lambda}{2}\sum_{t=1}^T (m_t-1).
% \]

% \medskip\noindent
% \textit{(d) Optimizing over $\lambda$.}
% Minimizing the right-hand side over $\lambda>0$ by choosing
% \[
% \lambda^\star
% \;=\;
% \sqrt{\frac{2\Big(
% \sum_{t=1}^T \mathbb E_{P_t\sim\mathcal P_t}\mathrm{KL}(Q_t\|P_t)
% +\sum_{t=1}^T \mathrm{KL}(\mathcal P_t\|\mathcal P_{t-1})
% +\log\tfrac1\delta\Big)}{\sum_{t=1}^T (m_t-1)}}
% \]
% gives
% \[
% \sum_{t=1}^T \Delta_t
% \;\le\;
% \sqrt{2 M_T\left(
% \sum_{t=1}^T \mathbb E_{P_t\sim\mathcal P_t}\mathrm{KL}(Q_t\|P_t)
% +\sum_{t=1}^T \mathrm{KL}(\mathcal P_t\|\mathcal P_{t-1})
% +\log\tfrac1\delta\right)}.
% \]
% Dividing by $T$ and expanding $\Delta_t$ yields the bound in \eqref{eq:main-bound}.
% \end{proof}

% \begin{remark}[Prior-weighted variant]
% If one prefers to keep the outer expectation with respect to $\mathcal P_{t-1}$ (to mirror the algorithmic flow $P_t\sim\mathcal P_{t-1}$), replace $\mathbb E_{P_t\sim\mathcal P_t}$ by $\mathbb E_{P_t\sim\mathcal P_{t-1}}$ on both sides of \eqref{eq:main-bound}.
% In the proof, reverse the hyper-level DV direction in Step~2.
% The bound retains the same structure and denominator, with the per-task complexity term becoming $\sum_t \mathbb E_{P_t\sim\mathcal P_{t-1}}\mathrm{KL}(Q_t\|P_t)$.
% \end{remark}

% \begin{corollary}[Gaussian smoothing / LoRA subspace]
% If $Q_t(\cdot\mid P_t,S_t)=\mathcal N(\hat W_t,\Sigma_t)$ (e.g., restricted to a LoRA subspace), then
% \[
% \mathbb E_{W\sim Q_t(\cdot\mid P_t,S_t)}\hat L_{S_t}(W)
% =
% \mathbb E_{\varepsilon\sim\mathcal N(0,\Sigma_t)}\hat L_{S_t}(\hat W_t+\varepsilon).
% \]
% In a frozen-backbone setting with $P_t^{\mathrm{full}}=\delta_{W_0}\otimes P_t^\Delta$ and $Q_t^{\mathrm{full}}=\delta_{W_0}\otimes Q_t^\Delta$,
% \(
% \mathrm{KL}(Q_t^{\mathrm{full}}\|P_t^{\mathrm{full}})
% =
% \mathrm{KL}(Q_t^\Delta\|P_t^\Delta),
% \)
% so both the per-task KL and the smoothing term are confined to the LoRA subspace.
% \end{corollary}





% \paragraph{Measurable structure and notation}
% Let $\mathcal W_\Delta\subseteq\mathbb R^d$ be a linear subspace and let $W+\mathcal W_\Delta$ be its affine translate. Equip both with the trace Borel $\sigma$-algebra induced from $\mathbb R^d$:
% \[
% \mathcal B(\mathcal W_\Delta)\ :=\ {B\cap \mathcal W_\Delta:\ B\in \mathcal B(\mathbb R^d)},\qquad
% \mathcal B(W+\mathcal W_\Delta)\ :=\ {B\cap (W+\mathcal W_\Delta):\ B\in \mathcal B(\mathbb R^d)}.
% \]
% These coincide with the Borel $\sigma$-algebras generated by the subspace topologies on $\mathcal W_\Delta$ and $W+\mathcal W_\Delta$, respectively. In particular, both are standard Borel spaces. The translation
% \[
% T_W:\ (\mathcal W_\Delta,\mathcal B(\mathcal W_\Delta))\ \to\ (W+\mathcal W_\Delta,\mathcal B(W+\mathcal W_\Delta)),\quad
% T_W(\Delta)=W+\Delta
% \]
% is a bimeasurable bijection (Borel isomorphism). 


% For the product view, endow ${W}$ with the trivial $\sigma$-algebra ${\varnothing,{W}}$ and ${W}\times\mathcal W_\Delta$ with the product $\sigma$-algebra ${\varnothing,{W}}\otimes\mathcal B(\mathcal W_\Delta)$. The map
% \[
% \Phi:\ {W}\times \mathcal W_\Delta \to W+\mathcal W_\Delta,\qquad \Phi(W,\Delta)=W+\Delta
% \]
% is a measurable isomorphism.  For any measurable $F$, the pushforward $F_{\text{\#}}\mu$ is defined by $(F_{\#}\mu)(A)=\mu(F^{-1}(A))$.


\begin{lemma}[KL decomposition with a frozen backbone]
\label{equation: subspace-full}
Let $W\in\mathbb R^d$ be fixed and let $P^{\Delta},Q^{\Delta}$ be probability measures on $(\mathcal W_\Delta,\mathcal B(\mathcal W_\Delta))$ with $Q^{\Delta}\ll P^{\Delta}$. Define their pushforwards onto the affine subspace by
\[
P^{\mathrm{full}}:=(T_W)_{\text{\#}} P^{\Delta},\qquad
Q^{\mathrm{full}}:=(T_W)_{\text{\#}} Q^{\Delta}.
\]
Then $Q^{\mathrm{full}}\ll P^{\mathrm{full}}$ and
\[
\mathrm{KL}\left(Q^{\mathrm{full}}\middle|P^{\mathrm{full}}\right)
=\mathrm{KL}\left(Q^{\Delta}\middle|P^{\Delta}\right).
\]
Equivalently, identifying ${W}\times\mathcal W_\Delta$ with $W+\mathcal W_\Delta$ via $\Phi$, one may write
$P^{\mathrm{full}}=\Phi_{\text{\#}}(\delta_W\otimes P^{\Delta})$ and
$Q^{\mathrm{full}}=\Phi_{\text{\#}}(\delta_W\otimes Q^{\Delta})$, and the same equality holds.
\end{lemma}

\begin{proof}
  Absolute continuity is preserved by pushforward under measurable maps, hence $Q^{\mathrm{full}}\ll P^{\mathrm{full}}$. Since $T_W$ is a Borel isomorphism, by the Radon–Nikodym chain rule
 \[
 \frac{dQ^{\mathrm{full}}}{dP^{\mathrm{full}}}(T_W\Delta)=\frac{dQ^{\Delta}}{dP^{\Delta}}(\Delta)\quad\text{for }P^{\Delta}\text{-a.e.\ }\Delta.
 \]

\begin{equation*}
     \begin{aligned}
 \mathrm{KL}\left(Q^{\mathrm{full}}\middle|P^{\mathrm{full}}\right)
 &= \int_{W+\mathcal W_\Delta}
 \log\Big(\frac{dQ^{\mathrm{full}}}{dP^{\mathrm{full}}}(y)\Big)
 Q^{\mathrm{full}}(dy)\\
 &= \int_{\mathcal W_\Delta}
 \log\Big(\frac{dQ^{\mathrm{full}}}{dP^{\mathrm{full}}}\big(T_W\Delta\big)\Big)
 Q^{\Delta}(d\Delta) \quad \text{(} Q^{\mathrm{full}}=(T_W)_{\text{\#}}Q^{\Delta}\text{)} 
 \\
 &= \int_{\mathcal W_\Delta}
 \log\Big(\frac{dQ^{\Delta}}{dP^{\Delta}}(\Delta)\Big)
 Q^{\Delta}(d\Delta) \\
 &= \mathrm{KL}\left(Q^{\Delta}\middle|P^{\Delta}\right).
 \end{aligned}
\end{equation*}
% The KL identity follows by applying the KL invariance fact to the Borel isomorphism $T_W$. 
\end{proof}



%  Let $\mathcal W_\Delta\subseteq\mathbb R^d$ be a linear subspace and let $W+\mathcal W_\Delta$ be its affine translate. We equip both with the trace Borel $\sigma$-algebra induced from $\mathbb R^d$, i.e.
%  \[
%  \mathcal B(\mathcal W_\Delta)\ :=\ {B\cap \mathcal W_\Delta:\ B\in \mathcal B(\mathbb R^d)},\qquad
%  \mathcal B(W+\mathcal W_\Delta)\ :=\ {B\cap (W+\mathcal W_\Delta):\ B\in \mathcal B(\mathbb R^d)}.
%  \]
%  This equals the Borel $\sigma$-algebra generated by the subspace topology on $\mathcal W_\Delta$ (and likewise for $W+\mathcal W_\Delta$). In particular, both are standard Borel spaces, and the translation
%  \[
%  T_W:\ (\mathcal W_\Delta,\mathcal B(\mathcal W_\Delta))\ \to\ (W+\mathcal W_\Delta,\mathcal B(W+\mathcal W_\Delta)),\quad
%  T_W(\Delta)=W+\Delta
%  \]
%  is a bimeasurable bijection (a Borel isomorphism).



% \begin{lemma}[KL decomposition on the parameter space with a frozen backbone]
% \label{lemma:kl-fullspace-frozen}
% Let $W\in\mathbb R^d$ be fixed and let $\mathcal W_\Delta\subseteq\mathbb R^d$ be a linear subspace
% (endowed with the trace Borel $\sigma$-algebra $\mathcal B(\mathcal W_\Delta)$) that represents trainable
% updates. The actual parameters corresponding to an update $\Delta\in\mathcal W_\Delta$ live in the
% \emph{affine subspace} $W+\mathcal W_\Delta$ rather than in $\mathcal W_\Delta$ itself. To canonically lift
% probability measures from the update subspace to the full parameter space, define the \emph{translation map}
% $T_W:\mathcal W_\Delta\to\mathbb R^d$ by $T_W(\Delta)=W+\Delta$ and set
% \[
% P^{\mathrm{full}}:=(T_W)_\# P^{\Delta},\qquad
% Q^{\mathrm{full}}:=(T_W)_\# Q^{\Delta},
% \]
% where $P^{\Delta},Q^{\Delta}$ are probability measures on $(\mathcal W_\Delta,\mathcal B(\mathcal W_\Delta))$
% with $Q^{\Delta}\ll P^{\Delta}$, and $F_\#\mu$ denotes the \emph{pushforward} of a measure $\mu$ by a measurable
% map $F$, i.e., $(F_\#\mu)(A)=\mu(F^{-1}(A))$ for all measurable sets $A$.
% This construction places all on the affine subspace $W+\mathcal W_\Delta$, so
% $\operatorname{supp}(P^{\mathrm{full}}),\operatorname{supp}(Q^{\mathrm{full}})\subseteq W+\mathcal W_\Delta$.
% Moreover, absolute continuity is preserved under pushforward: $Q^{\Delta}\ll P^{\Delta}$ implies
% $Q^{\mathrm{full}}\ll P^{\mathrm{full}}$.

% \emph{Equivalent product-form view.} One may equivalently write
% $P^{\mathrm{full}}=\delta_W\otimes P^{\Delta}$ and $Q^{\mathrm{full}}=\delta_W\otimes Q^{\Delta}$ on the
% product space $\{W\}\times\mathcal W_\Delta$, where $\delta_W$ is the Dirac measure at $W$ and
% $\otimes$ denotes the product measure, and then identify $\{W\}\times\mathcal W_\Delta$ with $W+\mathcal W_\Delta$
% via $(w,\Delta)\mapsto w+\Delta$.

% Then $Q^{\mathrm{full}}\ll P^{\mathrm{full}}$ and
% \[
% \mathrm{KL}\!\left(Q^{\mathrm{full}}\middle\|P^{\mathrm{full}}\right)
% =\mathrm{KL}\!\left(Q^{\Delta}\middle\|P^{\Delta}\right),
% \]
% where $\mathrm{KL}(Q\|P)=\int \log\!\big(\tfrac{dQ}{dP}\big)\,dQ$ whenever $Q\ll P$.

% \emph{Proof.}
% The map $T_W$ is a measurable bijection from $\mathcal W_\Delta$ onto the affine subspace $W+\mathcal W_\Delta$
% (with inverse $T_{-W}$ on $W+\mathcal W_\Delta$), hence absolute continuity is preserved by pushforward and
% $Q^{\mathrm{full}}\ll P^{\mathrm{full}}$.
% By the Radon--Nikodym chain rule,
% $\frac{dQ^{\mathrm{full}}}{dP^{\mathrm{full}}}(T_W\Delta)=\frac{dQ^{\Delta}}{dP^{\Delta}}(\Delta)$
% for $P^{\Delta}$-a.e.\ $\Delta$.
% Changing variables $y=T_W\Delta$ yields
% $\mathrm{KL}(Q^{\mathrm{full}}\|P^{\mathrm{full}})=\mathrm{KL}(Q^{\Delta}\|P^{\Delta})$.
% \qed
% \end{lemma}

% \begin{lemma}[KL decomposition on the parameter space with a frozen backbone)]
% \label{lem:pecl-kl-decomposition}
% Let $W=W_0+\Delta W$ with $W_0\in\mathbb R^d$ fixed (frozen backbone) and $\Delta W\in\mathcal W_\Delta$ trainable.
% For task $t$, let $P_t^\Delta$ and $Q_t^\Delta$ be probability measures on $(\mathcal W_\Delta,\mathcal B(\mathcal W_\Delta))$
% with $Q_t^\Delta\ll P_t^\Delta$.
% Define full-space measures on $\{W_0\}\times\mathcal W_\Delta$ by
% \[
% P_t^{\mathrm{full}}:=\delta_{W_0}\otimes P_t^\Delta,\qquad
% Q_t^{\mathrm{full}}:=\delta_{W_0}\otimes Q_t^\Delta,
% \]
% where $\delta_{W_0}$ is the \emph{Dirac measure} at $W_0$ \emph{(unit mass at $W_0$)}, and $\mu\otimes\nu$ denotes
% the \emph{product measure} on the product measurable space. Identifying
% $\{W_0\}\times\mathcal W_\Delta$ with $W_0+\mathcal W_\Delta$ via $(w,\Delta)\mapsto w+\Delta$, we have
% \[
% \mathrm{KL}\!\left(Q_t^{\mathrm{full}}\middle\|P_t^{\mathrm{full}}\right)
% =\mathrm{KL}\!\left(Q_t^\Delta\middle\|P_t^\Delta\right).
% \]
% The same identity holds for hyper-distributions provided
% $\mathcal P_t^\Delta \ll \mathcal P_{t-1}^\Delta$ and
% $\mathcal P_t^{\mathrm{full}}=\delta_{W_0}\otimes \mathcal P_t^\Delta$,
% $\mathcal P_{t-1}^{\mathrm{full}}=\delta_{W_0}\otimes \mathcal P_{t-1}^\Delta$.

% \emph{Proof.}
% For product measures sharing the same first factor,
% \[
% \frac{d(\delta_{W_0}\otimes Q_t^\Delta)}{d(\delta_{W_0}\otimes P_t^\Delta)}(w,\Delta)
% =\frac{dQ_t^\Delta}{dP_t^\Delta}(\Delta)\quad \text{on } \{W_0\}\times\mathcal W_\Delta.
% \]
% Therefore,
% $\mathrm{KL}(Q_t^{\mathrm{full}}\|P_t^{\mathrm{full}})
% =\int \log\!\big(\tfrac{dQ_t^\Delta}{dP_t^\Delta}(\Delta)\big)\,dQ_t^\Delta(\Delta)
% =\mathrm{KL}(Q_t^\Delta\|P_t^\Delta)$.
% The hyper-distribution case is identical with $P_t^\Delta,Q_t^\Delta$ replaced by $\mathcal P_t^\Delta,\mathcal P_{t-1}^\Delta$.
% \qed
% \end{lemma}



















% Let $W\in\mathbb R^d$ be a fixed parameter vector and let $\mathcal W_\Delta\subseteq\mathbb R^d$ be a linear subspace (equipped with its Borel $\sigma$-algebra $\mathcal B(\mathcal W_\Delta)$) that represents trainable updates.
% Let $P^{\Delta W}$ and $Q^{\Delta W}$ be probability measures on $(\mathcal W_\Delta,\mathcal B(\mathcal W_\Delta))$ with $Q^{\Delta W}\ll P^{\Delta W}$.
% % \emph{(absolute continuity: $P^{\Delta W}(A)=0\Rightarrow Q^{\Delta W}(A)=0$ for all measurable $A$)}
% Define the translation map $T_W:\mathcal W_\Delta\to\mathbb R^d$ by $T_W(\Delta)=W+\Delta$ and set
% \[
% P^{\mathrm{full}}:=(T_W)_\# P^{\Delta W},\qquad
% Q^{\mathrm{full}}:=(T_W)_\# Q^{\Delta W},
% \]
% where $F_\#\mu$ denotes the \emph{pushforward} of a measure $\mu$ by a measurable map $F$, i.e., $(F_\#\mu)(A)=\mu(F^{-1}(A))$ for all measurable $A$.

% \begin{lemma}[KL decomposition on the parameter space with a frozen backbone]
% \label{lemma:kl-fullspace-frozen}
% Let $W\in\mathbb R^d$ be fixed and let $\mathcal W_\Delta\subseteq\mathbb R^d$ be a linear subspace
% (endowed with the trace Borel $\sigma$-algebra $\mathcal B(\mathcal W_\Delta)$) that represents trainable
% updates. The actual parameters corresponding to an update $\Delta\in\mathcal W_\Delta$ live in the
% \emph{affine subspace} $W+\mathcal W_\Delta$ rather than in $\mathcal W_\Delta$ itself. To canonically lift
% probability measures from the update subspace to the full parameter space, define the \emph{translation map}
% $T_W:\mathcal W_\Delta\to\mathbb R^d$ by $T_W(\Delta)=W+\Delta$ and set
% \[
% P^{\mathrm{full}}:=(T_W)_\# P^{\Delta W},\qquad
% Q^{\mathrm{full}}:=(T_W)_\# Q^{\Delta W},
% \]
% where $F_\#\mu$ denotes the \emph{pushforward} of a measure $\mu$ by a measurable map $F$, i.e.,
% $(F_\#\mu)(A)=\mu(F^{-1}(A))$ for all measurable sets $A$.
% This construction places all mass on the affine subspace $W+\mathcal W_\Delta$, so
% $\operatorname{supp}(P^{\mathrm{full}}),\operatorname{supp}(Q^{\mathrm{full}})\subseteq W+\mathcal W_\Delta$.
% Moreover, absolute continuity is preserved under pushforward: if $Q^{\Delta W}\ll P^{\Delta W}$, then
% $Q^{\mathrm{full}}\ll P^{\mathrm{full}}$.

% \emph{Equivalent product-form view.} One may equivalently write
% $P^{\mathrm{full}}=\delta_W\otimes P^{\Delta W}$ and $Q^{\mathrm{full}}=\delta_W\otimes Q^{\Delta W}$ on the
% product space $\{W\}\times\mathcal W_\Delta$, where $\delta_W$ is the Dirac measure at $W$ and
% $\otimes$ denotes the product measure, and identify $\{W\}\times\mathcal W_\Delta$ with $W+\mathcal W_\Delta$
% via $(w,\Delta)\mapsto w+\Delta$.

% Then $Q^{\mathrm{full}}\ll P^{\mathrm{full}}$ and
% \[
% \mathrm{KL}\!\left(Q^{\mathrm{full}}\middle\|P^{\mathrm{full}}\right)
% =\mathrm{KL}\!\left(Q^{\Delta W}\middle\|P^{\Delta W}\right),
% \]
% where $\mathrm{KL}(Q\|P)=\int \log\!\big(\tfrac{dQ}{dP}\big)\,dQ$ whenever $Q\ll P$.

% \emph{Proof.}
% The map $T_W$ is a measurable bijection from $\mathcal W_\Delta$ onto the affine subspace $W+\mathcal W_\Delta$ with inverse $T_{-W}$, hence absolute continuity is preserved by pushforward and $Q^{\mathrm{full}}\ll P^{\mathrm{full}}$.
% By the Radon--Nikodym chain rule,
% $\frac{dQ^{\mathrm{full}}}{dP^{\mathrm{full}}}(T_W\Delta)=\frac{dQ^{\Delta W}}{dP^{\Delta W}}(\Delta)$ for $P^{\Delta W}$-a.e.\ $\Delta$.
% Changing variables $y=T_W\Delta$ yields
% $\mathrm{KL}(Q^{\mathrm{full}}\|P^{\mathrm{full}})=\mathrm{KL}(Q^{\Delta W}\|P^{\Delta W})$.
% \qed
% \end{lemma}

% \begin{lemma}[KL decomposition in PECL]
% \label{lem:pecl-kl-decomposition}
% Let $W=W_0+\Delta W$ with $W_0\in\mathbb R^d$ fixed (frozen backbone) and $\Delta W\in\mathcal W_\Delta$ trainable.
% For task $t$, let $P_t^\Delta$ and $Q_t^\Delta$ be probability measures on $(\mathcal W_\Delta,\mathcal B(\mathcal W_\Delta))$ with $Q_t^\Delta\ll P_t^\Delta$.
% Define full-space measures on $\{W_0\}\times\mathcal W_\Delta$ by
% \[
% P_t^{\mathrm{full}}:=\delta_{W_0}\otimes P_t^\Delta,\qquad
% Q_t^{\mathrm{full}}:=\delta_{W_0}\otimes Q_t^\Delta,
% \]
% where $\delta_{W_0}$ is the \emph{Dirac measure} at $W_0$ \emph{(unit mass at $W_0$)}, and $\mu\otimes\nu$ denotes the \emph{product measure} of measures $\mu$ and $\nu$ on the product space.
% Then
% \[
% \mathrm{KL}\!\left(Q_t^{\mathrm{full}}\middle\|P_t^{\mathrm{full}}\right)
% =\mathrm{KL}\!\left(Q_t^\Delta\middle\|P_t^\Delta\right).
% \]
% The same identity holds for hyper-distributions if
% $\mathcal P_t^{\mathrm{full}}=\delta_{W_0}\otimes \mathcal P_t^\Delta$ and
% $\mathcal P_{t-1}^{\mathrm{full}}=\delta_{W_0}\otimes \mathcal P_{t-1}^\Delta$.

% \emph{Proof.}
% For product measures sharing the same first factor,
% \[
% \frac{d(\delta_{W_0}\otimes Q_t^\Delta)}{d(\delta_{W_0}\otimes P_t^\Delta)}(w,\Delta)
% =\frac{dQ_t^\Delta}{dP_t^\Delta}(\Delta)\quad \text{on } \{W_0\}\times\mathcal W_\Delta.
% \]
% Therefore,
% $\mathrm{KL}(Q_t^{\mathrm{full}}\|P_t^{\mathrm{full}})
% =\int \log\!\big(\tfrac{dQ_t^\Delta}{dP_t^\Delta}(\Delta)\big)\,dQ_t^\Delta(\Delta)
% =\mathrm{KL}(Q_t^\Delta\|P_t^\Delta)$.
% The hyper-distribution case is identical with $P_t^\Delta,Q_t^\Delta$ replaced by $\mathcal P_t^\Delta,\mathcal P_{t-1}^\Delta$.
% \qed
% \end{lemma}



% \begin{lemma}[KL decomposition on the parameter space with a frozen backbone]
% \label{lemma:kl-fullspace-frozen}
% Let \(W \in \mathbb{R}^d\) be a deterministic parameter vector, and let \(\mathcal{W}_\Delta \subseteq \mathbb{R}^d\) be a subspace representing parameter updates (e.g., the low-rank adaptation space). Consider two probability measures, \(P^{\Delta W}\) and \(Q^{\Delta W}\), defined on the measurable space \((\mathcal{W}_\Delta, \mathcal{B}(\mathcal{W}_\Delta))\), \(Q^{\Delta W} \ll P^{\Delta W}\).
% Let \(P^{\mathrm{full}}\) and \(Q^{\mathrm{full}}\) be the probability measures on the full parameter space \((\mathbb{R}^d, \mathcal{B}(\mathbb{R}^d))\) via the translation map \(T_W\). $P^{\mathrm{full}} \coloneqq (T_W)_{\#} P^{\Delta W}, \quad Q^{\mathrm{full}} \coloneqq (T_W)_{\#} Q^{\Delta W}.$
% Then, \(Q^{\mathrm{full}} \ll P^{\mathrm{full}}\), and
% $$
% \mathrm{KL}\left(Q^{\mathrm{full}} \,\middle\|\, P^{\mathrm{full}}\right) = \mathrm{KL}\left(Q^{\Delta W} \,\middle\|\, P^{\Delta W}\right).
% $$
% \end{lemma}


% \begin{lemma}[KL Decomposition in PECL]
% \label{lem:pecl-kl-decomposition}
% Let the full parameter space decompose as $W = W_0 + \Delta W$, where $W_0 \in \mathbb{R}^d$ is frozen and $\Delta W \in \mathcal{W}_\Delta \subseteq \mathbb{R}^d$ is the trainable adapter subspace. Define the full-space probability measures as:
% \[
% P_t^{\mathrm{full}} = \delta_{W_0} \otimes P_t^{\Delta}, \quad Q_t^{\mathrm{full}} = \delta_{W_0} \otimes Q_t^{\Delta},
% \]
% where $\delta_{W_0}$ is the Dirac measure at $W_0$, and $P_t^{\Delta}, Q_t^{\Delta}$ are probability measures on $\mathcal{W}_\Delta$. Then:
% \[
% \mathrm{KL}(Q_t^{\mathrm{full}} \| P_t^{\mathrm{full}}) = \mathrm{KL}(Q_t^{\Delta} \| P_t^{\Delta}).
% \]
% Moreover, for the hyper-distributions:
% \[
% \mathrm{KL}(\mathcal{P}_t^{\mathrm{full}} \| \mathcal{P}_{t-1}^{\mathrm{full}}) = \mathrm{KL}(\mathcal{P}_t^{\Delta} \| \mathcal{P}_{t-1}^{\Delta}).
% \]
% \end{lemma}

% \begin{proof}
% For the parameter-level equality, since $P_t^{\mathrm{full}}$ and $Q_t^{\mathrm{full}}$ are product measures with identical Dirac components:
% \[
% \mathrm{KL}(Q_t^{\mathrm{full}} \| P_t^{\mathrm{full}}) = \mathrm{KL}(\delta_{W_0} \| \delta_{W_0}) + \mathrm{KL}(Q_t^{\Delta} \| P_t^{\Delta}) = \mathrm{KL}(Q_t^{\Delta} \| P_t^{\Delta}),
% \]
% where $\mathrm{KL}(\delta_{W_0} \| \delta_{W_0}) = 0$ because identical Dirac measures have zero KL divergence.

% For the hyper-distribution level, the same product structure applies: if $\mathcal{P}_t^{\mathrm{full}} = \delta_{W_0} \otimes \mathcal{P}_t^{\Delta}$ and $\mathcal{P}_{t-1}^{\mathrm{full}} = \delta_{W_0} \otimes \mathcal{P}_{t-1}^{\Delta}$, then:
% \[
% \mathrm{KL}(\mathcal{P}_t^{\mathrm{full}} \| \mathcal{P}_{t-1}^{\mathrm{full}}) = \mathrm{KL}(\delta_{W_0} \| \delta_{W_0}) + \mathrm{KL}(\mathcal{P}_t^{\Delta} \| \mathcal{P}_{t-1}^{\Delta}) = \mathrm{KL}(\mathcal{P}_t^{\Delta} \| \mathcal{P}_{t-1}^{\Delta}).
% \]
% \end{proof}

% \begin{lemma}[Sharpness Decomposition in PECL]
% \label{lem:pecl-sharpness-decomposition}
% Under the same PECL decomposition $W = W_0 + \Delta W$,
% the corresponding full-space posterior is:
% $Q_t^{\mathrm{full}} = \delta_{W_0} \otimes Q_t^{\Delta}$
% where 
%  a Gaussian posterior on the adapter subspace is 
% $Q_t^{\Delta} = \mathcal{N}(\hat{\Delta W}_t, \Sigma_t)$.
% Then the expected sharpness decomposes as
% \[
% \mathbb{E}_{W \sim Q_t^{\mathrm{full}}}[\hat{L}_{S_t}(W)] =\mathbb{E}_{\varepsilon \sim \mathcal{N}(0, \Sigma_t)}[\hat{L}_{S_t}(W_0 + \hat{\Delta W}_t + \varepsilon)= \hat{L}_{S_t}(W_0 + \hat{\Delta W}_t) + \mathrm{E\text{-}Sh}_t^{\Delta}(\hat{\Delta W}_t, \Sigma_t),
% \]
% where the subspace expected sharpness is:
% \[
% \mathrm{E\text{-}Sh}_t^{\Delta}(\hat{\Delta W}_t, \Sigma_t) = \mathbb{E}_{\varepsilon \sim \mathcal{N}(0, \Sigma_t)}[\hat{L}_{S_t}(W_0 + \hat{\Delta W}_t + \varepsilon) - \hat{L}_{S_t}(W_0 + \hat{\Delta W}_t)].
% \]
% \end{lemma}

% \begin{proof}
% Since $Q_t^{\mathrm{full}}$ is supported on the affine subspace $W_0 + \mathcal{W}_\Delta$, we have:
% \[
% \mathbb{E}_{W \sim Q_t^{\mathrm{full}}}[\hat{L}_{S_t}(W)] = \mathbb{E}_{\Delta W \sim Q_t^{\Delta}}[\hat{L}_{S_t}(W_0 + \Delta W)].
% \]

% For Gaussian $Q_t^{\Delta} = \mathcal{N}(\hat{\Delta W}_t, \Sigma_t)$, parameterize $\Delta W = \hat{\Delta W}_t + \varepsilon$ with $\varepsilon \sim \mathcal{N}(0, \Sigma_t)$:
% \[
% \mathbb{E}_{\Delta W \sim Q_t^{\Delta}}[\hat{L}_{S_t}(W_0 + \Delta W)] = \mathbb{E}_{\varepsilon \sim \mathcal{N}(0, \Sigma_t)}[\hat{L}_{S_t}(W_0 + \hat{\Delta W}_t + \varepsilon)].
% \]

% Adding and subtracting $\hat{L}_{S_t}(W_0 + \hat{\Delta W}_t)$:
% \[
% \mathbb{E}_{\varepsilon \sim \mathcal{N}(0, \Sigma_t)}[\hat{L}_{S_t}(W_0 + \hat{\Delta W}_t + \varepsilon)] = \hat{L}_{S_t}(W_0 + \hat{\Delta W}_t) + \mathrm{E\text{-}Sh}_t^{\Delta}(\hat{\Delta W}_t, \Sigma_t),
% \]
% which completes the proof.
% \end{proof}

\section{Appendix B: Proof of Lemma~\ref{lemma: subspace-full}}
\begin{lemma}[Sharpness decomposition in PECL]
\label{lemma: Sharpness in PECL}
Let $W_0\in\mathbb R^d$ be fixed and let the trainable update lie in the adapter subspace $\mathcal W_\Delta\subseteq\mathbb R^d$. Consider a Gaussian posterior on $(\mathcal W_\Delta,\mathcal B(\mathcal W_\Delta))$,
 $
 Q_t^{\Delta}=\mathcal N(\hat{\Delta W}_t,\Sigma_t),
 $
 with mean $\hat{\Delta W}_t\in\mathcal W_\Delta$ and covariance operator $\Sigma_t\succeq 0$ acting on $\mathcal W_\Delta$ (viewed in $\mathbb R^d$, this covariance is supported in $\mathcal W_\Delta$). Define the full-space posterior on ${W_0}\times\mathcal W_\Delta$ by
 $
 Q_t^{\mathrm{full}}:=\delta_{W_0}\otimes Q_t^{\Delta},
 $
and identify ${W_0}\times\mathcal W_\Delta$ with the affine subspace $W_0+\mathcal W_\Delta$ via $(w,\Delta)\mapsto w+\Delta$. Let $\hat L_{S_t}:\mathbb R^d\to\mathbb R$ be measurable and integrable under $Q_t^{\mathrm{full}}$. Then
 \[
 \mathbb{E}_{W\sim Q_t^{\mathrm{full}}}\big[\hat{L}_{S_t}(W)\big]
 = \mathbb{E}_{\varepsilon\sim\mathcal N(0,\Sigma_t)}
 \big[\hat{L}_{S_t}(W_0+\hat{\Delta W}_t+\varepsilon)\big]
 = \hat{L}_{S_t}(W_0+\hat{\Delta W}_t) +\mathrm{E\text{-}Sh}_t^{\Delta}(\hat{\Delta W}_t,\Sigma_t) \]
   where the subspace expected sharpness is
   \[
   \mathrm{E\text{-}Sh}_t^{\Delta}(\hat{\Delta W}_t,\Sigma_t)
   = \mathbb{E}_{\varepsilon\sim\mathcal N(0,\Sigma_t)}
   \big[\hat{L}_{S_t}(W_0+\hat{\Delta W}_t+\varepsilon)
  - \hat{L}_{S_t}(W_0+\hat{\Delta W}_t)\big].
     \]
\end{lemma}
\begin{proof}
    Sampling $W\sim Q_t^{\mathrm{full}}$ is equivalent to sampling $\varepsilon\sim\mathcal N(0,\Sigma_t)$ on $\mathcal W_\Delta$ and setting $W=W_0+\hat{\Delta W}_t+\varepsilon$. Taking expectations and adding–subtracting $\hat L_{S_t}(W_0+\hat{\Delta W}_t)$ inside the expectation gives the decomposition. 
\end{proof}


% \begin{lemma}[Sharpness decomposition in PECL]
% \label{lem:pecl-sharpness-decomposition}
% Let $W=W_0+\Delta W$ with $W_0\in\mathbb R^d$ fixed (frozen backbone) and $\Delta W\in\mathcal W_\Delta$
% lying in the trainable adapter subspace. Define the full-space posterior on $\{W_0\}\times\mathcal W_\Delta$ by
% \[
% Q_t^{\mathrm{full}}:=\delta_{W_0}\otimes Q_t^{\Delta},
% \]
% where $\delta_{W_0}$ is the Dirac measure at $W_0$, and $Q_t^{\Delta}=\mathcal N(\hat{\Delta W}_t,\Sigma_t)$
% is a Gaussian measure on $(\mathcal W_\Delta,\mathcal B(\mathcal W_\Delta))$ with mean
% $\hat{\Delta W}_t\in\mathcal W_\Delta$ and covariance $\Sigma_t\succeq 0$ acting on $\mathcal W_\Delta$.
% Here \(\Sigma_t \succeq 0\) denotes a positive semidefinite covariance operator on the adapter subspace \(\mathcal W_\Delta\) (i.e., \(v^\top\Sigma_t v\ge 0\) for all \(v\in\mathcal W_\Delta)\); when represented in the ambient \(\mathbb R^d\), its range lies in \(\mathcal W_\Delta\) and the resulting covariance is typically singular.
% Identifying $\{W_0\}\times\mathcal W_\Delta$ with the affine subspace $W_0+\mathcal W_\Delta$ via
% $(w,\Delta)\mapsto w+\Delta$, we equivalently have the random representation
% \[
% W \ \stackrel{\text{distribution}}{=}\ W_0+\hat{\Delta W}_t+\varepsilon,\qquad
% \varepsilon\sim\mathcal N(0,\Sigma_t)\ \text{on } \mathcal W_\Delta.
% \]
% Assume $\hat L_{S_t}:\mathbb R^d\to\mathbb R$ is measurable and integrable under $Q_t^{\mathrm{full}}$.
% Then the posterior expected empirical loss decomposes as
% \[
% \mathbb{E}_{W\sim Q_t^{\mathrm{full}}}\big[\hat{L}_{S_t}(W)\big]
% = \mathbb{E}_{\varepsilon\sim\mathcal N(0,\Sigma_t)}
%     \big[\hat{L}_{S_t}(W_0+\hat{\Delta W}_t+\varepsilon)\big]
% = \hat{L}_{S_t}(W_0+\hat{\Delta W}_t)
%   + \mathrm{E\text{-}Sh}_t^{\Delta}(\hat{\Delta W}_t,\Sigma_t),
% \]
% where the (subspace) expected sharpness is
% \[
% \mathrm{E\text{-}Sh}_t^{\Delta}(\hat{\Delta W}_t,\Sigma_t)
% := \mathbb{E}_{\varepsilon\sim\mathcal N(0,\Sigma_t)}
%    \big[\hat{L}_{S_t}(W_0+\hat{\Delta W}_t+\varepsilon)
%        - \hat{L}_{S_t}(W_0+\hat{\Delta W}_t)\big].
% \]
% \end{lemma}

% \begin{proof}
% By construction $Q_t^{\mathrm{full}}=\delta_{W_0}\otimes\mathcal N(\hat{\Delta W}_t,\Sigma_t)$ on
% $\{W_0\}\times\mathcal W_\Delta$ and the identification $(w,\Delta)\mapsto w+\Delta$,
% sampling $W\sim Q_t^{\mathrm{full}}$ is equivalent to sampling
% $\varepsilon\sim\mathcal N(0,\Sigma_t)$ on $\mathcal W_\Delta$ and setting
% $W=W_0+\hat{\Delta W}_t+\varepsilon$. Therefore
% \[
% \mathbb{E}_{W\sim Q_t^{\mathrm{full}}}[\hat L_{S_t}(W)]
% = \mathbb{E}_{\varepsilon}[\hat L_{S_t}(W_0+\hat{\Delta W}_t+\varepsilon)].
% \]
% Adding and subtracting $\hat L_{S_t}(W_0+\hat{\Delta W}_t)$ inside the expectation gives the claimed
% decomposition with $\mathrm{E\text{-}Sh}_t^{\Delta}$.
% \end{proof}


%==============================
% PECL / LoRA specialization
%==============================
%==============================
% Clarified construction: full-space hyperposterior → adapter reduction
%==============================

% \begin{lemma}[Full-to-adapter KL invariance under a frozen backbone]
% \label{lem:full-to-adapter-kl}
% Fix a backbone $W_0\in\mathbb{R}^d$ and adapter subspace $\mathcal W_\Delta\subseteq\mathbb{R}^d$.
% Define the measurable embedding on \emph{parameters}
% \[
% \iota_\theta:\ \mathcal W_\Delta \to \mathbb{R}^d,\qquad \iota_\theta(\Delta W):=W_0+\Delta W,
% \]
% and the induced embedding on \emph{probability measures}
% \[
% \iota:\ \mathcal P(\mathcal W_\Delta)\to \mathcal P(\mathbb{R}^d),\qquad
% \iota(P^\Delta):=(\iota_\theta)_{\#}P^\Delta=\delta_{W_0}\otimes P^\Delta.
% \]
% Let $Q^\Delta\ll P^\Delta$ be probabilities on $(\mathcal W_\Delta,\mathcal B(\mathcal W_\Delta))$. Then
% \[
% \mathrm{KL}(Q^{\mathrm{full}}\|P^{\mathrm{full}})
% \;=\;
% \mathrm{KL}\!\big(\,Q^\Delta\ \big\|\ P^\Delta\,\big).
% \]
% Moreover, if $\mathcal P_t^{\mathrm{full}}$ and $\mathcal P_{t-1}^{\mathrm{full}}$ are probability measures on the \emph{full prior space} whose supports lie in the embedded manifold
% $\mathcal M:=\{\iota(P^\Delta):\,P^\Delta\in\mathcal P(\mathcal W_\Delta)\}$, and we define the adapter-level hyperdistributions by the measurable inverse
% $\pi_\Delta:\mathcal M\to\mathcal P(\mathcal W_\Delta)$ as
% $\mathcal P_t^\Delta:=(\pi_\Delta)_{\#}\mathcal P_t^{\mathrm{full}}$ and
% $\mathcal P_{t-1}^\Delta:=(\pi_\Delta)_{\#}\mathcal P_{t-1}^{\mathrm{full}}$,
% then
% \[
% \mathrm{KL}\!\big(\,\mathcal P_t^{\mathrm{full}}\ \big\|\ \mathcal P_{t-1}^{\mathrm{full}}\,\big)
% \;=\;
% \mathrm{KL}\!\big(\,\mathcal P_t^\Delta\ \big\|\ \mathcal P_{t-1}^\Delta\,\big).
% \]
% \end{lemma}

% \begin{proof}[Proof]
% For the parameter-level statement, write $P^{\mathrm{full}}:=\iota(P^\Delta)=\delta_{W_0}\otimes P^\Delta$ and $Q^{\mathrm{full}}:=\iota(Q^\Delta)$. Since $\iota_\theta$ is a measurable bijection between $(\mathcal W_\Delta,\mathcal B(\mathcal W_\Delta))$ and its image, the Radon–Nikodym derivative transforms as
% \[
% \frac{dQ^{\mathrm{full}}}{dP^{\mathrm{full}}}(W_0+\Delta W)
% \;=\;
% \frac{dQ^\Delta}{dP^\Delta}(\Delta W)\quad\text{for $P^\Delta$-a.e.\ $\Delta W$}.
% \]
% Thus, by definition of KL divergence and change-of-variables,
% \[
% \mathrm{KL}(Q^{\mathrm{full}}\|P^{\mathrm{full}})
% =\int \log\!\Big(\frac{dQ^{\mathrm{full}}}{dP^{\mathrm{full}}}\Big)\,dQ^{\mathrm{full}}
% =\int \log\!\Big(\frac{dQ^\Delta}{dP^\Delta}\Big)\,dQ^\Delta
% =\mathrm{KL}(Q^\Delta\|P^\Delta).
% \]
% For the hyper-level statement, the same argument applies on the prior space: $\iota$ is a measurable isomorphism from $\mathcal P(\mathcal W_\Delta)$ onto its image $\mathcal M$, with inverse $\pi_\Delta$. KL divergence is invariant under such measurable bijections, hence
% $\mathrm{KL}(\mathcal P_t^{\mathrm{full}}\|\mathcal P_{t-1}^{\mathrm{full}})
% =\mathrm{KL}(\mathcal P_t^\Delta\|\mathcal P_{t-1}^\Delta)$.
% \end{proof}



% \paragraph{Full-space hyperposterior.}
% At step $t$, let the hyperposterior be a probability measure on full-model priors,
% \[
% \mathcal P_t^{\mathrm{full}} \quad\text{on}\quad \mathsf{Prior}^{\mathrm{full}}:=\mathcal P(\mathbb R^d),
% \]
% and let the task-$t$ parameter posterior be the conditional law
% $Q_t^{\mathrm{full}}(\,\cdot\,\mid P^{\mathrm{full}})$ on $\mathbb R^d$.
% Define the \emph{after-training} joint law and the \emph{reference} joint law by
% \[
% \widehat\Pi_t^{\mathrm{full}}(dP^{\mathrm{full}},dW^{\mathrm{full}})
% := \mathcal P_t^{\mathrm{full}}(dP^{\mathrm{full}})\,Q_t^{\mathrm{full}}(dW^{\mathrm{full}}\mid P^{\mathrm{full}}),
% \qquad
% \Pi_{t-1}^{\mathrm{full}}(dP^{\mathrm{full}},dW^{\mathrm{full}})
% := \mathcal P_{t-1}^{\mathrm{full}}(dP^{\mathrm{full}})\,P^{\mathrm{full}}(dW^{\mathrm{full}}).
% \]

% \paragraph{PECL support restriction.}
% In PECL, the full parameter decomposes as $W^{\mathrm{full}}=W_0+\Delta W_t$ with $W_0$ frozen.
% Consequently, both the \emph{prior} and the \emph{posterior} at task $t$ live on the embedded manifold
% \[
% \mathcal M_t \;:=\; \big\{\;\delta_{W_0}\otimes R^\Delta \;:\; R^\Delta\in\mathcal P(\mathcal W_t)\;\big\}\ \subset\ \mathcal P(\mathbb R^d),
% \]
% and factor as
% \[
% P^{\mathrm{full}}=\delta_{W_0}\otimes P^\Delta,\qquad
% Q_t^{\mathrm{full}}=\delta_{W_0}\otimes Q_t^\Delta,
% \]
% with $P^\Delta,Q_t^\Delta$ probability measures on the LoRA subspace $\mathcal W_t$.
% Let $\iota:\mathcal P(\mathcal W_t)\to \mathcal P(\mathbb R^d)$ denote the embedding
% $\iota(P^\Delta):=\delta_{W_0}\otimes P^\Delta$ and let
% $\pi_\Delta:\mathcal M_t\to \mathcal P(\mathcal W_t)$ be its (measurable) inverse on $\mathcal M_t$.

% \paragraph{Induced adapter-level hyperdistributions.}
% Define the adapter-level hyperdistributions as pushforwards of the full-space hyperdistributions:
% \[
% \mathcal P_t^\Delta \;:=\; (\pi_\Delta)_{\#}\,\mathcal P_t^{\mathrm{full}}\Big|_{\mathcal M_t},
% \qquad
% \mathcal P_{t-1}^\Delta \;:=\; (\pi_\Delta)_{\#}\,\mathcal P_{t-1}^{\mathrm{full}}\Big|_{\mathcal M_t}.
% \]
% Then the \emph{adapter-form} joint laws are
% \[
% \widehat\Pi_t(dP^\Delta,d\Delta W)
% := \mathcal P_t^\Delta(dP^\Delta)\,Q_t^\Delta(d\Delta W\mid P^\Delta),\qquad
% \Pi_{t-1}(dP^\Delta,d\Delta W)
% := \mathcal P_{t-1}^\Delta(dP^\Delta)\,P^\Delta(d\Delta W).
% \]

% \paragraph{KL reduction (hyper- and parameter-level).}
% Because $\iota$ is injective and $\pi_\Delta$ is its inverse on $\mathcal M_t$, Kullback–Leibler
% divergence is invariant under this measurable isomorphism onto its image. Hence:
% \begin{align*}
% \mathrm{KL}\!\big(\mathcal P_t^{\mathrm{full}}\big\|\mathcal P_{t-1}^{\mathrm{full}}\big)
% \;=\;
% \mathrm{KL}\!\big(\mathcal P_t^\Delta\big\|\mathcal P_{t-1}^\Delta\big),
% \qquad
% \mathrm{KL}\!\big(Q_t^{\mathrm{full}}\big\|P^{\mathrm{full}}\big)
% \;=\;
% \mathrm{KL}\!\big(Q_t^\Delta\big\|P^\Delta\big).
% \end{align*}
% Therefore, the joint KL chain rule on the \emph{full} space,
% \[
% \mathrm{KL}\!\big(\widehat\Pi_t^{\mathrm{full}}\big\|\Pi_{t-1}^{\mathrm{full}}\big)
% =
% \mathrm{KL}\!\big(\mathcal P_t^{\mathrm{full}}\big\|\mathcal P_{t-1}^{\mathrm{full}}\big)
% +\mathbb E_{P^{\mathrm{full}}\sim \mathcal P_t^{\mathrm{full}}}
% \mathrm{KL}\!\big(Q_t^{\mathrm{full}}\big\|P^{\mathrm{full}}\big),
% \]
% reduces exactly to the adapter-level identity
% \[
% \boxed{\;
% \mathrm{KL}\!\big(\widehat\Pi_t\big\|\Pi_{t-1}\big)
% =
% \mathrm{KL}\!\big(\mathcal P_t^\Delta\big\|\mathcal P_{t-1}^\Delta\big)
% +\mathbb E_{P^\Delta\sim \mathcal P_t^\Delta}
% \mathrm{KL}\!\big(Q_t^\Delta\big\|P^\Delta\big).
% \;}
% \]

% \paragraph{Final form used in the bound.}
% With the above reduction, all expectations “under the joint law at step $t$” can be taken as
% \[
% \widehat\Pi_t(dP,dW)\equiv \widehat\Pi_t(dP^\Delta,d\Delta W)
% =\mathcal P_t^\Delta(dP^\Delta)\,Q_t^\Delta(d\Delta W\mid P^\Delta),
% \]
% and the complexity term
% $\mathrm{KL}(\widehat\Pi_t\|\Pi_{t-1})$ appearing in the bound equals the sum of
% (i) the hyper-drift in the adapter space, and (ii) the adapter-level task KL.
% In particular, the parameter-level KL does \emph{not} pay for directions orthogonal to $\mathcal W_t$,
% reflecting the PECL constraint.





% Fix bounded loss $\ell\in[0,1]$. In parameter-efficient continual learning (PECL),
% let the full parameter at task $t$ decompose as
% \[
% W^{\mathrm{full}} \;=\; W_0 \;+\; \Delta W_t,\qquad W_0\ \text{frozen},\ \ \Delta W_t\in\mathcal W_t.
% \]
% Put joint laws on the full space by factoring the frozen and adaptive parts:
% \[
% P^{\mathrm{full}}_t \;:=\; \delta_{W_0}\otimes P_t^{\Delta},\qquad
% Q^{\mathrm{full}}_t \;:=\; \delta_{W_0}\otimes Q_t^{\Delta},
% \]
% and define the hyperposterior over adapter priors $\mathcal P_t$ so that the
% after-training joint law at step $t$ is
% $\widehat\Pi_t(dP^\Delta,d\Delta W):=\mathcal P_t(dP^\Delta)\,Q_t^\Delta(d\Delta W\mid P^\Delta)$,
% with reference joint $\Pi_{t-1}(dP^\Delta,d\Delta W):=\mathcal P_{t-1}(dP^\Delta)\,P^\Delta(d\Delta W)$.
% Assume independent task datasets $S_t=\{Z_{tj}\}_{j=1}^{m_t}\overset{\text{i.i.d.}}{\sim}\mu_t$ and set
% $M_T:=\sum_{t=1}^T(m_t-1)$. Let $Q_t^\Delta=\mathcal N(\hat\Delta W_t,\Sigma_t)$ with
% $\mathrm{range}(\Sigma_t)\subseteq\mathcal W_t$ and define the subspace sharpness
% \[
% \mathrm{E\text{-}Sh}_t(\hat W_t,\Sigma_t)
% := \mathbb{E}_{\varepsilon\sim\mathcal N(0,\Sigma_t)}
% \Big[\widehat L_{S_t}\big(W_0+\hat\Delta W_t+\varepsilon\big)-\widehat L_{S_t}(W_0+\hat\Delta W_t)\Big],
% \quad \hat W_t:=W_0+\hat\Delta W_t.
% \]
% Then 


% \begin{lemma}[Subspace restriction of expected sharpness under PECL]
% \label{lem:esh-subspace}
% Let the full parameter decompose as $W^{\mathrm{full}} = W_0 + \Delta W$ with a frozen backbone $W_0$ and an adapter subspace $\mathcal W_\Delta \subset \mathbb R^d$. Consider the affine manifold
% \[
% \mathcal M \;:=\; W_0 + \mathcal W_\Delta \;=\; \{\;W_0+\Delta W \;:\; \Delta W\in\mathcal W_\Delta\;\}.
% \]
% Assume the \emph{full prior} at task $t$ is supported on $\mathcal M$, i.e.,
% $P_t^{\mathrm{full}}=\delta_{W_0}\otimes P_t^\Delta$ with $P_t^\Delta\in\mathcal P(\mathcal W_\Delta)$.
% Let $\widehat L_{S_t}$ be any bounded empirical risk. Then:

% \smallskip
% \noindent\textbf{(i) Support constraint from finite KL.}
% If a full posterior $Q_t^{\mathrm{full}}$ has finite divergence
% $\mathrm{KL}(Q_t^{\mathrm{full}}\|P_t^{\mathrm{full}})<\infty$, then
% $Q_t^{\mathrm{full}}$ must also be supported on $\mathcal M$ and can be written as the pushforward
% \[
% Q_t^{\mathrm{full}} \;=\; (\iota_\theta)_{\#} Q_t^\Delta,
% \qquad
% \iota_\theta:\ \mathcal W_\Delta \to \mathcal M,\ \ \iota_\theta(\Delta W):=W_0+\Delta W,
% \]
% for some $Q_t^\Delta\in\mathcal P(\mathcal W_\Delta)$.

% \smallskip
% \noindent\textbf{(ii) Equality of full- and subspace-expected losses.}
% For any $Q_t^\Delta$ as in \textup{(i)},
% \[
% \mathbb E_{W\sim Q_t^{\mathrm{full}}}\!\left[\widehat L_{S_t}(W)\right]
% \;=\;
% \mathbb E_{\Delta W\sim Q_t^\Delta}\!\left[\widehat L_{S_t}(W_0+\Delta W)\right].
% \]

% \smallskip
% \noindent\textbf{(iii) Gaussian posterior and subspace E--Sh.}
% If $Q_t^\Delta=\mathcal N(\hat{\Delta W}_t,\Sigma_t)$ on $\mathcal W_\Delta$ with
% $\mathrm{range}(\Sigma_t)\subseteq\mathcal W_\Delta$, and
% $Q_t^{\mathrm{full}}=(\iota_\theta)_{\#}Q_t^\Delta$ (equivalently,
% $Q_t^{\mathrm{full}}=\mathcal N(W_0+\hat{\Delta W}_t,\Sigma_{\mathrm{full}})$ with
% $\Sigma_{\mathrm{full}}=\Pi_\Delta \Sigma_t \Pi_\Delta^\top$ and $\Pi_\Delta$ the orthogonal projector onto $\mathcal W_\Delta$), then
% \[
% \mathbb E_{W\sim Q_t^{\mathrm{full}}}\!\left[\widehat L_{S_t}(W)\right]
% =
% \widehat L_{S_t}(W_0+\hat{\Delta W}_t)
% +
% \mathrm{E\text{-}Sh}^{\Delta}_t(\hat{\Delta W}_t,\Sigma_t),
% \]
% where the \emph{subspace expected sharpness} is
% \[
% \mathrm{E\text{-}Sh}^{\Delta}_t(\hat{\Delta W}_t,\Sigma_t)
% :=
% \mathbb E_{\varepsilon\sim \mathcal N(0,\Sigma_t)}
% \!\Big[\widehat L_{S_t}(W_0+\hat{\Delta W}_t+\varepsilon)-\widehat L_{S_t}(W_0+\hat{\Delta W}_t)\Big].
% \]
% In particular, the usual full-space E--Sh with covariance $\Sigma_{\mathrm{full}}$ coincides \emph{exactly} with the LoRA-subspace E--Sh above.
% \end{lemma}

% \begin{proof}[Proof]
% \textit{(i)} Since $P_t^{\mathrm{full}}$ is supported on $\mathcal M$, any $Q_t^{\mathrm{full}}$ with finite $\mathrm{KL}(Q_t^{\mathrm{full}}\|P_t^{\mathrm{full}})$ must satisfy $Q_t^{\mathrm{full}}\ll P_t^{\mathrm{full}}$ and hence be supported on $\mathcal M$ as well. The map $\iota_\theta:\mathcal W_\Delta\to\mathcal M$, $\iota_\theta(\Delta W)=W_0+\Delta W$, is a measurable bijection with inverse $\pi_\Delta:\mathcal M\to\mathcal W_\Delta$, $\pi_\Delta(W_0+\Delta W)=\Delta W$. Define $Q_t^\Delta:=(\pi_\Delta)_{\#}Q_t^{\mathrm{full}}$ to get $Q_t^{\mathrm{full}}=(\iota_\theta)_{\#}Q_t^\Delta$.

% \textit{(ii)} For any measurable $h$ on $\mathcal M$, change of variables under pushforward yields
% $\mathbb E_{W\sim(\iota_\theta)_{\#}Q_t^\Delta}[h(W)]
% =\mathbb E_{\Delta W\sim Q_t^\Delta}[h(\iota_\theta(\Delta W))]$.
% Taking $h=\widehat L_{S_t}$ gives the claim.

% \textit{(iii)} If $Q_t^\Delta=\mathcal N(\hat{\Delta W}_t,\Sigma_t)$ on $\mathcal W_\Delta$, then
% $(\iota_\theta)_{\#}Q_t^\Delta$ is the degenerate Gaussian
% $\mathcal N(W_0+\hat{\Delta W}_t,\Sigma_{\mathrm{full}})$ with
% $\Sigma_{\mathrm{full}}=\Pi_\Delta \Sigma_t \Pi_\Delta^\top$ and zero variance on $\mathcal W_\Delta^\perp$.
% Applying \textit{(ii)} with $h=\widehat L_{S_t}$ and subtracting
% $\widehat L_{S_t}(W_0+\hat{\Delta W}_t)$ on both sides yields the E--Sh identity.
% \end{proof}








\section*{Appendix C: Proof of Theorem}
\begin{theorem}[Sharpness-aware PAC--Bayes bound for PECL]\label{thm:pecl-bound}
For any $\delta\in(0,1]$, with probability at least $1-\delta$ over $S_1,\ldots,S_T$,
\begin{equation}\label{eq:pecl-main}
\begin{aligned}
&\frac{1}{T}\sum_{t=1}^T
\mathbb{E}_{P_t\sim\mathcal P_{t-1}}\mathbb{E}_{W\sim Q_t}L_{\mathcal{D}_t}(W)
\ \le\
\ \frac{1}{T}\sum_{t=1}^T
\mathbb{E}_{P_t\sim\mathcal P_{t-1}}
\mathbb{E}_{\varepsilon \sim \mathcal{N}(0, \Sigma_t)}[\hat{L}_{S_t}(W_0 + \hat{\Delta W}_t + \varepsilon)
\\
&\quad +\sqrt{\frac{
\displaystyle \sum_{t=1}^T \mathbb{E}_{P_t^\Delta\sim\mathcal P_{t-1}}\!\Big[\mathrm{KL}\!\big(Q_t^\Delta\|P_t^\Delta\big)\Big]
\;+\;\displaystyle \sum_{t=1}^T \mathrm{KL}\!\big(\mathcal P_t\|\mathcal P_{t-1}\big)
\;+\;\log\!\frac{1}{\delta}
}{2\,\displaystyle\sum_{t=1}^T (m_t-1)}}.
\end{aligned}
\end{equation}
\end{theorem}



\begin{proof}


% Fix $\lambda>0$ and consider the $\{\alpha_t\}$-weighted generalization gap
% $\Delta:=\sum_{t=1}^T \alpha_t(\mathcal L_t-\hat{\mathcal L}_t)$ with
% $\alpha_t=(m_t-1)/M_T$, $M_T=\sum_{t=1}^T(m_t-1)$.
% Apply the variational KL inequality (change-of-measure) on the \emph{joint} measures
% $\widehat\Pi_t(dP,dW):=\mathcal P_t(dP)\,Q_t(dW\!\mid\!P)$ versus
% $\Pi_{t-1}(dP,dW):=\mathcal P_{t-1}(dP)\,P(dW)$ to get
% \[
% \lambda\,\Delta \;\le\; \sum_{t=1}^T \mathrm{KL}(\widehat\Pi_t\!\|\Pi_{t-1})
% \;+\;\sum_{t=1}^T\log \mathbb E_{\Pi_{t-1}}
% \exp\!\big\{\lambda\alpha_t(L_{\mathcal D_t}-\widehat L_{S_t})\big\}.
% \]
% Using Hoeffding’s lemma on each block $S_t$ (bounded loss) and independence across $t$ yields
% $\mathbb E\exp\{\sum_{t}\log(\cdots)\}\le \exp\{\lambda^2/(2M_T)\}$.
% By Markov’s inequality and optimizing $\lambda$ we obtain the joint bound
% \[
% \frac1T\sum_{t=1}^T \mathcal L_t \;\le\; \frac1T\sum_{t=1}^T \hat{\mathcal L}_t
% \;+\;\sqrt{\frac{\sum_{t=1}^T \mathrm{KL}(\widehat\Pi_t\!\|\Pi_{t-1})+\log(1/\delta)}
% {2\sum_{t=1}^T (m_t-1)}}.
% \]
% By the KL chain rule on product measures,
% $\mathrm{KL}(\widehat\Pi_t\!\|\Pi_{t-1})
% =\mathrm{KL}(\mathcal P_t\!\|\mathcal P_{t-1})
% +\mathbb E_{P_t\sim\mathcal P_t}\!\mathrm{KL}(Q_t\!\|\!P_t)$,
% giving Theorem~\ref{thm:dynamic-hyper-CL}（联合界）.

% KL is invariant under taking a Dirac product on the frozen factor. Apply Theorem~\ref{thm:dynamic-hyper-CL} with $(P_t,W_t)=(P_t^{\Delta},\Delta W_t)$.
The bound follows by specializing Theorem 2 to the PECL setting through two key decompositions:

By Lemma \ref{lemma:kl-fullspace-frozen}, the full-space KL terms reduce to their adapter-space counterparts:
\[
\mathrm{KL}(Q_t^{\mathrm{full}} \| P_t^{\mathrm{full}}) = \mathrm{KL}(Q_t^\Delta \| P_t^\Delta),
\quad
\mathrm{KL}(\mathcal{P}_t^{\mathrm{full}} \| \mathcal{P}_{t-1}^{\mathrm{full}}) = \mathrm{KL}(\mathcal{P}_t^\Delta \| \mathcal{P}_{t-1}^\Delta).
\]

% \textbf{Step 2: Sharpness decomposition via Lemma 6.}  
Under the Gaussian posterior $Q_t^\Delta = \mathcal{N}(\Delta\hat{W}_t, \Sigma_t)$, Lemma \ref{lemma: Sharpness in PECL} gives:
\[
\mathbb{E}_{W \sim Q_t^{\mathrm{full}}}[\hat{L}_{S_t}(W)] 
= \mathrm{E}_{\varepsilon \sim \mathcal{N}(0, \Sigma_t)}{L}_{S_t}(W_{t-1}+\Delta\hat{W}_t + \varepsilon)
 = \hat{L}_{S_t}(W_{t-1} + \Delta\hat{W}_t) + \mathrm{E\text{-}Sh}_t^\Delta(\Delta\hat{W}_t, \Sigma_t),
\]
where the perturbations are confined to the adapter subspace.

% \textbf{Step 3: Substitution into general bound.}
Substituting these decompositions into Theorem 2 yields the final PECL-specific bound, where all complexity and sharpness terms are now restricted to the trainable LoRA subspace $\mathcal{W}_\Delta$.
% Substituting these decompositions into Theorem 2 yields the final PECL-specific bound, where all complexity and sharpness terms are now restricted to the trainable LoRA subspace $\mathcal{W}_\Delta$.

\end{proof}

\begin{corollary}[SeqLoRA]\label{cor:seq-lora}
If a single adapter is updated sequentially and $P_t^{\Delta}:=Q_{t-1}^{\Delta}$ (conditioning on realized priors),
the hyper-term vanishes and, with probability at least $1-\delta$,
\[
\frac1T\sum_{t=1}^T \mathbb{E}_{w\sim Q_t}L_{\mathcal{D}_t}(W)
\ \le\
\frac1T\sum_{t=1}^T \mathbb{E}_{\epsilon \sim \mathcal{N}(0,\sum)}\widehat L_{S_t}(W_0+\Delta W+\epsilon)
\;+\;
\sqrt{\frac{
\sum_{t=1}^T \mathrm{KL}(Q_t^{\Delta}\|Q_{t-1}^{\Delta})
+\log(1/\delta)}
{2\,\sum_{t=1}^T (m_t-1)}}.
\]
\end{corollary}

\paragraph{What changes from full parameters to LoRA?}
(i) Task KL: $\mathrm{KL}(Q_t\|P_t)\;\to\;\mathrm{KL}(Q_t^{\Delta}\|P_t^{\Delta})$ (support restricted to the LoRA subspace);
(ii) Sharpness: $\mathrm{E\text{-}Sh}_t\;\to\;\mathrm{E\text{-}Sh}^{\Delta}_t$ (perturbations only along trainable directions);
(iii) Hyper-drift: $\sum_t \mathrm{KL}(\mathcal{P}_t\|\mathcal{P}_{t-1})$ persists, reflecting sequential knowledge accumulation at the hyper level.



%==============================
% Sharpness–risk decomposition and its use
%==============================

% \begin{lemma}[Sharpness decomposition under a Gaussian posterior]\label{lem:sharp-decomp}
% Fix a task $t$ and a bounded loss $\ell\in[0,1]$. Let the task-$t$ posterior be
% $Q_t=\mathcal N(\hat W_t,\Sigma_t)$, where $\hat W_t$ is the posterior mean
% and $\Sigma_t\succeq 0$. Then the following identities hold:
% \begin{align}
% \mathbb{E}_{W\sim Q_t}\widehat L_{S_t}(W)
% &= \widehat L_{S_t}(\hat W_t)
%    + \mathrm{E\text{-}Sh}_t(\hat W_t,\Sigma_t), \label{eq:emp-decomp}\\
% \mathbb{E}_{W\sim Q_t}      L_{\mu_t}(W)
% &=      L_{\mu_t}(\hat W_t)
%    + \mathrm{E\text{-}Sh}^{\,\mu}_t(\hat W_t,\Sigma_t), \label{eq:pop-decomp}
% \end{align}
% where
% \[
% \mathrm{E\text{-}Sh}_t(\hat W_t,\Sigma_t)
% := \mathbb{E}_{\varepsilon\sim\mathcal N(0,\Sigma_t)}
%    \!\big[\widehat L_{S_t}(\hat W_t+\varepsilon)-\widehat L_{S_t}(\hat W_t)\big],
% \quad
% \mathrm{E\text{-}Sh}^{\,\mu}_t(\hat W_t,\Sigma_t)
% := \mathbb{E}_{\varepsilon\sim\mathcal N(0,\Sigma_t)}
%    \!\big[L_{\mu_t}(\hat W_t+\varepsilon)-L_{\mu_t}(\hat W_t)\big].
% \]
% \end{lemma}

% \begin{proof}
% Since $W\sim Q_t$ is equivalent in law to $W=\hat W_t+\varepsilon$ with
% $\varepsilon\sim\mathcal N(0,\Sigma_t)$, take expectations and add–subtract
% $\widehat L_{S_t}(\hat W_t)$ (resp.\ $L_{\mu_t}(\hat W_t)$), which yields
% \eqref{eq:emp-decomp} and \eqref{eq:pop-decomp} by definition of
% $\mathrm{E\text{-}Sh}_t$ and $\mathrm{E\text{-}Sh}^{\,\mu}_t$.
% \end{proof}

% \begin{proposition}[Link to the pooled generalization gap]\label{prop:delta-sharp}
% Let $\alpha_t=(m_t-1)/\sum_{s=1}^T(m_s-1)$ and define the pooled gap
% \[
% \Delta:=\sum_{t=1}^T \alpha_t\Big(
% \mathbb{E}_{P_t\sim\mathcal P_{t-1}}\mathbb{E}_{W\sim Q_t}L_{\mu_t}(W)
% -
% \mathbb{E}_{P_t\sim\mathcal P_{t-1}}\mathbb{E}_{W\sim Q_t}\widehat L_{S_t}(W)
% \Big).
% \]
% Under the Gaussian posterior approximation from Lemma~\ref{lem:sharp-decomp},
% \[
% \Delta
% =\sum_{t=1}^T \alpha_t\,
% \mathbb{E}_{P_t\sim\mathcal P_{t-1}}\!
% \Big[
%   \big(L_{\mu_t}(\hat W_t)-\widehat L_{S_t}(\hat W_t)\big)
%  +\big(\mathrm{E\text{-}Sh}^{\,\mu}_t(\hat W_t,\Sigma_t)
%      -\mathrm{E\text{-}Sh}_t(\hat W_t,\Sigma_t)\big)
% \Big].
% \]
% \end{proposition}

% \begin{proof}
% Apply Lemma~\ref{lem:sharp-decomp} to each task and take the outer expectation
% over $P_t\sim\mathcal P_{t-1}$, then sum with weights $\alpha_t$.
% \end{proof}

% \begin{corollary}[Plug-in form of the joint PAC--Bayes bound]
% \label{cor:bound-with-sharpness}
% In Theorem~\textup{\ref{thm:dynamic-hyper-CL}}, replace each empirical term by
% \eqref{eq:emp-decomp}. Then, with probability at least $1-\delta$,
% \[
% \frac{1}{T}\sum_{t=1}^T
% \mathbb{E}_{P_t\sim\mathcal P_{t-1}}\mathbb{E}_{W\sim Q_t}L_{\mu_t}(W)
% \ \le\
% \frac{1}{T}\sum_{t=1}^T
% \mathbb{E}_{P_t\sim\mathcal P_{t-1}}
% \Big[\widehat L_{S_t}(\hat W_t)+\mathrm{E\text{-}Sh}_t(\hat W_t,\Sigma_t)\Big]
% \;+\;
% \sqrt{\frac{\Xi_T+\log(1/\delta)}{2\sum_{t=1}^T (m_t-1)}},
% \]
% where
% \[
% \Xi_T\ :=\
% \sum_{t=1}^T \mathbb{E}_{P_t\sim\mathcal P_{t-1}}\!\big[\mathrm{KL}(Q_t\|P_t)\big]
% \;+\;
% \sum_{t=1}^T \mathrm{KL}\!\big(\mathcal P_t\|\mathcal P_{t-1}\big).
% \]
% \end{corollary}

% \begin{remark}[PECL/LoRA subspace support]\label{rmk:lora-support}
% If PECL/LoRA is used, $W=W_0+\Delta W_t$ with $W_0$ frozen and
% $Q_t=\mathcal N(\hat W_t,\Sigma_t)$ supported on the adapter subspace
% $\mathcal W_t$ (i.e.\ $\mathrm{range}(\Sigma_t)\subseteq\mathcal W_t$). Then
% the KL terms in $\Xi_T$ reduce to the LoRA subspace by Dirac factorization,
% and $\mathrm{E\text{-}Sh}_t$ measures loss sensitivity to perturbations
% \emph{only within} $\mathcal W_t$.
% \end{remark}









% \begin{theorem}[Dynamic-hyperdistribution PAC-Bayes bound for sequential CL]\label{thm:dyn-hyper}
% Let tasks $t=1,\dots,T$ arrive sequentially. Before seeing $S_t$ of size $m_t$, draw a task prior $P_t$ from the previous hyperposterior $\mathcal{P}_{t-1}$. After observing $S_t$, a learning algorithm outputs a posterior $Q_t=Q_t(S_t,P_t)$ over hypotheses/parameters. Assume a bounded loss $\ell\in[0,1]$. Define the expected risk and the empirical risk on task \( t \) as
% $$
% \mathcal{L}_t := \mathbb{E}_{P_t \sim \mathcal{P}_{t-1}} \mathbb{E}_{W_t \sim Q_t(S_t,P_t)} L_{\mathcal{D}_t}(W_t),
% \quad
% \widehat{\mathcal{L}}_t := \mathbb{E}_{P_t \sim \mathcal{P}_{t-1}} \mathbb{E}_{W_t \sim Q_t(S_t,P_t)} \hat{L}_{S_t}(W_t).
% $$
% For any $\delta\in(0,1]$, with probability at least $1-\delta$ over $(S_1,\dots,S_T)$,
% \begin{equation}
% \begin{aligned}
%  &\frac{1}{T}\sum_{t=1}^T \mathcal{L}_t
%  \le
%  \frac{1}{T}\sum_{t=1}^T \widehat{\mathcal{L}}_t \\
%  &\quad+
%  \sqrt{
%  \underbrace{\frac{1}{2H_{m-1}T}
%  \sum_{t=1}^T \mathbb{E}_{P_t\sim\mathcal{P}_{t-1}}\big[\mathrm{KL}(Q_t|P_t)\big]}_{\text{task-level KL}}
%  +
%  \underbrace{\frac{1}{2H_{m-1}T}
%  \sum_{t=1}^T \mathrm{KL}(\mathcal{P}_t|\mathcal{P}_{t-1})}_{\text{hyper-level KL}}
%  +
%  \underbrace{\frac{1}{H_{m-1}}
%  \log\Big(\tfrac{T\bar m_A}{\delta}\Big)}_{\text{statistical term}}
%  }.
%  \end{aligned}
% \end{equation}

% \begin{equation}
%     \begin{aligned}
% \frac{1}{T}\sum_{t=1}^T \mathcal{L}_t
% &\le
% \frac{1}{T}\sum_{t=1}^T \widehat{\mathcal{L}}_t \\
% &+
% \sqrt{\frac{
% \displaystyle \sum_{t=1}^T \mathbb{E}_{P_t\sim\mathcal{P}_{t-1}}\!\big[\mathrm{KL}(Q_t\|P_t)\big]
% +
% \displaystyle \sum_{t=1}^T \mathrm{KL}(\mathcal{P}_t\|\mathcal{P}_{t-1})
% \;+\; 
% T\,\log\!\Big(\tfrac{T\,\bar m_A}{\delta}\Big)}
% {2\,T\,H_{m-1}}}\,.
%     \end{aligned}
% \end{equation}
% where  $H_{m-1}:=\Big(\frac{1}{T}\sum_{t=1}^T \frac{1}{m_t-1}\Big)^{-1}$ be the harmonic mean of $\{m_t-1\}$ and $\bar m_A:=\frac{1}{T}\sum_{t=1}^T m_t$ the arithmetic mean. 
% \end{theorem}

% \begin{proof}
% % \textbf{Step 1 (single-task, conditional).}
% Fix $t$ and condition on the past $\mathcal{F}_{t-1}=\sigma(S_1,\dots,S_{t-1})$. Since $P_t\sim\mathcal{P}_{t-1}$ is $\mathcal{F}_{t-1}$-measurable and fixed before $S_t$, the bounded-loss PAC-Bayes inequality gives, for any $\delta_t\in(0,1]$, with probability at least $1-\delta_t$ over $S_t$:
% \[
% \mathbb{E}_{P_t\sim\mathcal{P}_{t-1}}\mathbb{E}_{W\sim Q_t} L_{\mathcal D_t}(W)
% \le
% \mathbb{E}_{P_t\sim\mathcal{P}_{t-1}}\mathbb{E}_{W\sim Q_t} \hat L_{S_t}(W)
% +
% \sqrt{\frac{\mathbb{E}_{P_t\sim\mathcal{P}_{t-1}}\!\big[\mathrm{KL}(Q_t\|P_t)\big]+\log(m_t/\delta_t)}{2(m_t-1)}}.
% \tag{1}
% \]

% % \textbf{Step 2 (hyper-drift).}
% Accounting for the update $\mathcal{P}_{t-1}\mapsto \mathcal{P}_t$ via the KL chain rule at the hyper level adds $\mathrm{KL}(\mathcal{P}_t\|\mathcal{P}_{t-1})$ to the complexity term for task $t$:
% % \[
% % \sqrt{\frac{\mathbb{E}_{P_t \sim \mathcal{P}_t}\mathrm{KL}(Q_t\|P_t)+\mathrm{KL}(\mathcal{P}_t\|\mathcal{P}_{t-1})+\log(m_t/\delta_t)}{2(m_t-1)}}.
% % \tag{2}
% % \]
% \begin{equation}
% \begin{aligned}
%      \mathbb E_{P\sim {\mathcal P}_{t}}\mathbb E_{W\sim Q_t(\cdot\mid P)} L_{\mathcal D_t}(W)
%  &\le 
%  \mathbb E_{P\sim {\mathcal P}_{t}}\mathbb E_{W\sim Q_t}\hat L_{S_t}(W) \\
%  &+\sqrt{\frac{\mathbb {E}_{P\sim \mathcal{P}_t}\mathrm{KL}(Q_t|P)+\mathrm{KL}({\mathcal P}_t|l{\mathcal P}_{t-1})+\log(m_t/\delta_t)}{2(m_t-1)}}.
% \end{aligned}
% \end{equation}



% % \textbf{Step 3 (union bound with $1-\delta$).}
% Set $\delta_t=\delta/T$ and take a union bound over $t=1,\dots,T$. With probability at least $1-\delta$,
% \[
% \sum_{t=1}^T \mathcal{L}_t 
% \le 
% \sum_{t=1}^T \widehat{\mathcal{L}}_t 
% +
% \sum_{t=1}^T 
% \sqrt{\frac{A_t}{2(m_t-1)}},
% \quad
% A_t:=\mathbb{E}_{P_t}\mathrm{KL}(Q_t\|P_t)+\mathrm{KL}(\mathcal{P}_t\|\mathcal{P}_{t-1})+\log\!\Big(\tfrac{T\,m_t}{\delta}\Big).
% \tag{3}
% \]

% % \textbf{Step 4 (aggregate the square-root terms).}
% By Cauchy--Schwarz,
% \[
% \sum_{t=1}^T \sqrt{\frac{A_t}{2(m_t-1)}} 
% \le 
% \sqrt{\frac{\sum_{t=1}^T A_t}{2}\cdot \sum_{t=1}^T \frac{1}{m_t-1}}
% =
% \sqrt{\frac{\sum_{t=1}^T A_t}{2}\cdot \frac{T}{H_{m-1}} }.
% \tag{4}
% \]
% For the logarithms,
% \[
% \sum_{t=1}^T \log\!\Big(\tfrac{T\,m_t}{\delta}\Big)
% =
% T\log\!\Big(\tfrac{T}{\delta}\Big)+\sum_{t=1}^T \log m_t
% \le
% T\log\!\Big(\tfrac{T}{\delta}\Big)+T\log \bar m_A
% =
% T\,\log\!\Big(\tfrac{T\,\bar m_A}{\delta}\Big),
% \tag{5}
% \]
% where we used Jensen’s inequality for $\log$. Combining \((3)\)–\((5)\) and dividing by $T$ gives the claim.
% \end{proof}

% \paragraph{Comparison to fixed-hyperprior bounds.}
% Unlike transfer-style analyses that assume \emph{independent} task priors from a \emph{fixed} hyperprior and produce terms like $\mathrm{KL}(Q\|P)$ (environment level) plus averaged $\mathbb{E}_{P\sim Q}\mathrm{KL}(Q_i\|P)$, our bound contains an explicit \emph{sequential hyper-drift} penalty $\sum_t \mathrm{KL}(\mathcal{P}_t\|\mathcal{P}_{t-1})$ and conditions each task KL on $P_t\sim\mathcal{P}_{t-1}$, reflecting dynamic accumulation of information in CL.



% \begin{theorem}[Sharpness-aware PAC-Bayes bound under PECL]\label{thm:pecl}
% Under Lemma\ref{lemma:kl-fullspace-frozen} and the setting of Theorem~\ref{thm:dyn-hyper}, with probability at least $1-\delta$,


% \begin{equation}
%     \begin{aligned}
%         \frac{1}{T}\sum_{t=1}^T \mathcal{L}_t
% &\;\le\;  
% \frac{1}{T}\sum_{t=1}^T 
% \Big\{
% \hat L_{S_t}(W_0+\hat{\Delta W}_t)
% + \mathrm{E\text{-}Sh}^{\Delta}_t(\hat{\Delta W}_t,\Sigma_t)
% \Big\} \\
% &+
% \sqrt{\frac{
% \displaystyle \sum_{t=1}^T \mathbb{E}_{P_t}\!\big[\mathrm{KL}(Q_t^{\Delta}\|P_t^{\Delta})\big]
% \;+\;
% \displaystyle \sum_{t=1}^T \mathrm{KL}(\mathcal{P}_t\|\mathcal{P}_{t-1})
% \;+\; 
% T\,\log\!\Big(\tfrac{T\,\bar m_A}{\delta}\Big)}
% {2\,T\,H_{m-1}}}\,,
%     \end{aligned}
% \end{equation}
% where $Q_t^{\Delta}\approx\mathcal{N}(\hat{\Delta W}_t,\Sigma_t)$ on $\mathcal{W}_t$ and
% \[
% \mathrm{E\text{-}Sh}^{\Delta}_t(\hat{\Delta W}_t,\Sigma_t)
% :=
% \mathbb{E}_{\varepsilon\sim\mathcal{N}(0,\Sigma_t)}\!
% \Big[
% \hat L_{S_t}(W_0+\hat{\Delta W}_t+\varepsilon)-\hat L_{S_t}(W_0+\hat{\Delta W}_t)
% \Big].
% \]
% \end{theorem}

% \begin{proof}
% Apply Theorem~\ref{thm:dyn-hyper} to the full measures $(P_t^{\mathrm{full}},Q_t^{\mathrm{full}})$. By Lemma~\ref{lemma:kl-fullspace-frozen},
% $\mathrm{KL}(Q_t^{\mathrm{full}}\|P_t^{\mathrm{full}})=\mathrm{KL}(Q_t^{\Delta}\|P_t^{\Delta})$,
% so the task-level KL collapses to the LoRA subspace. Since $Q_t^{\mathrm{full}}=(T_{W_0})_{\#}Q_t^{\Delta}$ with $Q_t^{\Delta}=\mathcal{N}(\hat{\Delta W}_t,\Sigma_t)$,
% \[
% \mathbb{E}_{W\sim Q_t^{\mathrm{full}}}\hat L_{S_t}(W)
% =
% \hat L_{S_t}(W_0+\hat{\Delta W}_t)
% +
% \mathrm{E\text{-}Sh}^{\Delta}_t(\hat{\Delta W}_t,\Sigma_t),
% \]
% which restricts sharpness to perturbations in $\mathcal{W}_t$. The hyper-drift term $\mathrm{KL}(\mathcal{P}_t\|\mathcal{P}_{t-1})$ remains unchanged. Substitute into Theorem~\ref{thm:dyn-hyper}.
% \end{proof}

% \paragraph{What changes from full parameters to LoRA?}
% (i) Task KL: $\mathrm{KL}(Q_t\|P_t)\;\to\;\mathrm{KL}(Q_t^{\Delta}\|P_t^{\Delta})$ (support restricted to the LoRA subspace);
% (ii) Sharpness: $\mathrm{E\text{-}Sh}_t\;\to\;\mathrm{E\text{-}Sh}^{\Delta}_t$ (perturbations only along trainable directions);
% (iii) Hyper-drift: $\sum_t \mathrm{KL}(\mathcal{P}_t\|\mathcal{P}_{t-1})$ persists, reflecting sequential knowledge accumulation at the hyper level.










% % \begin{theorem}[Dynamic-hyperdistribution PAC-Bayes bound for sequential CL]
% % \label{thm:sharpness-pacbayes-cl}
% % Consider $T$ sequential tasks under the hierarchical PECL process. 
% % For each task $t$, let $P_t \sim \mathcal{P}_{t-1}$ be the task prior and 
% % $Q_t = Q_t(S_t,P_t)$ be the learned task posterior over subspace, where
% % $W_t \sim Q_t(S_t,P_t)$.
% % For any $\delta \in (0,1]$, with probability at least $1-\delta$ over the draw of the task datasets 
% % $S_1,\ldots,S_T$,

% % \begin{equation}
% % \begin{aligned}
% % \frac{1}{T}\,\mathcal{R}_{1:T}
% % &\le
% % \frac{1}{T}
% % \sum_{t=1}^{T}
% % \mathbb{E}_{P_t \sim \mathcal{P}_{t-1}}
% % \;
% % \mathbb{E}_{W_t \sim Q_t(S_t,P_t)}
% % \Big[
% %     L_{S_t}(W_t)
% %     \;+\;
% %     \mathrm{E\text{-}Sh}_t(W_t,\Sigma_t)
% % \Big] \\
% % &+
% % \frac{1}{T(\bar{m}_T - 1)}
% % \sqrt{
% % \frac{
% % \sum_{t=1}^{T}
% % \mathrm{KL}\!\big(
% %     Q_t
% %     \,\big\|\,
% %     P_t
% % \big)
% % +
% % \sum_{t=1}^{T}
% % \mathrm{KL}\!\big(
% %     \mathcal{P}_t
% %     \,\big\|\,
% %     \mathcal{P}_{t-1}
% % \big)
% % \;+\;
% % \log\!\big(\bar{m}_T / \delta\big)
% % }{
% % 2(\bar{m}_T - 1)
% % }
% % },
% % \end{aligned}
% % \end{equation}

% % where
% % $\bar{m}_T
% % :=
% % \frac{T}{\sum_{t=1}^{T} \frac{1}{m_t}}$
% % is the harmonic mean of the task sample sizes.
% % \end{theorem}

% % \begin{proof}[Proof of Theorem~\ref{thm:sharpness-pacbayes-cl}]
% % Fix $t\in\{1,\dots,T\}$. Condition on a task prior $P_t\sim\mathcal{P}_{t-1}$ and its dataset $S_t$ of size $m_t$. 
% % By Theorem~\ref{thm:pac-basic} (bounded loss in $[0,1]$), with probability at least $1-\delta_t$ over $S_t$,
% % \begin{equation}\label{eq:pb-single}
% % \mathbb{E}_{W\sim Q_t(S_t,P_t)}\!\big[L_{\mathcal{D}_t}(W)\big]
% % \;\le\;
% % \mathbb{E}_{W\sim Q_t(S_t,P_t)}\!\big[\hat L_{S_t}(W)\big]
% % +\sqrt{\frac{\mathrm{KL}(Q_t\|P_t)+\log\!\frac{m_t}{\delta_t}}{2(m_t-1)}}.
% % \end{equation}
% % Taking the expectation with respect to $P_t\sim\mathcal{P}_{t-1}$ and using the hyperdistribution KL decomposition (chain rule of KL at the hyper level),
% % \[
% % \mathrm{KL}\!
% % \Big(Q_t(\cdot\mid P_t)\mathcal{P}_{t-1}\,\Big\|\,P_t\mathcal{P}_{t-1}  \Big)
% % \;=\;
% % \mathbb{E}_{P_t\sim\mathcal{P}_{t-1}}
% % \mathrm{KL}(Q_t\|P_t)
% % \;+\;
% % \mathrm{KL}(\mathcal{P}_t\|\mathcal{P}_{t-1})
% % \]







% % Therefore, with probability at least $1-\delta_t$,
% % \begin{equation}\label{eq:pb-single-hyper}
% % \mathbb{E}_{P_t}\mathbb{E}_{W\sim Q_t}\!\big[L_{\mathcal{D}_t}(W)\big]
% % \;\le\;
% % \mathbb{E}_{P_t}\mathbb{E}_{W\sim Q_t}\!\big[\hat L_{S_t}(W)\big]
% % +\sqrt{\frac{\mathbb{E}_{P_t}[\mathrm{KL}(Q_t\|P_t)]+\mathrm{KL}(\mathcal{P}_t\|\mathcal{P}_{t-1})+\log\!\frac{m_t}{\delta_t}}{2(m_t-1)}}.
% % \end{equation}

% % Next, apply the identity induced by the Gaussian posterior $Q_t=\mathcal{N}(\hat W_t,\Sigma_t)$, Substitute into \eqref{eq:pb-single-hyper}, we obtain for each $t$,
% % \begin{equation}\label{eq:per-task}
% % \mathbb{E}_{P_t}\mathbb{E}_{W\sim Q_t}\!\big[L_{\mu_t}(W)\big]
% % \;\le\;
% % \mathbb{E}_{P_t}\!\Big[\hat L_{S_t}(\hat W_t)+\mathrm{E\text{-}Sh}_t(\hat W_t,\Sigma_t)\Big]
% % +\sqrt{\frac{A_t}{2(m_t-1)}},
% \end{equation}
% where we abbreviated
% \[
% A_t
% \;:=\;
% \mathbb{E}_{P_t\sim\mathcal{P}_{t-1}}\!\mathrm{KL}(Q_t\|P_t)
% \;+\;
% \mathrm{KL}(\mathcal{P}_t\|\mathcal{P}_{t-1})
% \;+\;
% \log\!\frac{m_t}{\delta_t}.
% \]

% Now average \eqref{eq:per-task} over $t=1,\dots,T$ and use the definitions 
% $\mathcal{L}_t=\mathbb{E}_{P_t}\mathbb{E}_{W\sim Q_t}[L_{\mu_t}(W)]$ and $\mathcal{R}_{1:T}=\sum_t\mathcal{L}_t$:
% \[
% \frac{1}{T}\,\mathcal{R}_{1:T}
% \;\le\;
% \frac{1}{T}\sum_{t=1}^T
% \mathbb{E}_{P_t}
% \Big[\hat L_{S_t}(\hat W_t)+\mathrm{E\text{-}Sh}_t(\hat W_t,\Sigma_t)\Big]
% \;+\;
% \frac{1}{T}\sum_{t=1}^T
% \sqrt{\frac{A_t}{2(m_t-1)}}.
% \]
% To aggregate the square-root terms with heterogeneous $m_t$, apply Cauchy–Schwarz:
% \[
% \frac{1}{T}\sum_{t=1}^T \sqrt{\frac{A_t}{2(m_t-1)}}
% \;\le\;
% \sqrt{
% \frac{\sum_{t=1}^T A_t}{2T^2}
% \cdot
% \sum_{t=1}^T \frac{1}{m_t-1}
% }
% \;=\;
% \sqrt{
% \frac{\sum_{t=1}^T A_t}{2T^2(\bar m_T-1)}
% },
% \]
% where $\bar m_T := \dfrac{T}{\sum_{t=1}^T \frac{1}{m_t}}$ is the harmonic mean, hence
% $\sum_{t=1}^T \frac{1}{m_t-1}\le \dfrac{T}{\bar m_T-1}$ (the inequality is tight when all $m_t$ are equal; the bound with $\bar m_T$ is standard and simplifies exposition).

% Finally, select a confidence split $\delta_t$ that yields an aggregated term 
% $\sum_{t=1}^T \log\!\frac{m_t}{\delta_t}
% \;\le\; 
% \log\!\frac{\bar m_T}{\delta}$ 
% (e.g., by taking $\delta_t$ proportional to $m_t$ and applying the AM–GM inequality). Combining the pieces gives
% \[
% \frac{1}{T}\,\mathcal{R}_{1:T}
% \;\le\;
% \frac{1}{T}\sum_{t=1}^T
% \mathbb{E}_{P_t}
% \Big[\hat L_{S_t}(\hat W_t)+\mathrm{E\text{-}Sh}_t(\hat W_t,\Sigma_t)\Big]
% +
% \sqrt{
% \frac{
% \sum_{t=1}^{T}
% \Big(
% \mathbb{E}_{P_t}\mathrm{KL}(Q_t\|P_t)
% +
% \mathrm{KL}(\mathcal{P}_t\|\mathcal{P}_{t-1})
% \Big)
% \;+\;
% \log\!\big(\bar{m}_T / \delta\big)
% }{
% 2\,T^2\,(\bar{m}_T - 1)
% }
% },
% \]
% which is exactly \eqref{eq:pecl-main-bound}. 
% Lastly, by Corollary~\ref{cor:lora-factorization}, in PECL the task-level KL term reduces to the LoRA-subspace KL, and the posterior covariance $\Sigma_t$ may be taken to have support in the LoRA subspace, turning $\mathrm{E\text{-}Sh}_t$ into a subspace-restricted sharpness term. This completes the proof.
% \end{proof}





% \begin{lemma}[]
% Let $Q=\delta_a\otimes Q_\Theta$ and $P=\delta_a\otimes P_\Theta$,  where $\delta_a$ is the Dirac measure at a point $a$. Then
% $$
% \mathrm{KL}(Q\,\|\,P)\ =\ \mathrm{KL}(Q_\Theta\,\|\,P_\Theta).
% $$
% \end{lemma}


 % And in the PECL, $P^{\Delta W}$ is the  initial parapmeter of the LoRA and  $Q^{\Delta W}$ is after trained. 

% \begin{lemma}
% For all $\rho>0$, $\theta\in\mathcal H_\Delta$, and $z\in\mathcal Z$,
% $$
% \ell(W(\theta),z)=
% \ell\big(W_0+\Delta W(\theta),z\big)\ \le\
% \max_{\|u\|_\Delta\le\rho}\ \ell\!\big(W_0+(\Delta W(\theta)+\epsilon),\ z\big)
% \ =\ \ell_{\mathrm{Sh}}^\rho(\theta,z).
% $$
% \end{lemma}





%------------------------------
% Theorem: LoRA–OPB bound with SAM
%------------------------------
% \begin{proof}
% We use the OPB inequality \cite{NEURIPS2022_onlineBayes} with the bounded loss
% $\ell_{\mathrm{Sh}}^\rho$.

% \textbf{(i) Regularity of \(\ell_{\mathrm{Sh}}^\rho\) (assumptions required by Online PAC-Bayes).}

% By construction,
%  $$
%  \ell_{\mathrm{Sh}}^\rho(\theta,z)
% :=
%  \max_{\epsilon\in\mathcal H_\Delta,\ |\epsilon|_F\le \rho}
%  \ \ell\big(W_0+\Delta W(\theta)+\epsilon,\ z\big).
%  $$
% Since the base loss is bounded, \(\ell(W,z)\in[0,K]\) for all \(W,z\), and the maximization defining \(\ell_{\mathrm{SAM}}^\rho\) is taken over a bounded perturbation set, it follows immediately that
%  $$
%  0 \le \ell_{\mathrm{Sh}}^\rho(\theta,z) \le K
%  \quad\text{for all }(\theta,z).
%  $$
%  This satisfies the “bounded loss” requirement of the Online PAC-Bayes reference.

% The map
%  $$
%  (\theta,\epsilon)\ \longmapsto\ \ell\big(W_0+\Delta W(\theta)+\epsilon,\ z\big)
%  $$
%  is measurable in \((\theta,\epsilon)\) (a basic property of the loss function and the model composition). The perturbation set
%  $$
%  U_\rho = \{\epsilon\in\mathcal H_\Delta:\ |\epsilon|_F\le \rho\}
%  $$
%  is compact (closed and bounded). By the measurable maximum theorem (e.g., Berge’s maximum theorem), the pointwise supremum of a measurable function over a compact feasible set is measurable in the remaining arguments. Hence
%  $$
%  \ell_{\mathrm{Sh}}^\rho(\theta,z)
%  =
%  \sup_{\epsilon\in U_\rho}\ \ell\big(W_0+(\Delta W(\theta)+\epsilon),\ z\big)
%  $$
%  is a measurable function of \((\theta,z)\). Therefore \(\ell_{\mathrm{SAM}}^\rho\) meets the regularity conditions (boundedness and measurability) required by the Online PAC-Bayes inequalities.




% % By boundedness of $\ell$ and compactness of $\{\epsilon:\|\epsilon\|_\Delta\le\rho\}$, we have
% % $0\le \ell_{\mathrm{SAM}}^\rho(\theta,z)\le K$ and measurability holds (maximum of a measurable
% % map over a compact feasible set).

% \textbf{(ii) Apply OPB to $\ell_{\mathrm{Sh}}^\rho$.}
% According to Thm.2.3 \cite{NEURIPS2022_onlineBayes}, we have the follow conclusion:
% for the bounded loss \(\ell_{\mathrm{SAM}}^\rho\) and an online-predictable prior sequence \({P_t}\), for any \(\lambda>0\) and \(\delta\in(0,1)\), with probability at least \(1-\delta\),
% \begin{equation}
%     \begin{aligned}
% \sum_{i=1}^T \mathbb E_{\theta\sim Q_{i,S}}
% \left[\mathbb E\!\left[\ell_{\mathrm{Sh}}^\rho(\theta,z_i)\ \big|\ \mathcal F_{i-1}\right]\right]
%  &\le
% \sum_{i=1}^m \mathbb E_{\theta\sim Q_{i,S}}\!\big[\ell_{\mathrm{SAM}}^\rho(\theta,z_i)\big] \\
% &+\frac{\mathrm{KL}(Q_{i,S}\,\|\,P_{i,S})}{\lambda}
% +\frac{\lambda m K^2}{2}
% +\frac{\log(1/\delta)}{\lambda}
% \end{aligned}
% \end{equation}


% \textbf{(iii) Compare the original loss with the SAM loss.}
% By feasibility of $\epsilon=0$ in the maximization,
% $$
% \ell\!\big(W(\theta),z\big)=\ell\!\big(W_0+\Delta W(\theta),z\big)
% \ \le\ \ell_{\mathrm{SAM}}^\rho(\theta,z)\qquad\text{for all }\ (\theta,z).
% $$
% Taking $\mathbb E[\cdot\mid\mathcal F_{i-1}]$, then $\mathbb E_{\theta\sim Q_{i,S}}$, and summing over $t$,
% $$
% \sum_{i=1}^m \mathbb E_{\theta\sim Q_{i,S}}
% \!\left[\mathbb E\!\left[\ell\!\big(W(\theta),z_i\big)\ \big|\ \mathcal F_{i-1}\right]\right]
% \ \le\
% \sum_{i=1}^m \mathbb E_{\theta\sim Q_{i,S}}
% \!\left[\mathbb E\!\left[\ell_{\mathrm{SAM}}^\rho(\theta,z_i)\ \big|\ \mathcal F_{i-1}\right]\right].
% $$
% Combining with $(\star)$ gives the claim.
% \end{proof}

%------------------------------
% Lemma A: loss <= SAM loss
%------------------------------


% \begin{proof}
% Choose the feasible perturbation $u=0$.
% \end{proof}

%------------------------------
% Lemma B: product-KL with a Dirac factor
%------------------------------
% \begin{lemma}[KL with a Dirac factor]
% Let $Q=\delta_a\otimes Q_\Theta$ and $P=\delta_a\otimes P_\Theta$. Then
% $$
% \mathrm{KL}(Q\,\|\,P)\ =\ \mathrm{KL}(Q_\Theta\,\|\,P_\Theta).
% $$
% \end{lemma}

% \begin{proof}
% By the Radon--Nikodym derivative and additivity of product measures,
% $$
% \mathrm{KL}(Q\,\|\,P)
% =\int \log\!\left(\frac{dQ_\Theta}{dP_\Theta}\right)\,dQ_\Theta
% +\underbrace{\int \log\!\left(\tfrac{d\delta_a}{d\delta_a}\right)\,d\delta_a}\_{=0}.
% $$
% \end{proof}

%------------------------------
% Corollary: Gibbs posterior
%------------------------------
% \begin{corollary}[Gibbs posterior minimizes the per-step bound]
% Fix $i$ and a prior $P_{i,S}$. Define
% $$
% \mathcal J(Q)\ :=\ \mathbb E_{\theta\sim Q}\!\big[\ell_{\mathrm{SAM}}^\rho(\theta,z_i)\big]
% +\frac{1}{\lambda}\,\mathrm{KL}(Q\,\|\,P_{i,S}).
% $$
% Over all $Q\ll P_{i,S}$, the unique minimizer is the Gibbs posterior
% $$
% \frac{dQ^\star_{i+1,S}}{dP_{i,S}}(\theta)\ \propto\
% \exp\!\left(-\lambda\,\ell_{\mathrm{SAM}}^\rho(\theta,z_i)\right).
% $$
% \end{corollary}

% \begin{proof}
% Apply the Donsker--Varadhan variational identity to $\mathcal J(Q)$.
% \end{proof}

%------------------------------
% Remarks (optional in paper body)
%------------------------------
% \paragraph{Remarks.}
% (i) The optimizable terms of the bound (SAM loss and KL) live entirely in the LoRA subspace $\mathcal H_\Delta$,
% hence promoting flatness where the updates occur is sufficient to tighten generalization---no backbone perturbation is needed.
% (ii) Writing $H=\nabla^2_W\ell(W(\theta),z)$ and $J=\partial W(\theta)/\partial\theta$, for small $\rho$,
% $$
% \ell_{\mathrm{SAM}}^\rho(\theta,z)\ \approx\ \ell\!\big(W(\theta),z\big)\ +\
% \frac{\rho^2}{2}\,\lambda_{\max}\!\big(J^\top HJ\big),
% $$
% so LoRA--SAM directly penalizes the projected Hessian spectrum in the learnable subspace.




% \begin{proof}
    
% \end{proof}



\end{document}
